\documentclass[10pt,a4paper]{article}
\usepackage[margin=2cm]{geometry}
\usepackage{xcolor}
\usepackage{fontspec}
\usepackage{titlesec}
\usepackage{fancyhdr}
\usepackage{hyperref}
\usepackage{parskip}
\usepackage{enumitem}

% Cyberpunk colors
\definecolor{bg}{HTML}{0D0D0D}
\definecolor{cyan}{HTML}{00FFFF}
\definecolor{magenta}{HTML}{FF00FF}
\definecolor{green}{HTML}{00FF00}
\definecolor{yellow}{HTML}{FFFF00}
\definecolor{red}{HTML}{FF0000}
\definecolor{gray}{HTML}{888888}
\definecolor{white}{HTML}{CCCCCC}
\definecolor{darkgray}{HTML}{1A1A2E}

% Page background
\pagecolor{bg}
\color{white}

% Monospace font
\setmonofont{Menlo}[Scale=0.85]
\setmainfont{Helvetica Neue}

% Section styling
\titleformat{\section}{\Large\bfseries\color{cyan}}{}{0em}{}[\color{gray}\titlerule]
\titleformat{\subsection}{\large\bfseries\color{magenta}}{}{0em}{}

% Header/Footer
\pagestyle{fancy}
\fancyhf{}
\fancyhead[L]{\color{gray}\small ZPWR ENCYCLOPEDIA}
\fancyhead[R]{\color{gray}\small\thepage}
\fancyfoot[C]{\color{gray}\small \texttt{zpwr wizard} // MenkeTechnologies}
\renewcommand{\headrulecolor}{\color{cyan}}
\renewcommand{\headrule}{\hbox to\headwidth{\color{cyan}\leaders\hrule height \headrulewidth\hfill}}

% Hyperlinks
\hypersetup{colorlinks=true, linkcolor=cyan, urlcolor=green}

% Tight spacing
\setlength{\parskip}{2pt}
\setlength{\parindent}{0pt}

\begin{document}

% Title page
\begin{titlepage}
\vspace*{3cm}
\begin{center}
{\Huge\bfseries\color{cyan} THE ZPWR ENCYCLOPEDIA}\\[1cm]
{\Large\color{magenta} A Hacker's Field Guide to the Terminal}\\[2cm]
{\large\color{green} 175 Pages // 43 Chapters // 410+ Verbs}\\[1cm]
{\color{gray} By MenkeTechnologies}\\[0.5cm]
{\color{gray} \url{https://github.com/MenkeTechnologies/zpwr}}\\[3cm]
{\color{yellow}\rule{10cm}{0.5pt}}\\[1cm]
{\color{gray}\small Generated from \texttt{zpwr wizard} // Cyberpunk Edition}
\end{center}
\end{titlepage}

\tableofcontents
\newpage


\section{CHAPTER 1: THE RED PILL}

\subsection{Page 1: Welcome to the Machine}

\vspace{4pt}\textbf{\color{cyan}=== WELCOME TO THE MACHINE ===}\vspace{2pt}\\
\color{white}You found it. The rabbit hole.\\
\color{white}ZPWR is not a dotfiles repo. It is not a \textbackslash{}"framework.\textbackslash{}"\\
\color{white}It is a fully integrated shell operating system that happens\\
\color{white}to run inside your terminal.\\
\vspace{4pt}\textbf{\color{magenta}WHAT IS ZPWR?}\\
\color{white}A 400+ verb command dispatcher built on top of zsh, tmux,\\
\color{white}vim/nvim, fzf, and git. Every verb is a surgical tool.\\
\color{white}Every keybinding is a weapon. Every alias is a shortcut\\
\color{white}through the matrix.\\
\vspace{4pt}\textbf{\color{magenta}THE PHILOSOPHY}\\
\texttt{\color{green}1. Verbs, not scripts. Everything is zpwr <verb>.}\\
\color{white}2. Discoverability.   Tab completion on everything.\\
\color{white}3. Symmetry.          vim* and emacs* verbs mirror each other.\\
\color{white}4. Speed.             Compiled zwc files. Benchmarked startup.\\
\color{white}5. No magic.          Read the source. It is all zsh.\\
\vspace{4pt}\textbf{\color{magenta}THE DIRECTORY TREE}\\
\color{white}\textbackslash{}PWR ........................ root of the kingdom\\
\color{white}\textbackslash{}PWR\_ENV .................... env vars and shell init\\
\color{white}\textbackslash{}PWR\_SCRIPTS ................ executable scripts\\
\color{white}\textbackslash{}PWR\_AUTOLOAD ............... autoloaded functions\\
\color{white}\textbackslash{}PWR\_LOCAL .................. local config (gitignored)\\
\color{white}\textbackslash{}PWR\_INSTALL ................ installer scripts\\
\color{white}\textbackslash{}PWR\_TMUX ................... tmux configuration\\
\color{white}\textbackslash{}PWR\_TEST ................... test suite\\
\vspace{4pt}\textbf{\color{magenta}YOUR FIRST COMMAND}\\
\texttt{\color{green}zpwr about    \# behold the banner}\\
\texttt{\color{green}zpwr help     \# list every verb in color}\\
\color{white}Pro tip: zpwr is aliased to many short forms.\\
\texttt{\color{green}Type zpwr <tab> and let the completions wash over you.}\\
\color{white}WARNING: Side effects may include mass alias adoption,\\
\color{white}involuntary tab-completing in other people's terminals,\\
\color{white}and an irrational hatred of vanilla bash.\\

\vspace{8pt}

\subsection{Page 2: Your First 5 Minutes}

\vspace{4pt}\textbf{\color{cyan}=== YOUR FIRST 5 MINUTES ===}\vspace{2pt}\\
\color{white}You just installed zpwr. The terminal blinks. What now?\\
\color{white}Here is your survival kit for the first 300 seconds.\\
\vspace{4pt}\textbf{\color{magenta}STEP 1: SEE THE BANNER}\\
\texttt{\color{green}zpwr about              \# the ASCII art welcome screen}\\
\texttt{\color{green}zpwr banner             \# same thing, different alias}\\
\texttt{\color{green}zpwr bannercounts       \# banner + environment statistics}\\
\vspace{4pt}\textbf{\color{magenta}STEP 2: COUNT YOUR ARSENAL}\\
\texttt{\color{green}zpwr environmentcounts  \# how many aliases, functions, etc.}\\
\color{white}This prints a breakdown of your loaded shell environment:\\
\color{white}aliases, functions, builtins, parameters, completions.\\
\color{white}If the numbers are big, congratulations. You are armed.\\
\vspace{4pt}\textbf{\color{magenta}STEP 3: DISCOVER THE VERBS}\\
\texttt{\color{green}zpwr help               \# colorized list of ALL verbs}\\
\texttt{\color{green}zpwr verbs     \# fzf fuzzy search through verbs}\\
\color{white}The verbsfzf verb is your panic button (alias: subcommandsfzf).\\
\color{white}Forgot what a command is called? Fuzzy search it.\\
\color{white}This pipes every verb and its description into fzf.\\
\color{white}Type a few characters. Watch the list collapse.\\
\color{white}Hit enter. Execute.\\
\vspace{4pt}\textbf{\color{magenta}STEP 4: VERIFY YOUR ENVIRONMENT}\\
\texttt{\color{green}zpwr doctor             \# run environment health check}\\
\texttt{\color{green}zpwr envvars            \# list all zpwr env vars}\\
\texttt{\color{green}zpwr colortest          \# verify your terminal colors}\\
\texttt{\color{green}zpwr colortest256       \# the full 256-color palette}\\
\vspace{4pt}\textbf{\color{magenta}STEP 5: TAB COMPLETE EVERYTHING}\\
\texttt{\color{green}Type zpwr and press <tab>. You will see 400+ verbs.}\\
\texttt{\color{green}Type zpwr git and press <tab>. Filtered to git verbs.}\\
\texttt{\color{green}Type zpwr vim and press <tab>. Filtered to vim verbs.}\\
\color{white}The completion system knows every verb. Trust it.\\
\color{white}If you only remember one command from this page:\\
\color{white}zpwr verbs\\
\color{white}It is the cheat code to everything else.\\

\vspace{8pt}

\subsection{Page 3: The Anatomy of a Verb}

\vspace{4pt}\textbf{\color{cyan}=== THE ANATOMY OF A VERB ===}\vspace{2pt}\\
\color{white}Every zpwr command follows one sacred pattern:\\
\color{white}zpwr <verb> [args...]\\
\vspace{4pt}\textbf{\color{magenta}HOW DISPATCH WORKS}\\
\texttt{\color{green}When you type zpwr filesearch, here is what happens:}\\
\color{white}1. The zpwr function receives \textbackslash{}"filesearch\textbackslash{}" as \textbackslash{}\$1\\
\color{white}2. It looks up ZPWR\_VERBS[filesearch] in the verb table\\
\color{white}3. The table entry maps to a function or alias\\
\color{white}4. That function is invoked with any remaining args\\
\vspace{4pt}\textbf{\color{magenta}THE VERB TABLE}\\
\color{white}The associative array ZPWR\_VERBS is the brain.\\
\color{white}Each entry has the format:\\
\color{white}ZPWR\_VERBS[verb]='function\_or\_alias=description'\\
\vspace{4pt}\textbf{\color{magenta}MANY NAMES, ONE FUNCTION}\\
\color{white}Verbs often have aliases for muscle memory:\\
\texttt{\color{green}zpwr forward  = zpwr f   = zpwr fwd}\\
\texttt{\color{green}zpwr reverse  = zpwr r   = zpwr rvs = zpwr cdup}\\
\texttt{\color{green}zpwr github   = zpwr gh  = zpwr mygithub}\\
\texttt{\color{green}zpwr repo     = zpwr zp  = zpwr zpwr}\\
\vspace{4pt}\textbf{\color{magenta}TAB COMPLETION IN ACTION}\\
\color{white}zpwr git<tab>     \# shows: gitcommit, gitconfig, gitsearch...\\
\color{white}zpwr vim<tab>     \# shows: vimall, vimrecent, vimctags...\\
\color{white}zpwr home<tab>    \# shows: home, homeenv, homelocal...\\
\color{white}zpwr forgit<tab>  \# shows: forgitadd, forgitlog, forgitdiff...\\
\vspace{4pt}\textbf{\color{magenta}INSPECT THE TABLE YOURSELF}\\
\texttt{\color{green}zpwr subcommands        \# dump all verbs to stdout}\\
\texttt{\color{green}zpwr verbs     \# fuzzy search verbs with fzf}\\
\texttt{\color{green}zpwr verbscount          \# how many verbs exist}\\
\vspace{4pt}\textbf{\color{magenta}THE EDIT SUFFIX PATTERN}\\
\color{white}Many verbs come in pairs:\\
\color{white}filesearch      -> open file after selection\\
\color{white}filesearchedit  -> type into command line for editing\\
\color{white}The \textbackslash{}"edit\textbackslash{}" variant puts the result on your command line\\
\color{white}instead of executing immediately. Power and precision.\\

\vspace{8pt}

\subsection{Page 4: Configuring Your Rig}

\vspace{4pt}\textbf{\color{cyan}=== CONFIGURING YOUR RIG ===}\vspace{2pt}\\
\color{white}Your terminal is a cockpit. These are the dials.\\
\vspace{4pt}\textbf{\color{magenta}THE FOUR CONFIG SESSIONS}\\
\color{white}Each of these opens a curated vim session with the\\
\color{white}relevant config files already loaded in buffers:\\
\texttt{\color{green}zpwr brc    \# shell aliases file vim session}\\
\color{white}opens \textbackslash{}PWR\_LOCAL/.tokens.sh and related files\\
\color{white}this is where YOUR personal aliases live\\
\texttt{\color{green}zpwr zrc    \# zshrc vim session}\\
\color{white}opens the zshrc and related zsh configs\\
\color{white}the source of all shell initialization\\
\texttt{\color{green}zpwr vrc    \# vimrc/init.vim vim session}\\
\color{white}opens vim/nvim config files\\
\color{white}your editor's soul laid bare\\
\texttt{\color{green}zpwr trc    \# tmux.conf vim session}\\
\color{white}opens tmux configuration\\
\color{white}pane and window management settings\\
\vspace{4pt}\textbf{\color{magenta}THE TOKENS FILE}\\
\texttt{\color{green}zpwr tokens       \# vim the .tokens.sh file}\\
\texttt{\color{green}zpwr vimtokens    \# same thing, explicit vim}\\
\texttt{\color{green}zpwr emacstokens  \# same thing, emacs edition}\\
\color{white}Located at \textbackslash{}PWR\_LOCAL/.tokens.sh\\
\color{white}This file is gitignored. Put secrets and custom aliases here.\\
\color{white}Never commit this file. It is YOUR local override.\\
\vspace{4pt}\textbf{\color{magenta}EDITOR COMMANDS}\\
\texttt{\color{green}zpwr edit      \# launch \textbackslash{}DITOR}\\
\texttt{\color{green}zpwr editor    \# same thing}\\
\texttt{\color{green}zpwr vimconfig \# edit all zpwr configs in vim}\\
\vspace{4pt}\textbf{\color{magenta}THE FULL CONFIG STACK}\\
\color{white}\textbackslash{}PWR\_LOCAL/.tokens.sh ... your aliases and secrets\\
\color{white}\textbackslash{}PWR\_ENV/.zpwr\_env.sh .. env vars and paths\\
\color{white}\textasciitilde{}/.zshrc .................. shell entry point\\
\color{white}\textasciitilde{}/.vimrc / init.vim ....... editor config\\
\color{white}\textasciitilde{}/.tmux.conf .............. tmux config\\
\color{white}Think of it as: brc for shell, zrc for zsh,\\
\color{white}vrc for vim, trc for tmux. Four letters, four kingdoms.\\

\vspace{8pt}


\section{CHAPTER 2: NAVIGATION PROTOCOL}

\subsection{Page 5: Directory Warfare}

\vspace{4pt}\textbf{\color{cyan}=== DIRECTORY WARFARE ===}\vspace{2pt}\\
\color{white}In the filesystem, speed is survival.\\
\color{white}Typing full paths is for people who have given up on life.\\
\vspace{4pt}\textbf{\color{magenta}FORWARD AND REVERSE}\\
\texttt{\color{green}zpwr forward <dir>  \# cd into a child directory}\\
\texttt{\color{green}zpwr f <dir>        \# same thing, shorter}\\
\texttt{\color{green}zpwr fwd <dir>      \# same thing, medium length}\\
\texttt{\color{green}zpwr reverse <dir>  \# cd UP the tree to a parent matching <dir>}\\
\texttt{\color{green}zpwr r <dir>        \# same thing, one letter}\\
\texttt{\color{green}zpwr rvs <dir>      \# same thing, three letters}\\
\texttt{\color{green}zpwr cdup <dir>     \# same thing, cd-up style}\\
\color{white}Example: you are in \textasciitilde{}/code/myproject/src/lib/utils\\
\texttt{\color{green}zpwr r myproject    \# jumps UP to \textasciitilde{}/code/myproject}\\
\color{white}No more ../../../../ chains. Navigate like a human.\\
\vspace{4pt}\textbf{\color{magenta}BASIC CD}\\
\texttt{\color{green}zpwr cd <path>      \# cd to directory arg}\\
\color{white}Seems redundant but integrates with zpwr's dispatch.\\
\vspace{4pt}\textbf{\color{magenta}Z: FRECENCY JUMPING}\\
\texttt{\color{green}zpwr z <query>      \# jump to frecency-ranked directory}\\
\color{white}z tracks where you go. The more you visit a directory,\\
\color{white}the higher it ranks. Type a fragment, land at the top match.\\
\texttt{\color{green}zpwr zcd            \# fzf list of z-ranked dirs, pick one}\\
\color{white}When you cannot remember the exact fragment,\\
\color{white}zcd shows you the full ranked list in fzf.\\
\vspace{4pt}\textbf{\color{magenta}PRACTICAL COMBOS}\\
\color{white}zpwr z proj         \# land in your most-visited \textbackslash{}"proj\textbackslash{}" dir\\
\color{white}zpwr r src          \# climb back up to the \textbackslash{}"src\textbackslash{}" parent\\
\color{white}zpwr f tests        \# dive into the \textbackslash{}"tests\textbackslash{}" child\\
\color{white}The pattern is: z for teleportation, f for descent,\\
\color{white}r for ascent. Three directions. Total coverage.\\
\color{white}PROTIP: If you are still typing cd ../../../,\\
\color{white}this is your intervention.\\

\vspace{8pt}

\subsection{Page 6: Home Base Operations}

\vspace{4pt}\textbf{\color{cyan}=== HOME BASE OPERATIONS ===}\vspace{2pt}\\
\color{white}Every zpwr directory has a one-word teleporter.\\
\color{white}No paths. No thinking. Just say where.\\
\vspace{4pt}\textbf{\color{magenta}THE HOME VERB FAMILY}\\
\texttt{\color{green}zpwr home                \# cd \textbackslash{}PWR (the root)}\\
\texttt{\color{green}zpwr homelocal           \# cd \textbackslash{}PWR\_LOCAL}\\
\texttt{\color{green}zpwr homescripts         \# cd \textbackslash{}PWR\_SCRIPTS}\\
\texttt{\color{green}zpwr homeenv             \# cd \textbackslash{}PWR\_ENV}\\
\texttt{\color{green}zpwr homeinstall         \# cd \textbackslash{}PWR\_INSTALL}\\
\texttt{\color{green}zpwr homeautoload        \# cd \textbackslash{}PWR\_AUTOLOAD}\\
\texttt{\color{green}zpwr homeautoloadcommon   \# cd \textbackslash{}PWR\_AUTOLOAD\_COMMON}\\
\texttt{\color{green}zpwr homecomps           \# cd \textbackslash{}PWR\_COMPS}\\
\texttt{\color{green}zpwr hometest            \# cd \textbackslash{}PWR\_TEST}\\
\texttt{\color{green}zpwr hometmux            \# cd \textbackslash{}PWR\_TMUX}\\
\texttt{\color{green}zpwr homeexercism        \# cd \textbackslash{}OME/Exercism}\\
\vspace{4pt}\textbf{\color{magenta}WHAT LIVES WHERE}\\
\color{white}home .............. the zpwr repo root, start of everything\\
\color{white}homelocal ........ gitignored personal configs and tokens\\
\color{white}homescripts ...... executable shell scripts\\
\color{white}homeenv .......... environment variables and init files\\
\color{white}homeinstall ...... installer scripts for dependencies\\
\color{white}homeautoload ..... autoloaded functions (the engine room)\\
\color{white}homeautoloadcommon  platform-independent autoloads\\
\color{white}homecomps ........ zsh completion definitions\\
\color{white}hometest ......... test suite files\\
\color{white}hometmux ......... tmux configuration files\\
\vspace{4pt}\textbf{\color{magenta}EXAMPLE WORKFLOW}\\
\color{white}You want to edit an autoload function:\\
\texttt{\color{green}zpwr homeautoloadcommon}\\
\color{white}\# now you are in the autoload directory\\
\color{white}\# ls, grep, edit to your heart's content\\
\color{white}You want to check your local tokens:\\
\texttt{\color{green}zpwr homelocal}\\
\color{white}\# lands you right next to .tokens.sh\\
\color{white}Every home* verb is a bookmark you never need to set.\\
\color{white}The bookmarks are built into the system.\\

\vspace{8pt}

\subsection{Page 7: Named Directories}

\vspace{4pt}\textbf{\color{cyan}=== NAMED DIRECTORIES ===}\vspace{2pt}\\
\color{white}Zsh has a secret weapon most people never learn:\\
\color{white}named directories. They turn env vars into tilde paths.\\
\vspace{4pt}\textbf{\color{magenta}HOW IT WORKS}\\
\color{white}When you do hash -d ZPWR=\textbackslash{}PWR, zsh registers a\\
\color{white}named directory. Now \textasciitilde{}ZPWR expands to \textbackslash{}PWR everywhere:\\
\color{white}in cd, in tab completion, in prompt, in arguments.\\
\vspace{4pt}\textbf{\color{magenta}ZPWR'S NAMED DIRECTORIES}\\
\color{white}\textasciitilde{}ZPWR ................. \textbackslash{}PWR root\\
\color{white}\textasciitilde{}ZPWR\_SCRIPTS ........ \textbackslash{}PWR\_SCRIPTS\\
\color{white}\textasciitilde{}ZPWR\_AUTOLOAD ....... \textbackslash{}PWR\_AUTOLOAD\\
\color{white}\textasciitilde{}ZPWR\_ENV ............ \textbackslash{}PWR\_ENV\\
\color{white}\textasciitilde{}ZPWR\_LOCAL .......... \textbackslash{}PWR\_LOCAL\\
\color{white}\textasciitilde{}ZPWR\_INSTALL ........ \textbackslash{}PWR\_INSTALL\\
\color{white}\textasciitilde{}ZPWR\_TMUX ........... \textbackslash{}PWR\_TMUX\\
\color{white}\textasciitilde{}ZPWR\_TEST ........... \textbackslash{}PWR\_TEST\\
\color{white}\textasciitilde{}ZPWR\_COMPS .......... \textbackslash{}PWR\_COMPS\\
\vspace{4pt}\textbf{\color{magenta}WHAT THIS MEANS IN PRACTICE}\\
\color{white}Your prompt might show:\\
\color{white}\textasciitilde{}ZPWR\_AUTOLOAD/common  instead of the full path\\
\color{white}You can type:\\
\color{white}cd \textasciitilde{}ZPWR\_SCRIPTS  and it just works\\
\texttt{\color{green}vim \textasciitilde{}ZPWR\_ENV/.zpwr\_env.sh  tilde paths everywhere}\\
\vspace{4pt}\textbf{\color{magenta}TAB COMPLETION WITH TILDE}\\
\color{white}Type cd \textasciitilde{}ZPWR<tab> and watch zsh expand the options:\\
\color{white}\textasciitilde{}ZPWR  \textasciitilde{}ZPWR\_SCRIPTS  \textasciitilde{}ZPWR\_ENV  \textasciitilde{}ZPWR\_LOCAL ...\\
\color{white}This works in any command, not just cd.\\
\vspace{4pt}\textbf{\color{magenta}THE MECHANISM}\\
\color{white}Zsh option AUTO\_NAME\_DIRS turns any exported scalar\\
\color{white}parameter into a named directory automatically.\\
\color{white}ZPWR exports its directory variables, so they become\\
\color{white}named directories for free.\\
\color{white}Think of named directories as human-readable symlinks\\
\color{white}that live in your shell's brain instead of the filesystem.\\
\color{white}No files created. No cleanup needed. Pure memory.\\

\vspace{8pt}

\subsection{Page 8: Quick Jumps}

\vspace{4pt}\textbf{\color{cyan}=== QUICK JUMPS ===}\vspace{2pt}\\
\texttt{\color{green}The home* verbs land you in zpwr directories.}\\
\color{white}These verbs land you in the ecosystem around zpwr.\\
\vspace{4pt}\textbf{\color{magenta}ZPWR INTERNALS}\\
\texttt{\color{green}zpwr scripts       \# cd to scripts directory}\\
\color{white}shortcut alias: zs\\
\texttt{\color{green}zpwr autoload      \# cd to autoload directory}\\
\color{white}shortcut alias: zal\\
\texttt{\color{green}zpwr completions   \# cd to completion directory}\\
\texttt{\color{green}zpwr comps         \# same thing}\\
\color{white}shortcut alias: zco\\
\vspace{4pt}\textbf{\color{magenta}PLUGIN ECOSYSTEM}\\
\texttt{\color{green}zpwr plugins       \# cd to \textbackslash{}SH\_CUSTOM/plugins}\\
\color{white}shortcut alias: zpl\\
\color{white}This is where oh-my-zsh custom plugins live.\\
\color{white}Your third-party zsh plugins are neighbors here.\\
\texttt{\color{green}zpwr forked        \# cd to \textbackslash{}PWR\_FORKED\_DIR}\\
\texttt{\color{green}zpwr fp            \# same thing}\\
\color{white}Forked repos that zpwr depends on live here.\\
\color{white}Useful when you need to hack on a dependency.\\
\vspace{4pt}\textbf{\color{magenta}SPECIAL DIRECTORIES}\\
\texttt{\color{green}zpwr hidden        \# cd to \textbackslash{}PWR\_HIDDEN\_DIR}\\
\color{white}The hidden directory. Dot-prefixed internal state.\\
\texttt{\color{green}zpwr repo          \# cd to \textbackslash{}PWR\_REPO\_NAME}\\
\texttt{\color{green}zpwr zp            \# same thing}\\
\texttt{\color{green}zpwr zpwr          \# same thing (yes, zpwr zpwr)}\\
\vspace{4pt}\textbf{\color{magenta}THE NAVIGATION CHEAT SHEET}\\
\color{white}Use case                      Verb\\
\color{white}--------                      ----\\
\texttt{\color{green}Edit a script .............. zpwr scripts}\\
\texttt{\color{green}Fix a completion ........... zpwr comps}\\
\texttt{\color{green}Add an autoload function .. zpwr autoload}\\
\texttt{\color{green}Debug a plugin ............ zpwr plugins}\\
\texttt{\color{green}Patch a fork .............. zpwr forked}\\
\color{white}Each of these is one verb, zero path typing,\\
\color{white}and instant teleportation to where the work is.\\

\vspace{8pt}


\section{CHAPTER 3: THE GIT DIMENSION}

\subsection{Page 9: Git Basics}

\vspace{4pt}\textbf{\color{cyan}=== GIT BASICS ===}\vspace{2pt}\\
\color{white}Git is already a language. ZPWR makes it a dialect\\
\color{white}you can actually speak without consulting the man page.\\
\vspace{4pt}\textbf{\color{magenta}THE COMMIT PIPELINE}\\
\texttt{\color{green}zpwr gitcommit \textbackslash{}"your message here\textbackslash{}"}\\
\color{white}This is not just git commit. It does:\\
\color{white}git add -> git commit -> git push\\
\color{white}All in one verb. The full gacp (git add commit push).\\
\color{white}Because who commits without pushing? Cowards.\\
\vspace{4pt}\textbf{\color{magenta}CONFIG AND IGNORES}\\
\texttt{\color{green}zpwr gitconfig         \# edit git config file}\\
\texttt{\color{green}zpwr gitglobalconfig   \# same thing, explicit name}\\
\texttt{\color{green}zpwr gitignore         \# edit .git/info/exclude (local)}\\
\texttt{\color{green}zpwr gitignores        \# edit global git ignores file}\\
\texttt{\color{green}zpwr gitglobalignores  \# same as gitignores}\\
\vspace{4pt}\textbf{\color{magenta}STASH AND DIFF (VIA FORGIT)}\\
\color{white}ZPWR delegates interactive stash and diff to forgit:\\
\texttt{\color{green}zpwr forgitdiff        \# fzf-powered git diff viewer}\\
\texttt{\color{green}zpwr forgitstash       \# fzf-powered stash browser}\\
\color{white}See page 13 for the full forgit arsenal.\\
\vspace{4pt}\textbf{\color{magenta}REMOTE OPERATIONS}\\
\texttt{\color{green}zpwr gitremotes        \# list all git remotes}\\
\texttt{\color{green}zpwr gitdir            \# check if pwd is a git dir}\\
\vspace{4pt}\textbf{\color{magenta}TAGS}\\
\texttt{\color{green}zpwr gitedittag        \# print tag to \textbackslash{}UFFER}\\
\texttt{\color{green}zpwr gitupdatetag      \# print tag and latest msg to \textbackslash{}UFFER}\\
\vspace{4pt}\textbf{\color{magenta}REAL-WORLD EXAMPLE}\\
\color{white}\# You just finished a feature. Ship it:\\
\texttt{\color{green}zpwr gitcommit \textbackslash{}"add quantum flux capacitor\textbackslash{}"}\\
\color{white}\# Stages everything, commits with message, pushes.\\
\color{white}\# One line. Done. Next feature.\\
\color{white}\# Need to check your git config:\\
\texttt{\color{green}zpwr gitconfig}\\
\color{white}\# Opens it in \textbackslash{}DITOR. Fix, save, done.\\
\color{white}Git verbs follow the pattern: git + action noun.\\
\color{white}If you can name it, zpwr probably has a verb for it.\\

\vspace{8pt}

\subsection{Page 10: GitHub Integration}

\vspace{4pt}\textbf{\color{cyan}=== GITHUB INTEGRATION ===}\vspace{2pt}\\
\color{white}GitHub is not a website. It is an API with a website.\\
\color{white}ZPWR treats it accordingly.\\
\vspace{4pt}\textbf{\color{magenta}CREATE AND DESTROY}\\
\texttt{\color{green}zpwr githubcreate      \# create a remote GitHub repo}\\
\texttt{\color{green}zpwr hc               \# same thing, two letters}\\
\color{white}Creates a repo on GitHub from your current directory.\\
\color{white}Sets up the remote. Ready to push.\\
\texttt{\color{green}zpwr githubdelete      \# delete a remote GitHub repo}\\
\texttt{\color{green}zpwr hd               \# same thing, two letters}\\
\color{white}This is permanent. It asks for confirmation.\\
\color{white}But once confirmed, the repo is gone. Poof.\\
\vspace{4pt}\textbf{\color{magenta}BROWSE YOUR GITHUB}\\
\texttt{\color{green}zpwr github            \# open your GitHub profile in browser}\\
\texttt{\color{green}zpwr gh               \# same thing}\\
\texttt{\color{green}zpwr mygithub          \# same thing}\\
\color{white}Opens your profile page. Quick ego check.\\
\texttt{\color{green}zpwr githubzpwr        \# open zpwr repo on GitHub}\\
\texttt{\color{green}zpwr ghz              \# same thing}\\
\color{white}Jump straight to the zpwr source. File issues.\\
\color{white}Star it. You know you want to.\\
\vspace{4pt}\textbf{\color{magenta}CONTRIBUTION STATS}\\
\texttt{\color{green}zpwr githubcontribcount \# count GitHub contribs in last year}\\
\color{white}Tells you how green your contribution graph is.\\
\color{white}The number that keeps you pushing commits at 2 AM.\\
\vspace{4pt}\textbf{\color{magenta}REVEAL: THE BROWSER SHORTCUT}\\
\texttt{\color{green}zpwr reveal            \# show current repo in browser}\\
\color{white}You are in a git repo. You want to see it on GitHub.\\
\color{white}One word. Browser opens. You are there.\\
\vspace{4pt}\textbf{\color{magenta}WORKFLOW EXAMPLE}\\
\color{white}\# Start a new project from scratch:\\
\color{white}mkdir new-project \&\& cd new-project\\
\color{white}git init\\
\texttt{\color{green}zpwr hc               \# repo created on GitHub}\\
\texttt{\color{green}zpwr gitcommit \textbackslash{}"initial commit\textbackslash{}"}\\
\texttt{\color{green}zpwr reveal            \# admire it in the browser}\\
\color{white}From zero to pushed repo in four commands.\\
\color{white}That is the GitHub integration promise.\\

\vspace{8pt}

\subsection{Page 11: Git Archaeology}

\vspace{4pt}\textbf{\color{cyan}=== GIT ARCHAEOLOGY ===}\vspace{2pt}\\
\color{white}The past is not dead. It is not even past.\\
\color{white}It is stored in .git/objects and zpwr can dig it up.\\
\vspace{4pt}\textbf{\color{magenta}SEARCHING THE LOG}\\
\texttt{\color{green}zpwr gitsearch             \# search for regex in git log}\\
\color{white}Uses fzf to grep through your entire commit history.\\
\color{white}That one commit where you \textbackslash{}"fixed the thing\textbackslash{}"? Found.\\
\texttt{\color{green}zpwr gitcommits            \# search git commits with fzf}\\
\color{white}Interactive commit browser. Preview diffs inline.\\
\color{white}Like git log but you can actually find things.\\
\vspace{4pt}\textbf{\color{magenta}COUNTING COMMITS}\\
\texttt{\color{green}zpwr gitcommitcount        \# total number of commits}\\
\color{white}Just a number. But what a number.\\
\color{white}The sum total of your decisions in this repo.\\
\vspace{4pt}\textbf{\color{magenta}CONTRIBUTION ANALYSIS}\\
\texttt{\color{green}zpwr gitcontribcount       \# contributions by author}\\
\color{white}Who committed the most? The leaderboard of blame.\\
\texttt{\color{green}zpwr gitcontribcountdirs   \# contributions across dirs}\\
\color{white}Same analysis but across multiple repositories.\\
\color{white}Useful for org-wide contribution audits.\\
\texttt{\color{green}zpwr gitcontribcountlines  \# lines contributed by author}\\
\color{white}Lines of code per author. The metric that lies to you.\\
\color{white}But it is a fun lie.\\
\vspace{4pt}\textbf{\color{magenta}SIZE ANALYSIS}\\
\texttt{\color{green}zpwr gittotallines         \# total line count of git files}\\
\color{white}How big is this repo, really? Now you know.\\
\texttt{\color{green}zpwr gitlargest            \# largest files in git history}\\
\color{white}Shows the biggest blobs stored by git, descending.\\
\color{white}That 50MB binary someone committed in 2019? There it is.\\
\color{white}Now you know what to purge (see page 12).\\
\vspace{4pt}\textbf{\color{magenta}DIG SESSION EXAMPLE}\\
\color{white}\# Who wrote the most code in this repo?\\
\texttt{\color{green}zpwr gitcontribcount}\\
\color{white}\# What is eating all the disk space?\\
\texttt{\color{green}zpwr gitlargest}\\
\color{white}\# Find that commit about the database migration:\\
\texttt{\color{green}zpwr gitsearch}\\
\color{white}\# Type \textbackslash{}"migration\textbackslash{}" in fzf. There it is.\\

\vspace{8pt}

\subsection{Page 12: Git Surgery}

\vspace{4pt}\textbf{\color{cyan}=== GIT SURGERY ===}\vspace{2pt}\\
\color{white}Sometimes you need to rewrite history.\\
\color{white}Not the geopolitical kind. The git kind.\\
\color{white}These are destructive operations. Use with intent.\\
\vspace{4pt}\textbf{\color{magenta}CLEARING THE CACHE}\\
\texttt{\color{green}zpwr gitclearcache     \# clear old git refs and objects}\\
\color{white}Runs garbage collection on git's internal state.\\
\color{white}Prunes unreachable objects, packs loose refs.\\
\color{white}Your .git directory gets a deep clean.\\
\vspace{4pt}\textbf{\color{magenta}REMOVING COMMITS}\\
\texttt{\color{green}zpwr gitclearcommit    \# remove matching commits from history}\\
\color{white}This is filter-branch territory.\\
\color{white}Rewrites history. Force push required after.\\
\color{white}Use when someone committed passwords, API keys,\\
\color{white}or that embarrassing commented-out code block.\\
\vspace{4pt}\textbf{\color{magenta}REMOVING FILES FROM ALL HISTORY}\\
\texttt{\color{green}zpwr gitclearfile      \# rm file from all git refs and objects}\\
\color{white}Nuclear option for files.\\
\color{white}That 200MB video someone committed in 2020?\\
\color{white}This removes it from every commit that ever existed.\\
\color{white}Your repo size drops. Your clone time drops.\\
\color{white}Your blood pressure drops.\\
\vspace{4pt}\textbf{\color{magenta}EMAIL SURGERY}\\
\texttt{\color{green}zpwr gitemail          \# change email with git filter-branch}\\
\color{white}Changes both author and committer email across history.\\
\texttt{\color{green}zpwr gitaemail         \# change author email only}\\
\color{white}When you committed with your personal email at work.\\
\texttt{\color{green}zpwr gitcemail         \# change committer email only}\\
\color{white}When the committer field has the wrong identity.\\
\vspace{4pt}\textbf{\color{magenta}WHEN TO USE THESE}\\
\color{white}Scenario: you committed with wrong@email.com\\
\texttt{\color{green}zpwr gitemail}\\
\color{white}\# Follow the prompts to fix every commit\\
\color{white}Scenario: .env file in git history with secrets\\
\texttt{\color{green}zpwr gitclearfile}\\
\color{white}\# Purges it from every ref and object\\
\color{white}REMEMBER: These rewrite history.\\
\color{white}Coordinate with your team before force pushing.\\
\color{white}Or do not coordinate and watch the chaos unfold.\\

\vspace{8pt}

\subsection{Page 13: Forgit: Git on Steroids}

\vspace{4pt}\textbf{\color{cyan}=== FORGIT: GIT ON STEROIDS ===}\vspace{2pt}\\
\color{white}Forgit wraps common git operations with fzf previews.\\
\color{white}Every action gets a live preview pane. Every choice\\
\color{white}is informed. No more blind git commands.\\
\vspace{4pt}\textbf{\color{magenta}STAGING AND UNSTAGING}\\
\texttt{\color{green}zpwr forgitadd            \# fzf-powered git add}\\
\color{white}Shows unstaged files with inline diffs in the preview.\\
\color{white}Select files with tab. Enter to stage. Visual staging.\\
\texttt{\color{green}zpwr forgitreset          \# fzf-powered git reset HEAD}\\
\color{white}The reverse of forgitadd. Unstage files visually.\\
\vspace{4pt}\textbf{\color{magenta}VIEWING CHANGES}\\
\texttt{\color{green}zpwr forgitdiff           \# fzf-powered git diff viewer}\\
\color{white}Browse changed files. Preview shows colorized diffs.\\
\color{white}Side-by-side comparison in your terminal.\\
\texttt{\color{green}zpwr forgitlog            \# fzf-powered git log}\\
\color{white}Commit history with preview. See the diff for each\\
\color{white}commit without leaving the log view. Beautiful.\\
\vspace{4pt}\textbf{\color{magenta}STASH OPERATIONS}\\
\texttt{\color{green}zpwr forgitstash          \# fzf-powered stash browser}\\
\color{white}Lists stashes. Preview shows what each stash contains.\\
\color{white}Pick one, apply it. No more \textbackslash{}"was it stash@\{3\}?\textbackslash{}"\\
\vspace{4pt}\textbf{\color{magenta}CLEANUP}\\
\texttt{\color{green}zpwr forgitclean          \# fzf-powered git clean}\\
\color{white}Shows untracked files. Select which to delete.\\
\color{white}Preview before you destroy. Civilized cleanup.\\
\texttt{\color{green}zpwr forgitrestore        \# fzf-powered git restore}\\
\color{white}Discard changes to selected files. With preview.\\
\vspace{4pt}\textbf{\color{magenta}GITIGNORE MANAGEMENT}\\
\texttt{\color{green}zpwr forgitignore         \# browse gitignore templates}\\
\texttt{\color{green}zpwr forgitignoreget      \# download a gitignore template}\\
\texttt{\color{green}zpwr forgitignorelist     \# list available templates}\\
\texttt{\color{green}zpwr forgitignoreupdate   \# update gitignore templates}\\
\texttt{\color{green}zpwr forgitignoreclean    \# clean gitignore cache}\\
\vspace{4pt}\textbf{\color{magenta}EXTRAS}\\
\texttt{\color{green}zpwr forgitinfo           \# show git repo information}\\
\texttt{\color{green}zpwr forgitwarn           \# show forgit warnings}\\
\color{white}Forgit turns git into a visual experience.\\
\color{white}Every forgit* verb is a regular git command with eyes.\\

\vspace{8pt}

\subsection{Page 14: Git Battalion Operations}

\vspace{4pt}\textbf{\color{cyan}=== GIT BATTALION OPERATIONS ===}\vspace{2pt}\\
\color{white}One repo is easy. Twenty repos is a war.\\
\color{white}These verbs let you run commands across all of them.\\
\vspace{4pt}\textbf{\color{magenta}THE GENERAL: FORALLDIR}\\
\texttt{\color{green}zpwr gitforalldir <cmd>   \# run cmd in ALL git dirs}\\
\color{white}Walks every git directory it can find and runs your\\
\color{white}command inside each one. The nuclear option.\\
\vspace{4pt}\textbf{\color{magenta}CACHED REPOS}\\
\texttt{\color{green}zpwr gitreposexec <cmd>    \# run cmd in cached git dirs}\\
\color{white}Uses the cached list of known repos. Faster than\\
\color{white}foralldir because it does not rescan the filesystem.\\
\vspace{4pt}\textbf{\color{magenta}CLEAN VS DIRTY TARGETING}\\
\texttt{\color{green}zpwr gitreposcleanexec <cmd>       \# refresh + run in clean repos}\\
\texttt{\color{green}zpwr gitreposcleancacheexec <cmd>  \# run in cached clean repos}\\
\texttt{\color{green}zpwr gitreposdirtyexec <cmd>       \# refresh + run in dirty repos}\\
\texttt{\color{green}zpwr gitreposdirtycacheexec <cmd>  \# run in cached dirty repos}\\
\color{white}Clean repos have no uncommitted changes.\\
\color{white}Dirty repos have pending work.\\
\color{white}Target your commands accordingly.\\
\vspace{4pt}\textbf{\color{magenta}MASS UPDATE}\\
\texttt{\color{green}zpwr gitupdatefordir      \# run git updates in all git dirs}\\
\color{white}Pulls latest, updates submodules, the full sync.\\
\color{white}Run this on a Monday morning. Coffee. Wait.\\
\vspace{4pt}\textbf{\color{magenta}BRANCH WIPE OPERATIONS}\\
\texttt{\color{green}zpwr gitfordirmain        \# wipe to main branch in git dirs}\\
\texttt{\color{green}zpwr gitfordirdevelop     \# wipe to develop branch in git dirs}\\
\texttt{\color{green}zpwr gitzfordir <branch> [z]       \# wipe to any branch + z jump}\\
\texttt{\color{green}zpwr gitzfordirmain [z]            \# z jump + wipe to main}\\
\texttt{\color{green}zpwr gitzfordirdevelop [z]         \# z jump + wipe to develop}\\
\vspace{4pt}\textbf{\color{magenta}THE REPOS FILE}\\
\texttt{\color{green}zpwr gitreposfile         \# edit the git repos file}\\
\color{white}This file lists all known repositories.\\
\color{white}Edit it to add or remove repos from batch operations.\\
\vspace{4pt}\textbf{\color{magenta}EXAMPLE: STATUS CHECK ALL REPOS}\\
\texttt{\color{green}zpwr gitreposexec git status -s}\\
\color{white}Runs git status -s in every cached repo.\\
\color{white}Instantly see which repos have uncommitted work.\\
\color{white}Battalion ops turn you from a developer\\
\color{white}into a fleet commander.\\

\vspace{8pt}


\section{CHAPTER 4: THE SEARCH ENGINE}

\subsection{Page 15: Finding Files}

\vspace{4pt}\textbf{\color{cyan}=== FINDING FILES ===}\vspace{2pt}\\
\color{white}The filesystem is vast. You are small.\\
\color{white}But you have fzf and a verb for every search strategy.\\
\vspace{4pt}\textbf{\color{magenta}FILESEARCH: LOCAL TREE}\\
\texttt{\color{green}zpwr filesearch          \# find and open a file in a sub dir}\\
\color{white}Recursively lists files under \textbackslash{}WD into fzf.\\
\color{white}Type fragments of the filename. Fuzzy match. Enter.\\
\color{white}Opens the file in your pager.\\
\texttt{\color{green}zpwr filesearchedit      \# same but puts path on command line}\\
\color{white}Instead of opening immediately, the selected path\\
\color{white}lands on your \textbackslash{}UFFER for further editing.\\
\color{white}Append commands, pipe it, whatever you want.\\
\vspace{4pt}\textbf{\color{magenta}DIRSEARCH: DIRECTORIES ONLY}\\
\texttt{\color{green}zpwr dirsearch           \# cd to a sub directory via fzf}\\
\color{white}Same idea but only directories. Pick one, land in it.\\
\color{white}Faster than cd-tab-tab-tab-tab.\\
\vspace{4pt}\textbf{\color{magenta}FINDSEARCH: FULL DISK}\\
\texttt{\color{green}zpwr findsearch          \# find across the entire drive}\\
\color{white}Uses find(1) to search the whole filesystem.\\
\color{white}Results feed into fzf. Pick one, open it.\\
\color{white}Slow on first run. Worth the wait.\\
\texttt{\color{green}zpwr findsearchedit      \# same but result goes to \textbackslash{}UFFER}\\
\vspace{4pt}\textbf{\color{magenta}LOCATESEARCH: INDEX-BASED}\\
\texttt{\color{green}zpwr locatesearch        \# search the locate database}\\
\color{white}Uses locate(1) which searches a pre-built index.\\
\color{white}Blazing fast. But only as fresh as your last updatedb.\\
\texttt{\color{green}zpwr locatesearchedit    \# same but result goes to \textbackslash{}UFFER}\\
\vspace{4pt}\textbf{\color{magenta}THE SEARCH HIERARCHY}\\
\color{white}Speed    Scope          Verb\\
\color{white}-----    -----          ----\\
\color{white}Fast ... current tree .. filesearch\\
\color{white}Fast ... dirs only ..... dirsearch\\
\color{white}Slow ... entire disk ... findsearch\\
\color{white}Instant  indexed disk .. locatesearch\\
\color{white}Start with filesearch for local work.\\
\color{white}Escalate to locatesearch when you need the whole disk.\\
\color{white}Fall back to findsearch when the index is stale.\\

\vspace{8pt}

\subsection{Page 16: Searching Content}

\vspace{4pt}\textbf{\color{cyan}=== SEARCHING CONTENT ===}\vspace{2pt}\\
\color{white}Finding files is one thing.\\
\color{white}Finding the line inside the file? That is the real game.\\
\vspace{4pt}\textbf{\color{magenta}WORDSEARCH: CONTENT IN LOCAL TREE}\\
\texttt{\color{green}zpwr wordsearch          \# find file by content in sub dir}\\
\color{white}Greps through files in the current tree.\\
\color{white}Results show matching lines with filenames.\\
\color{white}Fzf preview shows the file with match highlighted.\\
\color{white}Select one, jump straight to it.\\
\texttt{\color{green}zpwr wordsearchedit      \# same but result goes to \textbackslash{}UFFER}\\
\color{white}Puts the selected file path on your command line.\\
\color{white}Useful when you want to pipe or chain commands.\\
\vspace{4pt}\textbf{\color{magenta}GREP: AG INTO FZF}\\
\texttt{\color{green}zpwr grep               \# grep through pwd with ag into fzf}\\
\color{white}This is the heavy artillery. Uses the silver searcher\\
\color{white}(ag) piped into fzf for interactive filtering.\\
\color{white}Respects .gitignore. Skips binaries. Fast.\\
\vspace{4pt}\textbf{\color{magenta}FILESEARCH VS WORDSEARCH}\\
\color{white}Know the filename?     Use filesearch\\
\color{white}Know the content?      Use wordsearch\\
\color{white}Know nothing?          Use grep and improvise\\
\vspace{4pt}\textbf{\color{magenta}THE EDIT SUFFIX PATTERN}\\
\color{white}Every search verb has an edit variant:\\
\color{white}filesearch     -> filesearchedit\\
\color{white}wordsearch     -> wordsearchedit\\
\color{white}findsearch     -> findsearchedit\\
\color{white}locatesearch   -> locatesearchedit\\
\color{white}Without edit: opens the result immediately.\\
\color{white}With edit: puts the path on your command line.\\
\texttt{\color{green}This lets you compose: vim \textbackslash{}\$(zpwr filesearchedit)}\\
\vspace{4pt}\textbf{\color{magenta}REAL-WORLD EXAMPLES}\\
\color{white}\# Find where \textbackslash{}"ZPWR\_VERBS\textbackslash{}" is defined:\\
\texttt{\color{green}zpwr wordsearch}\\
\color{white}\# type ZPWR\_VERBS in fzf, browse the matches\\
\color{white}\# Quick grep for TODO comments:\\
\texttt{\color{green}zpwr grep}\\
\color{white}\# type TODO, see every file that has one\\
\color{white}The search verbs are your eyes inside the codebase.\\
\color{white}Use them before you reach for a GUI.\\

\vspace{8pt}

\subsection{Page 17: The FZF Dimension}

\vspace{4pt}\textbf{\color{cyan}=== THE FZF DIMENSION ===}\vspace{2pt}\\
\color{white}FZF is the nervous system of zpwr.\\
\color{white}It is wired into search, selection, navigation,\\
\color{white}git, vim, environment browsing, and verb dispatch.\\
\color{white}If zpwr has a UI, fzf IS that UI.\\
\vspace{4pt}\textbf{\color{magenta}VERB DISCOVERY}\\
\texttt{\color{green}zpwr verbs     \# fuzzy search all zpwr verbs}\\
\color{white}Every verb, every description, in one fzf window.\\
\color{white}This is the god-mode cheat sheet. Learn it first.\\
\vspace{4pt}\textbf{\color{magenta}ENVIRONMENT BROWSING}\\
\texttt{\color{green}zpwr envaccept          \# accept from fzf env}\\
\color{white}Browse all shell environment entries (aliases, functions,\\
\color{white}parameters, builtins) in fzf. Select one, accept it\\
\color{white}onto the command line.\\
\texttt{\color{green}zpwr envedit            \# edit from fzf env}\\
\color{white}Same browsing but puts the result on \textbackslash{}UFFER for editing.\\
\texttt{\color{green}zpwr envcachesearch     \# search cached aliases, functions, etc.}\\
\color{white}Searches the pre-built cache of your entire environment.\\
\color{white}Faster than envedit because it uses the cache file.\\
\vspace{4pt}\textbf{\color{magenta}FZF PATTERNS ACROSS ZPWR}\\
\color{white}FZF appears in dozens of verbs. The pattern is always:\\
\color{white}1. A data source generates lines (files, commits, dirs)\\
\color{white}2. Lines pipe into fzf with a preview command\\
\color{white}3. You fuzzy-filter, select, and take an action\\
\color{white}Verbs with fzf built in:\\
\color{white}filesearch     file list -> fzf -> open\\
\color{white}wordsearch     grep results -> fzf -> open\\
\color{white}dirsearch      dir list -> fzf -> cd\\
\color{white}gitcommits     git log -> fzf -> show diff\\
\color{white}forgitadd      unstaged files -> fzf -> stage\\
\color{white}forgitlog      commit log -> fzf -> preview\\
\color{white}zcd            z database -> fzf -> cd\\
\color{white}subcommandsfzf all verbs -> fzf -> execute\\
\vspace{4pt}\textbf{\color{magenta}THE FZF KEYBINDINGS}\\
\texttt{\color{green}zpwr vimsearch          \# search vim keybindings in fzf}\\
\texttt{\color{green}zpwr allsearch          \# search ALL keybindings in fzf}\\
\color{white}When you want to find a keybinding but do not remember\\
\color{white}where it is defined. Search them all.\\
\color{white}FZF is not optional in zpwr. It is the interface.\\
\color{white}Master fzf and you master zpwr.\\

\vspace{8pt}


\section{CHAPTER 5: EDITOR MASTERY}

\subsection{Page 18: Vim Integration Overview}

\vspace{4pt}\textbf{\color{cyan}=== VIM INTEGRATION OVERVIEW ===}\vspace{2pt}\\
\color{white}Vim is not just an editor in zpwr. It is a first-class\\
\color{white}citizen with its own verb family. Over 20 verbs,\\
\color{white}each opening vim in a specific context.\\
\vspace{4pt}\textbf{\color{magenta}THE PHILOSOPHY}\\
\color{white}Every vim verb answers one question:\\
\color{white}\textbackslash{}"What do I want to edit right now?\textbackslash{}"\\
\color{white}The answer is always one verb away.\\
\vspace{4pt}\textbf{\color{magenta}THE vim* VERB FAMILY}\\
\texttt{\color{green}vimall .............. open ALL zpwr files (for :argdo)}\\
\texttt{\color{green}vimalledit .......... pick and edit 1+ zpwr configs}\\
\color{white}vimautoload ......... open all autoloads (for :argdo)\\
\color{white}vimcd ............... vim edit, then cd to first dir\\
\texttt{\color{green}vimconfig ........... edit all zpwr config files}\\
\texttt{\color{green}vimctags ............ vim zpwr ctags}\\
\texttt{\color{green}vimemacsconfig ...... vim edit emacs zpwr configs}\\
\color{white}vimfilesearch ....... vim a file found via fzf\\
\color{white}vimfilesearchedit ... vim edit a file found via fzf\\
\color{white}vimfindsearch ....... vim a file from find(1) via fzf\\
\color{white}vimfindsearchedit ... vim edit find result\\
\texttt{\color{green}vimgtags ............ vim zpwr GNU global tags}\\
\texttt{\color{green}vimgtagsedit ........ vim edit zpwr GNU global tags}\\
\color{white}vimlocatesearch ..... vim a file from locate via fzf\\
\color{white}vimlocatesearchedit . vim edit locate result\\
\color{white}vimplugincount ...... count vim plugins\\
\color{white}vimpluginlist ....... list all vim plugins\\
\color{white}vimrecent ........... edit most recent files\\
\color{white}vimrecentcd ......... cd + edit most recent files\\
\color{white}vimrecentsudo ....... sudo edit most recent files\\
\color{white}vimrecentsudocd ..... cd + sudo edit recent files\\
\texttt{\color{green}vimscripts .......... open all zpwr scripts (for :argdo)}\\
\color{white}vimscriptedit ....... pick and edit 1+ scripts\\
\color{white}vimsearch ........... search vim keybindings via fzf\\
\texttt{\color{green}vimtests ............ edit all zpwr tests}\\
\color{white}vimtokens ........... vim the .tokens.sh file\\
\color{white}vimwordsearch ....... vim a file found by content\\
\color{white}vimwordsearchedit ... vim edit file found by content\\
\vspace{4pt}\textbf{\color{magenta}THE :argdo PATTERN}\\
\color{white}Verbs like vimall and vimscripts load many files\\
\color{white}into vim's arglist. Then you can:\\
\color{white}:argdo \%s/old/new/g | update\\
\color{white}Mass find-and-replace across the entire zpwr codebase.\\
\color{white}If it can be edited, there is a vim verb for it.\\

\vspace{8pt}

\subsection{Page 19: Vim File Operations}

\vspace{4pt}\textbf{\color{cyan}=== VIM FILE OPERATIONS ===}\vspace{2pt}\\
\color{white}Opening vim is easy. Opening vim with the right files\\
\color{white}already loaded? That is where zpwr earns its keep.\\
\vspace{4pt}\textbf{\color{magenta}THE BIG OPENS}\\
\texttt{\color{green}zpwr vimall             \# load ALL zpwr files into arglist}\\
\color{white}Every .zsh, .sh, .py, .pl file in the zpwr tree.\\
\color{white}Hundreds of files. Ready for :argdo mass operations.\\
\color{white}This is the tactical nuke of text editing.\\
\texttt{\color{green}zpwr vimalledit         \# fzf pick 1+ configs to edit}\\
\color{white}Like vimall but you choose which files to open.\\
\color{white}Surgical precision instead of carpet bombing.\\
\vspace{4pt}\textbf{\color{magenta}RECENT FILES}\\
\texttt{\color{green}zpwr vimrecent          \# edit most recently opened files}\\
\color{white}Pulls from vim's file history. Opens the files you\\
\color{white}were just working on. Resume where you left off.\\
\texttt{\color{green}zpwr vimrecentcd        \# cd to dir + open recent files}\\
\color{white}Same as vimrecent but also changes to the directory.\\
\color{white}Your shell and editor are now in the same place.\\
\texttt{\color{green}zpwr vimrecentsudo      \# sudo edit recent files}\\
\color{white}When you need root to edit those files.\\
\texttt{\color{green}zpwr vimrecentsudocd    \# cd + sudo edit recent files}\\
\color{white}The full combo: navigate, elevate, edit.\\
\vspace{4pt}\textbf{\color{magenta}ZPWR SUBSYSTEMS}\\
\texttt{\color{green}zpwr vimscripts         \# load all zpwr scripts into arglist}\\
\texttt{\color{green}zpwr vimscriptedit      \# pick and edit 1+ scripts}\\
\texttt{\color{green}zpwr vimautoload        \# load all autoloads into arglist}\\
\texttt{\color{green}zpwr vimtests           \# edit all zpwr test files}\\
\vspace{4pt}\textbf{\color{magenta}VIM CD}\\
\texttt{\color{green}zpwr vimcd              \# edit file, cd to its directory}\\
\color{white}Opens a file in vim, and when you exit, your shell\\
\color{white}is cd'd to the directory containing that file.\\
\color{white}Edit something, then be where it lives. Elegant.\\
\vspace{4pt}\textbf{\color{magenta}WORKFLOW EXAMPLE}\\
\color{white}\# Resume yesterday's work:\\
\texttt{\color{green}zpwr vimrecentcd}\\
\color{white}\# Opens recent files, lands you in the right dir\\
\color{white}\# Mass rename a function across all scripts:\\
\texttt{\color{green}zpwr vimscripts}\\
\color{white}:argdo \%s/oldFunc/newFunc/g | update\\
\color{white}\# Done. Every script updated.\\

\vspace{8pt}

\subsection{Page 20: Vim Search and Tags}

\vspace{4pt}\textbf{\color{cyan}=== VIM SEARCH AND TAGS ===}\vspace{2pt}\\
\color{white}Finding files is step one. Finding definitions\\
\color{white}inside those files is step two. Tags are step three:\\
\color{white}jump to any symbol by name.\\
\vspace{4pt}\textbf{\color{magenta}VIM FILE SEARCH}\\
\texttt{\color{green}zpwr vimfilesearch          \# fzf find file -> open in vim}\\
\color{white}Lists files in the current tree via fzf.\\
\color{white}Select one, it opens in vim. No path typing.\\
\texttt{\color{green}zpwr vimfilesearchedit      \# fzf find file -> path to \textbackslash{}UFFER}\\
\color{white}Same search but puts the path on your command line.\\
\vspace{4pt}\textbf{\color{magenta}VIM WORD SEARCH}\\
\texttt{\color{green}zpwr vimwordsearch          \# fzf grep content -> open in vim}\\
\color{white}Searches file contents, shows matches in fzf.\\
\color{white}Select a match, vim opens at that location.\\
\texttt{\color{green}zpwr vimwordsearchedit      \# same but path to \textbackslash{}UFFER}\\
\vspace{4pt}\textbf{\color{magenta}FULL DISK SEARCH INTO VIM}\\
\texttt{\color{green}zpwr vimfindsearch          \# find(1) -> fzf -> vim}\\
\texttt{\color{green}zpwr vimfindsearchedit      \# find(1) -> fzf -> \textbackslash{}UFFER}\\
\texttt{\color{green}zpwr vimlocatesearch        \# locate -> fzf -> vim}\\
\texttt{\color{green}zpwr vimlocatesearchedit    \# locate -> fzf -> \textbackslash{}UFFER}\\
\color{white}The search escalation ladder from page 15 but with\\
\color{white}vim as the final destination instead of a pager.\\
\vspace{4pt}\textbf{\color{magenta}CTAGS: JUMP TO DEFINITION}\\
\texttt{\color{green}zpwr vimctags               \# open zpwr ctags in vim}\\
\color{white}Ctags indexes function and variable definitions.\\
\color{white}With tags loaded, Ctrl-] on any symbol jumps to\\
\color{white}its definition. Ctrl-T jumps back.\\
\color{white}Code navigation at the speed of thought.\\
\vspace{4pt}\textbf{\color{magenta}GNU GLOBAL TAGS}\\
\texttt{\color{green}zpwr vimgtags               \# open zpwr GNU global tags in vim}\\
\texttt{\color{green}zpwr vimgtagsedit           \# edit zpwr GNU global tags}\\
\color{white}GNU Global is ctags on steroids. It indexes not just\\
\color{white}definitions but references too. Find every caller\\
\color{white}of a function, not just its definition.\\
\vspace{4pt}\textbf{\color{magenta}THE SEARCH MATRIX}\\
\color{white}What you know       Verb to use\\
\color{white}-------------       -----------\\
\color{white}Filename fragment   vimfilesearch\\
\color{white}String in a file    vimwordsearch\\
\color{white}File somewhere      vimlocatesearch\\
\color{white}Symbol name         vimctags then Ctrl-]\\
\color{white}All references      vimgtags\\

\vspace{8pt}

\subsection{Page 21: Vim Config and Sessions}

\vspace{4pt}\textbf{\color{cyan}=== VIM CONFIG AND SESSIONS ===}\vspace{2pt}\\
\color{white}Configuration is not a one-time event.\\
\color{white}It is a lifestyle. ZPWR gives you session commands\\
\color{white}that open the right config files as a group.\\
\vspace{4pt}\textbf{\color{magenta}THE SESSION VERBS}\\
\color{white}Each of these opens a curated set of config files\\
\color{white}in vim buffers, ready for synchronized editing:\\
\texttt{\color{green}zpwr vrc          \# vim/nvim config session}\\
\color{white}Opens your vimrc or init.vim along with related files.\\
\color{white}All your editor settings in one session.\\
\texttt{\color{green}zpwr brc          \# shell aliases session}\\
\color{white}Opens the tokens file and shell alias configs.\\
\color{white}Where you tune your personal shell experience.\\
\texttt{\color{green}zpwr zrc          \# zshrc session}\\
\color{white}Opens the zshrc and zsh init chain.\\
\color{white}The root of your shell's boot sequence.\\
\texttt{\color{green}zpwr trc          \# tmux.conf session}\\
\color{white}Opens tmux configuration files.\\
\color{white}Pane layouts, keybindings, status bar.\\
\vspace{4pt}\textbf{\color{magenta}CONFIG EDITING}\\
\texttt{\color{green}zpwr vimconfig    \# edit all zpwr configs in vim}\\
\color{white}Opens every zpwr configuration file. The master session.\\
\color{white}Use :bn and :bp to navigate between buffers.\\
\texttt{\color{green}zpwr vimtokens   \# vim the .tokens.sh file}\\
\texttt{\color{green}zpwr tokens      \# same thing}\\
\color{white}Quick jump to your personal aliases and secrets file.\\
\color{white}This is the file that is gitignored. Your sanctuary.\\
\vspace{4pt}\textbf{\color{magenta}PLUGIN INVENTORY}\\
\texttt{\color{green}zpwr vimplugincount  \# how many vim plugins installed}\\
\texttt{\color{green}zpwr vimpluginlist   \# list all vim plugins}\\
\color{white}Know thy plugins. Count them. Name them.\\
\color{white}If the count scares you, it should.\\
\vspace{4pt}\textbf{\color{magenta}TYPICAL CONFIG WORKFLOW}\\
\color{white}\# Add a new shell alias:\\
\texttt{\color{green}zpwr brc}\\
\color{white}\# Edit .tokens.sh, add your alias, :wq\\
\color{white}\# Source it or open a new terminal\\
\color{white}\# Tweak a vim keybinding:\\
\texttt{\color{green}zpwr vrc}\\
\color{white}\# Navigate to the mapping, change it, :wq\\
\color{white}Three letters to open any config.\\
\color{white}No file paths. No guessing. Just intent.\\

\vspace{8pt}

\subsection{Page 22: Emacs Mirror Universe}

\vspace{4pt}\textbf{\color{cyan}=== EMACS MIRROR UNIVERSE ===}\vspace{2pt}\\
\color{white}For every vim verb, there is an emacs verb.\\
\color{white}Identical semantics. Different editor. Perfect symmetry.\\
\color{white}ZPWR does not take sides in the editor war.\\
\color{white}It arms both factions equally.\\
\vspace{4pt}\textbf{\color{magenta}THE SYMMETRY TABLE}\\
\color{white}vim verb              emacs verb\\
\color{white}--------              ----------\\
\color{white}vimall .............. emacsall\\
\color{white}vimalledit .......... emacsalledit\\
\color{white}vimautoload ......... emacsautoload\\
\color{white}vimcd ............... emacscd\\
\color{white}vimconfig ........... emacsconfig\\
\color{white}vimctags ............ emacsctags\\
\color{white}vimfilesearch ....... emacsfilesearch\\
\color{white}vimfilesearchedit ... emacsfilesearchedit\\
\color{white}vimfindsearch ....... emacsfindsearch\\
\color{white}vimfindsearchedit ... emacsfindsearchedit\\
\color{white}vimgtags ............ emacsgtags\\
\color{white}vimgtagsedit ........ emacsgtagsedit\\
\color{white}vimlocatesearch ..... emacslocatesearch\\
\color{white}vimlocatesearchedit . emacslocatesearchedit\\
\color{white}vimrecent ........... emacsrecent\\
\color{white}vimrecentcd ......... emacsrecentcd\\
\color{white}vimrecentsudo ....... emacsrecentsudo\\
\color{white}vimrecentsudocd ..... emacsrecentsudocd\\
\color{white}vimscripts .......... emacsscripts\\
\color{white}vimscriptedit ....... emacsscriptedit\\
\color{white}vimtests ............ emacstests\\
\color{white}vimtokens ........... emacstokens\\
\color{white}vimwordsearch ....... emacswordsearch\\
\color{white}vimwordsearchedit ... emacswordsearchedit\\
\color{white}vimplugincount ...... emacsplugincount\\
\color{white}vimpluginlist ....... emacspluginlist\\
\vspace{4pt}\textbf{\color{magenta}EMACS-SPECIFIC}\\
\texttt{\color{green}zpwr emacsallserver       \# emacs all zpwr files (server mode)}\\
\texttt{\color{green}zpwr emacsemacsconfig     \# emacs edit emacs-specific configs}\\
\texttt{\color{green}zpwr emacszpwr            \# emacs the zpwr directory}\\
\color{white}The lesson: zpwr is editor-agnostic by design.\\
\color{white}Replace \textbackslash{}"vim\textbackslash{}" with \textbackslash{}"emacs\textbackslash{}" in any verb. It works.\\
\color{white}Both sides get the same firepower.\\

\vspace{8pt}


\section{CHAPTER 6: PERFORMANCE CULT}

\subsection{Page 23: Benchmarking Startup}

\vspace{4pt}\textbf{\color{cyan}=== BENCHMARKING STARTUP ===}\vspace{2pt}\\
\color{white}A shell that takes 2 seconds to start is a shell\\
\color{white}that has already lost. ZPWR takes startup time\\
\color{white}seriously enough to build a benchmarking verb.\\
\vspace{4pt}\textbf{\color{magenta}THE BENCH VERB}\\
\texttt{\color{green}zpwr bench}\\
\color{white}Runs multiple shell startup iterations and reports:\\
\color{white}- Mean startup time\\
\color{white}- Percentiles (p50, p95, p99)\\
\color{white}- Baseline diff (zpwr overhead vs bare zsh)\\
\vspace{4pt}\textbf{\color{magenta}WHAT THE NUMBERS MEAN}\\
\color{white}< 100ms .... imperceptible. You are golden.\\
\color{white}100-200ms .. noticeable if you open many terminals.\\
\color{white}200-500ms .. getting sluggish. Time to investigate.\\
\color{white}> 500ms .... something is wrong. Check your plugins.\\
\vspace{4pt}\textbf{\color{magenta}HOW IT WORKS}\\
\color{white}The bench verb spawns N fresh zsh instances,\\
\color{white}each running through the full init chain.\\
\color{white}It measures wall-clock time from fork to prompt.\\
\color{white}Then it does the same for bare zsh (no zpwr).\\
\color{white}The diff tells you exactly what zpwr costs.\\
\vspace{4pt}\textbf{\color{magenta}WHAT SLOWS YOU DOWN}\\
\color{white}Common culprits in order of likelihood:\\
\color{white}1. Too many vim/emacs plugins loading at shell init\\
\color{white}2. nvm/pyenv/rbenv hooking into every prompt\\
\texttt{\color{green}3. Uncompiled completion files (run zpwr compile)}\\
\color{white}4. Network calls in your init chain (stop that)\\
\color{white}5. Massive HISTFILE being loaded synchronously\\
\vspace{4pt}\textbf{\color{magenta}OPTIMIZATION CHECKLIST}\\
\texttt{\color{green}zpwr compile        \# recompile all cache and completions}\\
\texttt{\color{green}zpwr compilefpath    \# recompile all fpath files}\\
\texttt{\color{green}zpwr compilefiles    \# recompile all files}\\
\color{white}Compiled .zwc files load faster than source.\\
\color{white}This is the single biggest performance lever.\\
\vspace{4pt}\textbf{\color{magenta}BENCHMARK REGULARLY}\\
\texttt{\color{green}Run zpwr bench after any config change.}\\
\color{white}If the numbers go up, you broke something.\\
\color{white}If they go down, you are a performance wizard.\\
\color{white}Your shell starts thousands of times a year.\\
\color{white}Every millisecond is multiplied by every terminal open.\\
\color{white}Measure. Optimize. Measure again.\\

\vspace{8pt}

\subsection{Page 24: Flame Charts}

\vspace{4pt}\textbf{\color{cyan}=== FLAME CHARTS ===}\vspace{2pt}\\
\color{white}Bench tells you HOW LONG startup takes.\\
\color{white}Flame tells you WHERE the time goes.\\
\vspace{4pt}\textbf{\color{magenta}THE FLAME VERB}\\
\texttt{\color{green}zpwr flame}\\
\color{white}Generates a startup flame chart from zprof data.\\
\color{white}Shows you exactly which functions and files are\\
\color{white}consuming your startup time, in descending order.\\
\vspace{4pt}\textbf{\color{magenta}READING THE OUTPUT}\\
\color{white}The flame chart shows a ranked list of functions:\\
\color{white}Each entry shows:\\
\color{white}- Function/file name\\
\color{white}- Time consumed (milliseconds)\\
\color{white}- Percentage of total startup time\\
\color{white}The top of the list is where optimization effort\\
\color{white}pays off the most. If one function takes 40\% of\\
\color{white}startup, that is your target.\\
\vspace{4pt}\textbf{\color{magenta}HOW ZPROF WORKS}\\
\color{white}Zsh has a built-in profiler: zprof.\\
\color{white}When enabled, it instruments every function call\\
\color{white}during shell init and records the wall-clock time.\\
\color{white}The flame verb wraps this into a readable report.\\
\vspace{4pt}\textbf{\color{magenta}COMMON HOTSPOTS}\\
\color{white}What you will typically see at the top:\\
\color{white}compinit ........... completion system init\\
\color{white}compaudit .......... completion security check\\
\color{white}compdump ........... writing completion cache\\
\color{white}\_zpwr\_* ............ zpwr autoload functions\\
\color{white}nvm/pyenv/rbenv .... version manager hooks\\
\vspace{4pt}\textbf{\color{magenta}OPTIMIZATION TIPS}\\
\texttt{\color{green}1. If compinit is slow, run zpwr compile to cache it}\\
\color{white}2. If a version manager dominates, lazy-load it\\
\color{white}3. If an autoload is slow, check for disk I/O in it\\
\color{white}4. If compaudit is slow, fix the directory permissions\\
\texttt{\color{green}5. After changes, run zpwr bench to verify improvement}\\
\vspace{4pt}\textbf{\color{magenta}THE OPTIMIZATION LOOP}\\
\texttt{\color{green}zpwr flame      \# identify the bottleneck}\\
\color{white}... fix it ...\\
\texttt{\color{green}zpwr bench      \# measure the improvement}\\
\texttt{\color{green}zpwr flame      \# find the next bottleneck}\\
\color{white}Repeat until you hit diminishing returns.\\
\color{white}Flame charts turn gut feelings into data.\\
\color{white}Stop guessing. Start profiling.\\

\vspace{8pt}

\subsection{Page 25: Live Dashboard}

\vspace{4pt}\textbf{\color{cyan}=== LIVE DASHBOARD ===}\vspace{2pt}\\
\color{white}Bench is a snapshot. Flame is a profile.\\
\color{white}Top is a live feed of your shell's vital signs.\\
\vspace{4pt}\textbf{\color{magenta}THE TOP VERB}\\
\texttt{\color{green}zpwr top}\\
\color{white}Launches a live shell resource dashboard.\\
\color{white}Think htop, but for your zsh environment.\\
\vspace{4pt}\textbf{\color{magenta}WHAT EACH SECTION SHOWS}\\
\color{white}Shell Environment\\
\color{white}- Number of aliases currently defined\\
\color{white}- Number of functions loaded\\
\color{white}- Number of parameters (variables) in scope\\
\color{white}- Builtin and keyword counts\\
\color{white}Completion System\\
\color{white}- Completion functions registered\\
\color{white}- Cache file sizes and freshness\\
\color{white}Resource Usage\\
\color{white}- Memory footprint of the shell process\\
\color{white}- History file size and entry count\\
\color{white}- Autoload function count\\
\vspace{4pt}\textbf{\color{magenta}INTERPRETING THE DATA}\\
\color{white}High alias count?    Normal for zpwr. 1000+ is fine.\\
\color{white}High function count? Also normal. Autoloaded functions\\
\color{white}are lazy-loaded, they cost nothing\\
\color{white}until called.\\
\texttt{\color{green}Large HISTFILE?      Consider zpwr backuphistory and trim.}\\
\texttt{\color{green}Stale cache?         Run zpwr compile to refresh.}\\
\vspace{4pt}\textbf{\color{magenta}COMPANION VERBS}\\
\texttt{\color{green}zpwr environmentcounts   \# one-shot environment stats}\\
\texttt{\color{green}zpwr autoloadcount       \# total number of autoloads}\\
\texttt{\color{green}zpwr autoloadlist        \# list all autoloads}\\
\texttt{\color{green}zpwr vimplugincount      \# count vim plugins}\\
\texttt{\color{green}zpwr emacsplugincount    \# count emacs plugins}\\
\color{white}Use these for quick spot checks without the full dashboard.\\
\vspace{4pt}\textbf{\color{magenta}THE PERFORMANCE TRILOGY}\\
\texttt{\color{green}zpwr bench  -> how fast is startup?}\\
\texttt{\color{green}zpwr flame  -> where does the time go?}\\
\texttt{\color{green}zpwr top    -> what is loaded right now?}\\
\color{white}Three verbs. Three perspectives. Complete observability.\\
\color{white}Your shell is no longer a black box. It is an open book.\\

\vspace{8pt}


\section{CHAPTER 7: HEALTH \& DIAGNOSTICS}

\subsection{Page 26: The Doctor Is In}

\vspace{4pt}\textbf{\color{cyan}=== THE DOCTOR IS IN ===}\vspace{2pt}\\
\textit{\color{gray}// every system needs a physician. yours lives at zpwr doctor. //}\\
\color{white}Your shell is a living organism. It accumulates cruft, grows\\
\color{white}tumors in PATH, and occasionally develops broken symlinks.\\
\color{white}The doctor runs a full-body scan and tells you what hurts.\\
\vspace{4pt}\textbf{\color{magenta}INVOKE THE DOCTOR}\\
\texttt{\color{green}zpwr doctor              \# full diagnostic scan}\\
\texttt{\color{green}zpwr doctor -h           \# help with detailed check list}\\
\vspace{4pt}\textbf{\color{magenta}WHAT IT CHECKS}\\
\color{white}1. STALE .zwc FILES\\
\color{white}Compiled zsh files older than their source.\\
\color{white}Your shell is loading yesterday's code today.\\
\texttt{\color{green}Fix: zpwr recompile}\\
\color{white}2. ZCOMPDUMP FRESHNESS\\
\color{white}The completion dump file. If it is older than 1 week,\\
\color{white}tab completion may be missing new commands.\\
\texttt{\color{green}Fix: zpwr regenzsh}\\
\color{white}3. DUPLICATE PATH ENTRIES\\
\color{white}Every duplicate PATH entry is a wasted hash lookup.\\
\color{white}The doctor counts them and names the offenders.\\
\texttt{\color{green}Fix: zpwr deduppaths}\\
\color{white}4. DUPLICATE/MISSING FPATH ENTRIES\\
\color{white}Same disease, different organ. FPATH controls\\
\color{white}where autoloaded functions are found.\\
\color{white}5. BROKEN SYMLINKS\\
\color{white}Dangling symlinks in \textbackslash{}PWR. Every broken link is a\\
\color{white}source of cryptic errors at 3am.\\
\color{white}6. NAMED DIRECTORIES\\
\color{white}Invalid named dirs (\textasciitilde{}foo -> /nonexistent) and p10k\\
\color{white}pollution from auto\_name\_dirs leaking gitstatus vars.\\
\color{white}7. DEPENDENCIES\\
\color{white}Checks for: git, eza, fzf, tmux, perl, python3, zinit.\\
\color{white}Missing any? The doctor names names.\\
\color{white}8. HISTORY \& BACKUPS\\
\color{white}Is HISTFILE present? How many backup copies exist\\
\color{white}in \textbackslash{}PWR\_LOCAL/rcBackups? Zero backups = living dangerously.\\
\color{white}9. CONFIG SYMLINKS\\
\color{white}Checks .zshrc, .vimrc, .tmux.conf, .p10k.zsh.\\
\color{white}Are they symlinks into \textbackslash{}PWR? Or rogue files?\\
\vspace{4pt}\textbf{\color{magenta}READING THE RESULTS}\\
\color{white}green checkmark  = healthy, move along\\
\color{white}yellow warning   = degraded but functional\\
\color{white}red X            = broken, fix it now\\
\vspace{4pt}\textbf{\color{magenta}THE SIGNAL BAR}\\
\color{white}At the bottom: a 10-block signal meter.\\
\color{white}██████████  = HEALTHY. All green. You are a clean machine.\\
\color{white}██████░░░░  = WARNINGS. Stuff works but something is off.\\
\color{white}██░░░░░░░░  = CRITICAL. Your shell is on fire.\\
\vspace{4pt}\textbf{\color{magenta}PRO MOVE: CRON IT}\\
\texttt{\color{green}Add zpwr doctor to your weekly cron. Catch rot early.}\\
\texttt{\color{green}Pair it with zpwr recompile for a full tune-up.}\\

\vspace{8pt}

\subsection{Page 27: Dependency Cartography}

\vspace{4pt}\textbf{\color{cyan}=== DEPENDENCY CARTOGRAPHY ===}\vspace{2pt}\\
\textit{\color{gray}// mapping the call graph. who calls whom. who dies alone. //}\\
\color{white}zpwr has hundreds of functions. They call each other in a\\
\color{white}tangled web of dependencies. The deps verb maps that web.\\
\color{white}Think of it as Google Maps for your function call graph.\\
\vspace{4pt}\textbf{\color{magenta}BASIC USAGE}\\
\texttt{\color{green}zpwr deps                          \# summary overview}\\
\texttt{\color{green}zpwr deps zpwrClearList             \# what does zpwrClearList call?}\\
\texttt{\color{green}zpwr deps -h                        \# full help}\\
\vspace{4pt}\textbf{\color{magenta}THE FLAGS}\\
\color{white}-r, --reverse    Who calls this function? (reverse lookup)\\
\color{white}-o, --orphans    Functions called by nobody internally\\
\color{white}-s, --summary    Overview stats (default when no args)\\
\color{white}-d, --dot        Output dot format for graphviz rendering\\
\vspace{4pt}\textbf{\color{magenta}FORWARD LOOKUP: WHAT DOES X CALL?}\\
\texttt{\color{green}zpwr deps zpwrDoctor}\\
\color{white}Shows every function zpwrDoctor calls, plus depth-2 callees.\\
\color{white}Output:  zpwrDoctor -> zpwrPrettyPrint -> [zpwrLogDebug, ...]\\
\vspace{4pt}\textbf{\color{magenta}REVERSE LOOKUP: WHO CALLS X?}\\
\texttt{\color{green}zpwr deps -r zpwrPrettyPrint}\\
\color{white}Answers: \textbackslash{}"Who depends on zpwrPrettyPrint?\textbackslash{}"\\
\color{white}If the answer is 30 functions, that function is load-bearing.\\
\color{white}Touch it carefully.\\
\vspace{4pt}\textbf{\color{magenta}ORPHAN DETECTION}\\
\texttt{\color{green}zpwr deps --orphans}\\
\color{white}Lists functions that no other zpwr function calls.\\
\color{white}These are either entry points (user-invoked) or dead code.\\
\color{white}Orphans are not necessarily bad -- many are top-level verbs\\
\color{white}or zinit hooks. But if you wrote zpwrFoo and nobody calls it,\\
\color{white}maybe zpwrFoo was a fever dream.\\
\vspace{4pt}\textbf{\color{magenta}GRAPHVIZ EXPORT}\\
\texttt{\color{green}zpwr deps --dot > graph.dot}\\
\color{white}dot -Tsvg graph.dot -o graph.svg\\
\color{white}Produces a visual dependency graph. Nodes are functions,\\
\color{white}edges are call relationships. The dot output uses cyberpunk\\
\color{white}colors: dark background, cyan nodes, magenta edges.\\
\color{white}Open graph.svg in a browser and behold the nervous system.\\
\vspace{4pt}\textbf{\color{magenta}THE SUMMARY VIEW}\\
\color{white}With no args, deps shows two ranked lists:\\
\color{white}TOP 10 MOST DEPENDENCIES (functions that call the most others)\\
\color{white}TOP 10 MOST DEPENDED ON  (functions called by the most others)\\
\color{white}Plus stats: total functions, orphan count, total call edges.\\
\color{white}The bar charts use Unicode block characters. Very aesthetic.\\
\vspace{4pt}\textbf{\color{magenta}HOW IT WORKS}\\
\color{white}Scans autoload dirs: common, darwin, linux, comp\_utils.\\
\color{white}Reads each function body and string-matches against all\\
\color{white}other function names. Also scans \textbackslash{}PWR\_SCRIPTS/lib.sh.\\
\color{white}It is O(n\textasciicircum{}2) and proud of it. Your function count is not\\
\color{white}big enough to matter.\\

\vspace{8pt}

\subsection{Page 28: Command Autopsy}

\vspace{4pt}\textbf{\color{cyan}=== COMMAND AUTOPSY ===}\vspace{2pt}\\
\textit{\color{gray}// when you need to know why a command is slow, fat, or dead //}\\
\color{white}The trace verb is a command profiler. It wraps any command\\
\color{white}in instrumentation and reports wall time, CPU time, memory,\\
\color{white}context switches, page faults, and more.\\
\vspace{4pt}\textbf{\color{magenta}BASIC USAGE}\\
\texttt{\color{green}zpwr trace ls -la                   \# full trace of ls -la}\\
\texttt{\color{green}zpwr trace -q make                  \# quick: timing + resources only}\\
\texttt{\color{green}zpwr trace -h                       \# full help}\\
\vspace{4pt}\textbf{\color{magenta}THE FLAGS}\\
\color{white}-q, --quick      Timing + resource stats only (runs once)\\
\color{white}-s, --strace     Syscall counts only (Linux; no ltrace)\\
\color{white}-l, --ltrace     C library calls only (Linux; no strace)\\
\color{white}-a, --all        Both strace + ltrace (default behavior)\\
\vspace{4pt}\textbf{\color{magenta}WHAT YOU GET: THE RESOURCE CARD}\\
\color{white}exit code          green = 0, red = nonzero\\
\color{white}wall time          real clock time with rating:\\
\color{white}INSTANT (<100ms), FAST (<1s),\\
\color{white}NOMINAL (<5s), SLOW (<30s), HEAVY (30s+)\\
\color{white}user/sys time      CPU time in user/kernel space\\
\color{white}peak RSS           max resident memory (auto-scaled B/KB/MB/GB)\\
\color{white}ctx switches       voluntary + involuntary context switches\\
\color{white}page reclaims      soft faults (memory mapped in from cache)\\
\color{white}page faults        hard faults (actual disk I/O for memory)\\
\color{white}block I/O          filesystem reads and writes\\
\color{white}signals recv       number of signals delivered\\
\color{white}IPC messages       sent/received inter-process messages\\
\color{white}shell delta        memory change of the parent shell\\
\color{white}child procs        before/after child process count\\
\vspace{4pt}\textbf{\color{magenta}PLATFORM DIFFERENCES}\\
\color{white}macOS:  /usr/bin/time -l for resource stats.\\
\color{white}SIP blocks dtruss/dtrace on most binaries.\\
\color{white}You get timing + memory + kernel stats.\\
\color{white}Linux:  strace -c for syscall counts, ltrace -c for\\
\color{white}C library calls. Command re-runs for each tool.\\
\vspace{4pt}\textbf{\color{magenta}PRACTICAL EXAMPLES}\\
\texttt{\color{green}zpwr trace git status              \# is git status slow in this repo?}\\
\texttt{\color{green}zpwr trace -q rm -rf node\_modules  \# quick mode for destructive ops}\\
\texttt{\color{green}zpwr trace -s npm install          \# syscalls only (Linux)}\\
\texttt{\color{green}zpwr trace zsh -c 'source \textasciitilde{}/.zshrc' \# profile shell startup}\\
\vspace{4pt}\textbf{\color{magenta}THE SIGNAL BAR}\\
\color{white}Bottom of output: a 10-block performance signal.\\
\color{white}██████████  score 10 = under 100ms. Lightning.\\
\color{white}██████░░░░  score 6  = under 2s. Acceptable.\\
\color{white}██░░░░░░░░  score 2  = under 30s. Go get coffee.\\
\color{white}WARNING: Default mode re-runs the command for strace/ltrace.\\
\color{white}Use -q for destructive commands (rm, truncate, etc).\\

\vspace{8pt}

\subsection{Page 29: Usage Intelligence}

\vspace{4pt}\textbf{\color{cyan}=== USAGE INTELLIGENCE ===}\vspace{2pt}\\
\textit{\color{gray}// your shell keeps a diary. these verbs read it. //}\\
\color{white}You have 400+ verbs and thousands of aliases.\\
\color{white}Which ones do you actually use? funcrank and aliasrank\\
\color{white}mine your shell history to answer that question.\\
\vspace{4pt}\textbf{\color{magenta}FUNCRANK: FUNCTION \& VERB PROFILER}\\
\texttt{\color{green}zpwr funcrank                       \# top 20 functions and verbs}\\
\texttt{\color{green}zpwr funcrank 40                    \# top 40}\\
\texttt{\color{green}zpwr funcrank --reset               \# nuke cache, rebuild from scratch}\\
\texttt{\color{green}zpwr funcrank --reset 50            \# rebuild then show top 50}\\
\color{white}funcrank scans your HISTFILE for occurrences of known\\
\color{white}function names and zpwr verbs. It produces two ranked lists:\\
\color{white}1. Top functions by history occurrence count\\
\color{white}2. Top zpwr verbs by how often you type \textbackslash{}"zpwr <verb>\textbackslash{}"\\
\color{white}Each entry gets a Unicode bar chart. The biggest user wins.\\
\vspace{4pt}\textbf{\color{magenta}ALIASRANK: REVERSE EXPANSION PROFILER}\\
\texttt{\color{green}zpwr aliasrank                      \# top 20 by expanded usage}\\
\texttt{\color{green}zpwr aliasrank 50                   \# top 50}\\
\texttt{\color{green}zpwr aliasrank --reset              \# force cache rebuild}\\
\color{white}Here is the trick: zpwr uses zsh-expand to expand aliases\\
\color{white}before they hit history. So history contains the expanded\\
\color{white}form, not the alias name. aliasrank does a reverse lookup:\\
\color{white}it matches history lines against alias values to figure\\
\color{white}out which alias you would have typed.\\
\vspace{4pt}\textbf{\color{magenta}THE CACHE SYSTEM}\\
\color{white}\textbackslash{}PWR\_LOCAL/zpwr-aliasrank-cache.zsh   command counts\\
\color{white}\textbackslash{}PWR\_LOCAL/zpwr-verbrank-cache.zsh    verb counts\\
\color{white}\textbackslash{}PWR\_LOCAL/zpwr-rank-stamp            history line watermark\\
\color{white}The cache auto-rebuilds when history grows by 1000+ lines.\\
\color{white}The watermark file stores the last-seen history line count.\\
\color{white}No need to rebuild every time -- incremental is the way.\\
\color{white}Use --reset to force a full rebuild if cache seems stale.\\
\vspace{4pt}\textbf{\color{magenta}INTERPRETING THE RESULTS}\\
\color{white}If your top verb is help, you are still learning. Good.\\
\color{white}If your top verb is cleanall, you break things a lot.\\
\color{white}If your top function is cd, you need more named dirs.\\
\color{white}If an alias never appears, maybe it is too long to remember.\\
\vspace{4pt}\textbf{\color{magenta}PRACTICAL INTELLIGENCE}\\
\color{white}Run funcrank monthly. Notice patterns:\\
\color{white}- Verbs you never use might need better names or keybindings\\
\color{white}- Functions with high counts are candidates for optimization\\
\color{white}- Low-count zpwr verbs might benefit from shorter aliases\\
\color{white}Your history is a behavioral log. These verbs turn it into\\
\color{white}actionable intelligence. Know thyself, operative.\\

\vspace{8pt}


\section{CHAPTER 8: ENVIRONMENT MASTERY}

\subsection{Page 30: Know Your Shell}

\vspace{4pt}\textbf{\color{cyan}=== KNOW YOUR SHELL ===}\vspace{2pt}\\
\textit{\color{gray}// you cannot master what you cannot measure //}\\
\color{white}Three verbs to audit your environment. Each one reveals a\\
\color{white}different cross-section of the machine you live inside.\\
\vspace{4pt}\textbf{\color{magenta}ENVVARS: THE ZPWR VARIABLES}\\
\texttt{\color{green}zpwr envvars                        \# list all ZPWR\_* env vars}\\
\texttt{\color{green}zpwr ev                              \# same thing, shorter}\\
\color{white}Shows every environment variable starting with ZPWR\_.\\
\color{white}These are the knobs and dials of the framework:\\
\color{white}paths, flags, plugin manager config, colorscheme, etc.\\
\color{white}If something is misconfigured, this is where you look.\\
\vspace{4pt}\textbf{\color{magenta}ENVVARSALL: THE FULL PICTURE}\\
\texttt{\color{green}zpwr envvarsall                     \# every env var, all of them}\\
\texttt{\color{green}zpwr eva                             \# same thing, shorter}\\
\color{white}Dumps every exported variable. This includes ZPWR\_*,\\
\color{white}PATH, HOME, EDITOR, LANG, and everything else.\\
\color{white}Pipe it to fzf when you are hunting for a specific var:\\
\texttt{\color{green}zpwr eva | fzf}\\
\vspace{4pt}\textbf{\color{magenta}ENVIRONMENTCOUNTS: THE NUMBERS GAME}\\
\texttt{\color{green}zpwr environmentcounts              \# full environment statistics}\\
\texttt{\color{green}zpwr envcounts                      \# alias}\\
\texttt{\color{green}zpwr e                               \# shortest form}\\
\color{white}This is the census. It counts everything loaded in your\\
\color{white}current shell session:\\
\color{white}aliases         regular + global + suffix aliases\\
\color{white}functions       autoloaded + defined functions\\
\color{white}builtins        zsh builtin commands\\
\color{white}parameters      shell variables (scalar, array, assoc)\\
\color{white}completions     registered completion functions\\
\color{white}modules         loaded zsh modules\\
\color{white}named dirs      hash -d entries\\
\color{white}PATH entries    directories in PATH\\
\color{white}FPATH entries   directories in FPATH\\
\color{white}The numbers tell you how armed you are.\\
\color{white}A fresh zpwr install might show 3000+ aliases,\\
\color{white}500+ functions, and 100+ completions.\\
\color{white}If the numbers are low, something failed to load.\\
\vspace{4pt}\textbf{\color{magenta}PRO TIP: PAIR THEM}\\
\texttt{\color{green}zpwr doctor                   \# is anything broken?}\\
\texttt{\color{green}zpwr e                         \# how much loaded?}\\
\texttt{\color{green}zpwr ev                        \# are my zpwr vars right?}\\
\color{white}Run all three after installing or updating zpwr.\\
\color{white}They are your post-deployment smoke test.\\

\vspace{8pt}

\subsection{Page 31: The Option Matrix}

\vspace{4pt}\textbf{\color{cyan}=== THE OPTION MATRIX ===}\vspace{2pt}\\
\textit{\color{gray}// zsh has 200+ options. most people know 5. time to fix that. //}\\
\color{white}The zsh option system is a maze of boolean switches that\\
\color{white}control everything from globbing to history to error handling.\\
\color{white}These verbs help you navigate the labyrinth.\\
\vspace{4pt}\textbf{\color{magenta}OPTIONS: THE BOOLEAN SWITCHES}\\
\texttt{\color{green}zpwr options                        \# list all set options}\\
\color{white}Shows every zsh option that is currently ON.\\
\color{white}The important ones for zpwr:\\
\color{white}extendedglob      (\#, \textasciitilde{}, \textasciicircum{} in globs)  -- zpwr needs this\\
\color{white}autopushd         (cd pushes to dirstack)\\
\color{white}histignorealldups (no duplicate history)\\
\color{white}sharehistory      (history shared across sessions)\\
\color{white}completeinword    (tab complete mid-word)\\
\color{white}promptsubst       (parameter expansion in prompts)\\
\color{white}If extendedglob is off, half of zpwr's glob tricks die.\\
\color{white}If sharehistory is off, your sessions are islands.\\
\vspace{4pt}\textbf{\color{magenta}MODULES: THE LOADED EXTENSIONS}\\
\texttt{\color{green}zpwr modules                        \# list loaded zsh modules}\\
\color{white}Zsh modules are C libraries that add features:\\
\color{white}zsh/datetime      EPOCHREALTIME, strftime\\
\color{white}zsh/stat          zstat builtin for file metadata\\
\color{white}zsh/zutil         zparseopts for flag parsing\\
\color{white}zsh/complist       menuselect completion listing\\
\color{white}zsh/parameter      \textbackslash{}\$functions, \textbackslash{}\$aliases, \textbackslash{}\$parameters hashes\\
\color{white}zsh/terminfo       terminal capability strings\\
\color{white}If zsh/parameter is not loaded, zpwr env verbs break.\\
\color{white}If zsh/datetime is missing, timing-based verbs fail.\\
\color{white}Load a missing module: zmodload zsh/modulename\\
\vspace{4pt}\textbf{\color{magenta}ZSTYLE: THE FINE-TUNING KNOBS}\\
\texttt{\color{green}zpwr zstyle                         \# list all zstyle settings}\\
\color{white}zstyle is zsh's configuration registry. It controls:\\
\color{white}:completion:*     how tab completion behaves\\
\color{white}:vcs\_info:*       git/svn info in prompts\\
\color{white}:zpwr:*           zpwr-specific settings\\
\color{white}Example styles zpwr sets:\\
\color{white}zstyle ':completion:*' menu select\\
\color{white}zstyle ':completion:*' matcher-list 'm:\{a-z\}=\{A-Z\}'\\
\color{white}The first enables arrow-key menu selection in tab complete.\\
\color{white}The second makes completion case-insensitive.\\
\vspace{4pt}\textbf{\color{magenta}DETECTIVE WORK}\\
\color{white}Completions acting weird? Check:\\
\texttt{\color{green}zpwr options    \# is completeinword on?}\\
\texttt{\color{green}zpwr modules    \# is zsh/complist loaded?}\\
\texttt{\color{green}zpwr zstyle     \# what completion styles are set?}\\
\color{white}The answer is usually hiding in one of these three.\\

\vspace{8pt}

\subsection{Page 32: Cache Search}

\vspace{4pt}\textbf{\color{cyan}=== CACHE SEARCH ===}\vspace{2pt}\\
\textit{\color{gray}// the environment is a haystack. these are your magnets. //}\\
\color{white}zpwr maintains an environment cache for fast lookup.\\
\color{white}These verbs let you search, inspect, and edit that cache\\
\color{white}without grepping through raw variable dumps.\\
\vspace{4pt}\textbf{\color{magenta}ENVIRONMENTCACHESEARCH: FZF ENV NAVIGATION}\\
\texttt{\color{green}zpwr environmentcachesearch         \# fzf through cached env}\\
\color{white}Opens an fzf interface over the environment cache.\\
\color{white}Type a few characters to filter. The cache includes:\\
\color{white}- All exported environment variables\\
\color{white}- ZPWR\_* configuration variables\\
\color{white}- PATH and FPATH entries\\
\color{white}- Named directory mappings\\
\color{white}This is faster than env | grep because it uses the\\
\color{white}pre-built cache file rather than querying the live env.\\
\vspace{4pt}\textbf{\color{magenta}ENVACCEPT: ACCEPT A CACHE ENTRY}\\
\texttt{\color{green}zpwr envaccept                      \# accept/apply cached env entry}\\
\color{white}After finding a variable in the cache, this verb\\
\color{white}applies it to your current session. Useful when the\\
\color{white}cache has a value you want but the live env does not.\\
\vspace{4pt}\textbf{\color{magenta}ENVEDIT: EDIT THE ENVIRONMENT}\\
\texttt{\color{green}zpwr envedit                        \# open env cache in \textbackslash{}DITOR}\\
\color{white}Opens the environment cache file in your editor.\\
\color{white}Make changes, save, and the cache is updated.\\
\color{white}This is the surgical approach: when you know\\
\color{white}exactly which variable needs tweaking.\\
\vspace{4pt}\textbf{\color{magenta}THE CACHE LIFECYCLE}\\
\color{white}1. zpwr regenenvcache   builds the cache from live env\\
\color{white}2. environmentcachesearch fuzzy-searches the cache\\
\color{white}3. envaccept             applies an entry\\
\color{white}4. envedit               manual edits\\
\color{white}The cache lives at \textbackslash{}PWR\_LOCAL. It is regenerated during\\
\texttt{\color{green}shell startup and when you run zpwr regenenvcache.}\\
\vspace{4pt}\textbf{\color{magenta}WORKFLOW EXAMPLE}\\
\color{white}\textbackslash{}"What is ZPWR\_EDITOR set to?\textbackslash{}"\\
\texttt{\color{green}zpwr environmentcachesearch    \# type 'EDITOR' in fzf}\\
\color{white}Found it: ZPWR\_EDITOR=nvim\\
\color{white}\textbackslash{}"I want to change a bunch of ZPWR\_ vars.\textbackslash{}"\\
\texttt{\color{green}zpwr envedit                       \# opens in \textbackslash{}DITOR}\\
\color{white}Edit the vars, save, done.\\
\color{white}\textbackslash{}"Something is off, rebuild the cache.\textbackslash{}"\\
\texttt{\color{green}zpwr regenenvcache                 \# fresh cache from live env}\\
\vspace{4pt}\textbf{\color{magenta}PRO TIP}\\
\color{white}If you are debugging a missing env var, check two places:\\
\texttt{\color{green}zpwr environmentcachesearch (cached) and zpwr eva (live).}\\
\texttt{\color{green}If cache has it but live does not, run zpwr regenenvcache.}\\

\vspace{8pt}

\subsection{Page 33: Snapshots \& Time Travel}

\vspace{4pt}\textbf{\color{cyan}=== SNAPSHOTS \& TIME TRAVEL ===}\vspace{2pt}\\
\textit{\color{gray}// save your universe. restore it when everything goes wrong. //}\\
\color{white}The snapshot verb captures your entire terminal environment\\
\color{white}to a portable directory. The restore verb brings it back.\\
\color{white}This is git for your shell state.\\
\vspace{4pt}\textbf{\color{magenta}TAKING A SNAPSHOT}\\
\texttt{\color{green}zpwr snapshot                       \# timestamp-named snapshot}\\
\texttt{\color{green}zpwr snapshot myproject              \# named snapshot}\\
\texttt{\color{green}zpwr snapshot before-refactor        \# meaningful name}\\
\vspace{4pt}\textbf{\color{magenta}WHAT GETS CAPTURED}\\
\color{white}hash dirs         named directory mappings (hash -d)\\
\color{white}aliases           regular + global + suffix aliases\\
\color{white}env vars          all exported variables (secrets excluded)\\
\color{white}zpwr vars         ZPWR\_VARS associative array\\
\color{white}functions         user function count tracking\\
\color{white}history           full shell history dump\\
\color{white}tmux             sessions, windows, panes, resurrect state\\
\color{white}vim sessions      saved vim session files\\
\color{white}dir stack         current directory and pushd stack\\
\color{white}git status        branch, commit, dirty state of cwd\\
\color{white}Storage: \textbackslash{}PWR\_LOCAL/snapshots/<name>/\\
\color{white}Each snapshot is a directory of plain zsh files.\\
\color{white}Human-readable. Sourceable. Portable.\\
\vspace{4pt}\textbf{\color{magenta}RESTORING A SNAPSHOT}\\
\texttt{\color{green}zpwr restore                        \# list available snapshots}\\
\texttt{\color{green}zpwr restore myproject               \# restore everything}\\
\texttt{\color{green}zpwr restore myproject aliases env    \# restore only aliases+env}\\
\texttt{\color{green}zpwr restore myproject dirs           \# just cd to saved pwd}\\
\vspace{4pt}\textbf{\color{magenta}SELECTIVE RESTORE COMPONENTS}\\
\color{white}all       everything (default when no component given)\\
\color{white}hash      named directories only\\
\color{white}aliases   regular, global, suffix aliases\\
\color{white}env       environment variables\\
\color{white}tmux      stage tmux-resurrect file (press prefix+Ctrl-r)\\
\color{white}vim       vim session files to \textasciitilde{}/.vim/sessions\\
\color{white}history   merge snapshot history into current session\\
\color{white}dirs      cd to the saved working directory\\
\vspace{4pt}\textbf{\color{magenta}WHEN TO USE SNAPSHOTS}\\
\color{white}Before a big update:     zpwr snapshot before-update\\
\color{white}Before experiments:      zpwr snapshot stable-state\\
\color{white}Switching projects:      zpwr snapshot projectA\\
\color{white}New machine setup:       copy snapshot dir, restore on new host\\
\color{white}NOTE: Sensitive vars (PASSWORD, SECRET, TOKEN, AWS\_SECRET*,\\
\color{white}GITHUB\_TOKEN) are automatically excluded from env snapshots.\\
\color{white}Your secrets stay secret.\\
\color{white}History merge is additive -- it does not overwrite your\\
\color{white}current history. It uses fc -R to load entries.\\

\vspace{8pt}


\section{CHAPTER 9: THE JANITOR'S CLOSET}

\subsection{Page 34: Cleanup Operations}

\vspace{4pt}\textbf{\color{cyan}=== CLEANUP OPERATIONS ===}\vspace{2pt}\\
\textit{\color{gray}// entropy is the enemy. the janitor is your friend. //}\\
\color{white}Shell environments accumulate garbage like a hoarder's attic.\\
\color{white}Cache files grow stale. Temp dirs fill up. Logs balloon.\\
\color{white}These verbs are the Marie Kondo of your \textbackslash{}PWR.\\
\vspace{4pt}\textbf{\color{magenta}CLEANALL: THE THERMONUCLEAR BROOM}\\
\texttt{\color{green}zpwr cleanall                       \# clean everything at once}\\
\color{white}Calls all the individual clean verbs in sequence.\\
\color{white}This is the \textbackslash{}"I do not know what is wrong but I want\\
\color{white}a fresh start\textbackslash{}" button. Safe but thorough.\\
\vspace{4pt}\textbf{\color{magenta}INDIVIDUAL CLEAN VERBS}\\
\color{white}CLEANCACHE\\
\texttt{\color{green}zpwr cleancache                     \# clear general caches}\\
\color{white}Wipes cached data that zpwr generates during operation.\\
\color{white}Safe to run any time. Everything rebuilds on demand.\\
\color{white}CLEANCOMPCACHE\\
\texttt{\color{green}zpwr cleancompcache                 \# clear completion cache}\\
\color{white}Removes the zsh completion cache directory.\\
\color{white}Run this when tab completions are showing stale results\\
\color{white}or when a new command does not appear in completions.\\
\texttt{\color{green}Follow up with zpwr regenzsh to rebuild zcompdump.}\\
\color{white}CLEANDIRS\\
\texttt{\color{green}zpwr cleandirs                      \# clear cached dir listings}\\
\color{white}Removes directory cache files used by navigation commands.\\
\color{white}If directory completions point to dirs that no longer exist,\\
\color{white}this verb fixes that.\\
\color{white}CLEANTEMP\\
\texttt{\color{green}zpwr cleantemp                      \# clear temp files}\\
\color{white}Removes temporary files zpwr created during operation.\\
\color{white}These are transient by nature -- safe to nuke.\\
\color{white}CLEANDIRSANDTEMP\\
\texttt{\color{green}zpwr cleandirsandtemp               \# dirs + temp combined}\\
\color{white}Convenience combo of cleandirs and cleantemp.\\
\color{white}CLEANLOG\\
\texttt{\color{green}zpwr cleanlog                       \# truncate log files}\\
\color{white}The zpwr logfile grows forever if you never clean it.\\
\color{white}This verb truncates or removes old log entries.\\
\color{white}Check log size: ls -lh \textbackslash{}PWR\_LOGFILE\\
\vspace{4pt}\textbf{\color{magenta}THE CLEAN WORKFLOW}\\
\color{white}Light cleanup (weekly):\\
\texttt{\color{green}zpwr cleantemp \&\& zpwr cleanlog}\\
\color{white}Medium cleanup (monthly):\\
\texttt{\color{green}zpwr cleancache \&\& zpwr cleancompcache \&\& zpwr regenzsh}\\
\color{white}Full cleanup (when things feel sluggish):\\
\texttt{\color{green}zpwr cleanall}\\
\color{white}None of these delete your configuration or history.\\
\color{white}They only remove regenerable cache and temp data.\\

\vspace{8pt}

\subsection{Page 35: Git Cache Hygiene}

\vspace{4pt}\textbf{\color{cyan}=== GIT CACHE HYGIENE ===}\vspace{2pt}\\
\textit{\color{gray}// zpwr tracks your git repos. the cache needs a bath sometimes. //}\\
\color{white}zpwr maintains caches of git repository locations so that\\
\color{white}repo-aware verbs (navigation, status, batch operations) do not\\
\color{white}have to scan the filesystem every time. Three verbs manage\\
\color{white}these caches.\\
\vspace{4pt}\textbf{\color{magenta}CLEANGITCACHE: THE GENERAL REPO CACHE}\\
\texttt{\color{green}zpwr cleangitcache                  \# clear all git repo caches}\\
\color{white}Wipes the master list of known git repositories.\\
\color{white}This cache is built by scanning for .git directories\\
\texttt{\color{green}and is used by verbs like zpwr fordir and repo}\\
\color{white}navigation commands.\\
\color{white}When to use it:\\
\color{white}- You moved or deleted git repos\\
\color{white}- Navigation is suggesting repos that no longer exist\\
\color{white}- You cloned new repos and they are not showing up\\
\texttt{\color{green}After clearing, run zpwr regengitrepocache to rebuild.}\\
\vspace{4pt}\textbf{\color{magenta}CLEANGITCLEANCACHE: CLEAN REPOS ONLY}\\
\texttt{\color{green}zpwr cleangitcleancache             \# clear clean repo cache}\\
\color{white}zpwr separately tracks which repos have a clean working\\
\color{white}tree (no uncommitted changes). This cache speeds up\\
\color{white}operations that only care about clean repos.\\
\color{white}Clear this when:\\
\color{white}- The \textbackslash{}"clean repos\textbackslash{}" list seems wrong\\
\color{white}- You committed in a repo but it still shows as dirty\\
\vspace{4pt}\textbf{\color{magenta}CLEANGITDIRTYCACHE: DIRTY REPOS ONLY}\\
\texttt{\color{green}zpwr cleangitdirtycache             \# clear dirty repo cache}\\
\color{white}The mirror of the clean cache. Tracks repos with\\
\color{white}uncommitted changes. Clear this when:\\
\color{white}- A repo shows as dirty but you already committed\\
\color{white}- The dirty count in the banner seems wrong\\
\texttt{\color{green}Rebuild with: zpwr regengitdirtyrepocache}\\
\vspace{4pt}\textbf{\color{magenta}THE GIT CACHE ARCHITECTURE}\\
\color{white}regengitrepocache        scans for all git repos\\
\color{white}regengitdirtyrepocache   classifies clean vs dirty\\
\color{white}cleangitcache            wipes the master list\\
\color{white}cleangitcleancache       wipes the clean-only list\\
\color{white}cleangitdirtycache       wipes the dirty-only list\\
\color{white}The caches live in \textbackslash{}PWR\_LOCAL as plain text files.\\
\color{white}One repo path per line. Grep-friendly. Human-readable.\\
\vspace{4pt}\textbf{\color{magenta}THE FULL REFRESH}\\
\texttt{\color{green}zpwr cleangitcache}\\
\texttt{\color{green}zpwr regengitrepocache}\\
\texttt{\color{green}zpwr regengitdirtyrepocache}\\
\color{white}Three commands. Fresh git cache. Takes a few seconds\\
\color{white}depending on how many repos you have scattered around.\\

\vspace{8pt}

\subsection{Page 36: Path Deduplication}

\vspace{4pt}\textbf{\color{cyan}=== PATH DEDUPLICATION ===}\vspace{2pt}\\
\textit{\color{gray}// every duplicate entry is a tiny crime against performance //}\\
\color{white}PATH and FPATH grow duplicates like weeds. Every subshell,\\
\color{white}every sourced script, every tool that helpfully \textbackslash{}"adds itself\\
\color{white}to your PATH\textbackslash{}" can insert duplicates. The deduppaths verb\\
\color{white}removes them.\\
\vspace{4pt}\textbf{\color{magenta}THE VERB}\\
\texttt{\color{green}zpwr deduppaths                     \# deduplicate PATH and FPATH}\\
\color{white}Removes duplicate entries from both PATH and FPATH while\\
\color{white}preserving the order of first occurrence. If /usr/bin\\
\color{white}appears 3 times, only the first one survives.\\
\vspace{4pt}\textbf{\color{magenta}WHY DUPLICATES HAPPEN}\\
\color{white}1. Shell inheritance\\
\color{white}Parent shell sets PATH. Child shell re-sources .zshrc.\\
\color{white}PATH grows. Nested tmux sessions make it worse.\\
\color{white}2. Multiple init scripts\\
\color{white}Homebrew, nvm, pyenv, rbenv, cargo -- each one adds\\
\color{white}to PATH in their init hook. Source them twice, PATH grows.\\
\color{white}3. /etc/zprofile on macOS\\
\color{white}macOS ships with /etc/zprofile that runs path\_helper.\\
\color{white}path\_helper reads /etc/paths and /etc/paths.d/ and\\
\color{white}prepends entries. This runs on every login shell,\\
\color{white}BEFORE your .zshrc. Duplicates guaranteed.\\
\color{white}4. Manual PATH manipulation\\
\color{white}export PATH=\textbackslash{}"/my/dir:\textbackslash{}ATH\textbackslash{}" in .zshrc.\\
\color{white}Every time .zshrc is sourced, another copy appears.\\
\vspace{4pt}\textbf{\color{magenta}HOW TO PREVENT DUPLICATES}\\
\color{white}Zsh has a built-in mechanism: the typeset -U flag.\\
\color{white}typeset -U path fpath\\
\color{white}This marks the arrays as \textbackslash{}"unique\textbackslash{}" -- zsh automatically\\
\color{white}removes duplicates when entries are added. zpwr sets\\
\color{white}this early in startup. But if something adds to PATH\\
\color{white}before zpwr loads, duplicates can sneak in.\\
\vspace{4pt}\textbf{\color{magenta}DIAGNOSING THE PROBLEM}\\
\texttt{\color{green}zpwr doctor                         \# checks for duplicate PATH/FPATH}\\
\color{white}print -l \textbackslash{}\$path | sort | uniq -d     \# show only duplicated entries\\
\color{white}echo \textbackslash{}\$\#path                         \# count PATH entries\\
\color{white}A healthy PATH has 10-30 entries. If you see 50+,\\
\color{white}something is adding entries in a loop.\\
\vspace{4pt}\textbf{\color{magenta}THE FIX}\\
\texttt{\color{green}zpwr deduppaths                     \# instant dedup}\\
\color{white}Run it. Then check: echo \textbackslash{}\$\#path\\
\color{white}If the count dropped, you had duplicates.\\
\color{white}If it drops a lot, investigate the source.\\
\color{white}For a permanent fix, find the script that adds duplicates\\
\color{white}and guard it with a check:\\
\color{white}[[ -z \textbackslash{}\$\{path[(r)/my/dir]\} ]] \&\& path=(/my/dir \textbackslash{}\$path)\\

\vspace{8pt}

\subsection{Page 37: The Nuclear Options}

\vspace{4pt}\textbf{\color{cyan}=== THE NUCLEAR OPTIONS ===}\vspace{2pt}\\
\textit{\color{gray}// when gentle cleanup is not enough. when you need to nuke from orbit. //}\\
\color{white}Most clean verbs are surgical. These two are orbital strikes.\\
\color{white}They clean, regenerate, rebuild, and update in one shot.\\
\color{white}Use them when you want a fresh start without reinstalling.\\
\vspace{4pt}\textbf{\color{magenta}CLEANREFRESHUPDATE: THE BIG RED BUTTON}\\
\texttt{\color{green}zpwr cleanrefreshupdate             \# clean + regen + update}\\
\color{white}This is the full-cycle maintenance verb. It:\\
\color{white}1. Cleans all caches (equivalent to cleanall)\\
\color{white}2. Regenerates compiled files, completions, env cache\\
\color{white}3. Runs the update cycle for zpwr and plugins\\
\color{white}Think of it as \textbackslash{}"factory reset without losing data.\textbackslash{}"\\
\color{white}Your configuration, history, and local files survive.\\
\color{white}Everything else gets rebuilt from source.\\
\color{white}When to use it:\\
\color{white}- Shell startup feels slow and you have tried individual fixes\\
\color{white}- After a major zsh or macOS upgrade\\
\color{white}- When individual clean verbs did not solve the problem\\
\color{white}- When you want to feel powerful\\
\color{white}TIME COST: Several minutes. This is not instant.\\
\color{white}It downloads plugin updates and recompiles everything.\\
\vspace{4pt}\textbf{\color{magenta}CLEANUPDATEDEPS: THE DEPENDENCY PURGE}\\
\texttt{\color{green}zpwr cleanupdatedeps                \# clean deps + update them}\\
\color{white}Focused on external dependencies rather than caches.\\
\color{white}This verb:\\
\color{white}1. Cleans stale dependency state\\
\color{white}2. Updates all managed dependencies\\
\color{white}When to use it:\\
\color{white}- Plugin versions are stuck at old commits\\
\color{white}- A dependency update broke something\\
\color{white}- You want the latest of everything\\
\vspace{4pt}\textbf{\color{magenta}DECISION MATRIX}\\
\color{white}Symptom                          Verb\\
\texttt{\color{green}Stale completions                zpwr cleancompcache}\\
\texttt{\color{green}Slow startup                     zpwr recompile}\\
\texttt{\color{green}Missing new repos in nav         zpwr cleangitcache}\\
\texttt{\color{green}General weirdness                zpwr cleanall}\\
\texttt{\color{green}Everything is broken              zpwr cleanrefreshupdate}\\
\texttt{\color{green}Plugins are ancient               zpwr cleanupdatedeps}\\
\vspace{4pt}\textbf{\color{magenta}THE ESCALATION PATH}\\
\color{white}Always start small and escalate:\\
\texttt{\color{green}zpwr cleantemp                      \# level 1: temp files}\\
\texttt{\color{green}zpwr cleancache                     \# level 2: all caches}\\
\texttt{\color{green}zpwr cleanall                       \# level 3: everything}\\
\texttt{\color{green}zpwr cleanrefreshupdate             \# level 4: clean + rebuild + update}\\
\color{white}Do not start with level 4. That is like calling an airstrike\\
\color{white}because you saw a spider.\\

\vspace{8pt}


\section{CHAPTER 10: BUILD SYSTEM}

\subsection{Page 38: Compilation 101}

\vspace{4pt}\textbf{\color{cyan}=== COMPILATION 101 ===}\vspace{2pt}\\
\textit{\color{gray}// .zwc files: compiled zsh. faster loading. dark magic. //}\\
\color{white}Zsh can compile scripts to an internal bytecode format\\
\color{white}stored in .zwc (Zsh Word Code) files. Loading compiled\\
\color{white}files is faster than parsing source text every time.\\
\vspace{4pt}\textbf{\color{magenta}WHAT IS A .ZWC FILE?}\\
\color{white}A .zwc file is the compiled form of a .zsh script or an\\
\color{white}autoload function directory. When zsh finds foo.zwc next\\
\color{white}to foo.zsh, it loads the compiled version -- skipping\\
\color{white}the parsing step entirely.\\
\color{white}For individual files: zcompile foo.zsh -> foo.zsh.zwc\\
\color{white}For directories:      zcompile dir.zwc dir/* -> dir.zwc\\
\color{white}The speed difference is measurable on large files and\\
\color{white}during startup when many files are sourced.\\
\vspace{4pt}\textbf{\color{magenta}RECOMPILE: BUILD ALL .ZWC FILES}\\
\texttt{\color{green}zpwr recompile                      \# compile all zpwr files}\\
\texttt{\color{green}zpwr compile                        \# same thing}\\
\color{white}Scans all zpwr source files and compiles them to .zwc.\\
\color{white}This includes:\\
\color{white}- Autoload function files in \textbackslash{}PWR\_AUTOLOAD\\
\color{white}- Shell scripts in \textbackslash{}PWR\_ENV\\
\color{white}- Any .zsh file that has changed since last compile\\
\vspace{4pt}\textbf{\color{magenta}RELATED COMPILE VERBS}\\
\texttt{\color{green}zpwr recompilefiles                 \# compile specific file set}\\
\texttt{\color{green}zpwr compilefiles                   \# alias for above}\\
\texttt{\color{green}zpwr recompilefpath                 \# compile FPATH directories}\\
\texttt{\color{green}zpwr compilefpath                   \# alias for above}\\
\vspace{4pt}\textbf{\color{magenta}DECOMPILE: REMOVE .ZWC FILES}\\
\texttt{\color{green}zpwr decompile                      \# remove all .zwc files}\\
\texttt{\color{green}zpwr uncompile                      \# alias for above}\\
\color{white}Deletes all .zwc files in zpwr. After this, zsh falls\\
\color{white}back to parsing source files directly. Useful when:\\
\color{white}- Debugging a function that behaves differently than source\\
\color{white}- You suspect a stale .zwc file is loading old code\\
\color{white}- zpwr doctor reports stale .zwc files\\
\vspace{4pt}\textbf{\color{magenta}REFRESHZWC: SMART REBUILD}\\
\texttt{\color{green}zpwr refreshzwc                     \# refresh only stale .zwc files}\\
\color{white}Only recompiles .zwc files that are older than their\\
\color{white}source. Faster than a full recompile when only a few\\
\color{white}files changed.\\
\vspace{4pt}\textbf{\color{magenta}THE COMPILATION WORKFLOW}\\
\color{white}After editing zpwr source files:\\
\texttt{\color{green}zpwr recompile                    \# rebuild .zwc files}\\
\color{white}If something is weird after an update:\\
\texttt{\color{green}zpwr decompile \&\& zpwr recompile  \# clean rebuild}\\
\color{white}Quick check for stale files:\\
\texttt{\color{green}zpwr doctor                       \# reports stale .zwc files}\\
\color{white}The .zwc files are safe to delete. They are purely a\\
\color{white}performance cache. Source files are always the authority.\\

\vspace{8pt}

\subsection{Page 39: Compsys Cache}

\vspace{4pt}\textbf{\color{cyan}=== COMPSYS CACHE ===}\vspace{2pt}\\
\textit{\color{gray}// tab completion is a religion. the zcompdump is its scripture. //}\\
\color{white}The zsh completion system (compsys) is the most complex\\
\color{white}part of zsh. It maintains a dump file (zcompdump) that\\
\color{white}caches all registered completions. When it breaks, tab\\
\color{white}completion goes haywire. Here is how to fix it.\\
\vspace{4pt}\textbf{\color{magenta}WHAT IS ZCOMPDUMP?}\\
\color{white}A file (usually \textasciitilde{}/.zcompdump or in \textbackslash{}PWR\_LOCAL) that\\
\color{white}contains all compdef registrations. When zsh starts,\\
\color{white}it sources this file instead of re-scanning all completion\\
\color{white}functions. This saves 100-500ms of startup time.\\
\color{white}Location: \textbackslash{}SH\_COMPDUMP\\
\color{white}The file starts with \#files:TIMESTAMP and contains\\
\color{white}compdef calls for every registered completion.\\
\vspace{4pt}\textbf{\color{magenta}REGENZSH: THE COMPLETION RESET}\\
\texttt{\color{green}zpwr regenzsh                       \# regenerate zcompdump}\\
\color{white}Deletes the old zcompdump and forces compinit to rebuild it.\\
\color{white}This is the \#1 fix for completion problems:\\
\color{white}Symptom: New command installed but tab complete ignores it\\
\texttt{\color{green}Fix:     zpwr regenzsh}\\
\color{white}Symptom: Tab shows wrong completions for a command\\
\texttt{\color{green}Fix:     zpwr cleancompcache \&\& zpwr regenzsh}\\
\color{white}Symptom: \textbackslash{}"command not found: compdef\textbackslash{}" errors on startup\\
\texttt{\color{green}Fix:     zpwr regenzsh}\\
\color{white}Symptom: Startup is slow (500ms+ for completion init)\\
\texttt{\color{green}Fix:     zpwr regenzsh (forces fresh, optimized dump)}\\
\vspace{4pt}\textbf{\color{magenta}THE COMPLETION CACHE DIRECTORY}\\
\color{white}Separate from zcompdump, zsh also caches individual\\
\color{white}completion results (like package names, hostnames, etc.)\\
\color{white}in a cache directory.\\
\texttt{\color{green}zpwr cleancompcache                 \# wipe the cache directory}\\
\color{white}This forces completions to re-query their sources.\\
\color{white}Useful when package lists or hostnames changed.\\
\vspace{4pt}\textbf{\color{magenta}WHY COMPLETIONS BREAK}\\
\color{white}1. Stale zcompdump -- new compdef not registered\\
\color{white}2. Corrupted zcompdump -- syntax error in dump file\\
\color{white}3. Missing FPATH entry -- completion function not found\\
\color{white}4. Wrong function loaded -- .zwc has old version\\
\color{white}5. Plugin conflict -- two plugins register for same cmd\\
\vspace{4pt}\textbf{\color{magenta}THE NUCLEAR FIX}\\
\texttt{\color{green}zpwr cleancompcache}\\
\texttt{\color{green}zpwr decompile}\\
\texttt{\color{green}zpwr regenzsh}\\
\texttt{\color{green}zpwr recompile}\\
\color{white}Four commands: nuke cache, remove .zwc, regen dump, recompile.\\
\color{white}This fixes 99\% of completion problems. The remaining 1\%\\
\color{white}is probably a bug in the completion function itself.\\

\vspace{8pt}

\subsection{Page 40: Tag Generation}

\vspace{4pt}\textbf{\color{cyan}=== TAG GENERATION ===}\vspace{2pt}\\
\textit{\color{gray}// ctrl-] in vim. M-. in emacs. the tags file makes it work. //}\\
\color{white}Tags files are indexes of symbol definitions: functions,\\
\color{white}variables, classes, methods. Your editor uses them to jump\\
\color{white}to definitions with a single keystroke. zpwr generates them.\\
\vspace{4pt}\textbf{\color{magenta}CTAGS: THE CLASSIC}\\
\texttt{\color{green}zpwr regenctags                     \# regenerate ctags file}\\
\color{white}Runs exuberant-ctags or universal-ctags over the zpwr\\
\color{white}codebase. Produces a tags file that vim and emacs read.\\
\color{white}In vim: Ctrl-] jumps to the definition under cursor.\\
\color{white}In vim: Ctrl-T jumps back to where you were.\\
\texttt{\color{green}In vim: :tag zpwrDoctor jumps to zpwrDoctor.}\\
\color{white}ctags understands shell functions, C, Python, Ruby, and\\
\color{white}dozens of other languages out of the box.\\
\vspace{4pt}\textbf{\color{magenta}GTAGS: THE GNU GLOBAL SYSTEM}\\
\texttt{\color{green}zpwr regengtags                     \# regenerate GNU Global tags}\\
\color{white}GNU Global is a more powerful tagging system. It supports:\\
\color{white}- Forward references (where is this defined?)\\
\color{white}- Backward references (where is this called from?)\\
\color{white}- Symbol completion\\
\color{white}Produces GTAGS, GRTAGS, GPATH files in the project root.\\
\texttt{\color{green}Query from command line: global -x zpwrDoctor}\\
\texttt{\color{green}Query references: global -rx zpwrDoctor}\\
\vspace{4pt}\textbf{\color{magenta}GTAGS WITH PYGMENTS PARSER}\\
\texttt{\color{green}zpwr regengtagspygments             \# gtags with Pygments backend}\\
\color{white}Uses the Python Pygments parser instead of the built-in\\
\color{white}one. Pygments understands more languages and file types.\\
\color{white}Requires: pip install pygments\\
\vspace{4pt}\textbf{\color{magenta}GTAGS BY FILE TYPE}\\
\texttt{\color{green}zpwr regengtagstype                 \# gtags filtered by type}\\
\color{white}Generate tags for a specific file type only.\\
\color{white}Faster than a full regen when you only care about one language.\\
\vspace{4pt}\textbf{\color{magenta}HOW TAGS WORK WITH YOUR EDITOR}\\
\color{white}vim    reads 'tags' file. :set tags=./tags;,tags\\
\color{white}Ctrl-] to jump, Ctrl-T to return\\
\color{white}:tselect for multiple matches\\
\color{white}emacs  reads TAGS file (etags format) or uses gtags\\
\color{white}M-. to find definition, M-* to return\\
\color{white}ggtags-mode or helm-gtags for GNU Global\\
\color{white}nvim   same as vim, plus LSP which often replaces tags\\
\color{white}but tags still work as a fast fallback\\
\vspace{4pt}\textbf{\color{magenta}WHEN TO REGENERATE}\\
\color{white}- After adding or renaming functions\\
\color{white}- After a zpwr update that adds new verbs\\
\color{white}- When \textbackslash{}"tag not found\textbackslash{}" errors appear for known functions\\
\texttt{\color{green}- Use zpwr watch '*.zsh' 'zpwr regenctags' for auto-regen}\\

\vspace{8pt}

\subsection{Page 41: Cache Regeneration}

\vspace{4pt}\textbf{\color{cyan}=== CACHE REGENERATION ===}\vspace{2pt}\\
\textit{\color{gray}// the opposite of cleaning. building the caches back up. //}\\
\color{white}Cleaning removes caches. Regeneration rebuilds them.\\
\color{white}These verbs construct the indexes and precomputed data\\
\color{white}that make zpwr fast.\\
\vspace{4pt}\textbf{\color{magenta}REGEN: THE STANDARD REBUILD}\\
\texttt{\color{green}zpwr regen                          \# regenerate common caches}\\
\color{white}Rebuilds the most frequently used caches:\\
\color{white}environment cache, config links, keybindings.\\
\color{white}This is the verb you reach for after editing config files.\\
\vspace{4pt}\textbf{\color{magenta}REGENALL: THE FULL REBUILD}\\
\texttt{\color{green}zpwr regenall                       \# regenerate everything}\\
\color{white}Every cache. Every index. Every precomputed file.\\
\color{white}This is the thorough version of regen. Takes longer\\
\color{white}but leaves nothing stale.\\
\vspace{4pt}\textbf{\color{magenta}INDIVIDUAL REGEN VERBS}\\
\color{white}regenenvcache\\
\texttt{\color{green}zpwr regenenvcache                  \# rebuild environment cache}\\
\color{white}Snapshots all env vars to a cache file for fast lookup.\\
\color{white}Used by environmentcachesearch and other env verbs.\\
\color{white}regenkeybindings\\
\texttt{\color{green}zpwr regenkeybindings               \# rebuild keybinding index}\\
\color{white}Re-registers all zpwr keybindings. Run this after\\
\color{white}editing keybinding configuration or after a plugin\\
\color{white}clobbered your bindings.\\
\color{white}regenconfiglinks\\
\texttt{\color{green}zpwr regenconfiglinks               \# rebuild config symlinks}\\
\color{white}Re-creates the symlinks from \textasciitilde{}/.zshrc, \textasciitilde{}/.vimrc, etc.\\
\color{white}into the zpwr directory tree. Run this after a fresh\\
\color{white}clone or when zpwr doctor reports broken config links.\\
\color{white}regenhistory\\
\texttt{\color{green}zpwr regenhistory                   \# clean up history file}\\
\color{white}Processes the HISTFILE to remove corrupted entries.\\
\color{white}Zsh history files can accumulate broken lines from\\
\color{white}concurrent writes or improper shutdowns.\\
\color{white}regenpowerline\\
\texttt{\color{green}zpwr regenpowerline                 \# rebuild powerline config}\\
\color{white}Regenerates powerline/p10k configuration. Run when\\
\color{white}the prompt looks wrong after a terminal font change.\\
\color{white}regengitrepocache\\
\texttt{\color{green}zpwr regengitrepocache              \# scan and cache git repos}\\
\color{white}Walks the filesystem finding git repos and builds a\\
\color{white}master list. Used by repo navigation and batch verbs.\\
\vspace{4pt}\textbf{\color{magenta}THE REGEN CHEAT SHEET}\\
\texttt{\color{green}After editing config:     zpwr regen}\\
\texttt{\color{green}After zpwr update:        zpwr regenall}\\
\texttt{\color{green}After env var changes:    zpwr regenenvcache}\\
\texttt{\color{green}After repo changes:       zpwr regengitrepocache}\\
\texttt{\color{green}After completion issues:  zpwr regenzsh}\\

\vspace{8pt}

\subsection{Page 42: Updates \& Dependencies}

\vspace{4pt}\textbf{\color{cyan}=== UPDATES \& DEPENDENCIES ===}\vspace{2pt}\\
\textit{\color{gray}// staying current without breaking everything. a balancing act. //}\\
\color{white}zpwr is a living system that depends on external tools,\\
\color{white}plugins, and its own git repository. These verbs keep\\
\color{white}everything synchronized.\\
\vspace{4pt}\textbf{\color{magenta}UPDATE: THE CORE UPDATE}\\
\texttt{\color{green}zpwr update                         \# update zpwr itself}\\
\color{white}Pulls the latest zpwr from the git remote.\\
\color{white}This is a git pull on \textbackslash{}PWR.\\
\color{white}Your local config in \textbackslash{}PWR\_LOCAL is not touched.\\
\vspace{4pt}\textbf{\color{magenta}UPDATEALL: THE EVERYTHING UPDATE}\\
\texttt{\color{green}zpwr updateall                      \# update zpwr + all plugins}\\
\color{white}Updates zpwr, zinit plugins, and other managed components.\\
\color{white}This is the weekly maintenance verb. Run it when you want\\
\color{white}the latest of everything.\\
\vspace{4pt}\textbf{\color{magenta}UPDATEDEPS: DEPENDENCY UPDATES}\\
\texttt{\color{green}zpwr updatedeps                     \# update external dependencies}\\
\color{white}Updates the external tools zpwr depends on: zinit plugins,\\
\color{white}zsh-autosuggestions, zsh-syntax-highlighting, etc.\\
\color{white}Does not touch zpwr's own code.\\
\vspace{4pt}\textbf{\color{magenta}UPDATEDEPSCLEAN: CLEAN + UPDATE DEPS}\\
\texttt{\color{green}zpwr updatedepsclean                \# fresh dependency install}\\
\color{white}Cleans stale dependency state first, then updates.\\
\color{white}Use when updatedeps alone does not fix plugin issues.\\
\vspace{4pt}\textbf{\color{magenta}UPDATEZINIT: ZINIT ONLY}\\
\texttt{\color{green}zpwr updatezinit                    \# update zinit plugin manager}\\
\color{white}Updates just the zinit plugin manager binary, not the\\
\color{white}plugins it manages. Use when zinit itself has a bug or\\
\color{white}you need a newer zinit feature.\\
\vspace{4pt}\textbf{\color{magenta}UPDATEPULL: GIT PULL ALL REPOS}\\
\texttt{\color{green}zpwr updatepull                     \# pull all tracked git repos}\\
\color{white}Iterates over all cached git repos and runs git pull.\\
\color{white}Useful when you have many repos that track remotes and\\
\color{white}you want them all current.\\
\vspace{4pt}\textbf{\color{magenta}THE UPDATE HIERARCHY}\\
\color{white}update              just zpwr (git pull)\\
\color{white}updatedeps          just plugins and tools\\
\color{white}updatezinit         just the plugin manager\\
\color{white}updateall           zpwr + plugins + tools\\
\color{white}updatepull          all tracked git repos\\
\color{white}updatedepsclean     clean slate + update deps\\
\vspace{4pt}\textbf{\color{magenta}RECOMMENDED UPDATE SCHEDULE}\\
\texttt{\color{green}Weekly:      zpwr updateall}\\
\texttt{\color{green}After update: zpwr recompile}\\
\texttt{\color{green}If broken:   zpwr updatedepsclean}\\
\color{white}Always recompile after updating. New source files need\\
\color{white}new .zwc files. The old compiled versions will load\\
\color{white}stale code until you recompile.\\

\vspace{8pt}


\section{CHAPTER 11: MONITORING \& AUTOMATION}

\subsection{Page 43: File Watching}

\vspace{4pt}\textbf{\color{cyan}=== FILE WATCHING ===}\vspace{2pt}\\
\textit{\color{gray}// your files change. your command runs. automatically. //}\\
\color{white}The watch verb monitors files for changes and auto-executes\\
\color{white}a command when something is modified. It is the poor man's\\
\color{white}CI pipeline, running locally in your terminal.\\
\vspace{4pt}\textbf{\color{magenta}USAGE}\\
\texttt{\color{green}zpwr watch '<GLOB>' '<COMMAND>'}\\
\color{white}Two arguments: what to watch, and what to run.\\
\color{white}Quote both arguments to prevent premature expansion.\\
\vspace{4pt}\textbf{\color{magenta}PRACTICAL EXAMPLES}\\
\texttt{\color{green}zpwr watch '*.zsh' 'zpwr recompile'}\\
\color{white}Recompile zpwr whenever a .zsh file changes.\\
\color{white}Edit a function, save, and .zwc files auto-rebuild.\\
\texttt{\color{green}zpwr watch 'src/**/*.ts' 'npm run build'}\\
\color{white}Build your TypeScript project on every source change.\\
\texttt{\color{green}zpwr watch 'Makefile' 'make'}\\
\color{white}Run make whenever the Makefile changes.\\
\texttt{\color{green}zpwr watch '.' 'echo changed'}\\
\color{white}Watch the entire current directory.\\
\texttt{\color{green}zpwr watch '*.py' 'python -m pytest'}\\
\color{white}Run tests on every Python file change. TDD mode.\\
\vspace{4pt}\textbf{\color{magenta}PLATFORM BACKEND}\\
\color{white}macOS:   Uses fswatch (brew install fswatch)\\
\color{white}Linux:   Uses inotifywait (apt install inotify-tools)\\
\color{white}The verb auto-detects which tool is available.\\
\color{white}If neither is found, it tells you what to install.\\
\vspace{4pt}\textbf{\color{magenta}THE OUTPUT FORMAT}\\
\color{white}Each trigger shows a box with:\\
\color{white}- Trigger count (how many times it has fired)\\
\color{white}- Timestamp      (when the change was detected)\\
\color{white}- Filename       (which file changed)\\
\color{white}- Exit status    (green ok or red exit code)\\
\color{white}- Elapsed time   (how long the command took)\\
\vspace{4pt}\textbf{\color{magenta}HOW TO STOP}\\
\color{white}Press Ctrl-C. The watcher process terminates cleanly.\\
\vspace{4pt}\textbf{\color{magenta}GLOB HANDLING}\\
\color{white}The glob is split into directory and file filter:\\
\color{white}'src/**/*.ts' -> watch src/ recursively, filter *.ts\\
\color{white}'.'           -> watch current dir, no filter\\
\color{white}'*.zsh'       -> watch ./, filter *.zsh\\
\color{white}The glob-to-regex conversion handles * and ? patterns.\\
\color{white}For fswatch, it uses -e '.*' -i to exclude-then-include.\\
\color{white}For inotifywait, it uses zsh's glob matching on filenames.\\

\vspace{8pt}

\subsection{Page 44: Command Replay}

\vspace{4pt}\textbf{\color{cyan}=== COMMAND REPLAY ===}\vspace{2pt}\\
\textit{\color{gray}// your history is a script. replay it on demand. //}\\
\color{white}The replay verb plays back the last N commands from your\\
\color{white}shell history. Three modes: look, execute, or tmux pane.\\
\color{white}It is a VCR for your terminal.\\
\vspace{4pt}\textbf{\color{magenta}BASIC USAGE}\\
\texttt{\color{green}zpwr replay                         \# preview last 10 commands}\\
\texttt{\color{green}zpwr replay 20                      \# preview last 20}\\
\texttt{\color{green}zpwr replay 5 exec                  \# execute last 5 with timing}\\
\texttt{\color{green}zpwr replay 10 tmux                 \# replay in a tmux split pane}\\
\vspace{4pt}\textbf{\color{magenta}THE THREE MODES}\\
\color{white}PREVIEW MODE (default, safe)\\
\texttt{\color{green}zpwr replay}\\
\texttt{\color{green}zpwr replay 20 preview}\\
\color{white}Lists the commands with line numbers. Nothing executes.\\
\color{white}Use this to review what you did before replaying it.\\
\color{white}The output shows helpful hints for switching to exec/tmux.\\
\color{white}EXEC MODE (runs in current shell)\\
\texttt{\color{green}zpwr replay 5 exec}\\
\color{white}Executes each command sequentially in your current shell.\\
\color{white}Each command shows:\\
\color{white}- Command number [1/5]\\
\color{white}- The command being run\\
\color{white}- Exit status (green ok or red exit code)\\
\color{white}- Elapsed time in milliseconds\\
\color{white}WARNING: This runs real commands. If your last 5 commands\\
\color{white}included rm -rf something, it will do it again.\\
\color{white}Always preview first.\\
\color{white}TMUX MODE (runs in a new pane)\\
\texttt{\color{green}zpwr replay 10 tmux}\\
\color{white}Opens a new tmux split pane and replays commands there.\\
\color{white}Your current pane is untouched. You can watch the replay\\
\color{white}happen side-by-side with your current work.\\
\color{white}The pane waits for Enter when done, then cleans up.\\
\color{white}Requires an active tmux session.\\
\vspace{4pt}\textbf{\color{magenta}USE CASES}\\
\color{white}Reproduce a bug:       zpwr replay 3 exec\\
\color{white}Demo your workflow:     zpwr replay 20 tmux\\
\color{white}Audit what you did:     zpwr replay 50\\
\color{white}Re-run a build chain:   zpwr replay 4 exec\\
\vspace{4pt}\textbf{\color{magenta}HOW IT WORKS}\\
\color{white}Uses fc -rl -N to extract the last N history entries.\\
\color{white}Reverses them to chronological order. In exec mode,\\
\color{white}each command is eval'd with timing via EPOCHREALTIME.\\
\color{white}In tmux mode, a temp script is built and executed in\\
\color{white}a split pane with a 0.5s delay between commands for\\
\color{white}readability. The script self-destructs after completion.\\

\vspace{8pt}

\subsection{Page 45: Log Surveillance}

\vspace{4pt}\textbf{\color{cyan}=== LOG SURVEILLANCE ===}\vspace{2pt}\\
\textit{\color{gray}// watching logs in glorious technicolor. the ccze way. //}\\
\color{white}zpwr logs everything: startup, errors, plugin loads, debug\\
\color{white}messages. The taillog verb tails these logs with optional\\
\color{white}ccze colorization for maximum readability.\\
\vspace{4pt}\textbf{\color{magenta}BASIC USAGE}\\
\texttt{\color{green}zpwr taillog                        \# follow log with colors}\\
\texttt{\color{green}zpwr taillog -n 50                  \# last 50 lines, then follow}\\
\texttt{\color{green}zpwr taillog -f                     \# colorize without following}\\
\texttt{\color{green}zpwr taillog -r                     \# raw output with line numbers}\\
\vspace{4pt}\textbf{\color{magenta}THE FLAGS}\\
\color{white}-n, --lines N     Show last N lines before following\\
\color{white}-f, --no-follow   Just colorize, do not tail -f\\
\color{white}-r, --raw         No ccze, plain output with line numbers\\
\vspace{4pt}\textbf{\color{magenta}CUSTOM LOGFILE}\\
\texttt{\color{green}zpwr taillog /var/log/system.log}\\
\texttt{\color{green}zpwr taillog /var/log/syslog}\\
\color{white}The last argument can be any logfile path.\\
\color{white}Default: \textbackslash{}PWR\_LOGFILE (the zpwr logfile).\\
\vspace{4pt}\textbf{\color{magenta}CCZE: THE SECRET WEAPON}\\
\color{white}ccze is a log colorizer that understands dozens of log\\
\color{white}formats: syslog, Apache, Postfix, PHP, kernel messages,\\
\color{white}and more. It auto-detects the format and colorizes:\\
\color{white}timestamps     in one color\\
\color{white}warnings       in yellow\\
\color{white}errors         in red\\
\color{white}info           in cyan\\
\color{white}hostnames      in magenta\\
\color{white}Install: brew install ccze  or  apt install ccze\\
\color{white}If ccze is not found, taillog falls back to plain output\\
\color{white}with a helpful install hint.\\
\vspace{4pt}\textbf{\color{magenta}THE ZPWR LOGGING VERBS}\\
\color{white}taillog consumes logs. These verbs produce them:\\
\texttt{\color{green}zpwr logdebug <msg>                 \# debug level}\\
\texttt{\color{green}zpwr loginfo <msg>                  \# info level}\\
\texttt{\color{green}zpwr logerror <msg>                 \# error level}\\
\texttt{\color{green}zpwr logtrace <msg>                 \# trace level (most verbose)}\\
\color{white}The *console variants also print to stderr:\\
\texttt{\color{green}zpwr logdebugconsole <msg>}\\
\texttt{\color{green}zpwr loginfoconsole <msg>}\\
\texttt{\color{green}zpwr logerrorconsole <msg>}\\
\texttt{\color{green}zpwr logtraceconsole <msg>}\\
\vspace{4pt}\textbf{\color{magenta}LOG MAINTENANCE}\\
\color{white}Logs grow forever. Clean them periodically:\\
\texttt{\color{green}zpwr cleanlog                       \# truncate the logfile}\\
\color{white}Or check the size first: ls -lh \textbackslash{}PWR\_LOGFILE\\
\color{white}A 100MB logfile means your debug logging is too chatty.\\
\color{white}A 1KB logfile means nothing is being logged. Both are\\
\color{white}worth investigating.\\

\vspace{8pt}

\subsection{Page 46: Backup \& Restore}

\vspace{4pt}\textbf{\color{cyan}=== BACKUP \& RESTORE ===}\vspace{2pt}\\
\textit{\color{gray}// your history is priceless. these verbs are the insurance policy. //}\\
\color{white}Your shell history is years of accumulated knowledge:\\
\color{white}every command you ran, every fix you found at 2am, every\\
\color{white}pipeline you crafted. Losing it is losing institutional\\
\color{white}memory. These verbs prevent that.\\
\vspace{4pt}\textbf{\color{magenta}BACKUP: THE GENERAL BACKUP}\\
\texttt{\color{green}zpwr backup                         \# backup zpwr configuration}\\
\color{white}Creates a timestamped backup of zpwr configuration files.\\
\color{white}This captures your local settings, tokens, and any\\
\color{white}customizations in \textbackslash{}PWR\_LOCAL.\\
\color{white}Backup destination: \textbackslash{}PWR\_LOCAL/rcBackups/\\
\color{white}Each backup is timestamped so you can have many.\\
\vspace{4pt}\textbf{\color{magenta}BACKUPHISTORY: THE HISTORY VAULT}\\
\texttt{\color{green}zpwr backuphistory                  \# backup HISTFILE}\\
\color{white}Copies your HISTFILE to the backup directory with a\\
\color{white}timestamp. This is the most important backup in zpwr.\\
\color{white}Your HISTFILE can be corrupted by:\\
\color{white}- Power loss during write\\
\color{white}- Concurrent sessions fighting over the file\\
\color{white}- Disk full causing truncated writes\\
\color{white}- Bad .zsh\_history format from import/merge\\
\texttt{\color{green}A weekly zpwr backuphistory gives you a safety net.}\\
\vspace{4pt}\textbf{\color{magenta}RESTOREHISTORY: THE TIME MACHINE}\\
\texttt{\color{green}zpwr restorehistory                 \# restore from backup}\\
\color{white}Restores HISTFILE from the most recent backup.\\
\color{white}Use when your history is corrupted or accidentally\\
\color{white}truncated.\\
\color{white}The restore process:\\
\color{white}1. Finds the latest backup in \textbackslash{}PWR\_LOCAL/rcBackups\\
\color{white}2. Backs up the current (possibly broken) HISTFILE\\
\color{white}3. Copies the backup over HISTFILE\\
\color{white}4. Reloads history with fc -R\\
\vspace{4pt}\textbf{\color{magenta}SNAPSHOTS VS BACKUPS}\\
\color{white}snapshot          captures full env state (aliases, vars, tmux)\\
\color{white}backup            copies config files to safety\\
\color{white}backuphistory     specifically targets HISTFILE\\
\color{white}Use snapshots for environment state.\\
\color{white}Use backups for file safety.\\
\color{white}They complement each other.\\
\vspace{4pt}\textbf{\color{magenta}THE SAFETY PROTOCOL}\\
\color{white}Before a big change:\\
\texttt{\color{green}zpwr snapshot before-change}\\
\texttt{\color{green}zpwr backuphistory}\\
\texttt{\color{green}zpwr backup}\\
\color{white}Now you have three layers of protection:\\
\color{white}env snapshot, history backup, and config backup.\\
\color{white}If the change goes sideways, you can recover everything.\\
\vspace{4pt}\textbf{\color{magenta}CRON IT}\\
\color{white}Add to your crontab (crontab -e):\\
\texttt{\color{green}0 3 * * 0 zsh -c 'source \textasciitilde{}/.zshrc; zpwr backuphistory'}\\
\color{white}Every Sunday at 3am. Automatic insurance.\\

\vspace{8pt}


\section{CHAPTER 12: UTILITY ARSENAL}

\subsection{Page 47: File \& Process Operations}

\vspace{4pt}\textbf{\color{cyan}=== FILE \& PROCESS OPERATIONS ===}\vspace{2pt}\\
\textit{\color{gray}// the everyday tools. sharpened and color-coded. //}\\
\color{white}zpwr wraps common file and process operations with\\
\color{white}enhanced output, color, and fzf integration. These are\\
\color{white}the verbs you reach for twenty times a day.\\
\vspace{4pt}\textbf{\color{magenta}FILE LISTING}\\
\texttt{\color{green}zpwr clearlist                      \# clear screen + list files}\\
\texttt{\color{green}zpwr clearls                        \# alias for clearlist}\\
\texttt{\color{green}zpwr info                           \# alias for clearlist}\\
\texttt{\color{green}zpwr list                           \# list files (no clear)}\\
\texttt{\color{green}zpwr ls                             \# alias for list}\\
\color{white}clearlist is the most-used verb after help. It clears\\
\color{white}the terminal and shows a fresh directory listing.\\
\color{white}Uses eza if available, falls back to ls.\\
\vspace{4pt}\textbf{\color{magenta}FILE VIEWING}\\
\texttt{\color{green}zpwr cat <file>                     \# colorized cat}\\
\texttt{\color{green}zpwr c <file>                       \# shortest form}\\
\texttt{\color{green}zpwr catcd <dir>                    \# cat files in a directory}\\
\color{white}cat/c uses bat or pygmentize for syntax highlighting\\
\color{white}when available. Plain cat as fallback. catcd lists the\\
\color{white}contents of each file in a directory.\\
\vspace{4pt}\textbf{\color{magenta}PROCESS MANAGEMENT}\\
\texttt{\color{green}zpwr ps                             \# enhanced process listing}\\
\texttt{\color{green}zpwr kill                           \# fzf-powered process kill}\\
\texttt{\color{green}zpwr killedit                       \# kill then edit related file}\\
\texttt{\color{green}zpwr kill is the crown jewel here. It pipes process}\\
\color{white}listings through fzf, letting you fuzzy-search for the\\
\color{white}process you want to terminate. Select it, confirm, dead.\\
\color{white}No more guessing PIDs.\\
\vspace{4pt}\textbf{\color{magenta}FILE DESCRIPTORS}\\
\texttt{\color{green}zpwr lsof                           \# list open files (enhanced)}\\
\texttt{\color{green}zpwr lsofedit                       \# lsof then edit result}\\
\color{white}zpwr lsof wraps lsof with fzf filtering.\\
\color{white}Useful when tracking down which process holds a file\\
\color{white}or which files a process has open.\\
\vspace{4pt}\textbf{\color{magenta}FILE OPENING}\\
\texttt{\color{green}zpwr open                           \# open file/dir with system}\\
\color{white}Opens the target with the system handler: Finder on macOS,\\
\color{white}xdg-open on Linux. Works with files, directories, URLs.\\
\vspace{4pt}\textbf{\color{magenta}THE COMBO MOVES}\\
\color{white}Clear terminal, list files, check what is running:\\
\texttt{\color{green}zpwr clearlist \&\& zpwr ps}\\
\color{white}Find and kill a zombie process:\\
\texttt{\color{green}zpwr kill     \# fzf search -> select -> terminate}\\
\color{white}Find what is holding a port:\\
\texttt{\color{green}zpwr lsof     \# fzf through open files and ports}\\

\vspace{8pt}

\subsection{Page 48: Batch Execution}

\vspace{4pt}\textbf{\color{cyan}=== BATCH EXECUTION ===}\vspace{2pt}\\
\textit{\color{gray}// one command, many targets. the industrial strength loop. //}\\
\color{white}When you need to run the same operation across multiple\\
\color{white}files, directories, or git repos, these verbs replace\\
\color{white}hand-rolled for loops with tested, colorized execution.\\
\vspace{4pt}\textbf{\color{magenta}FOR: THE BASIC ITERATOR}\\
\texttt{\color{green}zpwr for <command>                  \# run command in current dir context}\\
\color{white}Executes a command with zpwr's enhanced output formatting.\\
\color{white}Shows timing, exit status, and colorized output.\\
\vspace{4pt}\textbf{\color{magenta}FOR10: THE STRESS TEST}\\
\texttt{\color{green}zpwr for10 <command>                \# run command 10 times}\\
\color{white}Executes a command 10 times. Useful for:\\
\color{white}- Benchmarking (is the timing consistent?)\\
\color{white}- Flaky test detection (does it fail intermittently?)\\
\color{white}- Load testing (can it handle rapid repeated execution?)\\
\vspace{4pt}\textbf{\color{magenta}FORDIR: ITERATE OVER GIT REPOS}\\
\texttt{\color{green}zpwr fordir <command>               \# run command in each git repo}\\
\color{white}Iterates over all cached git repos, cd's into each one,\\
\color{white}and runs the command. Extremely powerful for bulk operations:\\
\texttt{\color{green}zpwr fordir 'git status -s'         \# dirty files in all repos}\\
\texttt{\color{green}zpwr fordir 'git pull'              \# pull all repos}\\
\texttt{\color{green}zpwr fordir 'git stash list'        \# find forgotten stashes}\\
\color{white}Each repo's output is labeled with its path so you\\
\color{white}know which repo produced which output.\\
\vspace{4pt}\textbf{\color{magenta}FORDIRUPDATE: PULL ALL REPOS}\\
\texttt{\color{green}zpwr fordirupdate                   \# git pull in every repo}\\
\color{white}Specialized version of fordir that runs git pull in\\
\texttt{\color{green}every cached repo. Same as zpwr fordir 'git pull'}\\
\color{white}but optimized for the common case.\\
\vspace{4pt}\textbf{\color{magenta}EXECGLOB: GLOB-BASED EXECUTION}\\
\texttt{\color{green}zpwr execglob <glob> <command>      \# run command on glob matches}\\
\color{white}Finds files matching a glob and runs a command on each:\\
\texttt{\color{green}zpwr execglob '*.log' 'wc -l'      \# line count each log file}\\
\texttt{\color{green}zpwr execglob '*.bak' 'rm'         \# delete all .bak files}\\
\vspace{4pt}\textbf{\color{magenta}EXECGLOBPARALLEL: PARALLEL GLOB EXECUTION}\\
\texttt{\color{green}zpwr execglobparallel <glob> <cmd>  \# parallel execution}\\
\color{white}Same as execglob but runs commands in parallel.\\
\color{white}Useful when processing many files that are I/O-bound:\\
\texttt{\color{green}zpwr execglobparallel '*.png' 'optipng'  \# compress images}\\
\vspace{4pt}\textbf{\color{magenta}THE POWER COMBOS}\\
\color{white}Find all repos with uncommitted work:\\
\texttt{\color{green}zpwr fordir 'git diff --stat HEAD'}\\
\color{white}Count lines in all Python files:\\
\texttt{\color{green}zpwr execglob '**/*.py' 'wc -l'}\\
\color{white}Stress test a flaky command:\\
\texttt{\color{green}zpwr for10 'curl -s localhost:3000/health'}\\

\vspace{8pt}

\subsection{Page 49: Text Operations}

\vspace{4pt}\textbf{\color{cyan}=== TEXT OPERATIONS ===}\vspace{2pt}\\
\textit{\color{gray}// find and replace. rename. count. transform. //}\\
\color{white}Text manipulation is the bread and butter of shell work.\\
\color{white}These verbs handle the common cases: renaming files,\\
\color{white}replacing strings, counting lines, and transforming text.\\
\vspace{4pt}\textbf{\color{magenta}RENAME: FILE RENAMING}\\
\texttt{\color{green}zpwr rename                         \# rename files with patterns}\\
\color{white}Renames files using pattern matching. Safer than raw mv\\
\color{white}because it shows what will happen before doing it.\\
\vspace{4pt}\textbf{\color{magenta}REPLACER: FIND AND REPLACE IN FILES}\\
\texttt{\color{green}zpwr replacer                       \# search and replace in files}\\
\color{white}Performs find-and-replace across multiple files.\\
\color{white}The codebase refactoring tool. When you renamed a\\
\color{white}function and need to update every call site.\\
\vspace{4pt}\textbf{\color{magenta}CHANGENAME: RENAME WITH HISTORY}\\
\texttt{\color{green}zpwr changename                     \# rename with tracking}\\
\color{white}Renames files or directories with additional tracking.\\
\color{white}More sophisticated than basic rename for complex\\
\color{white}renaming operations.\\
\vspace{4pt}\textbf{\color{magenta}CHANGENAMEZPWR: ZPWR-SPECIFIC RENAME}\\
\texttt{\color{green}zpwr changenamezpwr                 \# rename within zpwr codebase}\\
\color{white}Specialized rename for zpwr function and file names.\\
\color{white}Handles the autoload naming convention (file name must\\
\color{white}match function name) and updates references.\\
\vspace{4pt}\textbf{\color{magenta}LINECOUNT: COUNT LINES}\\
\texttt{\color{green}zpwr linecount                      \# count lines in files}\\
\color{white}Enhanced line counting with formatting and totals.\\
\color{white}Shows per-file counts and a grand total.\\
\color{white}More informative than raw wc -l.\\
\vspace{4pt}\textbf{\color{magenta}PRE \& POST: TEXT INJECTION}\\
\texttt{\color{green}zpwr pre                            \# prepend text to input}\\
\texttt{\color{green}zpwr post                           \# append text to input}\\
\color{white}Pipe filters that add text before or after each line.\\
\color{white}Useful in pipelines for formatting:\\
\texttt{\color{green}ls | zpwr pre '  - '               \# bullet-point file list}\\
\texttt{\color{green}cat urls.txt | zpwr post '/api/v1'  \# append to each URL}\\
\vspace{4pt}\textbf{\color{magenta}URLSAFE: URL ENCODING}\\
\texttt{\color{green}zpwr urlsafe                        \# URL-encode a string}\\
\color{white}Converts special characters to percent-encoded form.\\
\color{white}Spaces become \%20, ampersands become \%26, etc.\\
\color{white}Essential when constructing URLs from user input.\\
\vspace{4pt}\textbf{\color{magenta}THE TEXT PIPELINE}\\
\color{white}zpwr text verbs compose naturally in pipelines:\\
\texttt{\color{green}cat data.txt | zpwr pre 'ROW: ' | zpwr post ' ;'}\\
\color{white}Result: ROW: line1 ;\\
\color{white}ROW: line2 ;\\
\color{white}Rename all .txt to .md:\\
\texttt{\color{green}zpwr rename                         \# interactive rename}\\
\color{white}Replace 'oldFunc' with 'newFunc' across the codebase:\\
\texttt{\color{green}zpwr replacer                       \# interactive replace}\\

\vspace{8pt}

\subsection{Page 50: Display \& Banners}

\vspace{4pt}\textbf{\color{cyan}=== DISPLAY \& BANNERS ===}\vspace{2pt}\\
\textit{\color{gray}// because aesthetics matter. even in a terminal. //}\\
\color{white}zpwr ships with a fleet of display verbs that produce\\
\color{white}ASCII art, colorized output, and terminal animations.\\
\color{white}Some are informational. Some are just cool to look at.\\
\vspace{4pt}\textbf{\color{magenta}THE BANNER FAMILY}\\
\texttt{\color{green}zpwr about                          \# the welcome banner}\\
\texttt{\color{green}zpwr banner                         \# same as about}\\
\texttt{\color{green}zpwr bannercounts                   \# banner + environment stats}\\
\color{white}The core banner shows the zpwr ASCII logo with version,\\
\color{white}author, and system info. bannercounts adds the environment\\
\color{white}census (alias count, function count, etc.).\\
\texttt{\color{green}zpwr bannerlolcat                   \# banner piped through lolcat}\\
\color{white}Rainbow gradient banner. Requires: gem install lolcat\\
\color{white}For when monochrome is not enough and you need to\\
\color{white}assert terminal dominance with a rainbow.\\
\texttt{\color{green}zpwr bannerpony                     \# banner with ponysay}\\
\color{white}Displays the banner inside an ASCII pony speech bubble.\\
\color{white}Requires: pip install ponysay\\
\color{white}Peak terminal culture.\\
\texttt{\color{green}zpwr bannernopony                   \# banner without the pony}\\
\color{white}For the pony-averse. All the info, none of the horse.\\
\vspace{4pt}\textbf{\color{magenta}PRETTY PRINTING}\\
\texttt{\color{green}zpwr prettyprint                    \# formatted text output}\\
\texttt{\color{green}zpwr altprettyprint                 \# alternate format}\\
\texttt{\color{green}zpwr printmap                       \# print associative array}\\
\color{white}prettyprint is used internally by many zpwr verbs to\\
\color{white}format output with consistent color and spacing.\\
\color{white}printmap formats associative arrays as aligned key-value\\
\color{white}tables. Useful for debugging hash tables.\\
\vspace{4pt}\textbf{\color{magenta}COLOR TESTING}\\
\texttt{\color{green}zpwr colortest                      \# test basic terminal colors}\\
\texttt{\color{green}zpwr colortest256                   \# full 256-color palette}\\
\color{white}Run these after changing terminals, fonts, or color\\
\color{white}schemes. colortest shows the 16 ANSI colors. colortest256\\
\color{white}shows the full 256-color grid. If colors look wrong,\\
\color{white}your TERM or terminal emulator config needs adjustment.\\
\vspace{4pt}\textbf{\color{magenta}ANIMATION}\\
\texttt{\color{green}zpwr animate                        \# terminal animation}\\
\color{white}Plays an ASCII animation in the terminal.\\
\color{white}Because sometimes you need to entertain yourself\\
\color{white}while waiting for a build.\\
\vspace{4pt}\textbf{\color{magenta}THE DISPLAY ARSENAL AT A GLANCE}\\
\texttt{\color{green}Impress a colleague:      zpwr bannerlolcat}\\
\texttt{\color{green}Debug your terminal:       zpwr colortest256}\\
\texttt{\color{green}Count your weapons:        zpwr bannercounts}\\
\texttt{\color{green}Maximum absurdity:         zpwr bannerpony}\\
\texttt{\color{green}Format a hash table:       zpwr printmap}\\
\texttt{\color{green}Kill time aesthetically:   zpwr animate}\\

\vspace{8pt}


\section{CHAPTER 13: EMACS DIMENSION}

\subsection{Page 51: Emacs Integration Overview}

\vspace{4pt}\textbf{\color{cyan}>>> ENTERING THE EMACS DIMENSION <<<}\vspace{2pt}\\
\color{white}where M-x is a way of life and pinky fingers go to die\\
\color{white}You thought there were only vim verbs?\\
\color{white}Think again, operative. ZPWR mirrors every vim operation\\
\color{white}with an emacs equivalent. Same power, different religion.\\
\vspace{4pt}\textbf{\color{magenta}THE 28 EMACS VERBS:}\\
\color{white}FILE OPS      emacsall, emacsalledit, emacsautoload\\
\color{white}RECENT        emacsrecent, emacsrecentcd, emacsrecentsudo\\
\color{white}SEARCH        emacsfilesearch, emacsfilesearchedit\\
\color{white}emacsfindsearch, emacsfindsearchedit\\
\color{white}emacslocatesearch, emacslocatesearchedit\\
\color{white}emacswordsearch, emacswordsearchedit\\
\color{white}TAGS          emacsctags, emacsgtags, emacsgtagsedit\\
\color{white}CONFIG        emacsconfig, emacsemacsconfig, emacstokens\\
\color{white}SERVER        emacsallserver\\
\color{white}SCRIPTS       emacsscripts, emacsscriptedit\\
\color{white}META          emacsplugincount, emacspluginlist, emacszpwr\\
\color{white}THE MIRROR PRINCIPLE:\\
\color{white}Every vim verb has an emacs twin. The architecture is\\
\color{white}symmetrical by design. You can swap editors without\\
\color{white}losing any ZPWR functionality.\\
\color{white}vim verb              emacs twin\\
\texttt{\color{green}zpwr vimall     -->   zpwr emacsall}\\
\texttt{\color{green}zpwr vimrecent  -->   zpwr emacsrecent}\\
\texttt{\color{green}zpwr vimscripts -->   zpwr emacsscripts}\\
\vspace{4pt}\textbf{\color{magenta}PHILOSOPHY:}\\
\color{white}ZPWR does not care about your editor war allegiance.\\
\color{white}It arms both sides equally. The real enemy is notepad.\\
\textit{\color{gray}// the emacs dimension runs parallel to vim-space}\\
\textit{\color{gray}// choose your weapon, the framework adapts}\\
\textit{\color{gray}// \textbackslash{}PWR\_EDITOR controls which universe you inhabit}\\

\vspace{8pt}

\subsection{Page 52: Emacs File Operations}

\vspace{4pt}\textbf{\color{cyan}>>> EMACS FILE OPS: OPENING THINGS SINCE 1976 <<<}\vspace{2pt}\\
\color{white}emacs: an operating system lacking only a decent editor\\
\vspace{4pt}\textbf{\color{magenta}zpwr emacsall}\\
\color{white}Opens all ZPWR framework files in emacs via fzf.\\
\color{white}Every autoload, script, config -- the whole kingdom.\\
\color{white}\textbackslash{}\$ zpwr emacsall              \# browse all files\\
\color{white}\textbackslash{}\$ zpwr emacsalledit          \# same but editable\\
\vspace{4pt}\textbf{\color{magenta}zpwr emacsrecent}\\
\color{white}Opens recently edited files. Emacs remembers where\\
\color{white}you have been, like a browser history for code.\\
\color{white}\textbackslash{}\$ zpwr emacsrecent            \# fzf recent files\\
\color{white}\textbackslash{}\$ zpwr emacsrecentcd          \# cd to recent dir\\
\color{white}\textbackslash{}\$ zpwr emacsrecentsudo        \# recent with sudo\\
\color{white}\textbackslash{}\$ zpwr emacsrecentsudocd      \# sudo + cd combo\\
\vspace{4pt}\textbf{\color{magenta}zpwr emacsscripts}\\
\color{white}Browse and open ZPWR scripts directory.\\
\color{white}Where the shell scripts live and breathe.\\
\color{white}\textbackslash{}\$ zpwr emacsscripts           \# browse scripts\\
\color{white}\textbackslash{}\$ zpwr emacsscriptedit        \# edit a script\\
\vspace{4pt}\textbf{\color{magenta}zpwr emacsautoload}\\
\color{white}Open autoloaded zsh functions in emacs.\\
\color{white}The beating heart of ZPWR's function library.\\
\color{white}\textbackslash{}\$ zpwr emacsautoload          \# browse autoloads\\
\textit{\color{gray}// protip: emacsrecent is your daily driver}\\
\textit{\color{gray}// it learns your habits and surfaces what matters}\\
\textit{\color{gray}// sudo variants for when /etc needs surgery}\\

\vspace{8pt}

\subsection{Page 53: Emacs Search \& Tags}

\vspace{4pt}\textbf{\color{cyan}>>> SEARCH \& TAGS: FIND ANYTHING, ANYWHERE <<<}\vspace{2pt}\\
\color{white}grep is for amateurs, we have fzf-powered search beams\\
\vspace{4pt}\textbf{\color{magenta}FILE SEARCH (by filename):}\\
\color{white}\textbackslash{}\$ zpwr emacsfilesearch        \# fzf filename search\\
\color{white}\textbackslash{}\$ zpwr emacsfilesearchedit    \# search + edit\\
\color{white}Searches file NAMES across the ZPWR tree.\\
\color{white}Fast, fuzzy, and feeds results into emacs.\\
\vspace{4pt}\textbf{\color{magenta}WORD SEARCH (by content):}\\
\color{white}\textbackslash{}\$ zpwr emacswordsearch        \# search file content\\
\color{white}\textbackslash{}\$ zpwr emacswordsearchedit    \# content search + edit\\
\color{white}Searches INSIDE files for patterns. Like ag/rg but\\
\color{white}with fzf preview and direct-to-emacs piping.\\
\vspace{4pt}\textbf{\color{magenta}FIND \& LOCATE SEARCH:}\\
\color{white}\textbackslash{}\$ zpwr emacsfindsearch        \# find-based search\\
\color{white}\textbackslash{}\$ zpwr emacsfindsearchedit    \# find + edit\\
\color{white}\textbackslash{}\$ zpwr emacslocatesearch      \# locate db search\\
\color{white}\textbackslash{}\$ zpwr emacslocatesearchedit  \# locate + edit\\
\vspace{4pt}\textbf{\color{magenta}TAGS -- CODE NAVIGATION:}\\
\color{white}\textbackslash{}\$ zpwr emacsctags             \# ctags generation\\
\color{white}\textbackslash{}\$ zpwr emacsgtags             \# GNU Global tags\\
\color{white}\textbackslash{}\$ zpwr emacsgtagsedit         \# gtags + jump to def\\
\color{white}Tags let emacs jump to function definitions instantly.\\
\color{white}ctags for basic tagging, gtags for cross-reference.\\
\textit{\color{gray}// the search hierarchy:}\\
\textit{\color{gray}//   wordsearch  = grep inside files (content)}\\
\textit{\color{gray}//   filesearch  = find by filename (path)}\\
\textit{\color{gray}//   locatesearch = system locate db (everywhere)}\\
\textit{\color{gray}//   findsearch  = find command (directory tree)}\\
\textit{\color{gray}//   tags        = jump to definitions (symbols)}\\

\vspace{8pt}

\subsection{Page 54: Emacs Config \& Server}

\vspace{4pt}\textbf{\color{cyan}>>> CONFIG \& SERVER: THE EMACS MOTHERSHIP <<<}\vspace{2pt}\\
\color{white}because starting emacs fresh every time is a war crime\\
\vspace{4pt}\textbf{\color{magenta}zpwr emacsconfig}\\
\color{white}Edit ZPWR's emacs configuration files.\\
\color{white}The bridge between the shell framework and emacs-land.\\
\color{white}\textbackslash{}\$ zpwr emacsconfig            \# edit zpwr emacs cfg\\
\vspace{4pt}\textbf{\color{magenta}zpwr emacsemacsconfig}\\
\color{white}Edit your personal emacs init files (init.el, etc).\\
\color{white}Deeper than emacsconfig -- this is YOUR emacs setup.\\
\color{white}\textbackslash{}\$ zpwr emacsemacsconfig       \# edit \textasciitilde{}/.emacs.d\\
\vspace{4pt}\textbf{\color{magenta}zpwr emacsallserver}\\
\color{white}Opens files using emacsclient, connecting to a running\\
\color{white}emacs daemon. Instant file opens, no startup penalty.\\
\color{white}\textbackslash{}\$ zpwr emacsallserver         \# open via daemon\\
\color{white}HOW THE SERVER WORKFLOW WORKS:\\
\color{white}1. Start emacs daemon:  emacs --daemon\\
\color{white}2. Use emacsallserver to open files via emacsclient\\
\color{white}3. Files open instantly in the running emacs instance\\
\color{white}4. No more 3-second startup tax per file\\
\vspace{4pt}\textbf{\color{magenta}zpwr emacstokens}\\
\color{white}Edit your tokens/secrets file in emacs.\\
\color{white}API keys, env vars, the classified stuff.\\
\color{white}\textbackslash{}\$ zpwr emacstokens            \# edit \textbackslash{}PWR\_LOCAL/.tokens.sh\\
\vspace{4pt}\textbf{\color{magenta}zpwr emacszpwr}\\
\color{white}Open the entire ZPWR directory tree in emacs.\\
\color{white}\textbackslash{}\$ zpwr emacszpwr              \# the whole framework\\
\textit{\color{gray}// emacsclient + daemon = the cheat code}\\
\textit{\color{gray}// startup once, open files forever}\\
\textit{\color{gray}// this is how emacs users actually survive}\\

\vspace{8pt}


\section{CHAPTER 14: TMUX WARFARE}

\subsection{Page 55: Tmux Session Management}

\vspace{4pt}\textbf{\color{cyan}>>> TMUX WARFARE: TERMINAL MULTIPLEXER COMBAT <<<}\vspace{2pt}\\
\color{white}one terminal to rule them all, split them, and in panes bind them\\
\vspace{4pt}\textbf{\color{magenta}zpwr attach}\\
\color{white}Attach to an existing tmux session.\\
\color{white}If no session exists, one will be summoned.\\
\color{white}\textbackslash{}\$ zpwr attach                 \# rejoin the matrix\\
\vspace{4pt}\textbf{\color{magenta}zpwr detach}\\
\color{white}Detach from current tmux session gracefully.\\
\color{white}Your processes keep running in the background.\\
\color{white}\textbackslash{}\$ zpwr detach                 \# leave but don't kill\\
\vspace{4pt}\textbf{\color{magenta}zpwr killmux}\\
\color{white}Kill the tmux server. Nuclear option. Everything dies.\\
\color{white}WARNING: all sessions, windows, and panes -- gone.\\
\color{white}\textbackslash{}\$ zpwr killmux                \# scorched earth\\
\vspace{4pt}\textbf{\color{magenta}zpwr start}\\
\color{white}Start a new tmux session with ZPWR's custom layout.\\
\color{white}Preconfigured panes, windows, and keybindings.\\
\color{white}\textbackslash{}\$ zpwr start                  \# bootstrap tmux\\
\vspace{4pt}\textbf{\color{magenta}zpwr starttabs}\\
\color{white}Start tmux with multiple pre-configured tabs.\\
\color{white}Each tab gets its own purpose -- dev, logs, git, etc.\\
\color{white}\textbackslash{}\$ zpwr starttabs              \# multi-tab setup\\
\textit{\color{gray}// the tmux lifecycle:}\\
\textit{\color{gray}//   start --> attach --> detach --> attach --> killmux}\\
\textit{\color{gray}//   sessions persist through SSH disconnects}\\
\textit{\color{gray}//   your terminal state is immortal}\\

\vspace{8pt}

\subsection{Page 56: Key Duplication -- Multi-Pane Workflow}

\vspace{4pt}\textbf{\color{cyan}>>> KEY DUPLICATION: TYPE ONCE, EXECUTE EVERYWHERE <<<}\vspace{2pt}\\
\color{white}the closest thing to being in two places at once\\
\vspace{4pt}\textbf{\color{magenta}zpwr startsend}\\
\color{white}Send keystrokes to ALL panes in the current window.\\
\color{white}Whatever you type goes to every pane simultaneously.\\
\color{white}\textbackslash{}\$ zpwr startsend              \# enable key broadcast\\
\color{white}USE CASE: You have 4 panes SSHed into 4 servers.\\
\color{white}You need to run the same command on all of them.\\
\color{white}startsend + type once = executed on all 4.\\
\vspace{4pt}\textbf{\color{magenta}zpwr startsendfull}\\
\color{white}Like startsend, but broadcasts across ALL windows\\
\color{white}and ALL panes in the entire tmux session.\\
\color{white}\textbackslash{}\$ zpwr startsendfull          \# global broadcast\\
\color{white}DANGER: every pane in every window receives input.\\
\color{white}Make sure you know what you are doing.\\
\vspace{4pt}\textbf{\color{magenta}zpwr stopsend}\\
\color{white}Turn off key broadcasting. Back to single-pane input.\\
\color{white}\textbackslash{}\$ zpwr stopsend               \# disable broadcast\\
\vspace{4pt}\textbf{\color{magenta}THE MULTI-PANE WORKFLOW:}\\
\color{white}1. \textbackslash{}\$ zpwr start                   \# launch tmux\\
\color{white}2. Split into 4 panes (Ctrl-b \%)\\
\color{white}3. SSH into different servers in each\\
\color{white}4. \textbackslash{}\$ zpwr startsend               \# enable broadcast\\
\color{white}5. Type: sudo apt update \&\& sudo apt upgrade -y\\
\color{white}6. All 4 servers update simultaneously\\
\color{white}7. \textbackslash{}\$ zpwr stopsend                \# back to normal\\
\textit{\color{gray}// sysadmins call this poor mans ansible}\\
\textit{\color{gray}// but honestly it is faster for quick ops}\\

\vspace{8pt}

\subsection{Page 57: Clipboard \& Buffer Operations}

\vspace{4pt}\textbf{\color{cyan}>>> CLIPBOARD OPS: COPY-PASTE ON STEROIDS <<<}\vspace{2pt}\\
\color{white}because Ctrl-C Ctrl-V is so last millennium\\
\vspace{4pt}\textbf{\color{magenta}zpwr pastebuffer}\\
\color{white}Paste from the tmux paste buffer into the terminal.\\
\color{white}Tmux has its own clipboard, separate from the system.\\
\color{white}\textbackslash{}\$ zpwr pastebuffer            \# paste tmux buffer\\
\vspace{4pt}\textbf{\color{magenta}zpwr pastestring}\\
\color{white}Paste a specific string into the active pane.\\
\color{white}Programmatic pasting -- useful in scripts.\\
\color{white}\textbackslash{}\$ zpwr pastestring 'echo hello' \# paste literal text\\
\vspace{4pt}\textbf{\color{magenta}zpwr pastecommand}\\
\color{white}Paste and execute a command in the active pane.\\
\color{white}Like pastestring but hits enter for you.\\
\color{white}\textbackslash{}\$ zpwr pastecommand 'ls -la'  \# paste + execute\\
\vspace{4pt}\textbf{\color{magenta}zpwr google}\\
\color{white}Open a Google search in your default browser.\\
\color{white}Because sometimes you need to leave the terminal.\\
\color{white}\textbackslash{}\$ zpwr google 'zsh completion' \# search the web\\
\vspace{4pt}\textbf{\color{magenta}zpwr openurl}\\
\color{white}Open any URL in the default browser.\\
\textit{\color{gray}\textbackslash{}\$ zpwr openurl https://github.com \# open URL}\\
\textit{\color{gray}// the clipboard pipeline:}\\
\textit{\color{gray}//   terminal output --> tmux buffer --> pastebuffer}\\
\textit{\color{gray}//   system clipboard --> pastestring --> pane}\\
\textit{\color{gray}//   brain --> google --> browser --> back to terminal}\\

\vspace{8pt}

\subsection{Page 58: Autocomplete Integration}

\vspace{4pt}\textbf{\color{cyan}>>> AUTOCOMPLETE: THE TAB KEY'S FINAL FORM <<<}\vspace{2pt}\\
\color{white}zsh-autocomplete + tmux = sentient tab completion\\
\vspace{4pt}\textbf{\color{magenta}zpwr startauto}\\
\color{white}Enable the ZPWR autocomplete integration.\\
\color{white}Completions appear as you type, no Tab needed.\\
\color{white}\textbackslash{}\$ zpwr startauto              \# enable live complete\\
\vspace{4pt}\textbf{\color{magenta}zpwr stopauto}\\
\color{white}Disable autocomplete. Sometimes you want silence.\\
\color{white}\textbackslash{}\$ zpwr stopauto               \# disable live complete\\
\vspace{4pt}\textbf{\color{magenta}HOW IT WORKS WITH TMUX:}\\
\color{white}The autocomplete system hooks into zsh's completion\\
\color{white}engine and renders suggestions in real-time.\\
\color{white}FEATURES:\\
\color{white}* Live suggestions as you type\\
\color{white}* History-based predictions\\
\color{white}* Path completion with preview\\
\color{white}* Command argument completion\\
\color{white}* Works across all tmux panes independently\\
\vspace{4pt}\textbf{\color{magenta}TMUX + AUTOCOMPLETE TIPS:}\\
\color{white}* Each tmux pane runs its own zsh instance\\
\color{white}* Autocomplete state is per-pane\\
\color{white}* startauto/stopauto affects current pane only\\
\color{white}* If completions feel slow, check \textbackslash{}PWR\_COMPS\_CACHE\\
\color{white}\textbackslash{}\$ zpwr stopauto               \# when it gets noisy\\
\color{white}\textbackslash{}\$ zpwr startauto              \# when you miss it\\
\textit{\color{gray}// autocomplete is like a copilot for your shell}\\
\textit{\color{gray}// it watches, it learns, it suggests}\\
\textit{\color{gray}// toggle it off when you need to think in silence}\\

\vspace{8pt}

\subsection{Page 59: Process Monitoring}

\vspace{4pt}\textbf{\color{cyan}>>> PROCESS MONITORING: WATCHING THE WATCHERS <<<}\vspace{2pt}\\
\color{white}every PID tells a story, most of them tragic\\
\vspace{4pt}\textbf{\color{magenta}zpwr pstreemonitor}\\
\color{white}Displays a live process tree. Watch your processes\\
\color{white}branch, fork, and occasionally die unexpectedly.\\
\color{white}\textbackslash{}\$ zpwr pstreemonitor          \# live process tree\\
\color{white}Uses pstree under the hood with continuous refresh.\\
\color{white}Like htop's tree view but from the ZPWR command line.\\
\vspace{4pt}\textbf{\color{magenta}zpwr tty}\\
\color{white}Show info about the current terminal/tty.\\
\color{white}\textbackslash{}\$ zpwr tty                    \# current tty info\\
\vspace{4pt}\textbf{\color{magenta}TMUX + MONITORING WORKFLOW:}\\
\color{white}The power move: dedicate a tmux pane to monitoring.\\
\color{white}LAYOUT EXAMPLE:\\
\color{white}+-------------------+--------------------+\\
\color{white}|                   |                    |\\
\color{white}|   your work       |  pstreemonitor     |\\
\color{white}|                   |                    |\\
\color{white}+-------------------+--------------------+\\
\color{white}|         zpwr taillog                   |\\
\color{white}+-----------------------------------------+\\
\vspace{4pt}\textbf{\color{magenta}PROCESS HUNTING:}\\
\color{white}ZPWR integrates fzf with ps and lsof for killing.\\
\color{white}Ctrl-Y Ctrl-K                  \# fzf kill by ps\\
\color{white}Ctrl-Y Ctrl-L                  \# fzf kill by lsof\\
\textit{\color{gray}// pstreemonitor in a side pane = instant situational awareness}\\
\textit{\color{gray}// you will see zombie processes before they see you}\\

\vspace{8pt}


\section{CHAPTER 15: NETWORK \& CONTAINERS}

\subsection{Page 60: Network Recon}

\vspace{4pt}\textbf{\color{cyan}>>> NETWORK RECON: SCANNING THE WIRE <<<}\vspace{2pt}\\
\color{white}packets don't lie, but DNS sometimes does\\
\vspace{4pt}\textbf{\color{magenta}zpwr digs}\\
\color{white}Enhanced DNS lookup. Wraps dig with cleaner output.\\
\color{white}When you need to know where a domain really points.\\
\color{white}\textbackslash{}\$ zpwr digs github.com         \# DNS lookup\\
\color{white}\textbackslash{}\$ zpwr digs example.com        \# check any domain\\
\color{white}Shows A records, CNAME chains, TTLs -- the works.\\
\color{white}Faster than remembering dig +short +noall +answer.\\
\vspace{4pt}\textbf{\color{magenta}zpwr pi / zpwr ping}\\
\color{white}Ping a host. Simple, reliable, essential.\\
\color{white}'pi' is the short alias because typing matters.\\
\color{white}\textbackslash{}\$ zpwr pi google.com           \# quick ping\\
\color{white}\textbackslash{}\$ zpwr ping 8.8.8.8           \# verbose ping\\
\vspace{4pt}\textbf{\color{magenta}zpwr pir}\\
\color{white}Ping with extended options. More detailed network\\
\color{white}diagnostics for when simple ping is not enough.\\
\color{white}\textbackslash{}\$ zpwr pir 192.168.1.1        \# extended ping\\
\vspace{4pt}\textbf{\color{magenta}RECON WORKFLOW:}\\
\color{white}1. \textbackslash{}\$ zpwr digs suspect-domain.com  \# where does it resolve?\\
\color{white}2. \textbackslash{}\$ zpwr pi <ip-from-dig>         \# is it alive?\\
\color{white}3. \textbackslash{}\$ zpwr pir <ip>                 \# deeper probe\\
\textit{\color{gray}// three verbs, complete network recon}\\
\textit{\color{gray}// digs for DNS, pi for heartbeat, pir for details}\\
\textit{\color{gray}// when in doubt, ping first -- dead hosts tell no tales}\\

\vspace{8pt}

\subsection{Page 61: Tor Integration}

\vspace{4pt}\textbf{\color{cyan}>>> TOR INTEGRATION: ANONYMOUS MODE ENGAGED <<<}\vspace{2pt}\\
\color{white}the onion router: because privacy is not a crime\\
\vspace{4pt}\textbf{\color{magenta}zpwr torip}\\
\color{white}Check your current Tor exit node IP address.\\
\color{white}Confirms you're routing through the onion network.\\
\color{white}\textbackslash{}\$ zpwr torip                  \# show tor exit IP\\
\color{white}If Tor is running, this shows the exit node's IP.\\
\color{white}If not, it'll tell you Tor isn't connected.\\
\vspace{4pt}\textbf{\color{magenta}zpwr toriprenew}\\
\color{white}Request a new Tor circuit. Get a fresh exit node\\
\color{white}and a new apparent IP address.\\
\color{white}\textbackslash{}\$ zpwr toriprenew             \# new identity\\
\color{white}BEFORE YOU START:\\
\color{white}Tor must be installed and running as a service.\\
\color{white}macOS:   brew install tor \&\& brew services start tor\\
\color{white}Linux:   sudo apt install tor \&\& sudo systemctl start tor\\
\vspace{4pt}\textbf{\color{magenta}THE ANONYMOUS WORKFLOW:}\\
\color{white}1. Start tor service\\
\color{white}2. \textbackslash{}\$ zpwr torip            \# confirm you're connected\\
\color{white}3. Do your thing\\
\color{white}4. \textbackslash{}\$ zpwr toriprenew       \# rotate identity\\
\color{white}5. \textbackslash{}\$ zpwr torip            \# verify new IP\\
\color{white}DISCLAIMER:\\
\color{white}Tor provides anonymity, not invincibility.\\
\color{white}DNS leaks, browser fingerprinting, and OPSEC\\
\color{white}failures can still expose you. Use responsibly.\\
\textit{\color{gray}// \textbackslash{}"I am become anonymous, destroyer of logs\textbackslash{}"}\\
\textit{\color{gray}//   -- J. Robert Oppenheimer, probably not}\\

\vspace{8pt}

\subsection{Page 62: Docker Management}

\vspace{4pt}\textbf{\color{cyan}>>> DOCKER MANAGEMENT: CONTAINER WARFARE <<<}\vspace{2pt}\\
\color{white}it works on my machine --> we ship the machine\\
\vspace{4pt}\textbf{\color{magenta}zpwr dockerwipe}\\
\color{white}The nuclear option for Docker. Removes EVERYTHING:\\
\color{white}containers, images, volumes, networks -- all of it.\\
\color{white}\textbackslash{}\$ zpwr dockerwipe             \# purge all docker\\
\color{white}WARNING: This is not gentle. It will:\\
\color{white}- Stop all running containers\\
\color{white}- Remove all containers\\
\color{white}- Remove all images\\
\color{white}- Remove all volumes\\
\color{white}- Prune everything\\
\color{white}WHEN TO USE:\\
\color{white}* Disk space emergency (docker images love eating disk)\\
\color{white}* Fresh start after a botched experiment\\
\color{white}* When 'docker system df' shows 50GB of shame\\
\vspace{4pt}\textbf{\color{magenta}zpwr dfimage}\\
\color{white}Reverse-engineer a Docker image to show its Dockerfile.\\
\color{white}See how an image was built without the source.\\
\color{white}\textbackslash{}\$ zpwr dfimage nginx:latest   \# reverse Dockerfile\\
\color{white}\textbackslash{}\$ zpwr dfimage myapp:v2       \# inspect any image\\
\color{white}Useful for auditing third-party images.\\
\color{white}Trust but verify what's in that container.\\
\vspace{4pt}\textbf{\color{magenta}DOCKER TRIAGE:}\\
\color{white}1. \textbackslash{}\$ docker system df              \# how bad is it?\\
\color{white}2. \textbackslash{}\$ zpwr dfimage suspect:latest   \# what's inside?\\
\color{white}3. \textbackslash{}\$ zpwr dockerwipe               \# burn it all down\\
\color{white}4. \textbackslash{}\$ docker system df              \# ahh, fresh disk\\
\textit{\color{gray}// dockerwipe: for when docker system prune isn't angry enough}\\
\textit{\color{gray}// dfimage: for when you don't trust the README}\\

\vspace{8pt}

\subsection{Page 63: Clone \& GitHub Workflow}

\vspace{4pt}\textbf{\color{cyan}>>> CLONE \& GITHUB: THE FULL GIT WORKFLOW <<<}\vspace{2pt}\\
\color{white}git clone is just the beginning of the relationship\\
\vspace{4pt}\textbf{\color{magenta}zpwr clone}\\
\color{white}Clone a git repository with ZPWR's enhanced setup.\\
\color{white}Handles auth, directory placement, and post-clone hooks.\\
\color{white}\textbackslash{}\$ zpwr clone user/repo         \# clone from GitHub\\
\textit{\color{gray}\textbackslash{}\$ zpwr clone https://...       \# clone any URL}\\
\vspace{4pt}\textbf{\color{magenta}zpwr reveal}\\
\color{white}Open the current git repo's GitHub page in browser.\\
\color{white}Instant context switch from terminal to web.\\
\color{white}\textbackslash{}\$ zpwr reveal                 \# open repo in browser\\
\vspace{4pt}\textbf{\color{magenta}zpwr revealrecurse}\\
\color{white}Open GitHub pages for the repo AND its submodules.\\
\color{white}When you need the full picture of a multi-repo project.\\
\color{white}\textbackslash{}\$ zpwr revealrecurse          \# open all repos\\
\vspace{4pt}\textbf{\color{magenta}THE GITHUB WORKFLOW:}\\
\color{white}1. \textbackslash{}\$ zpwr clone user/repo       \# get the code\\
\color{white}2. \textbackslash{}\$ cd repo \&\& hack hack hack\\
\color{white}3. \textbackslash{}\$ zpwr reveal                \# check CI/issues\\
\color{white}4. \textbackslash{}\$ git push                   \# ship it\\
\vspace{4pt}\textbf{\color{magenta}RELATED TOOLS:}\\
\color{white}\textbackslash{}\$ zpwr killremote             \# kill remote sessions\\
\color{white}killremote terminates remote SSH/tmux sessions that\\
\color{white}have gone stale. Cleanup for the disciplined operator.\\
\textit{\color{gray}// reveal is the bridge between terminal and browser}\\
\textit{\color{gray}// faster than typing github.com/user/repo}\\
\textit{\color{gray}// your fingers will thank you}\\

\vspace{8pt}


\section{CHAPTER 16: THE LOGGING SYSTEM}

\subsection{Page 64: Log Levels}

\vspace{4pt}\textbf{\color{cyan}>>> LOG LEVELS: SIGNAL VS NOISE <<<}\vspace{2pt}\\
\color{white}if a tree falls in the forest and nobody tails the log...\\
\color{white}ZPWR has a full structured logging system with four levels.\\
\color{white}Each level has a file variant and a console variant.\\
\vspace{4pt}\textbf{\color{magenta}LEVEL 1: TRACE}\\
\color{white}\textbackslash{}\$ zpwr logtrace 'entering func'  \# most verbose\\
\color{white}The firehose. Every function entry, variable assignment,\\
\color{white}and internal state change. Use for deep debugging only.\\
\vspace{4pt}\textbf{\color{magenta}LEVEL 2: DEBUG}\\
\color{white}\textbackslash{}\$ zpwr logdebug 'cache miss'    \# developer detail\\
\color{white}Detailed info for development. Cache behavior,\\
\color{white}branch decisions, intermediate values.\\
\vspace{4pt}\textbf{\color{magenta}LEVEL 3: INFO}\\
\color{white}\textbackslash{}\$ zpwr loginfo 'update done'    \# normal operations\\
\color{white}Standard operational messages. \textbackslash{}"Started\textbackslash{}", \textbackslash{}"Completed\textbackslash{}",\\
\color{white}\textbackslash{}"Loaded 47 plugins\textbackslash{}". The daily heartbeat.\\
\vspace{4pt}\textbf{\color{magenta}LEVEL 4: ERROR}\\
\color{white}\textbackslash{}\$ zpwr logerror 'file not found' \# something broke\\
\color{white}Things that should not have happened but did.\\
\color{white}Missing files, failed commands, broken assumptions.\\
\vspace{4pt}\textbf{\color{magenta}SEVERITY PYRAMID:}\\

\vspace{8pt}

\subsection{Page 65: Console vs File Logging}

\vspace{4pt}\textbf{\color{cyan}>>> CONSOLE VS FILE: TWO DESTINATIONS <<<}\vspace{2pt}\\
\color{white}to stderr or to disk, that is the question\\
\color{white}Every log verb comes in two flavors:\\
\vspace{4pt}\textbf{\color{magenta}FILE VARIANTS (write to \textbackslash{}PWR\_LOGFILE):}\\
\color{white}\textbackslash{}\$ zpwr logtrace 'msg'          \# --> logfile\\
\color{white}\textbackslash{}\$ zpwr logdebug 'msg'          \# --> logfile\\
\color{white}\textbackslash{}\$ zpwr loginfo 'msg'           \# --> logfile\\
\color{white}\textbackslash{}\$ zpwr logerror 'msg'          \# --> logfile\\
\color{white}These write silently to the log file. No terminal\\
\color{white}output. Perfect for background operations and scripts.\\
\vspace{4pt}\textbf{\color{magenta}CONSOLE VARIANTS (write to stderr AND logfile):}\\
\color{white}\textbackslash{}\$ zpwr logtraceconsole 'msg'   \# --> stderr + file\\
\color{white}\textbackslash{}\$ zpwr logdebugconsole 'msg'   \# --> stderr + file\\
\color{white}\textbackslash{}\$ zpwr loginfoconsole 'msg'    \# --> stderr + file\\
\color{white}\textbackslash{}\$ zpwr logerrorconsole 'msg'   \# --> stderr + file\\
\color{white}These print to your terminal AND write to the log.\\
\color{white}For interactive debugging when you need to see it live.\\
\vspace{4pt}\textbf{\color{magenta}THE LOGFILE:}\\
\texttt{\color{green}\textbackslash{}PWR\_LOGFILE = \textbackslash{}PWR\_LOCAL/zpwr.log}\\
\color{white}\textbackslash{}\$ cat \textbackslash{}PWR\_LOGFILE | tail -20   \# recent entries\\
\color{white}\textbackslash{}\$ wc -l \textbackslash{}PWR\_LOGFILE           \# how big is it?\\
\vspace{4pt}\textbf{\color{magenta}CHOOSING YOUR VARIANT:}\\
\color{white}SCENARIO                          USE\\
\color{white}Background script running          loginfo\\
\color{white}Debugging interactively             logdebugconsole\\
\color{white}Tracking a rare bug                 logtraceconsole\\
\color{white}Production error handling           logerror\\
\textit{\color{gray}// file = silent, persistent, searchable}\\
\textit{\color{gray}// console = loud, immediate, ephemeral}\\
\textit{\color{gray}// use both strategically}\\

\vspace{8pt}

\subsection{Page 66: Reading Logs}

\vspace{4pt}\textbf{\color{cyan}>>> READING LOGS: THE ART OF LOG DIVING <<<}\vspace{2pt}\\
\color{white}somewhere in that wall of text is the answer you seek\\
\vspace{4pt}\textbf{\color{magenta}zpwr taillog}\\
\color{white}Live tail of the ZPWR log file with ccze colorization.\\
\color{white}Watch events stream by in glorious technicolor.\\
\color{white}\textbackslash{}\$ zpwr taillog                \# live colorized tail\\
\color{white}ccze automatically highlights timestamps, levels, paths,\\
\color{white}and errors. Makes raw logs actually readable.\\
\vspace{4pt}\textbf{\color{magenta}WHAT TO LOOK FOR:}\\
\color{white}ERROR  lines = something broke, investigate immediately\\
\color{white}WARN   lines = potential issue, worth noting\\
\color{white}INFO   lines = normal operations, sanity check\\
\color{white}DEBUG  lines = details for troubleshooting\\
\color{white}TRACE  lines = deep internals, very verbose\\
\vspace{4pt}\textbf{\color{magenta}FILTERING TECHNIQUES:}\\
\color{white}\textbackslash{}\$ grep ERROR \textbackslash{}PWR\_LOGFILE     \# errors only\\
\color{white}\textbackslash{}\$ grep -c ERROR \textbackslash{}PWR\_LOGFILE  \# count errors\\
\color{white}\textbackslash{}\$ tail -100 \textbackslash{}PWR\_LOGFILE      \# last 100 lines\\
\vspace{4pt}\textbf{\color{magenta}TMUX LOG MONITORING SETUP:}\\
\color{white}Dedicate a tmux pane to zpwr taillog:\\
\color{white}1. Split pane: Ctrl-b \%\\
\color{white}2. In new pane: zpwr taillog\\
\color{white}3. Work in other pane, watch logs stream\\
\color{white}4. Errors glow red thanks to ccze\\
\vspace{4pt}\textbf{\color{magenta}LOG MAINTENANCE:}\\
\color{white}Logs grow forever if you let them. ZPWR handles\\
\color{white}rotation, but you can manually truncate:\\
\color{white}\textbackslash{}\$ : > \textbackslash{}PWR\_LOGFILE            \# truncate to zero\\
\textit{\color{gray}// taillog + ccze = the cyberpunk log experience}\\
\textit{\color{gray}// errors in red, info in green, timestamps in cyan}\\

\vspace{8pt}

\subsection{Page 67: Debug Mode}

\vspace{4pt}\textbf{\color{cyan}>>> DEBUG MODE: WHEN THINGS GO SIDEWAYS <<<}\vspace{2pt}\\
\color{white}enabling verbose output is how you stop guessing\\
\vspace{4pt}\textbf{\color{magenta}ZPWR\_DEBUG}\\
\color{white}Set this env var to enable debug-level logging.\\
\color{white}More detail in the log, more signal for your hunt.\\
\color{white}\textbackslash{}\$ export ZPWR\_DEBUG=true       \# enable debug\\
\color{white}\textbackslash{}\$ zpwr reload                  \# reload with debug\\
\color{white}\textbackslash{}\$ zpwr taillog                 \# watch debug output\\
\vspace{4pt}\textbf{\color{magenta}ZPWR\_TRACE}\\
\color{white}Maximum verbosity. Every function entry and exit,\\
\color{white}every variable, every decision point.\\
\color{white}\textbackslash{}\$ export ZPWR\_TRACE=true       \# enable trace\\
\color{white}\textbackslash{}\$ zpwr reload                  \# reload with trace\\
\color{white}WARNING: Trace mode generates MASSIVE log output.\\
\color{white}Your logfile will grow fast. Use sparingly.\\
\vspace{4pt}\textbf{\color{magenta}DEBUGGING WORKFLOW:}\\
\color{white}1. Reproduce the problem once to confirm it\\
\color{white}2. \textbackslash{}\$ export ZPWR\_DEBUG=true\\
\color{white}3. \textbackslash{}\$ zpwr reload\\
\color{white}4. Reproduce the problem again\\
\color{white}5. \textbackslash{}\$ zpwr taillog              \# find the smoking gun\\
\color{white}6. Still lost? Escalate:\\
\color{white}7. \textbackslash{}\$ export ZPWR\_TRACE=true\\
\color{white}8. \textbackslash{}\$ zpwr reload \&\& reproduce\\
\vspace{4pt}\textbf{\color{magenta}DISABLING DEBUG:}\\
\color{white}\textbackslash{}\$ unset ZPWR\_DEBUG ZPWR\_TRACE  \# back to normal\\
\color{white}\textbackslash{}\$ zpwr reload                  \# reload clean\\
\textit{\color{gray}// debug mode: \textbackslash{}"trust but verify\textbackslash{}" in action}\\
\textit{\color{gray}// trace mode: \textbackslash{}"verify everything, trust nothing\textbackslash{}"}\\
\textit{\color{gray}// always disable after debugging or your log eats disk}\\

\vspace{8pt}


\section{CHAPTER 17: INTROSPECTION}

\subsection{Page 68: Counting Everything}

\vspace{4pt}\textbf{\color{cyan}>>> COUNTING EVERYTHING: THE CENSUS BUREAU <<<}\vspace{2pt}\\
\color{white}if you can't measure it, you can't brag about it\\
\color{white}ZPWR tracks the size of every component. These verbs\\
\color{white}tell you exactly how big your installation is.\\
\vspace{4pt}\textbf{\color{magenta}zpwr autoloadcount}\\
\color{white}Count of autoloaded zsh functions in the framework.\\
\color{white}\textbackslash{}\$ zpwr autoloadcount           \# how many functions?\\
\vspace{4pt}\textbf{\color{magenta}zpwr scriptcount}\\
\color{white}Count of shell scripts in \textbackslash{}PWR/scripts.\\
\color{white}\textbackslash{}\$ zpwr scriptcount             \# how many scripts?\\
\vspace{4pt}\textbf{\color{magenta}zpwr zshplugincount}\\
\color{white}Count of installed zsh plugins (zinit, oh-my-zsh, etc).\\
\color{white}\textbackslash{}\$ zpwr zshplugincount          \# zsh plugins loaded\\
\vspace{4pt}\textbf{\color{magenta}zpwr vimplugincount}\\
\color{white}Count of installed vim/neovim plugins.\\
\color{white}\textbackslash{}\$ zpwr vimplugincount          \# vim plugins loaded\\
\vspace{4pt}\textbf{\color{magenta}zpwr emacsplugincount}\\
\color{white}Count of installed emacs packages.\\
\color{white}\textbackslash{}\$ zpwr emacsplugincount        \# emacs pkgs loaded\\
\vspace{4pt}\textbf{\color{magenta}THE FULL CENSUS:}\\
\color{white}\textbackslash{}\$ zpwr autoloadcount \&\& zpwr scriptcount \&\& \textbackslash{}\textbackslash{}\\
\color{white}zpwr zshplugincount \&\& zpwr vimplugincount \&\& \textbackslash{}\textbackslash{}\\
\color{white}zpwr emacsplugincount\\
\color{white}Run all five to get the full picture.\\
\color{white}ZPWR typically ships with 300+ autoloads, 50+ scripts,\\
\color{white}and dozens of plugins across all three editors.\\
\textit{\color{gray}// these counts tell you the scope of ZPWR}\\
\textit{\color{gray}// hundreds of functions working in concert}\\
\textit{\color{gray}// like a small city running on your laptop}\\

\vspace{8pt}

\subsection{Page 69: Listing Everything}

\vspace{4pt}\textbf{\color{cyan}>>> LISTING EVERYTHING: THE INVENTORY MANIFEST <<<}\vspace{2pt}\\
\color{white}counting is nice, but names are power\\
\color{white}Counts tell you how many. Lists tell you WHAT.\\
\color{white}Every count verb has a list counterpart.\\
\vspace{4pt}\textbf{\color{magenta}zpwr autoloadlist}\\
\color{white}List every autoloaded function by name.\\
\color{white}\textbackslash{}\$ zpwr autoloadlist            \# all function names\\
\color{white}\textbackslash{}\$ zpwr autoloadlist | grep fzf \# find fzf funcs\\
\vspace{4pt}\textbf{\color{magenta}zpwr scriptlist}\\
\color{white}List every script in the scripts directory.\\
\color{white}\textbackslash{}\$ zpwr scriptlist              \# all scripts\\
\vspace{4pt}\textbf{\color{magenta}zpwr zshpluginlist}\\
\color{white}List installed zsh plugins with their sources.\\
\color{white}\textbackslash{}\$ zpwr zshpluginlist           \# all zsh plugins\\
\vspace{4pt}\textbf{\color{magenta}zpwr vimpluginlist}\\
\color{white}List installed vim/neovim plugins.\\
\color{white}\textbackslash{}\$ zpwr vimpluginlist           \# all vim plugins\\
\vspace{4pt}\textbf{\color{magenta}zpwr emacspluginlist}\\
\color{white}List installed emacs packages.\\
\color{white}\textbackslash{}\$ zpwr emacspluginlist         \# all emacs packages\\
\vspace{4pt}\textbf{\color{magenta}POWER COMBOS:}\\
\color{white}\textbackslash{}\$ zpwr autoloadlist | wc -l    \# quick count\\
\color{white}\textbackslash{}\$ zpwr autoloadlist | fzf      \# fuzzy browse\\
\color{white}\textbackslash{}\$ zpwr scriptlist | sort       \# alphabetical\\
\textit{\color{gray}// lists are grep-able, fzf-able, pipe-able}\\
\textit{\color{gray}// the raw material for your own exploration}\\
\textit{\color{gray}// pipe into fzf for interactive discovery}\\

\vspace{8pt}

\subsection{Page 70: History Archaeology}

\vspace{4pt}\textbf{\color{cyan}>>> HISTORY ARCHAEOLOGY: DIGGING UP THE PAST <<<}\vspace{2pt}\\
\color{white}your terminal remembers everything, even the mistakes\\
\vspace{4pt}\textbf{\color{magenta}zpwr hist}\\
\color{white}Browse command history with fzf. Search, preview,\\
\color{white}and re-execute past commands interactively.\\
\color{white}\textbackslash{}\$ zpwr hist                   \# fzf history browser\\
\vspace{4pt}\textbf{\color{magenta}zpwr histedit}\\
\color{white}Edit the history file directly. Surgically remove\\
\color{white}embarrassing commands or fix typos in saved history.\\
\color{white}\textbackslash{}\$ zpwr histedit               \# edit history file\\
\color{white}Use with care. History surgery is delicate work.\\
\vspace{4pt}\textbf{\color{magenta}zpwr histfile}\\
\color{white}Show the path to the current history file.\\
\color{white}Useful when you're not sure which histfile is active.\\
\color{white}\textbackslash{}\$ zpwr histfile               \# show histfile path\\
\vspace{4pt}\textbf{\color{magenta}zpwr catviminfo}\\
\color{white}Display vim's recent file history from viminfo.\\
\color{white}What files has vim touched recently?\\
\color{white}\textbackslash{}\$ zpwr catviminfo             \# vim file history\\
\vspace{4pt}\textbf{\color{magenta}zpwr catrecentfviminfo / catnviminforecentf}\\
\color{white}Cross-reference vim/nvim recent files with the\\
\color{white}recentf list. Where vim and shell history intersect.\\
\color{white}\textbackslash{}\$ zpwr catrecentfviminfo      \# vim+recentf cross\\
\color{white}\textbackslash{}\$ zpwr catnviminforecentf     \# nvim+recentf cross\\
\textit{\color{gray}// your history is a goldmine of workflow patterns}\\
\textit{\color{gray}// hist + fzf = instant command recall}\\
\textit{\color{gray}// histedit = revisionist history (use responsibly)}\\

\vspace{8pt}

\subsection{Page 71: Verb Discovery}

\vspace{4pt}\textbf{\color{cyan}>>> VERB DISCOVERY: FINDING WHAT'S POSSIBLE <<<}\vspace{2pt}\\
\color{white}you can't use what you don't know exists\\
\vspace{4pt}\textbf{\color{magenta}zpwr subcommands}\\
\color{white}List all zpwr subcommands (verbs). The master index\\
\color{white}of everything zpwr can do.\\
\color{white}\textbackslash{}\$ zpwr subcommands             \# all subcommands\\
\vspace{4pt}\textbf{\color{magenta}zpwr subcommandscount}\\
\color{white}How many subcommands exist? The number keeps growing.\\
\color{white}\textbackslash{}\$ zpwr subcommandscount        \# total count\\
\vspace{4pt}\textbf{\color{magenta}zpwr verbs (alias: subcommandsfzf)}\\
\color{white}Browse subcommands interactively with fzf.\\
\color{white}The best way to explore when you're not sure what to try.\\
\color{white}\textbackslash{}\$ zpwr verbs         \# fzf subcommands\\
\vspace{4pt}\textbf{\color{magenta}zpwr verbslist / verbsfzf / verbscount}\\
\color{white}Alternative verb-centric views of available commands.\\
\color{white}\textbackslash{}\$ zpwr verbslist              \# list all verbs\\
\color{white}\textbackslash{}\$ zpwr verbs               \# fzf browse verbs\\
\color{white}\textbackslash{}\$ zpwr verbscount             \# count verbs\\
\vspace{4pt}\textbf{\color{magenta}zpwr wizard / manual / tutorial / docs}\\
\color{white}You're reading it right now. The interactive encyclopedia\\
\color{white}with vim-style navigation. Very meta. All four names work.\\
\color{white}\textbackslash{}\$ zpwr wizard                 \# start the wizard\\
\color{white}\textbackslash{}\$ zpwr wizard --toc           \# table of contents\\
\color{white}\textbackslash{}\$ zpwr wizard --search docker  \# search pages\\
\color{white}\textbackslash{}\$ zpwr wizard 42              \# jump to page 42\\
\textit{\color{gray}// verbsfzf is the fastest way to discover (alias: subcommandsfzf)}\\
\textit{\color{gray}// type a few letters, fzf narrows it down}\\
\textit{\color{gray}// when in doubt, fzf it out}\\

\vspace{8pt}

\subsection{Page 72: Shell State Inspection}

\vspace{4pt}\textbf{\color{cyan}>>> SHELL STATE: KNOW THY ZSH <<<}\vspace{2pt}\\
\color{white}introspection: when the shell looks in the mirror\\
\vspace{4pt}\textbf{\color{magenta}zpwr exists}\\
\color{white}Check if a ZPWR component exists (functions, files, etc).\\
\color{white}\textbackslash{}\$ zpwr exists zpwrReload       \# does this exist?\\
\vspace{4pt}\textbf{\color{magenta}zpwr existscommand}\\
\color{white}Check if a command exists in PATH.\\
\color{white}Cleaner than 'which' or 'command -v' in scripts.\\
\color{white}\textbackslash{}\$ zpwr existscommand git      \# is git installed?\\
\color{white}\textbackslash{}\$ zpwr existscommand rg       \# is ripgrep here?\\
\vspace{4pt}\textbf{\color{magenta}zpwr reload}\\
\color{white}Reload the entire ZPWR framework without restarting.\\
\color{white}Re-sources configs, rebuilds caches, refreshes state.\\
\color{white}\textbackslash{}\$ zpwr reload                 \# hot reload everything\\
\vspace{4pt}\textbf{\color{magenta}ZSH INSPECTION TOOLS:}\\
\color{white}These are built into zsh, but ZPWR makes them accessible:\\
\color{white}\textbackslash{}\$ setopt                      \# list active options\\
\color{white}\textbackslash{}\$ zmodload                    \# list loaded modules\\
\color{white}\textbackslash{}\$ zstyle -L                   \# list all zstyles\\
\vspace{4pt}\textbf{\color{magenta}zpwr copycommand}\\
\color{white}Copy a command's path or definition to clipboard.\\
\color{white}\textbackslash{}\$ zpwr copycommand ls         \# copy ls path\\
\vspace{4pt}\textbf{\color{magenta}zpwr opencommand}\\
\color{white}Open a command's source file in your editor.\\
\color{white}Works for functions, scripts, and binaries.\\
\color{white}\textbackslash{}\$ zpwr opencommand zpwrReload  \# view the source\\
\textit{\color{gray}// existscommand before you call, reload after you change}\\
\textit{\color{gray}// opencommand when you need to understand the internals}\\

\vspace{8pt}


\section{CHAPTER 18: INSTALLATION \& RESET}

\subsection{Page 73: Config Sessions}

\vspace{4pt}\textbf{\color{cyan}>>> CONFIG SESSIONS: EDIT EVERYTHING FROM HERE <<<}\vspace{2pt}\\
\color{white}one verb to open any config, no path memorization\\
\color{white}ZPWR provides shortcut verbs for every config file\\
\color{white}you'll ever need to edit. No more 'vim \textasciitilde{}/.zshrc'.\\
\vspace{4pt}\textbf{\color{magenta}zpwr brc  -- Bash RC}\\
\color{white}Edit your .bashrc. For when bash compatibility matters.\\
\color{white}\textbackslash{}\$ zpwr brc                    \# edit .bashrc\\
\vspace{4pt}\textbf{\color{magenta}zpwr zrc  -- Zsh RC}\\
\color{white}Edit your .zshrc. The most-touched config in ZPWR.\\
\color{white}\textbackslash{}\$ zpwr zrc                    \# edit .zshrc\\
\vspace{4pt}\textbf{\color{magenta}zpwr vrc  -- Vim RC}\\
\color{white}Edit your .vimrc or init.vim/init.lua.\\
\color{white}\textbackslash{}\$ zpwr vrc                    \# edit vim config\\
\vspace{4pt}\textbf{\color{magenta}zpwr trc  -- Tmux RC}\\
\color{white}Edit your .tmux.conf. Keybindings, statusbar, the works.\\
\color{white}\textbackslash{}\$ zpwr trc                    \# edit tmux config\\
\vspace{4pt}\textbf{\color{magenta}zpwr tokens  -- Secrets File}\\
\color{white}Edit \textbackslash{}PWR\_LOCAL/.tokens.sh. API keys, secrets, env vars.\\
\color{white}\textbackslash{}\$ zpwr tokens                 \# edit tokens/secrets\\
\color{white}\textbackslash{}\$ zpwr vimtokens              \# same, force vim\\
\color{white}\textbackslash{}\$ zpwr emacstokens            \# same, force emacs\\
\vspace{4pt}\textbf{\color{magenta}zpwr zcompdump  -- Completion Cache}\\
\color{white}Inspect or regenerate the zsh completion dump file.\\
\color{white}When completions act weird, this is where you look.\\
\color{white}\textbackslash{}\$ zpwr zcompdump              \# inspect comp cache\\
\textit{\color{gray}// brc, zrc, vrc, trc -- four letters, four configs}\\
\textit{\color{gray}// tokens for secrets, zcompdump for completion fixes}\\

\vspace{8pt}

\subsection{Page 74: Install \& Uninstall}

\vspace{4pt}\textbf{\color{cyan}>>> INSTALL \& UNINSTALL: SETUP AND TEARDOWN <<<}\vspace{2pt}\\
\color{white}easy to install, clean to remove -- no malware energy\\
\vspace{4pt}\textbf{\color{magenta}zpwr install}\\
\color{white}Run the full ZPWR installation process.\\
\color{white}Installs dependencies, links configs, sets up plugins.\\
\color{white}\textbackslash{}\$ zpwr install                \# full installation\\
\color{white}WHAT IT DOES:\\
\color{white}1. Installs system dependencies (brew/apt)\\
\color{white}2. Links config files to their proper locations\\
\color{white}3. Installs zsh, vim, and emacs plugins\\
\color{white}4. Generates caches and completion dumps\\
\color{white}5. Sets up tmux configuration\\
\vspace{4pt}\textbf{\color{magenta}zpwr uninstall}\\
\color{white}Remove ZPWR configuration. Restores your shell to\\
\color{white}its pre-ZPWR state. Clean breakup, no hard feelings.\\
\color{white}\textbackslash{}\$ zpwr uninstall              \# remove zpwr config\\
\vspace{4pt}\textbf{\color{magenta}CONFIG LINK MANAGEMENT:}\\
\color{white}ZPWR uses symlinks to place configs where tools expect them.\\
\color{white}How it works:\\
\color{white}\textbackslash{}PWR/configs/.zshrc  -->  \textasciitilde{}/.zshrc     (symlink)\\
\color{white}\textbackslash{}PWR/configs/.vimrc  -->  \textasciitilde{}/.vimrc     (symlink)\\
\color{white}\textbackslash{}PWR/configs/.tmux   -->  \textasciitilde{}/.tmux.conf (symlink)\\
\color{white}The source of truth lives in \textbackslash{}PWR, the symlinks\\
\color{white}just point there. Version controlled, portable, clean.\\
\vspace{4pt}\textbf{\color{magenta}PARTIAL OPERATIONS:}\\
\color{white}Sometimes you don't want the full install/uninstall.\\
\color{white}You just need to fix the symlinks:\\
\color{white}regenconfiglinks = recreate all config symlinks\\
\color{white}rmconfiglinks    = remove symlinks only\\
\textit{\color{gray}// install is idempotent -- run it again when things drift}\\
\textit{\color{gray}// uninstall is clean -- your system goes back to stock}\\

\vspace{8pt}

\subsection{Page 75: Reset \& Recovery}

\vspace{4pt}\textbf{\color{cyan}>>> RESET \& RECOVERY: WHEN THINGS GO WRONG <<<}\vspace{2pt}\\
\color{white}every system needs a panic button -- here are yours\\
\vspace{4pt}\textbf{\color{magenta}zpwr reset}\\
\color{white}Full framework reset. Clears caches, regenerates\\
\color{white}configs, rebuilds completion dumps, relinks everything.\\
\color{white}\textbackslash{}\$ zpwr reset                  \# the big red button\\
\color{white}WHEN TO USE:\\
\color{white}* Completions are broken or stale\\
\color{white}* Config changes aren't taking effect\\
\color{white}* Something feels \textbackslash{}"off\textbackslash{}" and you can't pin it down\\
\color{white}* After a major ZPWR update\\
\vspace{4pt}\textbf{\color{magenta}zpwr reload}\\
\color{white}Hot reload without resetting. Re-sources all configs\\
\color{white}and refreshes the environment. Faster than reset.\\
\color{white}\textbackslash{}\$ zpwr reload                 \# hot reload\\
\color{white}Use reload when you've edited a config and want to\\
\color{white}apply changes without nuking caches.\\
\vspace{4pt}\textbf{\color{magenta}ESCALATION LADDER:}\\
\color{white}Level 1: zpwr reload           \# re-source configs\\
\color{white}Level 2: zpwr reset            \# rebuild everything\\
\color{white}Level 3: zpwr install          \# full reinstall\\
\color{white}Level 4: rm -rf \textbackslash{}PWR \&\& clone \# scorched earth\\
\color{white}Start at Level 1. Only escalate if the problem persists.\\
\color{white}Most issues are fixed by reload. Reset handles the rest.\\
\vspace{4pt}\textbf{\color{magenta}COMMON RECOVERY SCENARIOS:}\\
\color{white}\textbackslash{}"completions are wrong\textbackslash{}"    --> zpwr reset\\
\color{white}\textbackslash{}"alias isn't working\textbackslash{}"      --> zpwr reload\\
\color{white}\textbackslash{}"new plugin not loading\textbackslash{}"   --> zpwr reload\\
\color{white}\textbackslash{}"everything is broken\textbackslash{}"     --> zpwr reset\\
\color{white}\textbackslash{}"still broken after reset\textbackslash{}" --> zpwr install\\
\textit{\color{gray}// reload is ctrl-r for your shell framework}\\
\textit{\color{gray}// reset is the factory reset you hope to never need}\\
\textit{\color{gray}// but when you do need it, you'll be glad it exists}\\

\vspace{8pt}


\section{CHAPTER 19: BATCH OPERATIONS}

\subsection{Page 76: Loop Warriors}

\vspace{4pt}\textbf{\color{cyan}<<< LOOP WARRIORS — THE REPETITION TOOLKIT >>>}\vspace{2pt}\\
\color{white}because doing things once is for civilians\\
\color{white}The batch operations verbs are the assault rifles\\
\color{white}of the ZPWR arsenal. Point them at a problem,\\
\color{white}pull the trigger, and watch every target fall.\\
\vspace{4pt}\textbf{\color{magenta}THE FOR LOOP — YOUR BASIC INFANTRY}\\
\texttt{\color{green}zpwr for <cmd> <args...>}\\
\color{white}Executes a command for each argument passed to it.\\
\color{white}Stdin becomes your ammo belt — pipe lines in, each one fires.\\
\color{white}Example: run a command against multiple targets\\
\texttt{\color{green}echo 'file1.txt\textbackslash{}nfile2.txt' | zpwr for wc -l}\\
\vspace{4pt}\textbf{\color{magenta}FOR10 — THE BURST-FIRE VARIANT}\\
\texttt{\color{green}zpwr for10 <cmd>}\\
\color{white}Runs your command 10 times. No arguments, no mercy.\\
\color{white}Perfect for stress testing, benchmarking, or just\\
\color{white}asserting dominance over your CPU.\\
\color{white}Example: benchmark a command 10 times\\
\texttt{\color{green}zpwr for10 zpwr bench}\\
\vspace{4pt}\textbf{\color{magenta}WHILETRUE — THE INFINITE LOOP}\\
\texttt{\color{green}zpwr whiletrue <cmd>}\\
\color{white}Runs your command forever. Ctrl-C is your only exit.\\
\color{white}For monitoring, watching, waiting. The sentinel loop.\\
\color{white}Example: watch disk usage forever\\
\texttt{\color{green}zpwr whiletrue df -h /}\\
\color{white}WARNING: there is no break condition. You ARE the break.\\
\vspace{4pt}\textbf{\color{magenta}WHILESLEEP — THE PATIENT SNIPER}\\
\texttt{\color{green}zpwr whilesleep <seconds> <cmd>}\\
\color{white}Like whiletrue, but sleeps between rounds.\\
\color{white}Polling with grace. Surveillance with manners.\\
\color{white}Example: check a URL every 30 seconds\\
\texttt{\color{green}zpwr whilesleep 30 curl -s -o /dev/null -w '\%\{http\_code\}' https://example.com}\\
\color{white}TACTICAL SUMMARY:\\
\color{white}for          iterate over stdin lines\\
\color{white}for10        run 10 times, no questions asked\\
\color{white}whiletrue    infinite loop, Ctrl-C to escape\\
\color{white}whilesleep   infinite loop with naptime between rounds\\
\textit{\color{gray}// in the loop, there is no time, only iterations}\\
\textit{\color{gray}// the loop does not end. the operator ends the loop.}\\

\vspace{8pt}

\subsection{Page 77: Directory Sweeps}

\vspace{4pt}\textbf{\color{cyan}<<< DIRECTORY SWEEPS — CARPET BOMBING YOUR FILESYSTEM >>>}\vspace{2pt}\\
\color{white}one directory at a time, or all of them at once\\
\color{white}Single-file loops are amateur hour.\\
\color{white}Real operators sweep directories. They glob patterns.\\
\color{white}They parallelize. These verbs are your air support.\\
\vspace{4pt}\textbf{\color{magenta}FORDIR — THE DIRECTORY CRAWLER}\\
\texttt{\color{green}zpwr fordir <cmd>}\\
\color{white}Iterates into each subdirectory of the current dir\\
\color{white}and executes the command inside it. cd in, run, cd out.\\
\color{white}Like a SWAT team clearing rooms one by one.\\
\color{white}Example: check git status in every project dir\\
\texttt{\color{green}cd \textasciitilde{}/repos \&\& zpwr fordir git status -s}\\
\vspace{4pt}\textbf{\color{magenta}FORDIRUPDATE — THE MASS UPDATER}\\
\texttt{\color{green}zpwr fordirupdate}\\
\color{white}Enters every subdirectory and runs git pull.\\
\color{white}Your entire project portfolio, synced in one command.\\
\color{white}Monday morning? Run this first. Coffee second.\\
\color{white}Example: update all repos under \textasciitilde{}/repos\\
\texttt{\color{green}cd \textasciitilde{}/repos \&\& zpwr fordirupdate}\\
\vspace{4pt}\textbf{\color{magenta}EXECGLOB — GLOB AND EXECUTE}\\
\texttt{\color{green}zpwr execglob <glob> <cmd>}\\
\color{white}Finds files matching a glob pattern and runs your command\\
\color{white}against each one. Pattern matching meets execution.\\
\color{white}The regex sniper rifle of file operations.\\
\color{white}Example: lint every Python file in the tree\\
\texttt{\color{green}zpwr execglob '*.py' pylint}\\
\vspace{4pt}\textbf{\color{magenta}EXECGLOBPARALLEL — CONCURRENT DEVASTATION}\\
\texttt{\color{green}zpwr execglobparallel <glob> <cmd>}\\
\color{white}Same as execglob, but fires all jobs in parallel.\\
\color{white}Every CPU core gets drafted into service.\\
\color{white}Your fan will spin up. This is normal. This is power.\\
\color{white}Example: format every .rs file in parallel\\
\texttt{\color{green}zpwr execglobparallel '*.rs' rustfmt}\\
\color{white}CHOOSING YOUR WEAPON:\\
\color{white}fordir             run cmd in each subdirectory\\
\color{white}fordirupdate       git pull in each subdirectory\\
\color{white}execglob           glob match + execute sequentially\\
\color{white}execglobparallel   glob match + execute in parallel\\
\textit{\color{gray}// sequential is predictable. parallel is fast.}\\
\textit{\color{gray}// choose based on whether output order matters.}\\
\textit{\color{gray}// when in doubt, go parallel. life is short.}\\

\vspace{8pt}

\subsection{Page 78: Polyglot Execution}

\vspace{4pt}\textbf{\color{cyan}<<< POLYGLOT EXECUTION — BEYOND THE SHELL >>>}\vspace{2pt}\\
\color{white}zsh is the conductor, but the orchestra speaks many tongues\\
\color{white}Sometimes zsh loops aren't enough. Sometimes you need\\
\color{white}Python. Sometimes you need background jobs. Sometimes\\
\color{white}you need to poll a service until it wakes up or dies.\\
\vspace{4pt}\textbf{\color{magenta}EXECPY — THE PYTHON BRIDGE}\\
\texttt{\color{green}zpwr execpy <python\_expression>}\\
\color{white}Fires a Python expression from the comfort of your shell.\\
\color{white}No need to open a REPL. No need for a .py file.\\
\color{white}Quick math, string ops, one-liners — python as a verb.\\
\color{white}Example: quick calculation\\
\texttt{\color{green}zpwr execpy '2**128'}\\
\color{white}=> 340282366920938463463374607431768211456\\
\vspace{4pt}\textbf{\color{magenta}BACKGROUND — THE SHADOW OPERATIVE}\\
\texttt{\color{green}zpwr background <cmd>}\\
\color{white}Sends a command into the background. Detached. Silent.\\
\color{white}It runs in the shadows while you keep working.\\
\color{white}Output goes to the void unless you redirect it.\\
\color{white}Example: start a long build in the background\\
\texttt{\color{green}zpwr background make -j8}\\
\vspace{4pt}\textbf{\color{magenta}TIMER — THE STOPWATCH}\\
\texttt{\color{green}zpwr timer <cmd>}\\
\color{white}Wraps any command in a high-resolution timer.\\
\color{white}How long did that take? Now you know. No guessing.\\
\color{white}Benchmarking made trivial.\\
\color{white}Example: time your shell startup\\
\texttt{\color{green}zpwr timer zsh -ic exit}\\
\vspace{4pt}\textbf{\color{magenta}POLL — THE WATCHDOG}\\
\texttt{\color{green}zpwr poll <cmd>}\\
\color{white}Runs a command repeatedly until it succeeds (exit 0).\\
\color{white}Waiting for a server to come up? A build to finish?\\
\color{white}Poll is your patient, relentless sentinel.\\
\color{white}Example: wait for a port to open\\
\texttt{\color{green}zpwr poll nc -z localhost 8080}\\
\color{white}THE BATCH OPS PHILOSOPHY:\\
\color{white}Loops        for, for10, whiletrue, whilesleep\\
\color{white}Sweeps       fordir, fordirupdate, execglob*\\
\color{white}Polyglot     execpy, background, timer, poll\\
\textit{\color{gray}// the terminal is not a single-threaded mind}\\
\textit{\color{gray}// it is a factory floor. these verbs are the machines.}\\
\textit{\color{gray}// feed them work. they do not tire. they do not complain.}\\

\vspace{8pt}


\section{CHAPTER 20: FORGIT DEEP DIVE}

\subsection{Page 79: Forgit Add \& Diff}

\vspace{4pt}\textbf{\color{cyan}<<< FORGIT ADD \& DIFF — INTERACTIVE STAGING >>>}\vspace{2pt}\\
\color{white}git add, but make it beautiful\\
\color{white}forgit is the fusion of fzf and git. It wraps the most\\
\color{white}common git operations in fuzzy, interactive, preview-enabled\\
\color{white}interfaces. You will never go back to raw git add.\\
\vspace{4pt}\textbf{\color{magenta}FORGITADD — INTERACTIVE STAGING}\\
\texttt{\color{green}zpwr forgitadd}\\
\color{white}Opens an fzf interface showing all unstaged changes.\\
\color{white}Each file gets a diff preview in the right pane.\\
\color{white}Tab to select multiple files. Enter to stage them.\\
\color{white}The workflow:\\
\texttt{\color{green}1. Run zpwr forgitadd}\\
\color{white}2. Browse files with arrow keys\\
\color{white}3. See the diff preview for each file\\
\color{white}4. Tab to multi-select\\
\color{white}5. Enter to stage selected files\\
\color{white}No more git add -p scrolling through 500-line diffs.\\
\color{white}No more accidentally staging that debug print.\\
\vspace{4pt}\textbf{\color{magenta}FORGITDIFF — VISUAL DIFF BROWSER}\\
\texttt{\color{green}zpwr forgitdiff}\\
\color{white}Opens an fzf browser of all changed files with\\
\color{white}a live diff preview. Read diffs without committing\\
\color{white}to anything. Pure reconnaissance.\\
\color{white}Example workflow:\\
\texttt{\color{green}zpwr forgitdiff           \# browse all diffs visually}\\
\texttt{\color{green}zpwr forgitadd            \# stage what you want}\\
\texttt{\color{green}zpwr gitcommit 'msg'      \# commit with message}\\
\color{white}WHY FORGIT > RAW GIT:\\
\color{white}Raw git add:     type filenames from memory like a caveman\\
\color{white}forgitadd:       browse, preview, multi-select, confirm\\
\color{white}Raw git diff:    wall of text scrolling past\\
\color{white}forgitdiff:      file-by-file preview with fuzzy search\\
\textit{\color{gray}// staging should not be a memory test}\\
\textit{\color{gray}// it should be a visual, interactive experience}\\
\textit{\color{gray}// forgit turns git into a GUI without leaving the terminal}\\

\vspace{8pt}

\subsection{Page 80: Forgit Log \& Stash}

\vspace{4pt}\textbf{\color{cyan}<<< FORGIT LOG \& STASH — VISUAL HISTORY >>>}\vspace{2pt}\\
\color{white}time travel with a preview window\\
\color{white}Git log is a river of hashes. Git stash is a pile of\\
\color{white}forgotten context. Forgit turns both into something\\
\color{white}you can actually navigate without losing your mind.\\
\vspace{4pt}\textbf{\color{magenta}FORGITLOG — THE COMMIT BROWSER}\\
\texttt{\color{green}zpwr forgitlog}\\
\color{white}Opens your git history in fzf with commit previews.\\
\color{white}Each commit shows its full diff in the preview pane.\\
\color{white}Search commits by message, author, or hash.\\
\color{white}What you see:\\
\color{white}Left pane:  commit graph with short messages\\
\color{white}Right pane: full diff of highlighted commit\\
\color{white}Search:     fuzzy match on any commit field\\
\color{white}\textbackslash{}"Who wrote this garbage?\textbackslash{}" — you, using forgitlog,\\
\color{white}discovering it was also you, three months ago.\\
\vspace{4pt}\textbf{\color{magenta}FORGITSTASH — THE STASH MANAGER}\\
\texttt{\color{green}zpwr forgitstash}\\
\color{white}Browse your stash entries with preview diffs.\\
\color{white}See what is in each stash before applying it.\\
\color{white}No more git stash pop roulette.\\
\color{white}The old way:\\
\color{white}git stash list              \# cryptic one-liners\\
\color{white}git stash show -p stash@\{2\}  \# which one was it again?\\
\color{white}git stash pop               \# pray\\
\color{white}The forgit way:\\
\texttt{\color{green}zpwr forgitstash            \# browse, preview, apply}\\
\color{white}LOG + STASH TOGETHER:\\
\color{white}These two verbs give you eyes into the past.\\
\color{white}forgitlog shows committed history. forgitstash shows\\
\color{white}uncommitted history. Between them, nothing is lost.\\
\textit{\color{gray}// the commit graph is not a wall of text}\\
\textit{\color{gray}// it is a timeline. forgit lets you walk it.}\\
\textit{\color{gray}// stashes are not forgotten. they are filed.}\\

\vspace{8pt}

\subsection{Page 81: Forgit Ignore \& Clean}

\vspace{4pt}\textbf{\color{cyan}<<< FORGIT IGNORE \& CLEAN — .GITIGNORE MANAGEMENT >>>}\vspace{2pt}\\
\color{white}because node\_modules should never see the inside of a commit\\
\color{white}The .gitignore is your repo's bouncer list. These verbs\\
\color{white}manage it with the precision of a nightclub velvet rope.\\
\vspace{4pt}\textbf{\color{magenta}FORGITIGNORE — THE IGNORE BROWSER}\\
\texttt{\color{green}zpwr forgitignore}\\
\color{white}Interactive gitignore template browser powered by\\
\color{white}gitignore.io. Search for languages, frameworks, IDEs.\\
\color{white}Never hand-write a .gitignore again.\\
\vspace{4pt}\textbf{\color{magenta}THE IGNORE FAMILY:}\\
\color{white}forgitignorelist     list available gitignore templates\\
\color{white}forgitignoreget      fetch a specific template by name\\
\color{white}forgitignoreupdate   update your local template cache\\
\color{white}forgitignoreclean    clean the template cache\\
\color{white}Example: generate a Python + macOS gitignore\\
\texttt{\color{green}zpwr forgitignoreget python,macos > .gitignore}\\
\color{white}Example: browse all available templates\\
\texttt{\color{green}zpwr forgitignorelist}\\
\vspace{4pt}\textbf{\color{magenta}FORGITCLEAN — THE UNTRACKED FILE PURGE}\\
\texttt{\color{green}zpwr forgitclean}\\
\color{white}Shows untracked files in an fzf interface with previews.\\
\color{white}Select which ones to delete. Surgical, not scorched earth.\\
\color{white}Unlike git clean -fd, you see before you destroy.\\
\color{white}The old way:\\
\color{white}git clean -fd    \# delete everything. hope for the best.\\
\color{white}The forgit way:\\
\texttt{\color{green}zpwr forgitclean  \# browse, select, confirm, then delete.}\\
\color{white}IGNORE WORKFLOW:\\
\texttt{\color{green}1. zpwr forgitignorelist    browse templates}\\
\texttt{\color{green}2. zpwr forgitignoreget     fetch what you need}\\
\texttt{\color{green}3. zpwr forgitclean         purge untracked cruft}\\
\texttt{\color{green}4. zpwr forgitadd           stage the clean repo}\\
\textit{\color{gray}// a clean repo is a happy repo}\\
\textit{\color{gray}// an ignored file is a file that cannot hurt you}\\
\textit{\color{gray}// let the templates do the thinking}\\

\vspace{8pt}

\subsection{Page 82: Forgit Reset \& Restore}

\vspace{4pt}\textbf{\color{cyan}<<< FORGIT RESET \& RESTORE — UNDO WITH EYES OPEN >>>}\vspace{2pt}\\
\color{white}the nuclear options, but with a safety preview\\
\color{white}Git reset and restore are the most terrifying commands\\
\color{white}in version control. Forgit wraps them in bubble wrap\\
\color{white}and gives you a preview before you pull the trigger.\\
\vspace{4pt}\textbf{\color{magenta}FORGITRESET — UNSTAGE WITH PREVIEW}\\
\texttt{\color{green}zpwr forgitreset}\\
\color{white}Shows staged files in fzf with diff previews.\\
\color{white}Select files to unstage. Like git reset HEAD but\\
\color{white}you can see exactly what you are unstaging first.\\
\color{white}Use case: you ran git add -A like a maniac\\
\color{white}and now you need to surgically remove files.\\
\texttt{\color{green}zpwr forgitreset    \# browse and unstage selectively}\\
\vspace{4pt}\textbf{\color{magenta}FORGITRESTORE — DISCARD CHANGES SAFELY}\\
\texttt{\color{green}zpwr forgitrestore}\\
\color{white}Shows modified files with diff previews. Select files\\
\color{white}to restore to their last committed state.\\
\color{white}This discards changes permanently. No undo.\\
\color{white}The preview saves you. You see the diff before\\
\color{white}the changes evaporate into the void.\\
\vspace{4pt}\textbf{\color{magenta}FORGITWARN — THE SAFETY NET}\\
\texttt{\color{green}zpwr forgitwarn}\\
\color{white}Displays warnings about uncommitted changes, diverged\\
\color{white}branches, and other hazardous states. Your pre-flight\\
\color{white}checklist before doing anything destructive.\\
\vspace{4pt}\textbf{\color{magenta}FORGITINFO — THE STATUS DASHBOARD}\\
\texttt{\color{green}zpwr forgitinfo}\\
\color{white}Comprehensive repo information at a glance.\\
\color{white}Branch, remotes, recent commits, dirty state.\\
\color{white}Everything you need before making decisions.\\
\color{white}THE FORGIT COMPLETE ARSENAL:\\
\color{white}forgitadd        interactive stage\\
\color{white}forgitdiff       browse diffs\\
\color{white}forgitlog        commit history browser\\
\color{white}forgitstash      stash manager\\
\color{white}forgitreset      interactive unstage\\
\color{white}forgitrestore    discard changes with preview\\
\color{white}forgitclean      purge untracked files\\
\color{white}forgitignore*    gitignore template management\\
\color{white}forgitwarn       hazard warnings\\
\color{white}forgitinfo       repo dashboard\\
\textit{\color{gray}// forgit is not a git replacement}\\
\textit{\color{gray}// it is git with eyes. git with a UI. git for humans.}\\

\vspace{8pt}


\section{CHAPTER 21: ZPWR ENVIRONMENT VARIABLES}

\subsection{Page 83: Core Paths}

\vspace{4pt}\textbf{\color{cyan}<<< CORE PATHS — THE SKELETON OF THE MACHINE >>>}\vspace{2pt}\\
\color{white}every system has bones. these are ZPWR's.\\
\color{white}ZPWR does not hardcode paths. Every directory is an\\
\color{white}environment variable. Move the root, and the entire\\
\color{white}system follows. This is not a dotfiles repo — it is\\
\color{white}a relocatable operating system.\\
\vspace{4pt}\textbf{\color{magenta}\textbackslash{}PWR — THE ROOT OF ALL THINGS}\\
\texttt{\color{green}echo \textbackslash{}PWR   \# typically \textasciitilde{}/.zpwr}\\
\color{white}The root directory. Every other path derives from this.\\
\color{white}Change this one variable and the whole tree moves.\\
\vspace{4pt}\textbf{\color{magenta}\textbackslash{}PWR\_LOCAL — YOUR PRIVATE SANCTUM}\\
\color{white}echo \textbackslash{}PWR\_LOCAL   \# \textbackslash{}PWR/local\\
\color{white}Gitignored. Your tokens, secrets, local overrides.\\
\color{white}Nothing in here ever touches a remote repository.\\
\color{white}Your .tokens.sh lives here. Guard it with your life.\\
\vspace{4pt}\textbf{\color{magenta}\textbackslash{}PWR\_SCRIPTS — THE ARMORY}\\
\color{white}echo \textbackslash{}PWR\_SCRIPTS   \# \textbackslash{}PWR/scripts\\
\color{white}Executable scripts. Bash, Python, Perl — whatever\\
\color{white}gets the job done. These are on your PATH.\\
\vspace{4pt}\textbf{\color{magenta}\textbackslash{}PWR\_TMUX — THE MULTIPLEXER CONFIG}\\
\color{white}echo \textbackslash{}PWR\_TMUX   \# \textbackslash{}PWR/tmux\\
\color{white}Tmux configuration, scripts, and status line logic.\\
\color{white}The control center for your terminal multiplexer.\\
\vspace{4pt}\textbf{\color{magenta}\textbackslash{}PWR\_INSTALL — THE INSTALLER}\\
\color{white}echo \textbackslash{}PWR\_INSTALL   \# \textbackslash{}PWR/install\\
\color{white}OS-specific installation scripts. macOS, Ubuntu,\\
\color{white}Fedora, Arch — each gets its own installer.\\
\texttt{\color{green}Run zpwr install and the right one fires.}\\
\color{white}THE PATH MAP:\\
\color{white}\textbackslash{}PWR                \textasciitilde{}/.zpwr (the root)\\
\color{white}\textbackslash{}PWR\_LOCAL          \textbackslash{}PWR/local (gitignored secrets)\\
\color{white}\textbackslash{}PWR\_SCRIPTS        \textbackslash{}PWR/scripts (executables)\\
\color{white}\textbackslash{}PWR\_TMUX           \textbackslash{}PWR/tmux (multiplexer)\\
\color{white}\textbackslash{}PWR\_INSTALL        \textbackslash{}PWR/install (OS installers)\\
\textit{\color{gray}// paths are not strings. they are load-bearing structures.}\\
\textit{\color{gray}// hardcode a path and you weld yourself to the floor.}\\
\textit{\color{gray}// use the variables. always use the variables.}\\

\vspace{8pt}

\subsection{Page 84: Autoload System}

\vspace{4pt}\textbf{\color{cyan}<<< THE AUTOLOAD SYSTEM — LAZY LOADING FOR HACKERS >>>}\vspace{2pt}\\
\color{white}why load everything when you can load nothing until needed?\\
\color{white}Zsh autoloading is the secret to ZPWR's fast startup.\\
\color{white}Functions are not loaded into memory until you call them.\\
\color{white}The shell knows they exist, but their code stays on disk\\
\color{white}until the moment of invocation. Maximum laziness. Maximum speed.\\
\vspace{4pt}\textbf{\color{magenta}\textbackslash{}PWR\_AUTOLOAD — THE FUNCTION WAREHOUSE}\\
\color{white}echo \textbackslash{}PWR\_AUTOLOAD   \# \textbackslash{}PWR/autoload\\
\color{white}The root of the autoload tree. Every file in this tree\\
\color{white}becomes a function name. The filename IS the function.\\
\vspace{4pt}\textbf{\color{magenta}\textbackslash{}PWR\_AUTOLOAD\_COMMON — CROSS-PLATFORM FUNCTIONS}\\
\color{white}echo \textbackslash{}PWR\_AUTOLOAD\_COMMON   \# \textbackslash{}PWR/autoload/common\\
\color{white}Functions that work on every OS. The common denominator.\\
\color{white}macOS, Linux, BSD — these functions do not care.\\
\vspace{4pt}\textbf{\color{magenta}HOW AUTOLOADING WORKS:}\\
\color{white}1. Zsh scans \textbackslash{}\$fpath directories at startup\\
\color{white}2. Each filename is registered as a function stub\\
\color{white}3. When you call the function, zsh reads the file\\
\color{white}4. The file contents replace the stub\\
\color{white}5. The function executes with the loaded code\\
\color{white}This means 400+ functions cost near-zero at startup.\\
\color{white}Only the ones you actually use consume memory.\\
\vspace{4pt}\textbf{\color{magenta}THE COMPILE ACCELERATION:}\\
\texttt{\color{green}zpwr compile       \# compile autoloads to .zwc}\\
\texttt{\color{green}zpwr recompile     \# recompile all}\\
\texttt{\color{green}zpwr uncompile     \# remove compiled files}\\
\color{white}Compiled .zwc files load faster than raw source.\\
\color{white}ZPWR compiles its autoload dirs into wordcode files.\\
\color{white}It is like bytecode for your shell functions.\\
\color{white}INSPECTION VERBS:\\
\texttt{\color{green}zpwr autoload          show the autoload system}\\
\texttt{\color{green}zpwr autoloadcount     count autoloaded functions}\\
\texttt{\color{green}zpwr autoloadlist      list all autoloaded functions}\\
\textit{\color{gray}// the fastest code is code that never runs}\\
\textit{\color{gray}// autoloading is the art of deferred execution}\\
\textit{\color{gray}// 400 functions, zero startup cost. that is the deal.}\\

\vspace{8pt}

\subsection{Page 85: Behavior Flags}

\vspace{4pt}\textbf{\color{cyan}<<< BEHAVIOR FLAGS — THE KNOBS AND DIALS >>>}\vspace{2pt}\\
\color{white}tune the machine without cracking open the engine\\
\color{white}ZPWR is not a take-it-or-leave-it monolith.\\
\color{white}These boolean and string flags let you control\\
\color{white}startup behavior, logging verbosity, and debug output\\
\color{white}without editing a single line of source code.\\
\vspace{4pt}\textbf{\color{magenta}\textbackslash{}PWR\_PROFILING — STARTUP PROFILER}\\
\color{white}export ZPWR\_PROFILING=true\\
\color{white}Enables zsh startup profiling. Shows exactly where\\
\color{white}your milliseconds go during shell init. Pair with\\
\texttt{\color{green}zpwr bench for the full picture.}\\
\color{white}Set it, open a new shell, and read the profiling data.\\
\color{white}Find the bottleneck. Kill the bottleneck. Repeat.\\
\vspace{4pt}\textbf{\color{magenta}\textbackslash{}PWR\_DEBUG — DEBUG MODE}\\
\color{white}export ZPWR\_DEBUG=true\\
\color{white}Turns on debug-level log messages. When something\\
\color{white}is broken and you need the machine to talk, flip this.\\
\texttt{\color{green}Check output with zpwr logdebug.}\\
\vspace{4pt}\textbf{\color{magenta}\textbackslash{}PWR\_TRACE — FULL TRACE MODE}\\
\color{white}export ZPWR\_TRACE=true\\
\color{white}The firehose. Every function call, every expansion,\\
\texttt{\color{green}traced to the log. Use zpwr trace to run commands}\\
\texttt{\color{green}with tracing, or zpwr logtrace to read trace logs.}\\
\color{white}WARNING: produces enormous output. Use surgically.\\
\vspace{4pt}\textbf{\color{magenta}\textbackslash{}PWR\_BANNER\_CLEARLIST — STARTUP BANNER CONTROL}\\
\color{white}export ZPWR\_BANNER\_CLEARLIST=true\\
\color{white}Controls whether the ZPWR banner and clearlist show\\
\color{white}at shell startup. The banner is glorious, but if\\
\color{white}you open 50 terminals a day, you might want silence.\\
\color{white}Related verbs:\\
\texttt{\color{green}zpwr banner          show the banner manually}\\
\texttt{\color{green}zpwr clearlist       show clearlist manually}\\
\texttt{\color{green}zpwr bannerlolcat    banner with rainbow colors}\\
\color{white}QUICK REFERENCE:\\
\color{white}ZPWR\_PROFILING        startup timing data\\
\color{white}ZPWR\_DEBUG            debug-level logging\\
\color{white}ZPWR\_TRACE            full execution trace\\
\color{white}ZPWR\_BANNER\_CLEARLIST startup banner toggle\\
\textit{\color{gray}// the flags do not change the code. they reveal it.}\\
\textit{\color{gray}// debugging is not failure. it is x-ray vision.}\\

\vspace{8pt}

\subsection{Page 86: Plugin System}

\vspace{4pt}\textbf{\color{cyan}<<< PLUGIN SYSTEM — THE TURBO LOADER >>>}\vspace{2pt}\\
\color{white}zinit: loading plugins at the speed of light\\
\color{white}ZPWR uses zinit (formerly zplugin) as its plugin\\
\color{white}manager. But not the basic way. The turbo way.\\
\color{white}Plugins load asynchronously after prompt appears.\\
\color{white}You see your prompt instantly. Plugins catch up.\\
\vspace{4pt}\textbf{\color{magenta}\textbackslash{}PWR\_PLUGIN\_MANAGER — CHOOSE YOUR LOADER}\\
\color{white}echo \textbackslash{}PWR\_PLUGIN\_MANAGER   \# zinit\\
\color{white}Controls which plugin manager ZPWR uses. The default\\
\color{white}and recommended value is zinit. Other managers exist\\
\color{white}but zinit's turbo mode is what makes ZPWR fast.\\
\vspace{4pt}\textbf{\color{magenta}TURBO LOADING EXPLAINED:}\\
\color{white}Normal plugin loading:\\
\color{white}source plugin1.zsh  \# blocks\\
\color{white}source plugin2.zsh  \# blocks\\
\color{white}source plugin3.zsh  \# blocks\\
\color{white}\# prompt finally appears after ALL plugins load\\
\color{white}Zinit turbo loading:\\
\color{white}prompt appears immediately\\
\color{white}plugin1 loads in background ... done\\
\color{white}plugin2 loads in background ... done\\
\color{white}plugin3 loads in background ... done\\
\color{white}\# user never notices the delay\\
\vspace{4pt}\textbf{\color{magenta}PLUGIN INSPECTION:}\\
\texttt{\color{green}zpwr zshplugincount     count loaded zsh plugins}\\
\texttt{\color{green}zpwr zshpluginlist      list all loaded zsh plugins}\\
\texttt{\color{green}zpwr vimplugincount     count vim plugins}\\
\texttt{\color{green}zpwr vimpluginlist      list all vim plugins}\\
\texttt{\color{green}zpwr emacsplugincount   count emacs plugins}\\
\texttt{\color{green}zpwr emacspluginlist    list all emacs plugins}\\
\vspace{4pt}\textbf{\color{magenta}UPDATE AND MAINTAIN:}\\
\texttt{\color{green}zpwr updatezinit        update zinit itself}\\
\texttt{\color{green}zpwr updatedeps          update all plugin deps}\\
\texttt{\color{green}zpwr plugins             plugin information}\\
\textit{\color{gray}// plugins are not bloat. they are force multipliers.}\\
\textit{\color{gray}// but only if they load in zero perceived time.}\\
\textit{\color{gray}// turbo mode makes 50 plugins feel like zero.}\\

\vspace{8pt}

\subsection{Page 87: Display \& Input}

\vspace{4pt}\textbf{\color{cyan}<<< DISPLAY \& INPUT — THE HUMAN INTERFACE LAYER >>>}\vspace{2pt}\\
\color{white}how the machine talks to you, and how you talk back\\
\color{white}These variables control what you see and how fast\\
\color{white}the shell responds to your keystrokes. Milliseconds\\
\color{white}matter when your fingers are moving at light speed.\\
\vspace{4pt}\textbf{\color{magenta}\textbackslash{}PWR\_EXA\_COMMAND — THE LS REPLACEMENT}\\
\color{white}echo \textbackslash{}PWR\_EXA\_COMMAND   \# exa or eza\\
\color{white}Controls which modern ls replacement ZPWR uses.\\
\color{white}exa is the original, eza is the maintained fork.\\
\color{white}Both give you colored, git-aware directory listings.\\
\texttt{\color{green}ZPWR verbs like zpwr ls and zpwr clearls use this.}\\
\color{white}Set it in your tokens file to override the default.\\
\vspace{4pt}\textbf{\color{magenta}\textbackslash{}PWR\_KEYTIMEOUT — KEYSTROKE LATENCY}\\
\color{white}echo \textbackslash{}PWR\_KEYTIMEOUT   \# in hundredths of a second\\
\color{white}How long zsh waits for the next key in a multi-key\\
\color{white}sequence. Lower = snappier. Too low = missed sequences.\\
\color{white}Default is tuned for fast typists.\\
\color{white}If vi-mode feels laggy on Esc, lower this.\\
\color{white}If key combos misfire, raise this.\\
\color{white}Sweet spot is personal. Experiment.\\
\vspace{4pt}\textbf{\color{magenta}\textbackslash{}PWR\_EXPAND\_PRE\_EXEC — ALIAS EXPANSION}\\
\color{white}echo \textbackslash{}PWR\_EXPAND\_PRE\_EXEC\\
\color{white}Controls whether ZPWR expands aliases before execution.\\
\color{white}When enabled, you type the alias and see the full\\
\color{white}command in your history. Transparency mode.\\
\color{white}With expansion:    you type ga  -> history shows git add\\
\color{white}Without expansion: you type ga  -> history shows ga\\
\vspace{4pt}\textbf{\color{magenta}\textbackslash{}PWR\_SEND\_KEYS\_FULL — TMUX KEY PASSTHROUGH}\\
\color{white}echo \textbackslash{}PWR\_SEND\_KEYS\_FULL\\
\color{white}Controls full send-keys behavior in tmux integration.\\
\color{white}When ZPWR manages tmux panes, this determines how\\
\color{white}keystrokes are forwarded to target panes.\\
\color{white}TUNING GUIDE:\\
\color{white}Feeling laggy?       Lower ZPWR\_KEYTIMEOUT\\
\color{white}Missing combos?      Raise ZPWR\_KEYTIMEOUT\\
\color{white}Want transparency?   Enable ZPWR\_EXPAND\_PRE\_EXEC\\
\color{white}Using eza?           Set ZPWR\_EXA\_COMMAND=eza\\
\textit{\color{gray}// the interface is not decoration. it is velocity.}\\
\textit{\color{gray}// every millisecond of latency is a papercut.}\\
\textit{\color{gray}// tune these until your shell feels like thought.}\\

\vspace{8pt}

\subsection{Page 88: Runtime State}

\vspace{4pt}\textbf{\color{cyan}<<< RUNTIME STATE — THE MACHINE'S SELF-AWARENESS >>>}\vspace{2pt}\\
\color{white}the shell knows what it is. do you?\\
\color{white}These variables hold runtime state — what OS you are on,\\
\color{white}what TTY you are using, where logs go. They are set\\
\color{white}at startup and consulted throughout the session.\\
\vspace{4pt}\textbf{\color{magenta}\textbackslash{}PWR\_TTY — YOUR TERMINAL IDENTITY}\\
\color{white}echo \textbackslash{}PWR\_TTY   \# e.g., /dev/ttys003\\
\color{white}The TTY device for the current shell session.\\
\color{white}Used by ZPWR to track which terminal is which,\\
\color{white}especially when managing multiple tmux panes.\\
\texttt{\color{green}The zpwr tty verb shows your current TTY info.}\\
\vspace{4pt}\textbf{\color{magenta}\textbackslash{}PWR\_OS\_TYPE — PLATFORM DETECTION}\\
\color{white}echo \textbackslash{}PWR\_OS\_TYPE   \# Darwin, Linux, etc.\\
\color{white}Detected at startup. Used throughout ZPWR to branch\\
\color{white}on platform-specific behavior. macOS vs Linux vs BSD,\\
\color{white}the right code path fires automatically.\\
\color{white}This is why ZPWR works on macOS and Linux without\\
\color{white}you changing a single setting. The OS type decides.\\
\vspace{4pt}\textbf{\color{magenta}\textbackslash{}PWR\_LOGFILE — THE BLACK BOX}\\
\color{white}echo \textbackslash{}PWR\_LOGFILE\\
\color{white}Where ZPWR writes its log output. Debug, info, error,\\
\color{white}trace — all levels go here when enabled.\\
\color{white}Read the logs:\\
\texttt{\color{green}zpwr loginfo         info-level messages}\\
\texttt{\color{green}zpwr logdebug        debug-level messages}\\
\texttt{\color{green}zpwr logerror        error-level messages}\\
\texttt{\color{green}zpwr logtrace        trace-level messages}\\
\texttt{\color{green}zpwr taillog         live tail with colorized output}\\
\vspace{4pt}\textbf{\color{magenta}\textbackslash{}PWR\_COMPSYS\_CACHE — COMPLETION CACHE}\\
\color{white}echo \textbackslash{}PWR\_COMPSYS\_CACHE\\
\color{white}Controls the completion system cache. Zsh completions\\
\color{white}are expensive to regenerate. This flag keeps them cached.\\
\vspace{4pt}\textbf{\color{magenta}\textbackslash{}PWR\_HISTFILE — HISTORY FILE}\\
\color{white}echo \textbackslash{}PWR\_HISTFILE\\
\color{white}Your shell history file location. ZPWR manages history\\
\color{white}with deduplication and backup. Your commands are sacred.\\
\texttt{\color{green}Related: zpwr backuphistory and zpwr restorehistory}\\
\color{white}THE FULL ENV VAR MAP:\\
\color{white}Core paths:      ZPWR, ZPWR\_LOCAL, ZPWR\_SCRIPTS, ZPWR\_TMUX\\
\color{white}Autoload:        ZPWR\_AUTOLOAD, ZPWR\_AUTOLOAD\_COMMON\\
\color{white}Flags:           ZPWR\_PROFILING, ZPWR\_DEBUG, ZPWR\_TRACE\\
\color{white}Display:         ZPWR\_EXA\_COMMAND, ZPWR\_KEYTIMEOUT\\
\color{white}Runtime:         ZPWR\_TTY, ZPWR\_OS\_TYPE, ZPWR\_LOGFILE\\
\textit{\color{gray}// the machine knows itself. every variable is a mirror.}\\
\textit{\color{gray}// read them to understand the current state of reality.}\\

\vspace{8pt}


\section{CHAPTER 22: CREATIVE TOOLS}

\subsection{Page 89: Media Creation}

\vspace{4pt}\textbf{\color{cyan}<<< MEDIA CREATION — ART FROM THE COMMAND LINE >>>}\vspace{2pt}\\
\color{white}who needs photoshop when you have a terminal?\\
\color{white}The terminal is not just for text. ZPWR includes verbs\\
\color{white}for generating GIFs, rendering ASCII art, and creating\\
\color{white}visual content without ever opening a GUI application.\\
\vspace{4pt}\textbf{\color{magenta}CREATEGIF — SCREEN TO GIF}\\
\texttt{\color{green}zpwr creategif}\\
\color{white}Creates animated GIFs from your terminal sessions.\\
\color{white}Record your workflow, your demos, your flex.\\
\color{white}Perfect for README demos and pull request previews.\\
\color{white}Your terminal is the camera. The GIF is the film.\\
\color{white}No screen recording software needed.\\
\vspace{4pt}\textbf{\color{magenta}FIGLETFONTS — THE ASCII ART GALLERY}\\
\texttt{\color{green}zpwr figletfonts}\\
\color{white}Previews all available figlet fonts on your system.\\
\color{white}figlet turns text into massive ASCII art banners.\\
\color{white}This verb shows you every font so you can pick one.\\
\color{white}Ever wondered what your name looks like in 47\\
\color{white}different ASCII art styles? Now you know.\\
\color{white}Example output (one of many fonts):\\
\color{white}\_\_\_\_\_\_\_\_  \_\_\_\_      \_\_\_\_  \_\_\_\_  \_\_\_\_\\
\color{white}|\_\_\_  \_\_\_||  \_ \textbackslash{}\textbackslash{}    |  \_ \textbackslash{}\textbackslash{}|  \_ \textbackslash{}\textbackslash{}|  \_ \textbackslash{}\textbackslash{}\\
\color{white}/ /    | |\_) |   | |\_) | |\_) | |\_) |\\
\color{white}/ /     |  \_\_/    |  \_\_/|  \_ <|  \_ <\\
\color{white}/ /      | |       | |   | | \textbackslash{}\textbackslash{} | | \textbackslash{}\textbackslash{} \textbackslash{}\textbackslash{}\\
\color{white}/\_/       |\_|       |\_|   |\_|  \textbackslash{}\textbackslash{}|\_|  \textbackslash{}\textbackslash{}\_\textbackslash{}\textbackslash{}\\
\vspace{4pt}\textbf{\color{magenta}ANIMATE — TERMINAL ANIMATION}\\
\texttt{\color{green}zpwr animate}\\
\color{white}Runs terminal animations. Because sometimes the\\
\color{white}command line needs to look as cool as it feels.\\
\color{white}Eye candy for the operator who has everything.\\
\color{white}THE CREATIVE TOOLKIT:\\
\color{white}creategif       record terminal to GIF\\
\color{white}figletfonts     preview all ASCII art fonts\\
\color{white}animate         terminal animations\\
\color{white}banner          the ZPWR startup banner\\
\color{white}bannerlolcat    rainbow-colored banner\\
\textit{\color{gray}// the terminal is a canvas. most people see a spreadsheet.}\\
\textit{\color{gray}// hackers see a gallery. these verbs are your brushes.}\\

\vspace{8pt}

\subsection{Page 90: Color Systems}

\vspace{4pt}\textbf{\color{cyan}<<< COLOR SYSTEMS — PAINTING THE VOID >>>}\vspace{2pt}\\
\color{white}256 colors. true color. the terminal is not monochrome.\\
\color{white}Color is not cosmetic. It is information density.\\
\color{white}A colored diff tells you more in one glance than a\\
\color{white}monochrome one does in ten seconds. ZPWR gives you\\
\color{white}tools to explore, test, and master terminal color.\\
\vspace{4pt}\textbf{\color{magenta}COLORTEST — THE 16-COLOR PALETTE}\\
\texttt{\color{green}zpwr colortest}\\
\color{white}Displays the basic 16-color ANSI palette. The foundation.\\
\color{white}Black, red, green, yellow, blue, magenta, cyan, white —\\
\color{white}plus their bright variants. 16 soldiers in formation.\\
\color{white}Use this to verify your terminal theme is sane.\\
\color{white}If red looks orange, your theme needs fixing.\\
\vspace{4pt}\textbf{\color{magenta}COLORTEST256 — THE FULL SPECTRUM}\\
\texttt{\color{green}zpwr colortest256}\\
\color{white}Displays all 256 extended colors. A grid of every\\
\color{white}color code your terminal supports. From 0 to 255.\\
\color{white}Useful for picking colors for prompts and scripts.\\
\color{white}\textbackslash{}"What number is that teal?\textbackslash{}" Run this and find out.\\
\vspace{4pt}\textbf{\color{magenta}PYGMENTCOLORS — SYNTAX HIGHLIGHTING PALETTE}\\
\texttt{\color{green}zpwr pygmentcolors}\\
\color{white}Shows available Pygments color schemes. Pygments is the\\
\color{white}syntax highlighting engine used by many ZPWR tools.\\
\color{white}Pick a scheme that matches your aesthetic.\\
\vspace{4pt}\textbf{\color{magenta}PRETTYPRINT \& ALTPRETTYPRINT — FORMATTED OUTPUT}\\
\texttt{\color{green}zpwr prettyprint <file>}\\
\texttt{\color{green}zpwr altprettyprint <file>}\\
\color{white}Syntax-highlighted file display. prettyprint uses one\\
\color{white}engine, altprettyprint uses another. Both turn raw\\
\color{white}source code into technicolor glory in your terminal.\\
\color{white}Like cat, but for people who respect their eyes.\\
\color{white}COLOR EXPLORATION WORKFLOW:\\
\texttt{\color{green}1. zpwr colortest        check basic 16 colors}\\
\texttt{\color{green}2. zpwr colortest256     explore full 256 palette}\\
\texttt{\color{green}3. zpwr pygmentcolors    pick a syntax theme}\\
\texttt{\color{green}4. zpwr prettyprint      view files in color}\\
\textit{\color{gray}// monochrome terminals are a war crime}\\
\textit{\color{gray}// color is data. learn to read it.}\\
\textit{\color{gray}// 256 colors and you are still using white-on-black?}\\

\vspace{8pt}

\subsection{Page 91: Script Management}

\vspace{4pt}\textbf{\color{cyan}<<< SCRIPT MANAGEMENT — THE FORGE >>>}\vspace{2pt}\\
\color{white}create scripts. count scripts. print scripts to PDF. ship scripts.\\
\color{white}ZPWR manages its own script collection and gives you\\
\color{white}tools to create new ones, inventory them, and even\\
\color{white}export them to PDF for documentation or review.\\
\vspace{4pt}\textbf{\color{magenta}SCRIPTNEW — THE SCRIPT GENERATOR}\\
\texttt{\color{green}zpwr scriptnew <name>}\\
\color{white}Creates a new script with proper boilerplate. Shebang,\\
\color{white}error handling, argument parsing — all pre-wired.\\
\color{white}Stop writing boilerplate. Start writing logic.\\
\color{white}The script lands in \textbackslash{}PWR\_SCRIPTS, already executable,\\
\color{white}already on your PATH. Zero friction from idea to tool.\\
\vspace{4pt}\textbf{\color{magenta}SCRIPTTOPDF — THE DOCUMENTER}\\
\texttt{\color{green}zpwr scripttopdf}\\
\color{white}Converts scripts to syntax-highlighted PDFs.\\
\color{white}For code reviews, documentation, or printing out\\
\color{white}and framing on your wall like the art it is.\\
\color{white}Send your code to a printer. Assert dominance.\\
\vspace{4pt}\textbf{\color{magenta}SCRIPTCOUNT — THE CENSUS}\\
\texttt{\color{green}zpwr scriptcount}\\
\color{white}How many scripts live in \textbackslash{}PWR\_SCRIPTS? This tells you.\\
\color{white}Track the growth of your personal armory over time.\\
\vspace{4pt}\textbf{\color{magenta}SCRIPTLIST — THE INVENTORY}\\
\texttt{\color{green}zpwr scriptlist}\\
\color{white}Lists every script in the collection. Names, paths,\\
\color{white}the full manifest of your executable arsenal.\\
\color{white}THE SCRIPT LIFECYCLE:\\
\texttt{\color{green}1. zpwr scriptnew myutil     create the script}\\
\texttt{\color{green}2. zpwr vimscriptedit        edit scripts with vim/fzf}\\
\texttt{\color{green}3. zpwr scriptcount          check your inventory}\\
\texttt{\color{green}4. zpwr scriptlist           browse the collection}\\
\texttt{\color{green}5. zpwr scripttopdf          export for documentation}\\
\vspace{4pt}\textbf{\color{magenta}ALSO RELATED:}\\
\texttt{\color{green}zpwr makedir <name>     create a directory}\\
\texttt{\color{green}zpwr makefile <name>    create a file}\\
\color{white}Simple creation verbs. No boilerplate, just mkdir/touch\\
\color{white}with ZPWR's logging and validation baked in.\\
\textit{\color{gray}// a script unwritten is a problem unsolved}\\
\textit{\color{gray}// a script uncounted is a tool forgotten}\\
\textit{\color{gray}// the forge never closes. the armory always grows.}\\

\vspace{8pt}


\section{CHAPTER 23: EXERCISM \& EXTERNAL}

\subsection{Page 92: Exercism Integration}

\vspace{4pt}\textbf{\color{cyan}<<< EXERCISM INTEGRATION — LEVEL UP YOUR CODE >>>}\vspace{2pt}\\
\color{white}practice makes perfect. ZPWR makes practice frictionless.\\
\color{white}Exercism.org is a free coding practice platform with\\
\color{white}mentor-reviewed exercises in 50+ languages. ZPWR wraps\\
\color{white}the exercism CLI so you never leave your terminal.\\
\vspace{4pt}\textbf{\color{magenta}EXERCISM — THE LAUNCHER}\\
\texttt{\color{green}zpwr exercism}\\
\color{white}The base exercism command integration. Manages your\\
\color{white}exercism workspace, configuration, and submission\\
\color{white}workflow from within ZPWR.\\
\color{white}Your exercism workspace lives under \textbackslash{}PWR and is\\
\texttt{\color{green}accessible via zpwr homeexercism for quick navigation.}\\
\vspace{4pt}\textbf{\color{magenta}EXERCISMDOWNLOAD — FETCH EXERCISES}\\
\texttt{\color{green}zpwr exercismdownload}\\
\color{white}Downloads an exercise to your local workspace.\\
\color{white}The exercise appears, ready to solve. Tests included.\\
\color{white}No browser clicking. No manual downloads.\\
\color{white}The workflow:\\
\color{white}1. Browse exercises on exercism.org\\
\texttt{\color{green}2. zpwr exercismdownload to fetch it}\\
\color{white}3. Solve the exercise in your editor\\
\color{white}4. Run the tests locally\\
\color{white}5. Submit when you pass\\
\vspace{4pt}\textbf{\color{magenta}EXERCISMEDIT — SOLVE IN YOUR EDITOR}\\
\texttt{\color{green}zpwr exercismedit}\\
\color{white}Opens exercism exercises directly in your configured\\
\color{white}editor. Download, edit, test, submit — all from the\\
\color{white}terminal. The browser is optional.\\
\color{white}WHY EXERCISM + ZPWR:\\
\color{white}The exercism CLI is good. ZPWR makes it seamless.\\
\color{white}No context switching. No browser tabs. No friction.\\
\color{white}Languages supported by exercism:\\
\color{white}Python, Go, Rust, Elixir, JavaScript, TypeScript,\\
\color{white}C, C++, Java, Ruby, Haskell, Clojure, Kotlin,\\
\color{white}Swift, Bash, and 35+ more.\\
\textit{\color{gray}// coding practice should be as easy as}\\
\textit{\color{gray}// opening a terminal and typing three words}\\
\textit{\color{gray}// zpwr exercismdownload. done. now solve.}\\

\vspace{8pt}

\subsection{Page 93: Travis CI \& External Services}

\vspace{4pt}\textbf{\color{cyan}<<< TRAVIS CI \& EXTERNAL SERVICES >>>}\vspace{2pt}\\
\color{white}CI/CD from your terminal, not your browser\\
\color{white}ZPWR integrates with external services so you can\\
\color{white}monitor builds, check statuses, and manage CI/CD\\
\color{white}pipelines without opening a single browser tab.\\
\vspace{4pt}\textbf{\color{magenta}TRAVIS — THE CI DASHBOARD}\\
\texttt{\color{green}zpwr travis}\\
\color{white}Travis CI integration. View build status, logs, and\\
\color{white}pipeline information from your terminal.\\
\vspace{4pt}\textbf{\color{magenta}TRAVISBUILD — BUILD STATUS}\\
\texttt{\color{green}zpwr travisbuild}\\
\color{white}Check the current build status. Green, red, or pending.\\
\color{white}No need to refresh a web page in a loop.\\
\vspace{4pt}\textbf{\color{magenta}TRAVISBRANCH — BRANCH BUILDS}\\
\texttt{\color{green}zpwr travisbranch}\\
\color{white}View build status for a specific branch. When you push\\
\color{white}to a feature branch and want to know if CI passed.\\
\vspace{4pt}\textbf{\color{magenta}TRAVISPR — PULL REQUEST BUILDS}\\
\texttt{\color{green}zpwr travispr}\\
\color{white}Check build status for pull requests. The merge button\\
\color{white}is only as good as the CI behind it.\\
\vspace{4pt}\textbf{\color{magenta}OTHER EXTERNAL INTEGRATIONS:}\\
\texttt{\color{green}zpwr goclean          clean Go build cache}\\
\texttt{\color{green}zpwr google <query>    search Google from terminal}\\
\texttt{\color{green}zpwr digs <domain>     DNS lookups}\\
\texttt{\color{green}zpwr torip            check your Tor IP address}\\
\texttt{\color{green}zpwr toriprenew       renew your Tor circuit}\\
\texttt{\color{green}zpwr curl             enhanced curl operations}\\
\color{white}THE EXTERNAL SERVICES PHILOSOPHY:\\
\color{white}The browser is a context switch. Every tab you open\\
\color{white}is a distraction vector. ZPWR brings the services\\
\color{white}to you, in the terminal, where you already are.\\
\textit{\color{gray}// the terminal is the cockpit}\\
\textit{\color{gray}// external services are instruments on the dashboard}\\
\textit{\color{gray}// you read them. you do not visit them.}\\

\vspace{8pt}


\section{CHAPTER 24: TIPS, TRICKS \& EASTER EGGS}

\subsection{Page 94: The One-Letter Verbs}

\vspace{4pt}\textbf{\color{cyan}<<< THE ONE-LETTER VERBS — MAXIMUM POWER, MINIMUM KEYSTROKES >>>}\vspace{2pt}\\
\color{white}six letters that do the work of six paragraphs\\
\color{white}ZPWR has 408 verbs. Six of them are single characters.\\
\color{white}These are not abbreviations — they are the distilled\\
\color{white}essence of the most common operations. One keystroke\\
\color{white}each. The speed limit of human-shell interaction.\\
\vspace{4pt}\textbf{\color{magenta}a — AUTOLOAD}\\
\texttt{\color{green}zpwr a}\\
\color{white}Quick access to the autoload system. When you need\\
\color{white}to check on autoloaded functions fast.\\
\vspace{4pt}\textbf{\color{magenta}c — CAT}\\
\texttt{\color{green}zpwr c}\\
\color{white}The cat operation. View file contents. The most basic\\
\color{white}read operation, compressed to one letter.\\
\vspace{4pt}\textbf{\color{magenta}e — EDIT}\\
\texttt{\color{green}zpwr e}\\
\color{white}Open your editor. The gateway to modification.\\
\color{white}Uses your configured \textbackslash{}DITOR.\\
\vspace{4pt}\textbf{\color{magenta}f — FILESEARCH}\\
\texttt{\color{green}zpwr f}\\
\color{white}File search. Find things. The single most common\\
\color{white}operation after listing and editing.\\
\vspace{4pt}\textbf{\color{magenta}r — REPLAY}\\
\texttt{\color{green}zpwr r}\\
\color{white}Replay operations. Redo what you did before.\\
\color{white}History is not just for reading.\\
\vspace{4pt}\textbf{\color{magenta}z — ZCD}\\
\texttt{\color{green}zpwr z}\\
\color{white}Z-jump. Frecency-based directory navigation.\\
\color{white}Type a partial path, z figures out the rest.\\
\color{white}The single most addictive feature in ZPWR.\\
\color{white}THE ONE-LETTER ARSENAL:\\
\color{white}a   autoload system\\
\color{white}c   cat / view files\\
\color{white}e   edit files\\
\color{white}f   file search\\
\color{white}r   replay\\
\color{white}z   z-jump directory navigation\\
\textit{\color{gray}// every extra character is latency between thought and action}\\
\textit{\color{gray}// one letter is the event horizon of efficiency}\\
\textit{\color{gray}// below this, there is only telepathy}\\

\vspace{8pt}

\subsection{Page 95: Alias Chains}

\vspace{4pt}\textbf{\color{cyan}<<< ALIAS CHAINS — COMPOUND OPERATIONS >>>}\vspace{2pt}\\
\color{white}one verb triggers another. the cascade begins.\\
\color{white}ZPWR verbs do not exist in isolation. They chain.\\
\color{white}They compose. They call each other. This is not an\\
\color{white}accident — it is the architecture.\\
\vspace{4pt}\textbf{\color{magenta}THE SHORTHAND ALIASES:}\\
\color{white}ZPWR is aliased to shorter forms for speed:\\
\texttt{\color{green}zpwr   the full command}\\
\color{white}zp     shorter\\
\color{white}zpz    alternate short form\\
\color{white}All three are equivalent:\\
\texttt{\color{green}zpwr help  ==  zp help  ==  zpz help}\\
\vspace{4pt}\textbf{\color{magenta}COMPOUND VERB PATTERNS:}\\
\color{white}Many verbs are compounds that do two things:\\
\texttt{\color{green}filesearch + edit     = zpwr filesearchedit}\\
\texttt{\color{green}findsearch + edit     = zpwr findsearchedit}\\
\texttt{\color{green}locatesearch + edit   = zpwr locatesearchedit}\\
\texttt{\color{green}wordsearch + edit     = zpwr wordsearchedit}\\
\texttt{\color{green}vimrecent + cd        = zpwr vimrecentcd}\\
\texttt{\color{green}vimrecent + sudo      = zpwr vimrecentsudo}\\
\color{white}The pattern: action + modifier = compound verb\\
\color{white}Search, then edit. Browse, then cd. Open, then sudo.\\
\vspace{4pt}\textbf{\color{magenta}THE CLEAN CHAIN:}\\
\texttt{\color{green}zpwr cleanall             clean everything}\\
\texttt{\color{green}zpwr cleancache           clean caches only}\\
\texttt{\color{green}zpwr cleancompcache       clean completion cache}\\
\texttt{\color{green}zpwr cleangitcache        clean git repo cache}\\
\texttt{\color{green}zpwr cleangitcleancache   clean git clean cache}\\
\texttt{\color{green}zpwr cleangitdirtycache   clean git dirty cache}\\
\texttt{\color{green}zpwr cleantemp            clean temp files}\\
\texttt{\color{green}zpwr cleanrefreshupdate   clean + refresh + update}\\
\vspace{4pt}\textbf{\color{magenta}THE GIT CHAIN:}\\
\texttt{\color{green}zpwr gitforalldir     run git cmd in all dirs}\\
\texttt{\color{green}zpwr gitfordirmain    run git cmd in dirs on main}\\
\texttt{\color{green}zpwr gitzfordir       z-style git for all dirs}\\
\textit{\color{gray}// verbs are words. chains are sentences.}\\
\textit{\color{gray}// the grammar of ZPWR is combinatorial.}\\
\textit{\color{gray}// learn the words, and the sentences write themselves.}\\

\vspace{8pt}

\subsection{Page 96: The Token System}

\vspace{4pt}\textbf{\color{cyan}<<< THE TOKEN SYSTEM — YOUR PRIVATE OVERRIDE LAYER >>>}\vspace{2pt}\\
\color{white}secrets, aliases, and local config that never touch git\\
\color{white}ZPWR has sensible defaults. But you are not sensible.\\
\color{white}You have API keys, custom aliases, machine-specific\\
\color{white}paths, and opinions. The token system is where all\\
\color{white}of that lives — safely outside version control.\\
\vspace{4pt}\textbf{\color{magenta}THE TOKEN FILES:}\\
\color{white}\textbackslash{}PWR\_LOCAL/.tokens.sh\\
\color{white}Your main tokens file. Aliases, exports, secrets.\\
\color{white}This is sourced during shell init. Gitignored.\\
\color{white}\textbackslash{}PWR\_LOCAL/.tokens-pre.sh\\
\color{white}Sourced BEFORE ZPWR's main configuration.\\
\color{white}Override defaults before they are set.\\
\color{white}Set ZPWR\_* variables here to change behavior.\\
\color{white}\textbackslash{}PWR\_LOCAL/.tokens-post.sh\\
\color{white}Sourced AFTER ZPWR's main configuration.\\
\color{white}Override anything ZPWR set. Final word.\\
\color{white}Your aliases here trump everything.\\
\vspace{4pt}\textbf{\color{magenta}THE OVERRIDE PATTERN:}\\
\color{white}Want to change a ZPWR default? Do not edit ZPWR source.\\
\color{white}Put it in tokens-pre.sh (before) or tokens-post.sh (after).\\
\color{white}Example: override in tokens-pre.sh\\
\color{white}export ZPWR\_EXA\_COMMAND=eza\\
\color{white}export ZPWR\_KEYTIMEOUT=1\\
\color{white}export ZPWR\_BANNER\_CLEARLIST=false\\
\color{white}Example: add custom aliases in tokens-post.sh\\
\color{white}alias k=kubectl\\
\color{white}alias tf=terraform\\
\color{white}export GITHUB\_TOKEN=ghp\_xxxxx\\
\vspace{4pt}\textbf{\color{magenta}MANAGING TOKENS:}\\
\texttt{\color{green}zpwr tokens           view your tokens files}\\
\texttt{\color{green}zpwr vimtokens        edit tokens in vim}\\
\texttt{\color{green}zpwr emacstokens      edit tokens in emacs}\\
\color{white}THE LOAD ORDER:\\
\color{white}1. tokens-pre.sh      (your early overrides)\\
\color{white}2. ZPWR core config    (the framework defaults)\\
\color{white}3. tokens-post.sh     (your final overrides)\\
\color{white}4. .tokens.sh         (your secrets and aliases)\\
\color{white}NEVER commit token files. They are gitignored for a reason.\\
\color{white}API keys in git history haunt you forever.\\
\textit{\color{gray}// the token system is your private dimension}\\
\textit{\color{gray}// ZPWR updates without touching your config}\\
\textit{\color{gray}// your overrides survive every git pull}\\

\vspace{8pt}

\subsection{Page 97: FZF Everywhere}

\vspace{4pt}\textbf{\color{cyan}<<< FZF EVERYWHERE — THE FUZZY FINDER CONSPIRACY >>>}\vspace{2pt}\\
\color{white}fzf is not a tool. it is a religion. ZPWR is its cathedral.\\
\color{white}fzf (fuzzy finder) is woven into the fabric of ZPWR.\\
\color{white}Dozens of verbs pipe their output through fzf for\\
\color{white}interactive selection. If a verb ends in 'fzf', it is\\
\color{white}the fuzzy-interactive variant. But many others use it too.\\
\vspace{4pt}\textbf{\color{magenta}DEDICATED FZF VERBS:}\\
\texttt{\color{green}zpwr verbs     browse subcommands in fzf}\\
\texttt{\color{green}zpwr verbs           browse all verbs in fzf}\\
\vspace{4pt}\textbf{\color{magenta}VERBS THAT USE FZF INTERNALLY:}\\
\color{white}forgit* verbs:\\
\color{white}forgitadd, forgitdiff, forgitlog, forgitstash,\\
\color{white}forgitreset, forgitrestore, forgitclean\\
\color{white}All use fzf for interactive selection with previews.\\
\color{white}Search + edit verbs:\\
\color{white}filesearchedit, findsearchedit, locatesearchedit,\\
\color{white}wordsearchedit — search with fzf, open selection.\\
\color{white}Recent file verbs:\\
\color{white}vimrecent, vimrecentcd, emacsrecent, emacsrecentcd\\
\color{white}Browse recent files with fzf preview.\\
\color{white}Editor verbs:\\
\color{white}vimalledit, emacsalledit — select files to edit.\\
\color{white}Git tag verbs:\\
\color{white}vimgtagsedit, emacsgtagsedit — browse tags in fzf.\\
\vspace{4pt}\textbf{\color{magenta}THE FZF PATTERN:}\\
\color{white}The pattern is always the same:\\
\color{white}1. Generate a list of candidates\\
\color{white}2. Pipe through fzf with a preview command\\
\color{white}3. Take the selection and act on it\\
\color{white}This pattern turns every list into a menu.\\
\color{white}Every menu into a decision. Every decision into action.\\
\vspace{4pt}\textbf{\color{magenta}FZF CUSTOMIZATION:}\\
\color{white}Set FZF\_DEFAULT\_OPTS in your tokens file:\\
\color{white}export FZF\_DEFAULT\_OPTS='--height 80\% --reverse --border'\\
\textit{\color{gray}// if you can list it, you can fzf it}\\
\textit{\color{gray}// if you can fzf it, you can act on it}\\
\textit{\color{gray}// fzf is the universal adapter between data and decisions}\\

\vspace{8pt}


\section{CHAPTER 25: THE COMPLETE REFERENCE}

\subsection{Page 98: All 408 Verbs A-G}

\vspace{4pt}\textbf{\color{cyan}<<< ALL 408 VERBS: A through G (1/3) >>>}\vspace{2pt}\\
\color{white}the complete arsenal. memorize at your own peril.\\
\color{white}a                          quick autoload access\\
\color{white}about                      display ZPWR about info\\
\color{white}alacritty                  alacritty terminal config\\
\color{white}aliasrank                  rank aliases by usage frequency\\
\color{white}all                        show all ZPWR info\\
\color{white}allsearch                  search across all sources\\
\color{white}altprettyprint             alternate syntax-highlighted print\\
\color{white}animate                    terminal animations\\
\color{white}attach                     attach to tmux session\\
\color{white}autoload                   show autoload system\\
\color{white}autoloadcount              count autoloaded functions\\
\color{white}autoloadlist               list all autoloaded functions\\
\color{white}background                 run command in background\\
\color{white}backup                     backup ZPWR configuration\\
\color{white}backuphistory              backup shell history\\
\color{white}banner                     display ZPWR banner\\
\color{white}bannercounts               banner with counts\\
\color{white}bannerlolcat               rainbow banner via lolcat\\
\color{white}bannernopony               banner without ponysay\\
\color{white}bannerpony                 banner with ponysay\\
\color{white}bench                      benchmark shell startup\\
\color{white}brc                        source bashrc\\
\color{white}c                          cat / view file\\
\color{white}cat                        enhanced cat\\
\color{white}catcd                      cat file and cd to its dir\\
\color{white}catnviminforecentf         cat nvim info recent files\\
\color{white}catrecentfviminfo          cat recent files from viminfo\\
\color{white}catviminfo                 cat viminfo file\\
\color{white}cd                         enhanced cd\\
\color{white}cdup                       cd up one directory\\
\color{white}cfasd                      cd via fasd frecency\\
\color{white}changename                 rename files interactively\\
\color{white}changenamezpwr             rename ZPWR components\\
\color{white}cleanall                   clean everything\\
\color{white}cleancache                 clean caches\\
\color{white}cleancompcache             clean completion cache\\
\color{white}... continued. Run zpwr verbs for the complete list.\\
\color{white}Or use zpwr wizard --search <term> to find specific verbs.\\

\vspace{8pt}

\subsection{Page 99: All 408 Verbs H-S}

\vspace{4pt}\textbf{\color{cyan}<<< ALL 408 VERBS: forgitdiff through recompilefpath (2/3) >>>}\vspace{2pt}\\
\color{white}the middle of the alphabet is where the real work lives\\
\color{white}forgitdiff                 interactive diff via fzf\\
\color{white}forgitignore               browse gitignore templates\\
\color{white}forgitignoreclean          clean gitignore cache\\
\color{white}forgitignoreget            fetch gitignore template\\
\color{white}forgitignorelist           list gitignore templates\\
\color{white}forgitignoreupdate         update gitignore cache\\
\color{white}forgitinfo                 repo info dashboard\\
\color{white}forgitlog                  interactive git log via fzf\\
\color{white}forgitreset                interactive unstage via fzf\\
\color{white}forgitrestore              interactive restore via fzf\\
\color{white}forgitstash                interactive stash via fzf\\
\color{white}forgitwarn                 warn about repo hazards\\
\color{white}forked                     manage forked repos\\
\color{white}forward                    port forwarding\\
\color{white}fp                         fpath info\\
\color{white}funcrank                   rank functions by usage\\
\color{white}fwd                        forward shortcut\\
\color{white}gh                         GitHub CLI integration\\
\color{white}ghcontribcount             GitHub contribution count\\
\color{white}ghz                        GitHub + z integration\\
\color{white}gitaemail                  set git author email\\
\color{white}gitcemail                  set git committer email\\
\color{white}gitclearcache              clear git cache\\
\color{white}gitclearcommit             clear a git commit\\
\color{white}gitclearfile               clear file from git\\
\color{white}gitcommit                  enhanced git commit\\
\color{white}gitcommitcount             count commits\\
\color{white}gitcommits                 browse commits\\
\color{white}gitconfig                  show git config\\
\color{white}gitcontribcount            count contributions\\
\color{white}gitcontribcountdirs        count contribs per dir\\
\color{white}gitcontribcountlines       count contributed lines\\
\color{white}gitdir                     show git directory\\
\color{white}gitedittag                 edit a git tag\\
\color{white}gitemail                   show git email\\
\color{white}gitforalldir               git cmd in all dirs\\
\color{white}... continued. Run zpwr verbs for the complete list.\\
\color{white}Or use zpwr wizard --search <term> to find specific verbs.\\

\vspace{8pt}

\subsection{Page 100: All 408 Verbs T-Z + Epilogue}

\vspace{4pt}\textbf{\color{cyan}<<< ALL 408 VERBS: refreshzwc through zstyle (3/3) >>>}\vspace{2pt}\\
\color{white}the final stretch. the end of the alphabet. the beginning of mastery.\\
\color{white}refreshzwc                 refresh compiled wordcode\\
\color{white}regen                      regenerate single component\\
\color{white}regenall                   regenerate everything\\
\color{white}regenconfiglinks           regenerate config symlinks\\
\color{white}regenctags                 regenerate ctags\\
\color{white}regenenvcache              regenerate env cache\\
\color{white}regengitdirtyrepocache     regen dirty repo cache\\
\color{white}regengitrepocache          regen git repo cache\\
\color{white}regengtags                 regenerate gtags\\
\color{white}regengtagspygments         regen gtags + pygments\\
\color{white}regengtagstype             regen gtags by type\\
\color{white}regenhistory               regenerate history\\
\color{white}regenkeybindings           regen keybindings\\
\color{white}regenpowerline             regenerate powerline\\
\color{white}regenzsh                   regenerate zsh config\\
\color{white}reload                     reload shell config\\
\color{white}relpath                    get relative path\\
\color{white}rename                     rename files\\
\color{white}replacer                   find and replace\\
\color{white}replay                     replay commands\\
\color{white}repo                       repository info\\
\color{white}reset                      reset shell state\\
\color{white}restart                    restart shell\\
\color{white}restore                    restore configuration\\
\color{white}restorehistory             restore history backup\\
\color{white}return2                    return to previous dir\\
\color{white}reveal                     open repo in browser\\
\color{white}revealrecurse              reveal recursively\\
\color{white}reverse                    reverse input\\
\color{white}rmconfiglinks              remove config symlinks\\
\color{white}rvs                        reverse shortcut\\
\color{white}scriptcount                count scripts\\
\color{white}scriptlist                 list all scripts\\
\color{white}scriptnew                  create new script\\
\color{white}scripts                    browse scripts\\
\color{white}scripttopdf                convert script to PDF\\
\color{white}... continued. Run zpwr verbs for the complete list.\\
\color{white}Or use zpwr wizard --search <term> to find specific verbs.\\

\vspace{8pt}


\section{APPENDIX A: KEYBINDING REFERENCE}

\subsection{Page 101: Tmux Prefix Keys: Navigation \& Panes}

\vspace{4pt}\textbf{\color{cyan}<<< TMUX PREFIX KEYS: NAVIGATION \& PANES >>>}\vspace{2pt}\\
\color{white}pane traversal, splitting, resizing. the grid of power.\\
\vspace{4pt}\textbf{\color{magenta}PANE SELECTION (hjkl repeatable):}\\
\color{white}prefix h              select pane left\\
\color{white}prefix j              select pane down\\
\color{white}prefix k              select pane up\\
\color{white}prefix l              select pane right\\
\color{white}prefix ;              last-pane (toggle to previous)\\
\color{white}prefix o              select next pane in order\\
\color{white}prefix q              display-panes (5s timeout)\\
\vspace{4pt}\textbf{\color{magenta}PANE RESIZE (HJKL repeatable, 5 cells):}\\
\color{white}prefix H              resize pane left 5\\
\color{white}prefix J              resize pane down 5\\
\color{white}prefix K              resize pane up 5\\
\color{white}prefix L              resize pane right 5\\
\color{white}prefix Ctrl-Up/Down   resize pane by 1\\
\color{white}prefix Alt-Up/Down    resize pane by 5\\
\vspace{4pt}\textbf{\color{magenta}PANE SPLITTING:}\\
\color{white}prefix -              split vertical (same dir)\\
\color{white}prefix \textbackslash{}\textbackslash{}\textbackslash{}\textbackslash{}              split horizontal (same dir)\\
\color{white}prefix \textbackslash{}"              split vertical (default dir)\\
\color{white}prefix \%              split horizontal (default dir)\\
\color{white}prefix |              split horizontal\\
\color{white}prefix \_              split vertical\\
\vspace{4pt}\textbf{\color{magenta}PANE OPERATIONS:}\\
\color{white}prefix z              zoom/unzoom current pane\\
\color{white}prefix !              break pane into own window\\
\color{white}prefix x              kill current pane\\
\color{white}prefix m              mark current pane\\
\color{white}prefix +              split vertical then close prev pane\\
\color{white}prefix Y              split vertical then close prev pane\\
\color{white}prefix \{              rotate window\\
\color{white}prefix \}              swap pane down\\
\vspace{4pt}\textbf{\color{magenta}ARROW KEY SELECTION:}\\
\color{white}prefix Up/Down        select pane up/down\\
\color{white}prefix Left/Right     select pane left/right\\
\color{white}Alt-Up/Down           select pane (root table, no prefix)\\
\color{white}Alt-Left/Right        select pane (root table, no prefix)\\
\color{white}All hjkl/HJKL bindings are repeatable (-r flag).\\
\color{white}Arrow keys in root table work without prefix for speed.\\
\color{white}// the wizard navigates panes by instinct\\
\color{white}// hjkl is muscle memory, arrows are for civilians\\

\vspace{8pt}

\subsection{Page 102: Tmux Prefix Keys: Windows \& Sessions}

\vspace{4pt}\textbf{\color{cyan}<<< TMUX PREFIX KEYS: WINDOWS \& SESSIONS >>>}\vspace{2pt}\\
\color{white}window management, session juggling, the multiplexer core.\\
\vspace{4pt}\textbf{\color{magenta}WINDOW CREATION \& SELECTION:}\\
\color{white}prefix c              new window\\
\color{white}prefix a              last-window (toggle to previous)\\
\color{white}prefix 0-9            select window by number\\
\color{white}prefix '              select window by index prompt\\
\color{white}prefix n              next window (repeatable)\\
\color{white}prefix p              previous window (repeatable)\\
\color{white}prefix Ctrl-n         next window\\
\color{white}prefix Ctrl-p         previous window\\
\vspace{4pt}\textbf{\color{magenta}WINDOW OPERATIONS:}\\
\color{white}prefix ,              rename current window\\
\color{white}prefix \&              kill window (with confirm)\\
\color{white}prefix .              move window to target\\
\color{white}prefix f              find window by name\\
\color{white}prefix w              choose-tree window view\\
\color{white}prefix i              display-message (window info)\\
\color{white}prefix t              clock-mode\\
\vspace{4pt}\textbf{\color{magenta}SESSION MANAGEMENT:}\\
\color{white}prefix \textbackslash{}\$              rename current session\\
\color{white}prefix (              switch to previous client (repeatable)\\
\color{white}prefix )              switch to next client (repeatable)\\
\color{white}prefix d              detach client\\
\color{white}prefix s              choose-tree session view\\
\color{white}prefix Ctrl-z         suspend client\\
\vspace{4pt}\textbf{\color{magenta}RELOAD \& PERSIST:}\\
\color{white}prefix r              reload tmux config + display confirmation\\
\color{white}prefix Ctrl-s         save session (tmux-resurrect)\\
\color{white}prefix Ctrl-r         restore session (tmux-resurrect)\\
\vspace{4pt}\textbf{\color{magenta}ROOT TABLE SHORTCUTS:}\\
\color{white}Ctrl-\textbackslash{}\textbackslash{}\textbackslash{}\textbackslash{}               switch to next client (no prefix)\\
\color{white}Ctrl-]               switch to prev client (no prefix)\\
\vspace{4pt}\textbf{\color{magenta}META:}\\
\color{white}prefix ?              list all keys\\
\color{white}prefix :              command prompt\\
\color{white}prefix /              describe a key binding\\
\color{white}prefix \textasciitilde{}              show messages\\
\color{white}// sessions are parallel universes\\
\color{white}// the wizard keeps many realities open at once\\

\vspace{8pt}

\subsection{Page 103: Tmux Prefix Keys: Layouts \& Pane Configs}

\vspace{4pt}\textbf{\color{cyan}<<< TMUX PREFIX KEYS: LAYOUTS \& PANE CONFIGS >>>}\vspace{2pt}\\
\color{white}preconfigured pane arrangements. instant workspaces.\\
\vspace{4pt}\textbf{\color{magenta}ZPWR LAYOUT PRESETS:}\\
\color{white}prefix D              control-window layout\\
\color{white}prefix E              four vertical panes\\
\color{white}prefix F              four panes layout\\
\color{white}prefix G              eight panes layout\\
\color{white}prefix O              sixteen panes layout\\
\color{white}prefix R              thirty-two panes REPL layout\\
\color{white}prefix T              config files layout\\
\color{white}prefix M              learn layout\\
\color{white}Each layout sources a conf file from \textbackslash{}PWR/tmux/:\\
\color{white}D -> control-window.conf\\
\color{white}E -> fourVertical.conf\\
\color{white}F -> four-panes.conf\\
\color{white}G -> eight-panes.conf\\
\color{white}O -> sixteen-panes.conf\\
\color{white}R -> thirtytwo-panes-repl.conf\\
\color{white}T -> config-files.conf\\
\color{white}M -> learn.conf\\
\vspace{4pt}\textbf{\color{magenta}SYNCHRONIZE \& CUSTOMIZE:}\\
\color{white}prefix S              toggle synchronize-panes (type in all)\\
\color{white}prefix C              customize-mode\\
\vspace{4pt}\textbf{\color{magenta}BUILT-IN LAYOUTS (Alt+number):}\\
\color{white}prefix Alt-1          even-horizontal\\
\color{white}prefix Alt-2          even-vertical\\
\color{white}prefix Alt-3          main-horizontal\\
\color{white}prefix Alt-4          main-vertical\\
\color{white}prefix Alt-5          tiled\\
\color{white}prefix Alt-6          main-horizontal-mirrored\\
\color{white}prefix Alt-7          main-vertical-mirrored\\
\vspace{4pt}\textbf{\color{magenta}WINDOW ROTATION:}\\
\color{white}prefix Ctrl-o         rotate window panes\\
\color{white}prefix Alt-o          rotate window panes (reverse)\\
\color{white}Synchronize-panes (prefix S) sends keystrokes to ALL panes.\\
\color{white}Perfect for running the same command across many servers.\\
\color{white}// sixteen panes is not chaos\\
\color{white}// sixteen panes is the wizard's dashboard\\

\vspace{8pt}

\subsection{Page 104: Tmux Prefix Keys: Special Operations}

\vspace{4pt}\textbf{\color{cyan}<<< TMUX PREFIX KEYS: SPECIAL OPERATIONS >>>}\vspace{2pt}\\
\color{white}multi-pane execution, URL picking, buffers, plugins.\\
\vspace{4pt}\textbf{\color{magenta}MULTI-PANE EXECUTION:}\\
\color{white}prefix Space          run command in all panes (single)\\
\color{white}prefix v              run command in all panes (multi)\\
\color{white}prefix b              open URL/path in all panes\\
\color{white}prefix g              google search in all panes\\
\color{white}These source \textbackslash{}PWR/scripts/allPanes.zsh with different modes.\\
\color{white}Space -> single mode (one command per pane)\\
\color{white}v     -> multi mode (batch operations)\\
\color{white}b     -> single open (open URL in each pane)\\
\color{white}g     -> single google (search in each pane)\\
\vspace{4pt}\textbf{\color{magenta}URL HANDLING:}\\
\color{white}prefix e              fzf URL picker -> open\\
\color{white}prefix u              fzf URL picker -> search\\
\color{white}Both use tmux-fzf-url plugin with 30s timeout.\\
\vspace{4pt}\textbf{\color{magenta}PASTE \& BUFFERS:}\\
\color{white}prefix P              paste buffer\\
\color{white}prefix ]              paste buffer (bracketed)\\
\color{white}prefix =              choose-buffer (fzf style)\\
\color{white}prefix \#              list buffers\\
\color{white}prefix Ctrl-v         paste from system clipboard\\
\vspace{4pt}\textbf{\color{magenta}COPY MODE:}\\
\color{white}prefix [              enter copy mode\\
\color{white}prefix PgUp           enter copy mode and page up\\
\vspace{4pt}\textbf{\color{magenta}PLUGIN MANAGEMENT (TPM):}\\
\color{white}prefix I              install plugins\\
\color{white}prefix U              update plugins\\
\color{white}prefix Alt-u          clean unused plugins\\
\vspace{4pt}\textbf{\color{magenta}SEND PREFIX:}\\
\color{white}prefix Ctrl-a         send prefix key to inner tmux\\
\color{white}prefix Ctrl-]         send literal Ctrl-]\\
\color{white}Use prefix Ctrl-a when nesting tmux sessions.\\
\color{white}The outer session captures prefix, Ctrl-a forwards it.\\
\color{white}// buffers are the wizard's memory\\
\color{white}// every copy is a spell stored for later\\

\vspace{8pt}

\subsection{Page 105: Tmux Copy Mode (vi)}

\vspace{4pt}\textbf{\color{cyan}<<< TMUX COPY MODE (VI BINDINGS) >>>}\vspace{2pt}\\
\color{white}scrollback navigation, selection, copy, search. pure vi.\\
\vspace{4pt}\textbf{\color{magenta}ENTER/EXIT:}\\
\color{white}prefix [              enter copy mode\\
\color{white}q                     cancel and exit copy mode\\
\color{white}Escape                clear selection\\
\vspace{4pt}\textbf{\color{magenta}SELECTION \& COPY:}\\
\color{white}v                     begin selection\\
\color{white}V                     select entire line\\
\color{white}Ctrl-v               toggle rectangle selection\\
\color{white}Space                 begin selection (alternate)\\
\color{white}y                     copy selection to clipboard (no clear)\\
\color{white}Enter                 copy selection to clipboard and cancel\\
\color{white}D                     copy to end of line and cancel\\
\color{white}A                     append selection and cancel\\
\vspace{4pt}\textbf{\color{magenta}NAVIGATION:}\\
\color{white}h/j/k/l              cursor movement\\
\color{white}Ctrl-h/j/k/l         move cursor 4 cells at a time\\
\color{white}w/b/e                word forward/backward/end\\
\color{white}W/B/E                WORD (space-delimited) movement\\
\color{white}0 / \textbackslash{}\$                start/end of line\\
\color{white}\textasciicircum{}                     first non-blank character\\
\color{white}g / G                top / bottom of history\\
\color{white}H / M / L            top / middle / bottom of screen\\
\vspace{4pt}\textbf{\color{magenta}SCROLLING:}\\
\color{white}Ctrl-u / Ctrl-d      half page up / down\\
\color{white}Ctrl-b / Ctrl-f      full page up / down\\
\color{white}Ctrl-y / Ctrl-e      scroll up / down one line\\
\color{white}J / K                scroll down / up one line\\
\vspace{4pt}\textbf{\color{magenta}SEARCH:}\\
\color{white}/                     search forward (incremental)\\
\color{white}?                     search backward (incremental)\\
\color{white}n / N                next / previous match\\
\color{white}*                     search forward for word under cursor\\
\color{white}\#                     search backward for word under cursor\\
\vspace{4pt}\textbf{\color{magenta}ZPWR EXTENSIONS:}\\
\color{white}s                     copy + google search selection\\
\color{white}x                     copy + open URL selection\\
\color{white}z                     copy + google search selection\\
\color{white}RightClick            copy + google search selection\\
\color{white}Copy goes to system clipboard via reattach-to-user-namespace pbcopy.\\
\color{white}ZPWR extends vi copy mode with instant google/open actions.\\
\color{white}// in copy mode the wizard moves through time\\
\color{white}// scrollback is archaeology, search is divination\\

\vspace{8pt}

\subsection{Page 106: Zsh Insert Mode: Basic \& Movement}

\vspace{4pt}\textbf{\color{cyan}<<< ZSH INSERT MODE: BASIC \& MOVEMENT >>>}\vspace{2pt}\\
\color{white}the primary editing mode. where commands are born.\\
\vspace{4pt}\textbf{\color{magenta}LINE MOVEMENT:}\\
\color{white}Ctrl-A                beginning-of-line\\
\color{white}Ctrl-E                end-of-line\\
\color{white}Ctrl-O                EOL or next tab stop\\
\vspace{4pt}\textbf{\color{magenta}EDITING:}\\
\color{white}Ctrl-U                clear entire line\\
\color{white}Ctrl-W                delete word backward (autopair)\\
\color{white}Ctrl-H                delete char backward (autopair)\\
\color{white}Ctrl-Y                change surrounding quotes\\
\color{white}Ctrl-Z                undo\\
\color{white}Ctrl-R                redo\\
\color{white}Ctrl-Q                duplicate last word\\
\vspace{4pt}\textbf{\color{magenta}ZPWR SPECIAL:}\\
\color{white}Space                 zpwr supernatural expand (aliases+globs)\\
\color{white}Ctrl-@                zpwr expand + terminate space\\
\color{white}Ctrl-M (Enter)        zpwr magic enter (show status on empty)\\
\color{white}.                     zpwr rationalize dot (.. -> ../..)\\
\color{white}Ctrl-B                zpwr clipboard paste\\
\color{white}Ctrl-S                git add commit push (no check)\\
\color{white}Ctrl-K                zsh-learn-Learn (flashcard)\\
\vspace{4pt}\textbf{\color{magenta}COMPLETION \& FZF:}\\
\color{white}Tab (Ctrl-I)          fzf-completion\\
\color{white}Ctrl-D                list-choices\\
\color{white}Ctrl-L                clear-screen\\
\color{white}Ctrl-T                fzf file widget\\
\color{white}Alt-c                 fzf cd widget\\
\vspace{4pt}\textbf{\color{magenta}HISTORY:}\\
\color{white}Ctrl-N                down-history\\
\color{white}Ctrl-P                up-history\\
\color{white}Up/Down arrows        history substring search up/down\\
\vspace{4pt}\textbf{\color{magenta}MODE SWITCHING:}\\
\color{white}Escape                enter vi command mode\\
\color{white}jf / fj              enter vi command mode (chord escape)\\
\vspace{4pt}\textbf{\color{magenta}SURROUND \& PAIRS:}\\
\color{white}\textbackslash{}" ' ( [ \{ \textbackslash{}`          auto-pair/surround insert\\
\color{white}) ] \}                 autopair close\\
\color{white}Backspace             autopair-aware delete\\
\color{white}Space is the most powerful key: it expands aliases,\\
\color{white}global aliases, and abbreviations before inserting.\\
\color{white}// insert mode is where spells are typed\\
\color{white}// every keystroke is an incantation\\

\vspace{8pt}

\subsection{Page 107: Zsh Insert Mode: Ctrl-F \& Ctrl-V Chords}

\vspace{4pt}\textbf{\color{cyan}<<< ZSH INSERT MODE: CTRL-F \& CTRL-V CHORDS >>>}\vspace{2pt}\\
\color{white}two-key combos that unlock the full zpwr arsenal.\\
\vspace{4pt}\textbf{\color{magenta}CTRL-F CHORDS (fzf \& actions):}\\
\color{white}Ctrl-F Ctrl-D         pipe buffer into fzf\\
\color{white}Ctrl-F Ctrl-F         cd + vim fzf file search (accept)\\
\color{white}Ctrl-F Ctrl-G         google search\\
\color{white}Ctrl-F Ctrl-H         lsof into fzf\\
\color{white}Ctrl-F Ctrl-I         ag (silver searcher) into fzf\\
\color{white}Ctrl-F Ctrl-J         zpwr verbs widget (accept)\\
\color{white}Ctrl-F Ctrl-K         alternate quotes around buffer\\
\color{white}Ctrl-F Ctrl-L         list choices\\
\color{white}Ctrl-F Ctrl-M         zzcomplete\\
\color{white}Ctrl-F Ctrl-N         zpwr verbs widget (browse)\\
\color{white}Ctrl-F Ctrl-O         open URL from buffer\\
\color{white}Ctrl-F Ctrl-P         basic sed substitution\\
\color{white}Ctrl-F Ctrl-R         wrap buffer as variable\\
\color{white}Ctrl-F Ctrl-S         git diff check (gacp)\\
\color{white}Ctrl-F Ctrl-V         edit command line in \textbackslash{}DITOR\\
\color{white}Ctrl-F j              zpwr verbs widget (accept)\\
\color{white}Ctrl-F n              zpwr verbs widget (browse)\\
\vspace{4pt}\textbf{\color{magenta}CTRL-V CHORDS (vim/fzf/navigation):}\\
\color{white}Ctrl-V Ctrl-@         vim fzf (open file in vim)\\
\color{white}Ctrl-V Space          vim fzf (open file in vim)\\
\color{white}Ctrl-V Ctrl-F         fasd fzf (frecency file open)\\
\color{white}Ctrl-V Ctrl-G         fzf cd widget\\
\color{white}Ctrl-V Ctrl-K         emacs fzf\\
\color{white}Ctrl-V Ctrl-L         z fzf (frecency dirs)\\
\color{white}Ctrl-V Ctrl-N         vim fzf sudo\\
\color{white}Ctrl-V Ctrl-O         fzf tab complete\\
\color{white}Ctrl-V Ctrl-P         sudo command line toggle\\
\color{white}Ctrl-V Ctrl-R         history search multi-word\\
\color{white}Ctrl-V Ctrl-S         z fzf (alternate)\\
\color{white}Ctrl-V Ctrl-V         vim all files widget (accept)\\
\color{white}Ctrl-V Ctrl-Z         fzf history widget\\
\color{white}Ctrl-V ,              fzf environment variables\\
\color{white}Ctrl-V .              fzf all keybindings\\
\textit{\color{gray}Ctrl-V //             locate fzf}\\
\color{white}Ctrl-V /.             locate fzf + edit\\
\color{white}Ctrl-V ;              fzf surround\\
\color{white}Ctrl-V c              fzf git commits\\
\color{white}Ctrl-V k              fzf vim keybindings\\
\color{white}Ctrl-V m              vim all files widget (browse)\\
\color{white}Ctrl-V v              vim all files widget (accept)\\
\color{white}Ctrl-F is the action chord: search, edit, google, diff.\\
\color{white}Ctrl-V is the vim/fzf chord: open files, browse, navigate.\\
\color{white}// chords are the wizard's combos\\
\color{white}// two keys unlock what single keys cannot reach\\

\vspace{8pt}

\subsection{Page 108: Zsh Visual \& Menu Select Modes}

\vspace{4pt}\textbf{\color{cyan}<<< ZSH VISUAL \& MENU SELECT MODES >>>}\vspace{2pt}\\
\color{white}text objects, selection, and completion menu navigation.\\
\vspace{4pt}\textbf{\color{magenta}VISUAL MODE (enter with v/V in vicmd):}\\
\color{white}Escape                deactivate region\\
\color{white}j / k                move selection down / up\\
\color{white}o                     exchange point and mark\\
\color{white}p                     put replace selection\\
\color{white}x                     delete selection\\
\color{white}u                     downcase selection\\
\color{white}U                     upcase selection\\
\color{white}\textasciitilde{}                     swap case of selection\\
\vspace{4pt}\textbf{\color{magenta}TEXT OBJECTS (a = around, i = inner):}\\
\color{white}a\textbackslash{}" / i\textbackslash{}"              select around/inner double quotes\\
\color{white}a' / i'              select around/inner single quotes\\
\color{white}a\textbackslash{}` / i\textbackslash{}`              select around/inner backticks\\
\color{white}a( / i(              select around/inner parentheses\\
\color{white}a[ / i[              select around/inner brackets\\
\color{white}a\{ / i\{              select around/inner braces\\
\color{white}a< / i<              select around/inner angle brackets\\
\color{white}aB / iB              select around/inner block\\
\color{white}aw / iw              select a word / in word\\
\color{white}aW / iW              select a blank word / in blank word\\
\color{white}aa / ia              select a shell word / in shell word\\
\color{white}Extended quote chars also work as text objects:\\
\color{white}+ , - . / : ; = @ | \textbackslash{}\textbackslash{}\textbackslash{}\textbackslash{} are all valid quote delimiters.\\
\vspace{4pt}\textbf{\color{magenta}MENU SELECT MODE (during completion):}\\
\color{white}Ctrl-J / Ctrl-K       move down / up in menu\\
\color{white}Ctrl-H / Ctrl-L       move left / right in menu\\
\color{white}Tab / Ctrl-I          move forward (next item)\\
\color{white}Ctrl-S / Shift-Tab    reverse menu complete\\
\color{white}Ctrl-@ (Ctrl-Space)   accept line\\
\color{white}Ctrl-D                accept and menu complete (stay open)\\
\color{white}Ctrl-F                accept and infer next\\
\color{white}Ctrl-M (Enter)        accept line (.accept-line)\\
\color{white}Ctrl-V                insert mode within menu\\
\color{white}Ctrl-N / Ctrl-P       forward / backward word in menu\\
\color{white}Ctrl-X                incremental search forward in menu\\
\color{white}?                     incremental search backward in menu\\
\color{white}.                     self-insert (exit menu + type dot)\\
\color{white}Backspace             undo last menu selection\\
\color{white}Menu select uses hjkl-style navigation with Ctrl prefix.\\
\color{white}Ctrl-D is key: accept current match but keep menu open.\\
\color{white}// visual mode is precision surgery\\
\color{white}// select only what you need, transform only what you touch\\

\vspace{8pt}

\subsection{Page 109: Vim Normal Mode Highlights}

\vspace{4pt}\textbf{\color{cyan}<<< VIM NORMAL MODE: ZPWR HIGHLIGHTS >>>}\vspace{2pt}\\
\color{white}leader mappings, fzf integration, tmux synergy.\\
\color{white}leader key is Space. only zpwr-specific bindings listed.\\
\vspace{4pt}\textbf{\color{magenta}FILE \& BUFFER (Space prefix):}\\
\color{white}Space w               :w! (force write)\\
\color{white}Space e               :q! (force quit)\\
\color{white}Space q               :qa! (quit all)\\
\color{white}Space wq              :wq! (write and quit)\\
\color{white}Space n / Space p     next / prev buffer\\
\color{white}Space an / Space ap   next / prev in arglist\\
\color{white}Space Tab             alternate or next buffer\\
\color{white}Space t               new tab\\
\color{white}Space s / Space v     horizontal / vertical split\\
\vspace{4pt}\textbf{\color{magenta}FZF INTEGRATION (Space prefix):}\\
\color{white}Space f               :Files! (fzf files)\\
\color{white}Space b               :Buffers! (fzf buffers)\\
\color{white}Space j               :Lines! (fzf lines)\\
\color{white}Space ag              :Ag (silver searcher)\\
\color{white}Space aa              :Agg (ag variant)\\
\color{white}Space rg              :Rg (ripgrep)\\
\color{white}Space z               :Commits! (fzf git commits)\\
\color{white}Space ;               :Commands! (fzf commands)\\
\color{white}Space ,               :FZFMaps (fzf all mappings)\\
\color{white}Space ke / , / m / i  fzf keys / maps / nmap / imap\\
\color{white}Space ta / ma / env  fzf tags / marks / env vars\\
\color{white}Space / / .           locate all / colorscheme picker\\
\vspace{4pt}\textbf{\color{magenta}HISTORY (Space prefix):}\\
\color{white}Space hh              :History! (fzf file history)\\
\color{white}Space h/              :History/! (fzf search history)\\
\color{white}Space h:              :History:! (fzf command history)\\
\vspace{4pt}\textbf{\color{magenta}SEARCH \& REPLACE (Space prefix):}\\
\color{white}Space r               replace word under cursor (whole file)\\
\color{white}Space g               global substitute word under cursor\\
\color{white}Space x               repeat last substitution on line\\
\color{white}n / N                 next / prev search (centered)\\
\color{white}z/                    toggle auto-highlight word under cursor\\
\vspace{4pt}\textbf{\color{magenta}GIT \& TMUX:}\\
\color{white}Space hs              GitGutter stage hunk\\
\color{white}Space hu              GitGutter undo hunk\\
\color{white}Space hp              GitGutter preview hunk\\
\color{white}[c / ]c              prev / next git hunk\\
\color{white}Space vj              save + TmuxRepeat file\\
\color{white}Space os / ox / oc   google / open URL / copy word\\
\vspace{4pt}\textbf{\color{magenta}MINIMAP \& COMMENTS:}\\
\color{white}Space mm              open minimap\\
\color{white}Space mt              toggle minimap\\
\color{white}Space cc              NERDCommenter comment\\
\color{white}Space c Space         NERDCommenter toggle\\
\color{white}Space cu              NERDCommenter uncomment\\
\color{white}// the leader key is Space because power needs room\\
\color{white}// every binding is one thought away from execution\\

\vspace{8pt}

\subsection{Page 110: Vim Insert \& Visual Mode + Keybinding Search}

\vspace{4pt}\textbf{\color{cyan}<<< VIM INSERT \& VISUAL MODE + KEYBINDING SEARCH >>>}\vspace{2pt}\\
\color{white}vim mode highlights and how to find any keybinding.\\
\vspace{4pt}\textbf{\color{magenta}VIM INSERT MODE HIGHLIGHTS:}\\
\color{white}Ctrl-A                jump to first non-blank (like \textasciicircum{})\\
\color{white}Ctrl-E                end of line or dismiss popup\\
\color{white}Ctrl-D                delete char forward or default\\
\color{white}Ctrl-F                move right or scroll\\
\color{white}Ctrl-S                isurround (surround insert)\\
\color{white}Ctrl-G s / S          isurround / ISurround (line)\\
\color{white}Ctrl-Tab              list ultisnips snippets\\
\color{white}Shift-Tab             SuperTab backward completion\\
\color{white}Alt-d                 delete word forward\\
\color{white}\textbackslash{}" ' ( ) [ ] \textbackslash{}` \{ \}    delimitMate auto-pairing\\
\vspace{4pt}\textbf{\color{magenta}VIM VISUAL MODE HIGHLIGHTS:}\\
\color{white}Space r               replace word in selection range\\
\color{white}Space g               global substitute in selection range\\
\color{white}Space x               repeat last sub on selected lines\\
\color{white}Space \textbackslash{}" ' \textbackslash{}` ( [ \{    surround selection with quote/bracket\\
\color{white}Space b               copy selection to tmux buffer\\
\color{white}Space = / Space -     grow / shrink selection (4 lines)\\
\color{white}J / K                move selected lines down / up\\
\color{white}Y                     yank and move to end\\
\color{white}< / >                indent / outdent (stay in visual)\\
\color{white}S                     vim-surround visual surround\\
\color{white}go                    copy selection + open as URL\\
\color{white}gs                    copy selection + google search\\
\color{white}Space vl / vk / vj   TmuxRepeat repl / visual / file\\
\color{white}Space em / ev / ef   extract method / variable / fold\\
\vspace{4pt}\textbf{\color{magenta}SNEAK \& MOTION:}\\
\color{white}s / S                sneak forward / backward (2-char)\\
\color{white}f / F / t / T        sneak-enhanced find char\\
\color{white}; / ,                repeat / reverse last sneak/find\\
\color{white}w / b / e            CamelCaseMotion word movement\\
\vspace{4pt}\textbf{\color{cyan}========================================}\vspace{2pt}\\
\vspace{4pt}\textbf{\color{magenta}HOW TO SEARCH KEYBINDINGS:}\\
\vspace{4pt}\textbf{\color{cyan}========================================}\vspace{2pt}\\
\color{white}zpwr allsearch        search ALL keybindings in fzf\\
\color{white}zpwr vimsearch        search vim keybindings in fzf\\
\color{white}zpwr zshsearch        search zsh keybindings in fzf\\
\vspace{4pt}\textbf{\color{magenta}KEYBINDING CACHE FILES:}\\
\color{white}\textbackslash{}PWR\_LOCAL/zpwrAllKeybindings.txt\\
\color{white}\textbackslash{}PWR\_LOCAL/zpwr\{Vim,Zsh,Tmux\}Keybindings.txt\\
\vspace{4pt}\textbf{\color{magenta}REGENERATE:}\\
\color{white}zpwr regenkeybindings  regenerate all keybinding caches\\
\color{white}Cache files are auto-generated. Run regenkeybindings after\\
\color{white}adding new bindings to refresh the searchable index.\\
\color{white}// the wizard does not memorize 4000 keybindings\\
\color{white}// the wizard searches them at the speed of thought\\
\color{white}// fzf is the oracle, the cache is the grimoire\\

\vspace{8pt}


\section{APPENDIX B: COMPLETE KEYBINDING DUMP}

\subsection{Page 111: ALL Tmux Prefix Keys (1/4)}

\vspace{4pt}\textbf{\color{cyan}=== ALL TMUX PREFIX KEYS (1/4) ===}\vspace{2pt}\\
\color{white}bind-key  -T prefix  Space  run-shell -b 'ZPWR/scripts/allPanes.zsh single\\
\color{white}bind-key  -T prefix  .  break-pane\\
\color{white}bind-key  -T prefix  \textbackslash{}\textbackslash{}'  split-window\\
\color{white}bind-key  -T prefix  \textbackslash{}\textbackslash{}\#  list-buffers\\
\color{white}bind-key  -T prefix  \textbackslash{}\textbackslash{}  command-prompt -I '\#S' \{ rename-session '\%\%' \}\\
\color{white}bind-key  -T prefix  \textbackslash{}\textbackslash{}\%  split-window -h\\
\color{white}bind-key  -T prefix  \&  confirm-before -p 'kill-window \#W? (y/n)' kill-win\\
\color{white}bind-key  -T prefix  \textbackslash{}\textbackslash{}'  command-prompt -T window-target -p index \{ selec\\
\color{white}bind-key -r -T prefix  (  switch-client -p\\
\color{white}bind-key -r -T prefix  )  switch-client -n\\
\color{white}bind-key  -T prefix  +  split-window -v -c '\#\{pane\_current\_path\}' \textbackslash{}\textbackslash{}; sele\\
\color{white}bind-key  -T prefix  ,  command-prompt -I '\#W' \{ rename-window '\%\%' \}\\
\color{white}bind-key  -T prefix  -  split-window -v -c '\#\{pane\_current\_path\}'\\
\color{white}bind-key  -T prefix  .  command-prompt -T target \{ move-window -t '\%\%' \}\\
\color{white}bind-key  -T prefix  /  command-prompt -k -p key \{ list-keys -1N '\%\%' \}\\
\color{white}bind-key  -T prefix  0  select-window -t :=0\\
\color{white}bind-key  -T prefix  1  select-window -t :=1\\
\color{white}bind-key  -T prefix  2  select-window -t :=2\\
\color{white}bind-key  -T prefix  3  select-window -t :=3\\
\color{white}bind-key  -T prefix  4  select-window -t :=4\\
\color{white}bind-key  -T prefix  5  select-window -t :=5\\
\color{white}bind-key  -T prefix  6  select-window -t :=6\\
\color{white}bind-key  -T prefix  7  select-window -t :=7\\
\color{white}bind-key  -T prefix  8  select-window -t :=8\\
\color{white}bind-key  -T prefix  9  select-window -t :=9\\
\color{white}bind-key  -T prefix  :  command-prompt\\
\color{white}bind-key  -T prefix  \textbackslash{}\textbackslash{};  last-pane\\
\color{white}bind-key  -T prefix  <  display-menu -T '\#[align=centre]\#\{window\_index\}:\#\{\\
\color{white}bind-key  -T prefix  =  choose-buffer -Z\\
\color{white}bind-key  -T prefix  >  display-menu -T '\#[align=centre]\#\{pane\_index\} (\#\{p\\
\color{white}bind-key  -T prefix  ?  list-keys -N\\
\color{white}bind-key  -T prefix  C  customize-mode -Z\\
\color{white}bind-key  -T prefix  D  source-file ZPWR/tmux/control-window.conf\\
\color{white}bind-key  -T prefix  E  source-file ZPWR/tmux/fourVertical.conf\\
\color{white}bind-key  -T prefix  F  source-file ZPWR/tmux/four-panes.conf\\
\color{white}bind-key  -T prefix  G  source-file ZPWR/tmux/eight-panes.conf\\
\color{white}bind-key -r -T prefix  H  resize-pane -L 5\\
\color{white}bind-key  -T prefix  I  run-shell HOME/.tmux/plugins/tpm/bindings/install\_\\
\color{white}bind-key -r -T prefix  J  resize-pane -D 5\\
\color{white}bind-key -r -T prefix  K  resize-pane -U 5\\
\color{white}bind-key -r -T prefix  L  resize-pane -R 5\\
\color{white}bind-key  -T prefix  M  source-file ZPWR/tmux/learn.conf\\
\color{white}bind-key  -T prefix  O  source-file ZPWR/tmux/sixteen-panes.conf\\
\color{white}bind-key  -T prefix  P  paste-buffer\\
\color{white}bind-key  -T prefix  R  source-file ZPWR/tmux/thirtytwo-panes-repl.conf\\
\color{white}bind-key  -T prefix  S  set-window-option synchronize-panes\\
\color{white}bind-key  -T prefix  T  source-file ZPWR/tmux/config-files.conf\\
\color{white}bind-key  -T prefix  U  run-shell HOME/.tmux/plugins/tpm/bindings/update\_p\\

\vspace{8pt}

\subsection{Page 112: ALL Tmux Prefix Keys (2/4)}

\vspace{4pt}\textbf{\color{cyan}=== ALL TMUX PREFIX KEYS (2/4) ===}\vspace{2pt}\\
\color{white}bind-key  -T prefix  Y  split-window -v -c '\#\{pane\_current\_path\}' \textbackslash{}\textbackslash{}; sele\\
\color{white}bind-key  -T prefix  [  copy-mode\\
\color{white}bind-key  -T prefix  \textbackslash{}\textbackslash{}\textbackslash{}\textbackslash{}  split-window -h -c '\#\{pane\_current\_path\}'\\
\color{white}bind-key  -T prefix  ]  paste-buffer -p\\
\color{white}bind-key  -T prefix  \_  split-window -v\\
\color{white}bind-key  -T prefix  a  last-window\\
\color{white}bind-key  -T prefix  b  run-shell -b 'ZPWR/scripts/allPanes.zsh single ope\\
\color{white}bind-key  -T prefix  c  new-window\\
\color{white}bind-key  -T prefix  d  detach-client\\
\color{white}bind-key  -T prefix  e  run-shell -b 'HOME/.tmux/plugins/tmux-fzf-url/fzf-\\
\color{white}bind-key  -T prefix  f  command-prompt \{ find-window -Z '\%\%' \}\\
\color{white}bind-key  -T prefix  g  run-shell -b 'ZPWR/scripts/allPanes.zsh single goo\\
\color{white}bind-key -r -T prefix  h  select-pane -L\\
\color{white}bind-key  -T prefix  i  display-message\\
\color{white}bind-key -r -T prefix  j  select-pane -D\\
\color{white}bind-key -r -T prefix  k  select-pane -U\\
\color{white}bind-key -r -T prefix  l  select-pane -R\\
\color{white}bind-key  -T prefix  m  select-pane -m\\
\color{white}bind-key -r -T prefix  n  next-window\\
\color{white}bind-key  -T prefix  o  select-pane -t :.+\\
\color{white}bind-key -r -T prefix  p  previous-window\\
\color{white}bind-key  -T prefix  q  display-panes -d 5000\\
\color{white}bind-key  -T prefix  r  source-file ZPWR/tmux/init.conf \textbackslash{}\textbackslash{}; display-messag\\
\color{white}bind-key  -T prefix  s  choose-tree -Zs\\
\color{white}bind-key  -T prefix  t  clock-mode\\
\color{white}bind-key  -T prefix  u  run-shell -b 'HOME/.tmux/plugins/tmux-fzf-url/fzf-\\
\color{white}bind-key  -T prefix  v  run-shell -b 'ZPWR/scripts/allPanes.zsh multi'\\
\color{white}bind-key  -T prefix  w  choose-tree -Zw\\
\color{white}bind-key  -T prefix  x  kill-pane\\
\color{white}bind-key  -T prefix  z  resize-pane -Z\\
\color{white}bind-key -r -T prefix  \textbackslash{}\textbackslash{}\{  rotate-window\\
\color{white}bind-key  -T prefix  |  split-window -h\\
\color{white}bind-key  -T prefix  \textbackslash{}\textbackslash{}\}  swap-pane -D\\
\color{white}bind-key  -T prefix  \textbackslash{}\textbackslash{}\textasciitilde{}  show-messages\\
\color{white}bind-key -r -T prefix  DC  refresh-client -c\\
\color{white}bind-key  -T prefix  PPage  copy-mode -u\\
\color{white}bind-key -r -T prefix  Up  select-pane -U\\
\color{white}bind-key -r -T prefix  Down  select-pane -D\\
\color{white}bind-key -r -T prefix  Left  select-pane -L\\
\color{white}bind-key -r -T prefix  Right  select-pane -R\\
\color{white}bind-key  -T prefix  M-1  select-layout even-horizontal\\
\color{white}bind-key  -T prefix  M-2  select-layout even-vertical\\
\color{white}bind-key  -T prefix  M-3  select-layout main-horizontal\\
\color{white}bind-key  -T prefix  M-4  select-layout main-vertical\\
\color{white}bind-key  -T prefix  M-5  select-layout tiled\\
\color{white}bind-key  -T prefix  M-6  select-layout main-horizontal-mirrored\\
\color{white}bind-key  -T prefix  M-7  select-layout main-vertical-mirrored\\
\color{white}bind-key  -T prefix  M-n  next-window -a\\

\vspace{8pt}

\subsection{Page 113: ALL Tmux Prefix Keys (3/4)}

\vspace{4pt}\textbf{\color{cyan}=== ALL TMUX PREFIX KEYS (3/4) ===}\vspace{2pt}\\
\color{white}bind-key  -T prefix  M-o  rotate-window -D\\
\color{white}bind-key  -T prefix  M-p  previous-window -a\\
\color{white}bind-key  -T prefix  M-u  run-shell HOME/.tmux/plugins/tpm/bindings/clean\_\\
\color{white}bind-key -r -T prefix  M-Up  resize-pane -U 5\\
\color{white}bind-key -r -T prefix  M-Down  resize-pane -D 5\\
\color{white}bind-key -r -T prefix  M-Left  resize-pane -L 5\\
\color{white}bind-key -r -T prefix  M-Right  resize-pane -R 5\\
\color{white}bind-key  -T prefix  C-]  send-keys \textasciicircum{}]\\
\color{white}bind-key  -T prefix  C-a  send-prefix\\
\color{white}bind-key  -T prefix  C-n  next-window\\
\color{white}bind-key  -T prefix  C-o  rotate-window\\
\color{white}bind-key  -T prefix  C-p  previous-window\\
\color{white}bind-key  -T prefix  C-r  run-shell HOME/.tmux/plugins/tmux-resurrect/scri\\
\color{white}bind-key  -T prefix  C-s  run-shell HOME/.tmux/plugins/tmux-resurrect/scri\\
\color{white}bind-key  -T prefix  C-v  run-shell 'tmux set buffer \textbackslash{}\textbackslash{}'(reattach-to-user-\\
\color{white}bind-key  -T prefix  C-z  suspend-client\\
\color{white}bind-key -r -T prefix  C-Up  resize-pane -U\\
\color{white}bind-key -r -T prefix  C-Down  resize-pane -D\\
\color{white}bind-key -r -T prefix  C-Left  resize-pane -L\\
\color{white}bind-key -r -T prefix  C-Right  resize-pane -R\\
\color{white}bind-key -r -T prefix  S-Up  refresh-client -U 10\\
\color{white}bind-key -r -T prefix  S-Down  refresh-client -D 10\\
\color{white}bind-key -r -T prefix  S-Left  refresh-client -L 10\\
\color{white}bind-key -r -T prefix  S-Right  refresh-client -R 10\\
\color{white}i  <C-X><C-L>  * fzf\#vim\#complete(fzf\#wrap(\{ 'prefix': '\textasciicircum{}.*', 'source': 'r\\
\color{white}<Space><Space>  <Plug>(easymotion-prefix)\\
\color{white}<Plug>(easymotion-prefix)N  <Plug>(easymotion-N)\\
\color{white}<Plug>(easymotion-prefix)n  <Plug>(easymotion-n)\\
\color{white}<Plug>(easymotion-prefix)k  <Plug>(easymotion-k)\\
\color{white}<Plug>(easymotion-prefix)j  <Plug>(easymotion-j)\\
\color{white}<Plug>(easymotion-prefix)gE  <Plug>(easymotion-gE)\\
\color{white}<Plug>(easymotion-prefix)ge  <Plug>(easymotion-ge)\\
\color{white}<Plug>(easymotion-prefix)E  <Plug>(easymotion-E)\\
\color{white}<Plug>(easymotion-prefix)e  <Plug>(easymotion-e)\\
\color{white}<Plug>(easymotion-prefix)B  <Plug>(easymotion-B)\\
\color{white}<Plug>(easymotion-prefix)b  <Plug>(easymotion-b)\\
\color{white}<Plug>(easymotion-prefix)W  <Plug>(easymotion-W)\\
\color{white}<Plug>(easymotion-prefix)w  <Plug>(easymotion-w)\\
\color{white}<Plug>(easymotion-prefix)T  <Plug>(easymotion-T)\\
\color{white}<Plug>(easymotion-prefix)t  <Plug>(easymotion-t)\\
\color{white}<Plug>(easymotion-prefix)s  <Plug>(easymotion-s)\\
\color{white}<Plug>(easymotion-prefix)F  <Plug>(easymotion-F)\\
\color{white}<Plug>(easymotion-prefix)f  <Plug>(easymotion-f)\\
\color{white}<Space><Space>  <Plug>(easymotion-prefix)\\
\color{white}<Plug>(easymotion-prefix)N  <Plug>(easymotion-N)\\
\color{white}<Plug>(easymotion-prefix)n  <Plug>(easymotion-n)\\
\color{white}<Plug>(easymotion-prefix)k  <Plug>(easymotion-k)\\
\color{white}<Plug>(easymotion-prefix)j  <Plug>(easymotion-j)\\

\vspace{8pt}

\subsection{Page 114: ALL Tmux Prefix Keys (4/4)}

\vspace{4pt}\textbf{\color{cyan}=== ALL TMUX PREFIX KEYS (4/4) ===}\vspace{2pt}\\
\color{white}<Plug>(easymotion-prefix)gE  <Plug>(easymotion-gE)\\
\color{white}<Plug>(easymotion-prefix)ge  <Plug>(easymotion-ge)\\
\color{white}<Plug>(easymotion-prefix)E  <Plug>(easymotion-E)\\
\color{white}<Plug>(easymotion-prefix)e  <Plug>(easymotion-e)\\
\color{white}<Plug>(easymotion-prefix)B  <Plug>(easymotion-B)\\
\color{white}<Plug>(easymotion-prefix)b  <Plug>(easymotion-b)\\
\color{white}<Plug>(easymotion-prefix)W  <Plug>(easymotion-W)\\
\color{white}<Plug>(easymotion-prefix)w  <Plug>(easymotion-w)\\
\color{white}<Plug>(easymotion-prefix)T  <Plug>(easymotion-T)\\
\color{white}<Plug>(easymotion-prefix)t  <Plug>(easymotion-t)\\
\color{white}<Plug>(easymotion-prefix)s  <Plug>(easymotion-s)\\
\color{white}<Plug>(easymotion-prefix)F  <Plug>(easymotion-F)\\
\color{white}<Plug>(easymotion-prefix)f  <Plug>(easymotion-f)\\

\vspace{8pt}

\subsection{Page 115: ALL Tmux Copy-Mode-Vi (1/2)}

\vspace{4pt}\textbf{\color{cyan}=== ALL TMUX COPY-MODE-VI (1/2) ===}\vspace{2pt}\\
\color{white}bind-key  -T copy-mode-vi Enter  send-keys -X copy-pipe-and-cancel 'reatta\\
\color{white}bind-key  -T copy-mode-vi Escape  send-keys -X clear-selection\\
\color{white}bind-key  -T copy-mode-vi Space  send-keys -X begin-selection\\
\color{white}bind-key  -T copy-mode-vi \textbackslash{}\textbackslash{}\#  send-keys -FX search-backward -- '\#\{copy\_cu\\
\color{white}bind-key  -T copy-mode-vi \textbackslash{}\textbackslash{}  send-keys -X end-of-line\\
\color{white}bind-key  -T copy-mode-vi \textbackslash{}\textbackslash{}\%  send-keys -X next-matching-bracket\\
\color{white}bind-key  -T copy-mode-vi *  send-keys -FX search-forward -- '\#\{copy\_curso\\
\color{white}bind-key  -T copy-mode-vi ,  send-keys -X jump-reverse\\
\color{white}bind-key  -T copy-mode-vi /  command-prompt -i -p '(search down)' 'send -X\\
\color{white}bind-key  -T copy-mode-vi 0  send-keys -X start-of-line\\
\color{white}bind-key  -T copy-mode-vi 1  command-prompt -N -I 1 -p (repeat) \{ send-key\\
\color{white}bind-key  -T copy-mode-vi 2  command-prompt -N -I 2 -p (repeat) \{ send-key\\
\color{white}bind-key  -T copy-mode-vi 3  command-prompt -N -I 3 -p (repeat) \{ send-key\\
\color{white}bind-key  -T copy-mode-vi 4  command-prompt -N -I 4 -p (repeat) \{ send-key\\
\color{white}bind-key  -T copy-mode-vi 5  command-prompt -N -I 5 -p (repeat) \{ send-key\\
\color{white}bind-key  -T copy-mode-vi 6  command-prompt -N -I 6 -p (repeat) \{ send-key\\
\color{white}bind-key  -T copy-mode-vi 7  command-prompt -N -I 7 -p (repeat) \{ send-key\\
\color{white}bind-key  -T copy-mode-vi 8  command-prompt -N -I 8 -p (repeat) \{ send-key\\
\color{white}bind-key  -T copy-mode-vi 9  command-prompt -N -I 9 -p (repeat) \{ send-key\\
\color{white}bind-key  -T copy-mode-vi :  command-prompt -p '(goto line)' \{ send-keys -\\
\color{white}bind-key  -T copy-mode-vi \textbackslash{}\textbackslash{};  send-keys -X jump-again\\
\color{white}bind-key  -T copy-mode-vi ?  command-prompt -i -p '(search up)' 'send -X s\\
\color{white}bind-key  -T copy-mode-vi A  send-keys -X append-selection-and-cancel\\
\color{white}bind-key  -T copy-mode-vi B  send-keys -X previous-space\\
\color{white}bind-key  -T copy-mode-vi D  send-keys -X copy-pipe-end-of-line-and-cancel\\
\color{white}bind-key  -T copy-mode-vi E  send-keys -X next-space-end\\
\color{white}bind-key  -T copy-mode-vi F  command-prompt -1 -p '(jump backward)' \{ send\\
\color{white}bind-key  -T copy-mode-vi G  send-keys -X history-bottom\\
\color{white}bind-key  -T copy-mode-vi H  send-keys -X top-line\\
\color{white}bind-key  -T copy-mode-vi J  send-keys -X scroll-down\\
\color{white}bind-key  -T copy-mode-vi K  send-keys -X scroll-up\\
\color{white}bind-key  -T copy-mode-vi L  send-keys -X bottom-line\\
\color{white}bind-key  -T copy-mode-vi M  send-keys -X middle-line\\
\color{white}bind-key  -T copy-mode-vi N  send-keys -X search-reverse\\
\color{white}bind-key  -T copy-mode-vi P  send-keys -X toggle-position\\
\color{white}bind-key  -T copy-mode-vi T  command-prompt -1 -p '(jump to backward)' \{ s\\
\color{white}bind-key  -T copy-mode-vi V  send-keys -X select-line\\
\color{white}bind-key  -T copy-mode-vi W  send-keys -X next-space\\
\color{white}bind-key  -T copy-mode-vi X  send-keys -X set-mark\\
\color{white}bind-key  -T copy-mode-vi \textasciicircum{}  send-keys -X back-to-indentation\\
\color{white}bind-key  -T copy-mode-vi b  send-keys -X previous-word\\
\color{white}bind-key  -T copy-mode-vi e  send-keys -X next-word-end\\
\color{white}bind-key  -T copy-mode-vi f  command-prompt -1 -p '(jump forward)' \{ send-\\
\color{white}bind-key  -T copy-mode-vi g  send-keys -X history-top\\
\color{white}bind-key  -T copy-mode-vi h  send-keys -X cursor-left\\
\color{white}bind-key  -T copy-mode-vi j  send-keys -X cursor-down\\
\color{white}bind-key  -T copy-mode-vi k  send-keys -X cursor-up\\
\color{white}bind-key  -T copy-mode-vi l  send-keys -X cursor-right\\

\vspace{8pt}

\subsection{Page 116: ALL Tmux Copy-Mode-Vi (2/2)}

\vspace{4pt}\textbf{\color{cyan}=== ALL TMUX COPY-MODE-VI (2/2) ===}\vspace{2pt}\\
\color{white}bind-key  -T copy-mode-vi n  send-keys -X search-again\\
\color{white}bind-key  -T copy-mode-vi o  send-keys -X other-end\\
\color{white}bind-key  -T copy-mode-vi q  send-keys -X cancel\\
\color{white}bind-key  -T copy-mode-vi r  send-keys -X refresh-from-pane\\
\color{white}bind-key  -T copy-mode-vi s  send-keys -X copy-pipe-no-clear 'reattach-to-\\
\color{white}bind-key  -T copy-mode-vi t  command-prompt -1 -p '(jump to forward)' \{ se\\
\color{white}bind-key  -T copy-mode-vi v  send-keys -X begin-selection\\
\color{white}bind-key  -T copy-mode-vi w  send-keys -X next-word\\
\color{white}bind-key  -T copy-mode-vi x  send-keys -X copy-pipe-no-clear 'reattach-to-\\
\color{white}bind-key  -T copy-mode-vi y  send-keys -X copy-pipe-no-clear 'reattach-to-\\
\color{white}bind-key  -T copy-mode-vi z  send-keys -X copy-pipe-no-clear 'reattach-to-\\
\color{white}bind-key  -T copy-mode-vi \textbackslash{}\textbackslash{}\{  send-keys -X previous-paragraph\\
\color{white}bind-key  -T copy-mode-vi \textbackslash{}\textbackslash{}\}  send-keys -X next-paragraph\\
\color{white}bind-key  -T copy-mode-vi WheelUpPane  select-pane \textbackslash{}\textbackslash{}; send-keys -X -N 5 s\\
\color{white}bind-key  -T copy-mode-vi WheelDownPane  select-pane \textbackslash{}\textbackslash{}; send-keys -X -N 5\\
\color{white}bind-key  -T copy-mode-vi DoubleClick1Pane  select-pane \textbackslash{}\textbackslash{}; send-keys -X s\\
\color{white}bind-key  -T copy-mode-vi TripleClick1Pane  select-pane \textbackslash{}\textbackslash{}; send-keys -X s\\
\color{white}bind-key  -T copy-mode-vi BSpace  send-keys -X cursor-left\\
\color{white}bind-key  -T copy-mode-vi Home  send-keys -X start-of-line\\
\color{white}bind-key  -T copy-mode-vi End  send-keys -X end-of-line\\
\color{white}bind-key  -T copy-mode-vi NPage  send-keys -X page-down\\
\color{white}bind-key  -T copy-mode-vi PPage  send-keys -X page-up\\
\color{white}bind-key  -T copy-mode-vi Up  send-keys -X cursor-up\\
\color{white}bind-key  -T copy-mode-vi Down  send-keys -X cursor-down\\
\color{white}bind-key  -T copy-mode-vi Left  send-keys -X cursor-left\\
\color{white}bind-key  -T copy-mode-vi Right  send-keys -X cursor-right\\
\color{white}bind-key  -T copy-mode-vi M-x  send-keys -X jump-to-mark\\
\color{white}bind-key  -T copy-mode-vi C-b  send-keys -X page-up\\
\color{white}bind-key  -T copy-mode-vi C-c  send-keys -X cancel\\
\color{white}bind-key  -T copy-mode-vi C-d  send-keys -X halfpage-down\\
\color{white}bind-key  -T copy-mode-vi C-e  send-keys -X scroll-down\\
\color{white}bind-key  -T copy-mode-vi C-f  send-keys -X page-down\\
\color{white}bind-key  -T copy-mode-vi C-h  send-keys -X -N 4 cursor-left\\
\color{white}bind-key  -T copy-mode-vi C-j  send-keys -X -N 4 cursor-down\\
\color{white}bind-key  -T copy-mode-vi C-k  send-keys -X -N 4 cursor-up\\
\color{white}bind-key  -T copy-mode-vi C-l  send-keys -X -N 4 cursor-right\\
\color{white}bind-key  -T copy-mode-vi C-n  send-keys -X page-down\\
\color{white}bind-key  -T copy-mode-vi C-p  send-keys -X page-up\\
\color{white}bind-key  -T copy-mode-vi C-u  send-keys -X halfpage-up\\
\color{white}bind-key  -T copy-mode-vi C-v  send-keys -X rectangle-toggle\\
\color{white}bind-key  -T copy-mode-vi C-y  send-keys -X scroll-up\\
\color{white}bind-key  -T copy-mode-vi C-Up  send-keys -X scroll-up\\
\color{white}bind-key  -T copy-mode-vi C-Down  send-keys -X scroll-down\\

\vspace{8pt}

\subsection{Page 117: ALL Zsh Insert Mode (1/3)}

\vspace{4pt}\textbf{\color{cyan}=== ALL ZSH INSERT MODE (1/3) ===}\vspace{2pt}\\
\color{white}bindkey -M viins '\textasciicircum{}@' zpwrExpandTerminateSpace\\
\color{white}bindkey -M viins '\textasciicircum{}A' beginning-of-line\\
\color{white}bindkey -M viins '\textasciicircum{}B' zpwrClipboard\\
\color{white}bindkey -M viins '\textasciicircum{}C' self-insert\\
\color{white}bindkey -M viins '\textasciicircum{}D' list-choices\\
\color{white}bindkey -M viins '\textasciicircum{}E' end-of-line\\
\color{white}bindkey -M viins '\textasciicircum{}F\textasciicircum{}D' zpwrIntoFzf\\
\color{white}bindkey -M viins '\textasciicircum{}F\textasciicircum{}F' zpwrCdVimFzfFilesearchWidgetAccept\\
\color{white}bindkey -M viins '\textasciicircum{}F\textasciicircum{}G' zpwrGoogle\\
\color{white}bindkey -M viins '\textasciicircum{}F\textasciicircum{}H' zpwrLsoffzf\\
\color{white}bindkey -M viins '\textasciicircum{}F\textasciicircum{}I' zpwrIntoFzfAg\\
\color{white}bindkey -M viins '\textasciicircum{}F\textasciicircum{}J' zpwrVerbsWidgetAccept\\
\color{white}bindkey -M viins '\textasciicircum{}F\textasciicircum{}K' zpwrAlternateQuotes\\
\color{white}bindkey -M viins '\textasciicircum{}F\textasciicircum{}L' list-choices\\
\color{white}bindkey -M viins '\textasciicircum{}F\textasciicircum{}M' zzcomplete\\
\color{white}bindkey -M viins '\textasciicircum{}F\textasciicircum{}N' zpwrVerbsWidget\\
\color{white}bindkey -M viins '\textasciicircum{}F\textasciicircum{}O' zpwrOpen\\
\color{white}bindkey -M viins '\textasciicircum{}F\textasciicircum{}P' basicSedSub\\
\color{white}bindkey -M viins '\textasciicircum{}F\textasciicircum{}R' zpwrAsVar\\
\color{white}bindkey -M viins '\textasciicircum{}F\textasciicircum{}S' zsh-gacp-CheckDiff\\
\color{white}bindkey -M viins '\textasciicircum{}F\textasciicircum{}V' edit-command-line\\
\color{white}bindkey -M viins '\textasciicircum{}Fj' zpwrVerbsWidgetAccept\\
\color{white}bindkey -M viins '\textasciicircum{}Fn' zpwrVerbsWidget\\
\color{white}bindkey -M viins '\textasciicircum{}G' what-cursor-position\\
\color{white}bindkey -M viins '\textasciicircum{}H' autopair-delete\\
\color{white}bindkey -M viins '\textasciicircum{}I' fzf-completion\\
\color{white}bindkey -M viins '\textasciicircum{}J' accept-line\\
\color{white}bindkey -M viins '\textasciicircum{}K' zsh-learn-Learn\\
\color{white}bindkey -M viins '\textasciicircum{}L' clear-screen\\
\color{white}bindkey -M viins '\textasciicircum{}M' zpwrMagicEnter\\
\color{white}bindkey -M viins '\textasciicircum{}N' down-history\\
\color{white}bindkey -M viins '\textasciicircum{}O' zpwrEOLorNextTabStop\\
\color{white}bindkey -M viins '\textasciicircum{}O\textasciicircum{}I' zc-logo\\
\color{white}bindkey -M viins '\textasciicircum{}O\textasciicircum{}Z' zui-demo-various\\
\color{white}bindkey -M viins '\textasciicircum{}P' up-history\\
\color{white}bindkey -M viins '\textasciicircum{}Q' zpwrLastWordDouble\\
\color{white}bindkey -M viins '\textasciicircum{}R' redo\\
\color{white}bindkey -M viins '\textasciicircum{}S' zsh-gacp-NoCheck\\
\color{white}bindkey -M viins '\textasciicircum{}T' fzf-file-widget\\
\color{white}bindkey -M viins '\textasciicircum{}U' zpwrClearLine\\
\color{white}bindkey -M viins '\textasciicircum{}V\textasciicircum{}@' zpwrVimFzf\\
\color{white}bindkey -M viins '\textasciicircum{}V\textasciicircum{}F' zpwrFasdFZF\\
\color{white}bindkey -M viins '\textasciicircum{}V\textasciicircum{}G' fzf-cd-widget\\
\color{white}bindkey -M viins '\textasciicircum{}V\textasciicircum{}K' zpwrEmacsFzf\\
\color{white}bindkey -M viins '\textasciicircum{}V\textasciicircum{}L' zpwrZFZF\\
\color{white}bindkey -M viins '\textasciicircum{}V\textasciicircum{}N' zpwrVimFzfSudo\\
\color{white}bindkey -M viins '\textasciicircum{}V\textasciicircum{}O' fzf-tab-complete\\
\color{white}bindkey -M viins '\textasciicircum{}V\textasciicircum{}P' sudo-command-line\\

\vspace{8pt}

\subsection{Page 118: ALL Zsh Insert Mode (2/3)}

\vspace{4pt}\textbf{\color{cyan}=== ALL ZSH INSERT MODE (2/3) ===}\vspace{2pt}\\
\color{white}bindkey -M viins '\textasciicircum{}V\textasciicircum{}R' history-search-multi-word\\
\color{white}bindkey -M viins '\textasciicircum{}V\textasciicircum{}S' zpwrZFZF\\
\color{white}bindkey -M viins '\textasciicircum{}V\textasciicircum{}V' zpwrVimAllWidgetAccept\\
\color{white}bindkey -M viins '\textasciicircum{}V\textasciicircum{}Z' fzf-history-widget\\
\color{white}bindkey -M viins '\textasciicircum{}V ' zpwrVimFzf\\
\color{white}bindkey -M viins '\textasciicircum{}V,' zpwrFzfEnv\\
\color{white}bindkey -M viins '\textasciicircum{}V.' zpwrFzfAllKeybind\\
\color{white}bindkey -M viins '\textasciicircum{}V/.' zpwrLocateFzfEdit\\
\textit{\color{gray}bindkey -M viins '\textasciicircum{}V//' zpwrLocateFzf}\\
\color{white}bindkey -M viins '\textasciicircum{}V;' zpwrFzfSurround\\
\color{white}bindkey -M viins '\textasciicircum{}Vc' zpwrFzfCommits\\
\color{white}bindkey -M viins '\textasciicircum{}Vk' zpwrFzfVimKeybind\\
\color{white}bindkey -M viins '\textasciicircum{}Vl' zpwrZFZF\\
\color{white}bindkey -M viins '\textasciicircum{}Vm' zpwrVimAllWidget\\
\color{white}bindkey -M viins '\textasciicircum{}Vv' zpwrVimAllWidgetAccept\\
\color{white}bindkey -M viins '\textasciicircum{}W' autopair-delete-word\\
\color{white}bindkey -M viins '\textasciicircum{}X\textasciicircum{}R' \_read\_comp\\
\color{white}bindkey -M viins '\textasciicircum{}X\textasciicircum{}X' zpwrBufferXtrace\\
\color{white}bindkey -M viins '\textasciicircum{}X?' \_complete\_debug\\
\color{white}bindkey -M viins '\textasciicircum{}XC' \_correct\_filename\\
\color{white}bindkey -M viins '\textasciicircum{}Xa' \_expand\_alias\\
\color{white}bindkey -M viins '\textasciicircum{}Xc' \_correct\_word\\
\color{white}bindkey -M viins '\textasciicircum{}Xd' \_list\_expansions\\
\color{white}bindkey -M viins '\textasciicircum{}Xe' \_expand\_word\\
\color{white}bindkey -M viins '\textasciicircum{}Xh' \_complete\_help\\
\color{white}bindkey -M viins '\textasciicircum{}Xm' \_most\_recent\_file\\
\color{white}bindkey -M viins '\textasciicircum{}Xn' \_next\_tags\\
\color{white}bindkey -M viins '\textasciicircum{}Xt' \_complete\_tag\\
\color{white}bindkey -M viins '\textasciicircum{}X\textasciitilde{}' \_bash\_list-choices\\
\color{white}bindkey -M viins '\textasciicircum{}Y' zpwrChangeQuotes\\
\color{white}bindkey -M viins '\textasciicircum{}Z' undo\\
\color{white}bindkey -M viins '\textasciicircum{}[' vi-cmd-mode\\
\color{white}bindkey -M viins '\textasciicircum{}[\textasciicircum{}D' capitalize-word\\
\color{white}bindkey -M viins '\textasciicircum{}[\textasciicircum{}E' zpwrExpandGlobalAliases\\
\color{white}bindkey -M viins '\textasciicircum{}[\textasciicircum{}F' sub\\
\color{white}bindkey -M viins '\textasciicircum{}[\textasciicircum{}L' down-case-word\\
\color{white}bindkey -M viins '\textasciicircum{}[\textasciicircum{}M' self-insert-unmeta\\
\color{white}bindkey -M viins '\textasciicircum{}[\textasciicircum{}O' zpwrRunner\\
\color{white}bindkey -M viins '\textasciicircum{}[\textasciicircum{}P' zpwrUp8widget\\
\color{white}bindkey -M viins '\textasciicircum{}[\textasciicircum{}T' transpose-words\\
\color{white}bindkey -M viins '\textasciicircum{}[\textasciicircum{}U' up-case-word\\
\color{white}bindkey -M viins '\textasciicircum{}[ ' zpwrSshRegain\\
\color{white}bindkey -M viins '\textasciicircum{}[,' \_history-complete-newer\\
\color{white}bindkey -M viins '\textasciicircum{}[/' \_history-complete-older\\
\color{white}bindkey -M viins '\textasciicircum{}[OA' history-substring-search-up\\
\color{white}bindkey -M viins '\textasciicircum{}[OB' history-substring-search-down\\
\color{white}bindkey -M viins '\textasciicircum{}[OC' vi-forward-char\\
\color{white}bindkey -M viins '\textasciicircum{}[OD' vi-backward-char\\

\vspace{8pt}

\subsection{Page 119: ALL Zsh Insert Mode (3/3)}

\vspace{4pt}\textbf{\color{cyan}=== ALL ZSH INSERT MODE (3/3) ===}\vspace{2pt}\\
\color{white}bindkey -M viins '\textasciicircum{}[OP' zpwrUp8widget\\
\color{white}bindkey -M viins '\textasciicircum{}[OQ' sub\\
\color{white}bindkey -M viins '\textasciicircum{}[OQ\textasciicircum{}[OQ' npm\_toggle\_install\_uninstall\\
\color{white}bindkey -M viins '\textasciicircum{}[OR' zpwrGetrcWidget\\
\color{white}bindkey -M viins '\textasciicircum{}[Q' zi-browse-symbol\\
\color{white}bindkey -M viins '\textasciicircum{}[[1;2D' sub\\
\color{white}bindkey -M viins '\textasciicircum{}[[1;5A' zsh-gacp-CheckDiff\\
\color{white}bindkey -M viins '\textasciicircum{}[[1;5B' zpwrUpdater\\
\color{white}bindkey -M viins '\textasciicircum{}[[1;5C' zpwrTutsUpdate\\
\color{white}bindkey -M viins '\textasciicircum{}[[1;5D' dbz\\
\color{white}bindkey -M viins '\textasciicircum{}[[200\textasciitilde{}' bracketed-paste\\
\color{white}bindkey -M viins '\textasciicircum{}[[5\textasciitilde{}' zpwrClipboard\\
\color{white}bindkey -M viins '\textasciicircum{}[[A' history-substring-search-up\\
\color{white}bindkey -M viins '\textasciicircum{}[[B' history-substring-search-down\\
\color{white}bindkey -M viins '\textasciicircum{}[[C' vi-forward-char\\
\color{white}bindkey -M viins '\textasciicircum{}[[D' vi-backward-char\\
\color{white}bindkey -M viins '\textasciicircum{}[[Z' zpwrClipboard\\
\color{white}bindkey -M viins '\textasciicircum{}[c' fzf-cd-widget\\
\color{white}bindkey -M viins '\textasciicircum{}[man' man-command-line\\
\color{white}bindkey -M viins '\textasciicircum{}[\textasciitilde{}' \_bash\_complete-word\\
\color{white}bindkey -R -M viins '\textasciicircum{}\textbackslash{}\textbackslash{}\textbackslash{}\textbackslash{}\textbackslash{}\textbackslash{}\textbackslash{}\textbackslash{}'-'\textasciicircum{}\textbackslash{}\textbackslash{}\textasciicircum{}' self-insert\\
\color{white}bindkey -M viins '\textasciicircum{}\_' zpwrAsVar\\
\color{white}bindkey -M viins ' ' zpwrExpandSupernaturalSpace\\
\color{white}bindkey -M viins '.' self-insert\\
\color{white}bindkey -M viins '\textbackslash{}\textbackslash{}'' zpwrInterceptSurround\\
\color{white}bindkey -R -M viins '\#'-'\&' self-insert\\
\color{white}bindkey -R -M viins '''-'(' zpwrInterceptSurround\\
\color{white}bindkey -M viins ')' autopair-close\\
\color{white}bindkey -R -M viins '*'-'-' self-insert\\
\color{white}bindkey -M viins '.' zpwrRationalizeDot\\
\color{white}bindkey -R -M viins '/'-'Z' self-insert\\
\color{white}bindkey -M viins '[' zpwrInterceptSurround\\
\color{white}bindkey -M viins '\textbackslash{}\textbackslash{}\textbackslash{}\textbackslash{}\textbackslash{}\textbackslash{}\textbackslash{}\textbackslash{}' self-insert\\
\color{white}bindkey -M viins ']' autopair-close\\
\color{white}bindkey -R -M viins '\textbackslash{}\textbackslash{}\textasciicircum{}'-'\_' self-insert\\
\color{white}bindkey -M viins '\textbackslash{}\textbackslash{}'' zpwrInterceptSurround\\
\color{white}bindkey -R -M viins 'a'-'f' self-insert\\
\color{white}bindkey -M viins 'fj' vi-cmd-mode\\
\color{white}bindkey -R -M viins 'g'-'j' self-insert\\
\color{white}bindkey -M viins 'jf' vi-cmd-mode\\
\color{white}bindkey -R -M viins 'k'-'z' self-insert\\
\color{white}bindkey -M viins '\{' zpwrInterceptSurround\\
\color{white}bindkey -M viins '|' self-insert\\
\color{white}bindkey -M viins '\}' autopair-close\\
\color{white}bindkey -M viins '\textasciitilde{}' self-insert\\
\color{white}bindkey -M viins '\textasciicircum{}?' zpwrInterceptDelete\\
\color{white}bindkey -R -M viins '\textbackslash{}\textbackslash{}M-\textasciicircum{}@'-'\textbackslash{}\textbackslash{}M-\textasciicircum{}?' self-insert\\

\vspace{8pt}

\subsection{Page 120: ALL Zsh Visual Mode (1/2)}

\vspace{4pt}\textbf{\color{cyan}=== ALL ZSH VISUAL MODE (1/2) ===}\vspace{2pt}\\
\color{white}bindkey -a 'V' visual-line-mode\\
\color{white}bindkey -a 'v' visual-mode\\
\color{white}bindkey -M visual '\textasciicircum{}[' deactivate-region\\
\color{white}bindkey -M visual '\textasciicircum{}[OA' up-line\\
\color{white}bindkey -M visual '\textasciicircum{}[OB' down-line\\
\color{white}bindkey -M visual '\textasciicircum{}[[A' up-line\\
\color{white}bindkey -M visual '\textasciicircum{}[[B' down-line\\
\color{white}bindkey -M visual '-' vi-up-line-or-history\\
\color{white}bindkey -M visual 'U' vi-up-case\\
\color{white}bindkey -M visual 'a\textbackslash{}\textbackslash{}'' select-quoted\\
\color{white}bindkey -M visual 'a'' select-quoted\\
\color{white}bindkey -M visual 'a(' select-bracketed\\
\color{white}bindkey -M visual 'a)' select-bracketed\\
\color{white}bindkey -M visual 'a+' select-quoted\\
\color{white}bindkey -M visual 'a,' select-quoted\\
\color{white}bindkey -M visual 'a-' select-quoted\\
\color{white}bindkey -M visual 'a.' select-quoted\\
\color{white}bindkey -M visual 'a/' select-quoted\\
\color{white}bindkey -M visual 'a:' select-quoted\\
\color{white}bindkey -M visual 'a;' select-quoted\\
\color{white}bindkey -M visual 'a<' select-bracketed\\
\color{white}bindkey -M visual 'a=' select-quoted\\
\color{white}bindkey -M visual 'a>' select-bracketed\\
\color{white}bindkey -M visual 'a@' select-quoted\\
\color{white}bindkey -M visual 'aB' select-bracketed\\
\color{white}bindkey -M visual 'aW' select-a-blank-word\\
\color{white}bindkey -M visual 'a[' select-bracketed\\
\color{white}bindkey -M visual 'a\textbackslash{}\textbackslash{}\textbackslash{}\textbackslash{}\textbackslash{}\textbackslash{}\textbackslash{}\textbackslash{}' select-quoted\\
\color{white}bindkey -M visual 'a]' select-bracketed\\
\color{white}bindkey -M visual 'a\textbackslash{}\textbackslash{}'' select-quoted\\
\color{white}bindkey -M visual 'aa' select-a-shell-word\\
\color{white}bindkey -M visual 'ab' select-bracketed\\
\color{white}bindkey -M visual 'aw' select-a-word\\
\color{white}bindkey -M visual 'a\{' select-bracketed\\
\color{white}bindkey -M visual 'a|' select-quoted\\
\color{white}bindkey -M visual 'a\}' select-bracketed\\
\color{white}bindkey -M visual 'i\textbackslash{}\textbackslash{}'' select-quoted\\
\color{white}bindkey -M visual 'i'' select-quoted\\
\color{white}bindkey -M visual 'i(' select-bracketed\\
\color{white}bindkey -M visual 'i)' select-bracketed\\
\color{white}bindkey -M visual 'i+' select-quoted\\
\color{white}bindkey -M visual 'i,' select-quoted\\
\color{white}bindkey -M visual 'i-' select-quoted\\
\color{white}bindkey -M visual 'i.' select-quoted\\
\color{white}bindkey -M visual 'i/' select-quoted\\
\color{white}bindkey -M visual 'i:' select-quoted\\
\color{white}bindkey -M visual 'i;' select-quoted\\
\color{white}bindkey -M visual 'i<' select-bracketed\\

\vspace{8pt}

\subsection{Page 121: ALL Zsh Visual Mode (2/2)}

\vspace{4pt}\textbf{\color{cyan}=== ALL ZSH VISUAL MODE (2/2) ===}\vspace{2pt}\\
\color{white}bindkey -M visual 'i=' select-quoted\\
\color{white}bindkey -M visual 'i>' select-bracketed\\
\color{white}bindkey -M visual 'i@' select-quoted\\
\color{white}bindkey -M visual 'iB' select-bracketed\\
\color{white}bindkey -M visual 'iW' select-in-blank-word\\
\color{white}bindkey -M visual 'i[' select-bracketed\\
\color{white}bindkey -M visual 'i\textbackslash{}\textbackslash{}\textbackslash{}\textbackslash{}\textbackslash{}\textbackslash{}\textbackslash{}\textbackslash{}' select-quoted\\
\color{white}bindkey -M visual 'i]' select-bracketed\\
\color{white}bindkey -M visual 'i\textbackslash{}\textbackslash{}'' select-quoted\\
\color{white}bindkey -M visual 'ia' select-in-shell-word\\
\color{white}bindkey -M visual 'ib' select-bracketed\\
\color{white}bindkey -M visual 'iw' select-in-word\\
\color{white}bindkey -M visual 'i\{' select-bracketed\\
\color{white}bindkey -M visual 'i|' select-quoted\\
\color{white}bindkey -M visual 'i\}' select-bracketed\\
\color{white}bindkey -M visual 'j' down-line\\
\color{white}bindkey -M visual 'k' up-line\\
\color{white}bindkey -M visual 'o' exchange-point-and-mark\\
\color{white}bindkey -M visual 'p' put-replace-selection\\
\color{white}bindkey -M visual 'u' vi-down-case\\
\color{white}bindkey -M visual 'x' vi-delete\\
\color{white}bindkey -M visual '\textasciitilde{}' vi-oper-swap-case\\

\vspace{8pt}

\subsection{Page 122: ALL Vim Normal Mode (1/3)}

\vspace{4pt}\textbf{\color{cyan}=== ALL VIM NORMAL MODE (1/3) ===}\vspace{2pt}\\
\color{white}n  <Esc><C-C>  * wvU\\
\color{white}n  <Esc><C-T>  * :call TransposeWords()<CR>\\
\color{white}n  <Space>mt  * :MinimapToggle<CR>\\
\color{white}n  <Space>mc  * :MinimapClose<CR>\\
\color{white}n  <Space>mu  * :MinimapUpdate<CR>\\
\color{white}n  <Space>mm  * :Minimap<CR>\\
\color{white}n  <Space><Tab>  :call AltOrNextBuffer()<CR>\\
\color{white}n  <Space>/  * :LocateAll<CR>\\
\color{white}n  <Space>ma  * :Marks.<CR>\\
\color{white}n  <Space>]  * <C-W>\}<CR>\\
\color{white}n  <Space>ta  * :Tags.<CR>\\
\color{white}n  <Space>h:  * :History:.<CR>\\
\color{white}n  <Space>h/  * :History/.<CR>\\
\color{white}n  <Space>hh  * :History.<CR>\\
\color{white}n  <Space>rg  * :Rg<CR>\\
\color{white}n  <Space>oa  * :ALEToggle<CR>\\
\color{white}n  <Space>m  * :Map.<CR>\\
\color{white}n  <Space>j  * :Lines.<CR>\\
\color{white}n  <Space>ke  * :FZFKeys<CR>\\
\color{white}n  <Space>,  * :FZFMaps<CR>\\
\color{white}n  <Space>aa  * :Agg<CR>\\
\color{white}n  <Space>i  * :Imap<CR>\\
\color{white}n  <Space>f  * :Files.<CR>\\
\color{white}n  <Space>env  * :FZFEnv<CR>\\
\color{white}n  <Space>;  * :Commands.<CR>\\
\color{white}n  <Space>.  * :Colors<CR>\\
\color{white}n  <Space>b  * :Buffers.<CR>\\
\color{white}n  <Space>ag  * :Ag<CR>\\
\color{white}n  <Space>vj  * :w.<CR>:call TmuxRepeat('file')<CR>\\
\color{white}n  <Space>ev  :call ExtractVariable()<CR>\\
\color{white}n  <Space>os  * :call CopyCWordClip()<CR> :call system('bash ZPWR\_TMUX/goo\\
\color{white}n  <Space>ox  * :call CopyCWordClip()<CR> :call system('bash ZPWR\_TMUX/goo\\
\color{white}n  <Space>oc  * :call CopyCWordClip()<CR>\\
\color{white}n  <Space>x  * :normal mzg\&'zzz<CR>\\
\color{white}n  <Space>p  * :bprev<CR>\\
\color{white}n  <Space>n  * :bnext<CR>\\
\color{white}n  <Space>ap  * :prev<CR>\\
\color{white}n  <Space>an  * :next<CR>\\
\color{white}n  <Space>lo  * :lopen<CR>\\
\color{white}n  <Space>lc  * :lclose<CR>\\
\color{white}n  <Space>lp  * :lprev<CR>\\
\color{white}n  <Space>ln  * :lnext<CR>\\
\color{white}n  <Space>[  * :call Quoter('bracket')<CR>\\
\color{white}n  <Space>'  * :call Quoter('back')<CR>\\
\color{white}n  <Space>'  * :call Quoter('single')<CR>\\
\color{white}n  <Space>'  * :call Quoter('double')<CR>\\
\color{white}n  <Space>z  * :Commits.<CR>\\
\color{white}n  <Space>t  * :tabnew<CR>\\

\vspace{8pt}

\subsection{Page 123: ALL Vim Normal Mode (2/3)}

\vspace{4pt}\textbf{\color{cyan}=== ALL VIM NORMAL MODE (2/3) ===}\vspace{2pt}\\
\color{white}n  <Space>s  * :split<CR>\\
\color{white}n  <Space>v  * :vsplit<CR>\\
\color{white}n  <Space>w  * :w.<CR>\\
\color{white}n  <Space>e  * :q.<CR>\\
\color{white}n  <Space>wq  * :wq.<CR>\\
\color{white}n  <Space>q  * :qa.<CR>\\
\color{white}n  <Space>r  * :\%s@\textbackslash{}\textbackslash{}C\textbackslash{}\textbackslash{}<<C-R><C-W>\textbackslash{}\textbackslash{}>@<C-R><C-W>@g<Left><Left>\\
\color{white}n  <Space>g  * :\%s@\textbackslash{}\textbackslash{}C\textbackslash{}\textbackslash{}<<C-R><C-W>\textbackslash{}\textbackslash{}>@@g<Left><Left>\\
\color{white}n  <Space>=  * 4+\\
\color{white}n  <Space>-  * 4-\\
\color{white}n  \&\&  * :normal mzg\&'zzz<CR>\\
\color{white}n  \&  * :\&\&<CR>\\
\color{white}n  <s<Esc>  \& <Nop>\\
\color{white}n  =s<Esc>  \& <Nop>\\
\color{white}n  >s<Esc>  \& <Nop>\\
\color{white}n  N  * :call GoToLastSearch('?')<CR>\\
\color{white}n  Y  * yy'>\\
\color{white}n  [o<Esc>  \& <Nop>\\
\color{white}n  []  k][\%?\}<CR>]\}]]\}]]\\
\color{white}n  [[  ?\{<CR>w99[\{\\
\color{white}n  \textbackslash{}\textbackslash{}K  * :OnlineThesaurusCurrentWord<CR>\\
\color{white}n  ]o<Esc>  \& <Nop>\\
\color{white}n  ]]  j0[[\%/\{<CR>\\
\color{white}n  ][  /\}<CR>b99]\}\\
\color{white}n  n  * :call GoToLastSearch('/')<CR>\\
\color{white}n  yo<Esc>  \& <Nop>\\
\color{white}n  y<C-G>  \& :<C-U>call setreg(v:register, fugitive\#Object(@\%))<CR>\\
\color{white}n  z/  * :if AutoHighlightToggle()|set hls|endif<CR>\\
\color{white}n  <SNR>198\_:  * :<C-U><C-R>=v:count ? v:count : ''<CR>\\
\color{white}n  <C-G>  * :call multiple\_cursors\#new('n', 1)<CR>\\
\color{white}n  <F11>  * :call conque\_term\#exec\_file()<CR>\\
\color{white}n  <C-W>\textbackslash{}\textbackslash{}  * :vsplit<CR>\\
\color{white}n  <C-W>-  * :split<CR>\\
\color{white}n  <C-D>,  * :FZFEnv<CR>\\
\color{white}n  <C-D>.  * :FZFKeys<CR>\\
\color{white}n  <C-D>/  * :LocateAll<CR>\\
\color{white}n  <C-D>z  * :TlistAddFiles *<CR>:TlistToggle<CR>\\
\color{white}n  <C-D>y  * :update<CR>:SyntasticCheck<CR>\\
\color{white}n  <C-D>x  * :Marks<CR>\\
\color{white}n  <C-D>w  * :update<CR>\\
\color{white}n  <C-D>v  * :w.<CR>:call TmuxRepeatGeneric()<CR>\\
\color{white}n  <C-D>u  * :History:<CR>\\
\color{white}n  <C-D>]  * <C-W>\}<CR>\\
\color{white}n  <C-D>t  * :Tags<CR>\\
\color{white}n  <C-D>s  * :History/<CR>\\
\color{white}n  <C-D>rr  * :Rg<CR>\\
\color{white}n  <C-D>rq  * :silent .open -t \%:p:h<CR>:redraw.<CR>\\
\color{white}n  <C-D>q  * :SaveSession.<CR><Tab>\\

\vspace{8pt}

\subsection{Page 124: ALL Vim Normal Mode (3/3)}

\vspace{4pt}\textbf{\color{cyan}=== ALL VIM NORMAL MODE (3/3) ===}\vspace{2pt}\\
\color{white}n  <C-D>p  * :call PasteClip()<CR>\\
\color{white}n  <C-D>r  * :call GetRef()<CR>\\
\color{white}n  <C-D>o  * :ALEToggle<CR>\\
\color{white}n  <C-D>n  * :Snippets<CR>\\
\color{white}n  <C-D>m  * :Map.<CR>\\
\color{white}n  <C-D>l  * :Lines.<CR>\\
\color{white}n  <C-D>k  * :ALEFix<CR>\\
\color{white}n  <C-D>j  * :Agg<CR>\\
\color{white}n  <C-D>i  * :Imap<CR>\\
\color{white}n  <C-D>h  * :HistoryFiles<CR>\\
\color{white}n  <C-D>g  * :Commits.<CR>\\
\color{white}n  <C-D>f  * :Files.<CR>\\
\color{white}n  <C-D>e  * :ALEInfo<CR>\\
\color{white}n  <C-D>d  * :Commands<CR>\\
\color{white}n  <C-D>c  * :Colors<CR>\\
\color{white}n  <C-D>b  * :Buffers<CR>\\
\color{white}n  <C-D>a  * :Ag<CR>\\
\color{white}n  <C-D><C-D>  * :GitGutterUndoHunk<CR>\\
\color{white}n  <C-V>  * :w.<CR>:call TmuxRepeat('file')<CR>\\
\color{white}n  <C-Bslash>  * +\\
\color{white}n  <F7>  * :TTags<CR>\\
\color{white}n  <F6>  * :SyntasticToggleMode<CR>\\
\color{white}n  <F5>  * :LOTRToggle<CR>\\
\color{white}n  <F4>  * :MinimapToggle<CR>\\
\color{white}n  <F3>  * :TlistAddFiles *<CR>:TlistToggle<CR>\\
\color{white}n  <F2>  * :UndotreeToggle<CR>\\
\color{white}n  <F1>  * :NERDTreeToggle<CR>\\
\color{white}n  <F8>  * :\%s@@@g<Left><Left><Left>\\
\color{white}n  <C-Up>  * :<C-U>call GoToNextMarker('\{\{\{',1)<CR>\\
\color{white}n  <C-Down>  * :<C-U>call GoToNextMarker('\{\{\{',0)<CR>\\
\color{white}n  <End>  * G\\
\color{white}n  <Home>  * gg\\
\color{white}n  <C-F>  * :q.<CR>\\
\color{white}n  <C-C>  * :wq.<CR>:qa.<CR>\\
\color{white}n  <C-T>  * xp\\
\color{white}n  <C-L>  * 4l\\
\color{white}n  <C-H>  * 4h\\
\color{white}n  <C-K>  * 4k\\
\color{white}n  <C-J>  * 4j\\
\color{white}n  <D-v>  '*P\\

\vspace{8pt}

\subsection{Page 125: ALL Vim Insert Mode (1/2)}

\vspace{4pt}\textbf{\color{cyan}=== ALL VIM INSERT MODE (1/2) ===}\vspace{2pt}\\
\color{white}i  <M-d>  * <C-O>dw\\
\color{white}i  <F31>  * <C-O>dw\\
\color{white}i  <C-F>  * col('.')>strlen(getline('.'))?'\textbackslash{}\textbackslash{}<C-F>':'\textbackslash{}\textbackslash{}<Right>'\\
\color{white}i  <C-E>  * col('.')>strlen(getline('.'))||pumvisible()?'\textbackslash{}\textbackslash{}<C-E>':'\textbackslash{}\textbackslash{}<End>\\
\color{white}i  <C-D>  * col('.')>strlen(getline('.'))?'\textbackslash{}\textbackslash{}<C-D>':'\textbackslash{}\textbackslash{}<Del>'\\
\color{white}i  <C-X><C-A>  * <C-A>\\
\color{white}i  <C-A>  * <C-O>\textasciicircum{}\\
\color{white}i  <C-X>  * <C-R>=<SNR>69\_ManualCompletionEnter()<CR>\\
\color{white}i  <C-Tab>  * <C-R>=UltiSnips\#ListSnippets()<CR>\\
\color{white}i  <C-X><C-L>  * fzf\#vim\#complete(fzf\#wrap(\{ 'prefix': '\textasciicircum{}.*', 'source': 'r\\
\color{white}i  <C-X><C-K>  * fzf\#vim\#complete\#word(\{'left': '15\%'\})\\
\color{white}i  <F11>  * <C-X><C-T>\\
\color{white}i  <F10>  * <C-X><C-K>\\
\color{white}i  <C-D>n  * <C-X><C-O>\\
\color{white}i  <C-D>\textbackslash{}\textbackslash{}  * <C-X><C-L>\\
\color{white}i  <C-D>,  * <C-O>:FZFEnv<CR>\\
\color{white}i  <C-D>.  * <C-O>:FZFKeys<CR>\\
\color{white}i  <C-D>/  * <C-O>:LocateAll<CR>\\
\color{white}i  <C-D>z  * <Esc>:TlistAddFiles * <CR> :TlistToggle<CR>i\\
\color{white}i  <C-D>y  * <Esc>:update<CR>:SyntasticCheck<CR>a\\
\color{white}i  <C-D>x  * <C-O>:Marks<CR>\\
\color{white}i  <C-D>w  * <C-O>:update<CR>\\
\color{white}i  <C-D>v  * <Esc>:w.<CR>:call TmuxRepeatGeneric()<CR>a\\
\color{white}i  <C-D>u  * <C-O>:History:<CR>\\
\color{white}i  <C-D>]  * <C-W>\}<CR>\\
\color{white}i  <C-D>t  * <C-O>:Tags<CR>\\
\color{white}i  <C-D>s  * <C-O>:History/<CR>\\
\color{white}i  <C-D>rr  * <Esc>:Rg<CR>\\
\color{white}i  <C-D>rq  * <Esc>:silent .open -t \%:p:h<CR>:redraw.<CR>a\\
\color{white}i  <C-D>q  * <C-O>:SaveSession.<CR><Tab>\\
\color{white}i  <C-D>p  * <C-O>:call PasteClip()<CR>\\
\color{white}i  <C-D>r  * <C-O>:call GetRef()<CR>\\
\color{white}i  <C-D>o  * <C-O>:ALEToggle<CR>\\
\color{white}i  <C-D>m  * <C-O>:Map<CR>\\
\color{white}i  <C-D>l  * <C-O>:Lines<CR>\\
\color{white}i  <C-D>k  * <C-O>:ALEFix<CR>\\
\color{white}i  <C-D>j  * <C-O>:Agg<CR>\\
\color{white}i  <C-D>i  * <C-O>:Imap<CR>\\
\color{white}i  <C-D>h  * <C-O>:HistoryFiles<CR>\\
\color{white}i  <C-D>g  * <C-O>:Commits.<CR>\\
\color{white}i  <C-D>f  * <C-O>:Files<CR>\\
\color{white}i  <C-D>e  * <C-O>:ALEInfo<CR>\\
\color{white}i  <C-D>d  * <C-O>:Commands<CR>\\
\color{white}i  <C-D>c  * <C-O>:Colors<CR>\\
\color{white}i  <C-D>b  * <C-O>:Buffers<CR>\\
\color{white}i  <C-D>a  * <C-O>:Ag<CR>\\
\color{white}i  <C-D><C-D>  * <C-O>:GitGutterUndoHunk<CR>\\
\color{white}i  <C-Bslash>  * <Esc>+\\

\vspace{8pt}

\subsection{Page 126: ALL Vim Insert Mode (2/2)}

\vspace{4pt}\textbf{\color{cyan}=== ALL VIM INSERT MODE (2/2) ===}\vspace{2pt}\\
\color{white}i  <F7>  * <Esc>:TTags<CR>\\
\color{white}i  <F6>  * <Esc>:SyntasticToggleMode<CR>\\
\color{white}i  <F5>  * <Esc>:LOTRToggle<CR>\\
\color{white}i  <F4>  * <Esc>:MinimapToggle<CR>\\
\color{white}i  <F3>  * <Esc>:TlistAddFiles *<CR>:TlistToggle<CR>\\
\color{white}i  <F2>  * <Esc>:UndotreeToggle<CR>\\
\color{white}i  <F1>  * <Esc>:NERDTreeToggle<CR>\\
\color{white}i  <F8>  * <Esc>:\%s@@@g<Left><Left><Left>\\
\color{white}i  <C-Up>  * <Esc>:<C-U>call GoToNextMarker('\{\{\{',1)<CR>i\\
\color{white}i  <C-Down>  * <Esc>:<C-U>call GoToNextMarker('\{\{\{',0)<CR>i\\
\color{white}i  <C-B><C-N>  * <Esc>\textasciicircum{}2xji\\
\color{white}i  <C-B>  * getline('.')=\textasciitilde{}'\textasciicircum{}\textbackslash{}\textbackslash{}s*'\&\&col('.')>strlen(getline('.'))?'0\textbackslash{}\textbackslash{}<C-D>\\
\color{white}i  <End>  * <Esc>Gi\\
\color{white}i  <Home>  * <Esc>ggi\\
\color{white}i  <C-C>  * <Esc>:wq.<CR>:qa.<CR>\\
\color{white}i  <C-T>  * i<BS><C-O>:silent. undojoin | normal. xp<CR>\\
\color{white}i  <C-D><C-T>  * <C-O>:call TransposeWords()<CR>\\
\color{white}i  <C-L>  * <Esc>mbgg=G'bzza\\
\color{white}i  <Tab>  * <C-R>=UltiSnips\#ExpandSnippetOrJump()<CR>\\
\color{white}i  fj  <Esc>\\
\color{white}i  jf  <Esc>\\

\vspace{8pt}

\subsection{Page 127: ALL Vim Visual Mode (1/1)}

\vspace{4pt}\textbf{\color{cyan}=== ALL VIM VISUAL MODE (1/1) ===}\vspace{2pt}\\
\color{white}v  <Space>vl  * <Esc>:call TmuxRepeat('repl')<CR>gv\\
\color{white}v  <Space>vk  * <Esc>:call TmuxRepeat('visual')<CR>gv\\
\color{white}v  <Space>vj  * <Esc>:call TmuxRepeat('file')<CR>gv\\
\color{white}v  <Space>(  * :call InsertQuoteVisualMode('paren')<CR>\\
\color{white}v  <Space>\{  * :call InsertQuoteVisualMode('curlybracket')<CR>\\
\color{white}v  <Space>[  * :call InsertQuoteVisualMode('bracket')<CR>\\
\color{white}v  <Space>'  * :call InsertQuoteVisualMode('back')<CR>\\
\color{white}v  <Space>'  * :call InsertQuoteVisualMode('single')<CR<CR>\\
\color{white}v  <Space>'  * :call InsertQuoteVisualMode('double')<CR>\\
\color{white}v  <Space>r  * :'<,'>s@\textbackslash{}\textbackslash{}C\textbackslash{}\textbackslash{}<<C-R><C-W>\textbackslash{}\textbackslash{}>@<C-R><C-W>@g<Left><Left>\\
\color{white}v  <Space>g  * :'<,'>\%s@\textbackslash{}\textbackslash{}C\textbackslash{}\textbackslash{}<<C-R><C-W>\textbackslash{}\textbackslash{}>@@g<Left><Left>\\
\color{white}v  <Space>b  * :w .tmux set-buffer '(cat)'<CR><CR>\\
\color{white}v  <Space>=  * 4+\\
\color{white}v  <Space>-  * 4-\\
\color{white}v  <  * <gv\\
\color{white}v  >  * >gv\\
\color{white}v  J  * :m '> + <CR> gv\\
\color{white}v  K  * :m '< -- <CR> gv\\
\color{white}v  Y  * y'>j\\
\color{white}v  \textbackslash{}\textbackslash{}K  * y:Thesaurus <C-R>'<CR>\\
\color{white}v  go  * :call CopyClip()<CR> :call system('bash ZPWR\_TMUX/google.sh open'\\
\color{white}v  gs  * :call CopyClip()<CR> :call system('bash ZPWR\_TMUX/google.sh googl\\
\color{white}v  <C-D>d  * :<C-C>:update<CR>\\
\color{white}v  <C-D>y  * :<C-C>:update<CR>:SyntasticCheck<CR>\\
\color{white}v  <C-D>,  * :call NERDComment('x','Toggle')<CR>'>\\
\color{white}v  <RightMouse> * :call CopyClip()<CR>'>\\
\color{white}v  <C-B>  * :call CopyClip()<CR>'>\\
\color{white}v  <C-Up>  * :m '< -- <CR> gv\\
\color{white}v  <C-Down>  * :m '> + <CR> gv\\
\color{white}v  <C-Left>  * <gv\\
\color{white}v  <C-Right>  * >gv\\
\color{white}v  <C-F>  * :<C-C>:q.<CR>\\
\color{white}v  <C-L>  * 4l\\
\color{white}v  <C-K>  * 4k\\
\color{white}v  <C-J>  * 4j\\
\color{white}v  <D-x>  '*d\\
\color{white}v  <D-c>  '*y\\
\color{white}v  <D-v>  '-d'*P\\

\vspace{8pt}


\section{CHAPTER 26: POWERLEVEL10K \& INSTANT PROMPT}

\subsection{Page 128: Instant Prompt: Zero Latency}

\vspace{4pt}\textbf{\color{cyan}=== INSTANT PROMPT: ZERO LATENCY ===}\vspace{2pt}\\
\textit{\color{gray}// your prompt loads before .zshrc finishes. time is a flat circle. //}\\
\color{white}The single best UX trick in zpwr: the prompt appears\\
\color{white}INSTANTLY while .zshrc is still loading behind the curtain.\\
\color{white}Zero perceived latency. The illusion of speed IS speed.\\
\vspace{4pt}\textbf{\color{magenta}HOW IT WORKS}\\
\color{white}On previous shell exit, p10k cached the prompt to a file:\\
\color{white}\textbackslash{}\$\{XDG\_CACHE\_HOME:-\textbackslash{}OME/.cache\}/p10k-instant-prompt-\textbackslash{}\$\{(\%):-\%n\}.zsh\\
\color{white}This file contains the last rendered prompt as raw ANSI escape\\
\color{white}sequences. On next shell start, it is sourced BEFORE .zshrc loads.\\
\color{white}Result: the prompt appears in <10ms. The user can start typing.\\
\color{white}.zshrc continues loading in the background. Nobody notices.\\
\vspace{4pt}\textbf{\color{magenta}THE LOADING SEQUENCE IN .zshrc}\\
\color{white}1. zmodload zsh/datetime (timestamp)\\
\color{white}2. Source zpwr-hash-dirs.zsh (cached hash -d entries)\\
\color{white}3. Source p10k-instant-prompt-*.zsh (the magic file)\\
\color{white}4. promptTimestamp=\textbackslash{}POCHREALTIME (measure prompt time)\\
\color{white}5. ...rest of .zshrc loads (plugins, aliases, etc)...\\
\color{white}6. p10k finalize (at very end of .zshrc)\\
\color{white}Step 3 must come before ANY console output. If something\\
\color{white}prints to stdout before the instant prompt, p10k aborts it.\\
\vspace{4pt}\textbf{\color{magenta}THE HASH DIR CACHE}\\
\color{white}The prompt uses \%\textasciitilde{} to show named dirs like \textasciitilde{}ZPWR instead of\\
\color{white}/Users/you/.zpwr. But hash -d entries are set in zpwrBindDirs,\\
\color{white}which runs AFTER instant prompt. Chicken-and-egg problem.\\
\color{white}Solution: zpwrBindDirs writes a cache file:\\
\texttt{\color{green}\textbackslash{}PWR\_LOCAL/zpwr-hash-dirs.zsh}\\
\color{white}This cache is sourced at line 73 of .zshrc, BEFORE instant\\
\color{white}prompt, so \%\textasciitilde{} resolves \textasciitilde{}ZPWR in the cached prompt correctly.\\
\vspace{4pt}\textbf{\color{magenta}ZPWR\_BANNER\_CLEARLIST SUPPRESSION}\\
\color{white}zpwrClearList prints a banner + file listing on login. But if\\
\color{white}it runs before instant prompt finalizes, it corrupts the display.\\
\color{white}Fix: the banner is deferred until AFTER p10k finalize:\\
\color{white}if [[ -n \textbackslash{}"\textbackslash{}\$\_\_p9k\_instant\_prompt\_active\textbackslash{}" \&\& \textbackslash{}"\textbackslash{}PWR\_BANNER\_CLEARLIST\textbackslash{}" == true ]]; then\\
\texttt{\color{green}zpwrClearList || true}\\
\color{white}fi\\
\vspace{4pt}\textbf{\color{magenta}THE promptTimestamp MEASUREMENT}\\
\color{white}ZPWR\_VARS[promptTimeMs] records how fast the prompt appeared.\\
\texttt{\color{green}Run zpwr environmentcounts to see it. With instant prompt,}\\
\color{white}this is typically 0.010-0.050 seconds. Without it, 0.5-2.0s.\\
\vspace{4pt}\textbf{\color{magenta}p10k finalize}\\
\color{white}The last line of .zshrc:\\
\color{white}(( ! \textbackslash{}\$\{+functions[p10k]\} )) || p10k finalize\\
\color{white}This tells p10k that .zshrc is done. The real prompt replaces\\
\color{white}the cached one. Any deferred output can now safely print.\\

\vspace{8pt}

\subsection{Page 129: Powerlevel10k Configuration}

\vspace{4pt}\textbf{\color{cyan}=== POWERLEVEL10K CONFIGURATION ===}\vspace{2pt}\\
\textit{\color{gray}// the prompt is your cockpit HUD. configure it or fly blind. //}\\
\color{white}The p10k config lives at \textbackslash{}PWR\_ENV/.p10k.zsh -- a 1700+ line\\
\color{white}file that controls every pixel of your prompt. The wizard\\
\color{white}generated it, but you own it now.\\
\vspace{4pt}\textbf{\color{magenta}LEFT PROMPT SEGMENTS}\\
\color{white}These appear on the left side, top to bottom across 4 lines:\\
\color{white}Line 1: custom\_env, custom\_env\_bat, custom\_env\_editor\\
\color{white}Line 2: custom\_env\_term, ip\\
\color{white}Line 3: context, dir, history, ssh, root\_indicator, status, vcs\\
\color{white}Line 4: custom\_tty, prompt\_char\\
\color{white}The dir segment shows your current directory using \%\textasciitilde{} expansion.\\
\color{white}The vcs segment shows git branch, dirty status, ahead/behind.\\
\color{white}The prompt\_char shows green on success, red on failure.\\
\vspace{4pt}\textbf{\color{magenta}RIGHT PROMPT SEGMENTS}\\
\color{white}The right side is information-dense. Across 3 lines:\\
\color{white}Line 1: dir\_writable, background\_jobs, virtualenv, node\_version,\\
\color{white}go\_version, rust\_version, java\_version, package, os\_icon...\\
\color{white}Line 2: load, disk\_usage, ram, swap, wifi\\
\color{white}Line 3: battery, time, command\_execution\_time, vi\_mode, custom\_pid\\
\vspace{4pt}\textbf{\color{magenta}THE DIR SEGMENT AND \%\textasciitilde{} EXPANSION}\\
\color{white}Zsh's \%\textasciitilde{} prompt expansion replaces known paths with named dirs.\\
\texttt{\color{green}If you hash -d ZPWR=/Users/you/.zpwr, then /Users/you/.zpwr}\\
\color{white}shows as \textasciitilde{}ZPWR in your prompt. Shorter. Cleaner. Professional.\\
\color{white}POWERLEVEL9K\_DIR\_MAX\_LENGTH=80 caps the dir segment width.\\
\color{white}POWERLEVEL9K\_DIR\_MIN\_COMMAND\_COLUMNS=40 reserves typing space.\\
\vspace{4pt}\textbf{\color{magenta}HOW zpwrBindDirs POPULATES NAMED DIRECTORIES}\\
\color{white}The function runs hash -d for \textasciitilde{}20 directories:\\
\color{white}hash -d ZPWR=\textbackslash{}"\textbackslash{}PWR\textbackslash{}"\\
\color{white}hash -d ZPWR\_TMUX=\textbackslash{}"\textbackslash{}PWR\_TMUX\textbackslash{}"\\
\color{white}hash -d ZPWR\_SCRIPTS=\textbackslash{}"\textbackslash{}PWR\_SCRIPTS\textbackslash{}"\\
\color{white}hash -d TMUXH=\textbackslash{}"\textbackslash{}MUX\_HOME\textbackslash{}"\\
\color{white}Then it writes them all to \textbackslash{}PWR\_LOCAL/zpwr-hash-dirs.zsh\\
\color{white}for instant prompt to use on next startup. The cache iterates\\
\color{white}\textbackslash{}\$\{(kv)nameddirs\} to capture every registered named directory.\\
\vspace{4pt}\textbf{\color{magenta}THE MULTILINE FRAME}\\
\color{white}The prompt uses box-drawing characters for the frame:\\
\color{white}First line:  \%244F-----  (top connector)\\
\color{white}Middle:      \%244F-----  (mid connector)\\
\color{white}Last line:   \%244F-----  (bottom connector)\\
\color{white}Gap fill character is '---' in grey (color 244).\\
\color{white}Mode: nerdfont-complete. Separators use Powerline glyphs.\\
\vspace{4pt}\textbf{\color{magenta}prompt\_char SYMBOLS}\\
\color{white}Insert mode: ---  Command mode: ---  Visual: V  Overwrite: ---\\

\vspace{8pt}

\subsection{Page 130: P10k Prompt Customization}

\vspace{4pt}\textbf{\color{cyan}=== P10K PROMPT CUSTOMIZATION ===}\vspace{2pt}\\
\textit{\color{gray}// your prompt is a canvas. paint it with ANSI and Unicode. //}\\
\color{white}ZPWR ships custom p10k segments beyond the defaults.\\
\color{white}These are defined with POWERLEVEL9K\_CUSTOM\_* variables\\
\color{white}in \textbackslash{}PWR\_ENV/.p10k.zsh.\\
\vspace{4pt}\textbf{\color{magenta}CUSTOM SEGMENTS}\\
\color{white}custom\_tty       Shows \textbackslash{}PWR\_TTY (e.g. /dev/ttys003)\\
\color{white}Defined as:\\
\color{white}POWERLEVEL9K\_CUSTOM\_TTY='echo \textbackslash{}"\%K\{blue\}\%F\{white\} \textbackslash{}PWR\_TTY \%f\%k\%F\{blue\}\%f\textbackslash{}"'\\
\color{white}custom\_env       Environment info (custom shell, OS, etc.)\\
\color{white}custom\_env\_bat   Shows bat (syntax highlighter) version/status\\
\color{white}custom\_env\_editor Shows \textbackslash{}DITOR\\
\color{white}custom\_env\_term  Shows \textbackslash{}ERM\_PROGRAM\\
\color{white}custom\_pid       Shows shell PID (right prompt, bottom line)\\
\vspace{4pt}\textbf{\color{magenta}ADDING YOUR OWN SEGMENT}\\
\color{white}1. Define the command that generates output:\\
\color{white}typeset -g POWERLEVEL9K\_CUSTOM\_MYWIDGET='echo \textbackslash{}"hello\textbackslash{}"'\\
\color{white}2. Set colors:\\
\color{white}typeset -g POWERLEVEL9K\_CUSTOM\_MYWIDGET\_FOREGROUND=255\\
\color{white}typeset -g POWERLEVEL9K\_CUSTOM\_MYWIDGET\_BACKGROUND=24\\
\color{white}3. Add 'custom\_mywidget' to LEFT or RIGHT PROMPT\_ELEMENTS array.\\
\color{white}4. Run p10k reload or source \textbackslash{}PWR\_ENV/.p10k.zsh to apply.\\
\vspace{4pt}\textbf{\color{magenta}COLORS AND ICONS}\\
\color{white}p10k uses 256-color indices. Print the color map:\\
\color{white}for i in \{0..255\}; do print -Pn \textbackslash{}"\%K\{\textbackslash{}\$i\}  \%k\%F\{\textbackslash{}\$i\}\textbackslash{}\$\{(l:3::0:)i\}\%f \textbackslash{}"; done\\
\color{white}Foreground/background for each segment and state:\\
\color{white}POWERLEVEL9K\_DIR\_BACKGROUND=4        (blue)\\
\color{white}POWERLEVEL9K\_DIR\_FOREGROUND=254      (near-white)\\
\color{white}Anchor dirs (like .git roots) get bold + brighter color (255).\\
\vspace{4pt}\textbf{\color{magenta}TRUNCATION}\\
\color{white}POWERLEVEL9K\_DIR\_MAX\_LENGTH=80 caps directory width.\\
\color{white}POWERLEVEL9K\_DIR\_MIN\_COMMAND\_COLUMNS\_PCT=50 means at least\\
\color{white}half the terminal width is reserved for your command input.\\
\color{white}SHORTEN\_STRATEGY can be set to truncate\_to\_unique for smart\\
\color{white}shortening (commented out by default in zpwr).\\
\vspace{4pt}\textbf{\color{magenta}THE .p10k.zsh FILE STRUCTURE}\\
\color{white}The file is a single anonymous function () \{ ... \} that:\\
\color{white}- Unsets all POWERLEVEL9K\_* vars (clean slate)\\
\color{white}- Defines LEFT/RIGHT\_PROMPT\_ELEMENTS arrays\\
\color{white}- Sets mode, icons, colors, truncation per segment\\
\color{white}- Configures transient prompt, instant prompt, hot reload\\
\vspace{4pt}\textbf{\color{magenta}SEEING PROMPT TIMING}\\
\texttt{\color{green}zpwr environmentcounts}\\
\color{white}Shows promptTimeMs, startupTimeMs, and per-phase timestamps.\\
\color{white}This is how you prove your prompt loaded in 0.015 seconds.\\

\vspace{8pt}

\subsection{Page 131: Transient Prompt \& Visual Tweaks}

\vspace{4pt}\textbf{\color{cyan}=== TRANSIENT PROMPT \& VISUAL TWEAKS ===}\vspace{2pt}\\
\textit{\color{gray}// old prompts fade into ghosts. only the present matters. //}\\
\color{white}Transient prompt collapses previous prompts to a minimal\\
\color{white}form, keeping your scrollback clean and focused. In zpwr\\
\color{white}it is currently set to 'off' but easy to enable.\\
\vspace{4pt}\textbf{\color{magenta}TRANSIENT PROMPT}\\
\color{white}When enabled, after you press Enter, the full multi-line\\
\color{white}prompt collapses to just the prompt\_char symbol (--- or ---).\\
\color{white}All segments, git status, directory -- gone from old lines.\\
\color{white}Only the current prompt shows full detail.\\
\color{white}typeset -g POWERLEVEL9K\_TRANSIENT\_PROMPT=off\\
\color{white}Change to always or same-dir to enable.\\
\color{white}'same-dir' only collapses when you stay in the same directory.\\
\vspace{4pt}\textbf{\color{magenta}HOT RELOAD}\\
\color{white}Edit .p10k.zsh and apply changes without restarting zsh:\\
\color{white}p10k reload\\
\color{white}This re-sources .p10k.zsh and rebuilds the prompt in-place.\\
\color{white}No exec zsh needed. No losing your command history buffer.\\
\vspace{4pt}\textbf{\color{magenta}p10k configure}\\
\color{white}The interactive configuration wizard. Walks you through:\\
\color{white}- Font detection (powerline glyphs, nerdfont icons)\\
\color{white}- Style selection (rainbow, lean, classic, pure)\\
\color{white}- Separator shapes, line count, frame options\\
\color{white}- Instant prompt mode (verbose, quiet, off)\\
\color{white}It generates a fresh .p10k.zsh. ZPWR ships a pre-configured\\
\color{white}one in \textbackslash{}PWR\_ENV/.p10k.zsh (rainbow style, nerdfont-complete).\\
\color{white}Running the wizard overwrites it, so back up yours first.\\
\vspace{4pt}\textbf{\color{magenta}INSTANT PROMPT vs zpwrClearList}\\
\color{white}The banner (zpwrClearList) prints colorized file listings on\\
\color{white}login. Instant prompt redirects all stdout/stderr during init\\
\color{white}to a buffer. If zpwrClearList runs during this window, the\\
\color{white}output gets swallowed or garbled. ZPWR's fix:\\
\color{white}1. Check if \_\_p9k\_instant\_prompt\_active is set\\
\color{white}2. If so, defer zpwrClearList to AFTER p10k finalize\\
\color{white}3. The banner appears cleanly after the prompt settles\\
\vspace{4pt}\textbf{\color{magenta}THE eza --classify=always FIX}\\
\color{white}Under instant prompt, fd/stdout gets redirected. eza's\\
\color{white}--classify flag (append */ to dirs) auto-detects terminal\\
\color{white}capabilities. When stdout is redirected, it thinks there's\\
\color{white}no terminal and drops the indicators.\\
\color{white}Fix: use --classify=always to force classification regardless\\
\color{white}of fd state. ZPWR\_EXA\_COMMAND uses this:\\
\color{white}command eza --git -il --classify=always -H --color-scale -g -a\\
\vspace{4pt}\textbf{\color{magenta}\textbackslash{}TY IMMUNE TO FD REDIRECTS}\\
\color{white}Zsh sets \textbackslash{}TY before rc files load. Instant prompt redirects\\
\color{white}fds 0/1/2, but \textbackslash{}TY still points to the real terminal.\\
\color{white}ZPWR captures it: export ZPWR\_TTY=\textbackslash{}TY\\
\color{white}This is used by the custom\_tty prompt segment.\\

\vspace{8pt}


\section{CHAPTER 27: ZINIT \& PLUGIN SYSTEM}

\subsection{Page 132: Zinit: The Plugin Engine}

\vspace{4pt}\textbf{\color{cyan}=== ZINIT: THE PLUGIN ENGINE ===}\vspace{2pt}\\
\textit{\color{gray}// the package manager for your shell. apt-get for zsh plugins. //}\\
\color{white}Zinit (formerly zplugin) is a flexible zsh plugin manager\\
\color{white}with turbo mode, ice modifiers, and completion management.\\
\color{white}ZPWR uses it as the default: ZPWR\_PLUGIN\_MANAGER=zinit\\
\vspace{4pt}\textbf{\color{magenta}WHAT ZINIT DOES}\\
\color{white}- Downloads GitHub plugins to \textbackslash{}PWR\_PLUGIN\_MANAGER\_HOME\\
\color{white}- Sources them into your shell (or loads as completions)\\
\color{white}- Manages load order with turbo mode (deferred loading)\\
\color{white}- Handles OMZ snippets (individual files from oh-my-zsh)\\
\color{white}- Compiles plugins to .zwc bytecode for speed\\
\vspace{4pt}\textbf{\color{magenta}ZINIT LOAD vs ZINIT SNIPPET}\\
\color{white}zinit load user/repo      Clone full GitHub repo, source .plugin.zsh\\
\color{white}zinit snippet OMZP::git    Download single file from OMZ plugin\\
\color{white}zinit snippet OMZL::git.zsh Download single OMZ lib file\\
\color{white}load = full repo clone. snippet = single file download.\\
\color{white}ZPWR uses both: load for GitHub plugins, snippet for OMZ.\\
\vspace{4pt}\textbf{\color{magenta}THE PLUGIN ARRAYS}\\
\color{white}ZPWR organizes plugins into arrays in .zshrc:\\
\color{white}ZPWR\_GH\_PLUGINS   28+ GitHub repos (fzf-tab, zsh-z, forgit...)\\
\color{white}ZPWR\_OMZ\_PLUGINS  25+ OMZ snippets (git, python, golang...)\\
\color{white}ZPWR\_OMZ\_LIBS     4 OMZ lib files (git.zsh, directories.zsh...)\\
\color{white}ZPWR\_OMZ\_COMPS    12 OMZ completion files (ng, coffee, scala...)\\
\color{white}Conditional plugins are appended if the command exists:\\
\texttt{\color{green}zpwrCommandExists docker \&\& ZPWR\_GH\_PLUGINS+=( zsh-docker-aliases )}\\
\vspace{4pt}\textbf{\color{magenta}PLUGIN STORAGE}\\
\color{white}Everything lives under \textbackslash{}PWR\_PLUGIN\_MANAGER\_HOME:\\
\color{white}\textbackslash{}PWR\_PLUGIN\_MANAGER\_HOME/bin/         zinit itself\\
\color{white}\textbackslash{}PWR\_PLUGIN\_MANAGER\_HOME/plugins/     cloned repos\\
\color{white}\textbackslash{}PWR\_PLUGIN\_MANAGER\_HOME/snippets/    OMZ snippets\\
\color{white}\textbackslash{}PWR\_PLUGIN\_MANAGER\_HOME/completions/  symlinked completions\\
\vspace{4pt}\textbf{\color{magenta}THE 48+ PLUGIN ROSTER}\\
\color{white}Total plugins managed by zpwr (from .zshplugins):\\
\color{white}MenkeTechnologies forks: fzf-tab, forgit, zsh-expand,\\
\color{white}zsh-autocomplete, zsh-z, fasd-simple, revolver, zunit...\\
\color{white}zdharma-continuum: fast-syntax-highlighting, zbrowse, zui...\\
\color{white}zsh-users: completions, autosuggestions, history-substring\\
\color{white}Plus: hlissner/zsh-autopair, bigH/git-fuzzy, and more.\\
\vspace{4pt}\textbf{\color{magenta}FIRST RUN BOOTSTRAP}\\
\color{white}If \textbackslash{}PWR\_PLUGIN\_MANAGER\_HOME/bin doesn't exist, .zshrc\\
\color{white}auto-clones zinit from GitHub. After that, OPTIMIZE\_OUT\_DISK\_ACCESSES=1\\
\color{white}skips unnecessary stat() calls on subsequent starts.\\

\vspace{8pt}

\subsection{Page 133: Turbo Mode: Deferred Loading}

\vspace{4pt}\textbf{\color{cyan}=== TURBO MODE: DEFERRED LOADING ===}\vspace{2pt}\\
\textit{\color{gray}// why load 48 plugins at startup when you can load 1? //}\\
\color{white}Turbo mode is zinit's killer feature: plugins load AFTER\\
\color{white}the prompt appears, using zsh's sched mechanism. Your shell\\
\color{white}is interactive instantly; plugins trickle in behind the scenes.\\
\vspace{4pt}\textbf{\color{magenta}THE wait ICE MODIFIER}\\
\color{white}The 'wait' ice defers loading until after the first prompt:\\
\color{white}zinit ice wait            Load after prompt (default priority)\\
\color{white}zinit ice wait'0'         Same as wait (priority 0)\\
\color{white}zinit ice wait'0a'        Priority 0, sub-order a (loads first)\\
\color{white}zinit ice wait'0b'        Priority 0, sub-order b (loads second)\\
\color{white}zinit ice wait'0e'        Priority 0, sub-order e (loads later)\\
\color{white}The letter suffix controls load order WITHIN the same priority.\\
\color{white}a loads before b loads before c, etc. This is how ZPWR\\
\color{white}orchestrates 48 plugins into a precise loading sequence.\\
\vspace{4pt}\textbf{\color{magenta}THE GRADUATED LOADING SEQUENCE}\\
\color{white}Synchronous:  powerlevel10k (must load before prompt)\\
\color{white}wait:          OMZ completions, libs, plugins, GH plugins\\
\color{white}wait'0':       zsh-autopair (needs ZLE ready)\\
\color{white}wait'0a':      zsh-expand + late aliases + alias cache\\
\color{white}wait'0b':      zsh-autocomplete (if ZPWR\_AUTO\_COMPLETE=true)\\
\color{white}wait'0c':      zsh-history-substring-search\\
\color{white}wait'0d':      zsh-autosuggestions + FZF late bindings + verbs\\
\color{white}wait'0e':      zsh-zinit-final (penultimate, final keybindings)\\
\color{white}wait'0f':      zsh-more-completions (low priority, last resort)\\
\color{white}wait'\textbackslash{}\$\{ZPWR\_ZINIT\_COMPINIT\_DELAY\}g':  compinit + cdreplay\\
\color{white}wait'\textbackslash{}\$\{ZPWR\_ZINIT\_COMPINIT\_DELAY\}h':  fast-syntax-highlighting\\
\vspace{4pt}\textbf{\color{magenta}WHY STARTUP IS FAST}\\
\color{white}Only ONE plugin loads synchronously: powerlevel10k.\\
\color{white}Everything else is turbo-deferred. Combined with instant\\
\color{white}prompt, you see the prompt in \textasciitilde{}15ms. The other 47 plugins\\
\color{white}load in the next \textasciitilde{}200ms while you're still thinking.\\
\vspace{4pt}\textbf{\color{magenta}atinit AND atload HOOKS}\\
\color{white}These run code at precise moments during plugin loading:\\
\color{white}atinit   Runs BEFORE the plugin is sourced\\
\color{white}atload   Runs AFTER the plugin is sourced\\
\color{white}ZPWR uses these extensively to bind keymaps and aliases:\\
\texttt{\color{green}zinit ice wait'0a' atload'zpwrBindAliasesLate; zpwrCreateAliasCache'}\\
\color{white}The atload fires after zsh-expand is sourced, then late\\
\color{white}aliases override any OMZ defaults. Precise surgical timing.\\
\vspace{4pt}\textbf{\color{magenta}ZPWR\_ZINIT\_COMPINIT\_DELAY}\\
\color{white}Default: 0. Controls how long to wait before running compinit.\\
\color{white}Higher values defer completion system init even further.\\
\color{white}Set in \textbackslash{}PWR\_ENV/.zpwr\_env.sh. The 'g' and 'h' suffixes on\\
\color{white}the wait ice ensure compinit and syntax highlighting load last.\\

\vspace{8pt}

\subsection{Page 134: Ice Modifiers Deep Dive}

\vspace{4pt}\textbf{\color{cyan}=== ICE MODIFIERS DEEP DIVE ===}\vspace{2pt}\\
\textit{\color{gray}// ice modifiers are the spice of zinit. season generously. //}\\
\color{white}Every zinit load/snippet can be preceded by 'zinit ice'\\
\color{white}with modifiers that control HOW the plugin loads.\\
\color{white}Ice melts after one use -- it applies only to the next command.\\
\vspace{4pt}\textbf{\color{magenta}CORE ICE MODIFIERS USED BY ZPWR}\\
\color{white}lucid          Suppress the 'Loaded plugin' message\\
\color{white}ZPWR uses this on every plugin. No spam.\\
\color{white}nocompile      Skip .zwc compilation of the plugin\\
\color{white}Useful when zpwrRecompileFiles handles it globally.\\
\color{white}nocd           Don't cd to plugin dir before sourcing\\
\color{white}Prevents p10k internals from polluting named dirs.\\
\color{white}nocompletions  Don't install completions from this plugin\\
\color{white}Used for zsh-more-completions (fpath-only approach).\\
\color{white}as'completion' Load the file as a completion, not a plugin\\
\color{white}Used for OMZ completion snippets like \_ng, \_coffee.\\
\color{white}as'program'    Add the pick'd file to \textbackslash{}ATH as executable\\
\color{white}Used for fzf binary and git-fuzzy.\\
\color{white}as'null'       Don't source anything -- just run hooks\\
\color{white}Used for the compinit phase (atinit only).\\
\color{white}pick'null'     Don't source any file from the plugin dir\\
\color{white}Used with as'completion' for OMZ comp snippets.\\
\color{white}pick'bin/fzf'  Source or add this specific file\\
\color{white}With as'program', adds bin/fzf to PATH.\\
\vspace{4pt}\textbf{\color{magenta}HOOK ICE MODIFIERS}\\
\color{white}atinit         Code to run BEFORE the plugin sources\\
\color{white}atload         Code to run AFTER the plugin sources\\
\color{white}Example from ZPWR:\\
\texttt{\color{green}zinit ice lucid nocompile wait atinit='zpwrBindOverrideOMZ;zpwrBindForGit'}\\
\color{white}zinit load MenkeTechnologies/forgit\\
\color{white}Before forgit loads, zpwrBindOverrideOMZ runs to clear OMZ\\
\color{white}aliases that would conflict with forgit's git shortcuts.\\
\vspace{4pt}\textbf{\color{magenta}HOW ZPWR USES ICE BY CATEGORY}\\
\color{white}OMZ completions: lucid nocompile as'completion' pick'null' wait\\
\texttt{\color{green}OMZ libs:        lucid nocompile wait atload='zpwrOmzOverrides'}\\
\color{white}OMZ plugins:     lucid nocompile wait\\
\color{white}GH plugins:      lucid nocompile wait\\
\color{white}Programs:        as'program' lucid nocompile pick'bin/fzf' wait\\
\color{white}Syntax highlight:lucid nocompile wait\textbackslash{}"\textbackslash{}\$\{ZPWR\_ZINIT\_COMPINIT\_DELAY\}h\textbackslash{}"\\
\vspace{4pt}\textbf{\color{magenta}ICE STACKING}\\
\color{white}Multiple ices combine. Order doesn't matter (mostly).\\
\color{white}The = syntax for atload/atinit works with single quotes:\\
\texttt{\color{green}atload='zpwrBindAliasesLate; zpwrCreateAliasCache'}\\
\color{white}Or with quoted strings:\\
\texttt{\color{green}atinit'zpwrBindZdharma' atload'zpwrBindZdharmaPost'}\\

\vspace{8pt}

\subsection{Page 135: Zinit Completions \& Compinit}

\vspace{4pt}\textbf{\color{cyan}=== ZINIT COMPLETIONS \& COMPINIT ===}\vspace{2pt}\\
\textit{\color{gray}// compinit is the most expensive call in zsh startup. defer it. //}\\
\color{white}The completion system (compsys) needs compinit to initialize.\\
\color{white}But compinit scans all fpath entries and parses completion\\
\color{white}files. On a system with 400+ completions, this is slow.\\
\color{white}ZPWR defers it behind turbo mode.\\
\vspace{4pt}\textbf{\color{magenta}ZINIT COMPINIT OPTIONS}\\
\color{white}ZINIT[COMPINIT\_OPTS]='-C -u'\\
\color{white}-C skips security checks (don't re-verify file ownership).\\
\color{white}-u suppresses warnings about insecure directories.\\
\color{white}These flags make compinit fast but trust your fpath.\\
\color{white}ZINIT[ZCOMPDUMP\_PATH]=\textbackslash{}"\textbackslash{}SH\_COMPDUMP\textbackslash{}"\\
\color{white}Points compinit to zpwr's zcompdump location.\\
\vspace{4pt}\textbf{\color{magenta}THE DEFERRED COMPINIT}\\
\color{white}Compinit runs in the LAST turbo phase:\\
\color{white}zinit ice lucid nocompile nocd as'null'\\
\color{white}wait\textbackslash{}"\textbackslash{}\$\{ZPWR\_ZINIT\_COMPINIT\_DELAY\}g\textbackslash{}"\\
\texttt{\color{green}atinit'zicompinit; zicdreplay; zpwrBindOverrideOMZCompdefs'}\\
\color{white}zicompinit: runs compinit with the configured opts.\\
\color{white}zicdreplay: replays all compdef calls that were recorded\\
\color{white}while compinit wasn't loaded yet. Plugins that called compdef\\
\color{white}during turbo load had their calls queued -- now they fire.\\
\vspace{4pt}\textbf{\color{magenta}THE OMZ COMPLETION SYMLINK HACK}\\
\color{white}OMZ completion snippets land in:\\
\color{white}\textbackslash{}SH/snippets/OMZP::coffee/\_coffee\\
\color{white}But compinit needs them in the completions dir. ZPWR\\
\color{white}symlinks them using zsh/files builtin:\\
\color{white}zmodload -F zsh/files b:zf\_ln\\
\color{white}zf\_ln -sfn \textbackslash{}SH/snippets/OMZP::\textbackslash{}\$\{p\%/*\}/... \textbackslash{}SH/completions/...\\
\color{white}zf\_ln is a builtin ln(1), no fork needed. Fast.\\
\vspace{4pt}\textbf{\color{magenta}zpwrStaleZcompdump}\\
\color{white}The zcompdump cache gets stale when plugins update.\\
\color{white}zpwrStaleZcompdump checks the cache age. If it is older\\
\color{white}than \textasciitilde{}1 week, it triggers a full compinit rebuild.\\
\color{white}This runs automatically every shell startup (cheap check).\\
\vspace{4pt}\textbf{\color{magenta}FORCE REBUILD}\\
\texttt{\color{green}zpwr regenzsh}\\
\color{white}Nukes the zcompdump, re-runs compinit from scratch, and\\
\color{white}recompiles everything. Use when completions go missing or\\
\color{white}after adding new completion files to fpath.\\
\vspace{4pt}\textbf{\color{magenta}ZPWR\_ZINIT\_COMPINIT\_DELAY}\\
\color{white}Default: 0. Set in \textbackslash{}PWR\_ENV/.zpwr\_env.sh.\\
\color{white}Controls the numeric prefix of the wait ice for compinit.\\
\color{white}0 = load immediately after prompt. Higher = more delay.\\
\color{white}On fast machines, 0 is fine. On slow machines, 1 or 2\\
\color{white}lets other plugins finish first.\\

\vspace{8pt}

\subsection{Page 136: Plugin Troubleshooting}

\vspace{4pt}\textbf{\color{cyan}=== PLUGIN TROUBLESHOOTING ===}\vspace{2pt}\\
\textit{\color{gray}// when plugins misbehave, you need diagnostic tools. //}\\
\color{white}Zinit ships with built-in debugging commands.\\
\color{white}When something breaks, these are your first responders.\\
\vspace{4pt}\textbf{\color{magenta}DIAGNOSTIC COMMANDS}\\
\color{white}zinit report <plugin>      Show load status, errors, files\\
\color{white}zinit report zsh-users/zsh-autosuggestions\\
\color{white}Shows: source file, ice modifiers used, compdef calls,\\
\color{white}functions defined, aliases created, PATH changes.\\
\color{white}zinit times               Show plugin load times\\
\color{white}Sorted by load time. Find the slowest plugin.\\
\color{white}If something takes >50ms, investigate.\\
\color{white}zinit loaded              List all loaded plugins\\
\color{white}zinit unload <plugin>      Unload a plugin from current session\\
\color{white}Useful for debugging conflicts. Does not persist.\\
\color{white}zinit list                List all registered plugins\\
\color{white}zinit cd <plugin>          cd into plugin directory\\
\vspace{4pt}\textbf{\color{magenta}COMMON ISSUES}\\
\color{white}Plugin not loading?\\
\color{white}- Check if 'wait' ice is set (turbo mode defers loading)\\
\color{white}- Run zinit report to see if it errored silently\\
\color{white}- The 'lucid' ice suppresses messages -- remove it to debug\\
\color{white}Completion missing after plugin install?\\
\texttt{\color{green}- Run zpwr regenzsh to rebuild zcompdump}\\
\color{white}- Check fpath: print -l \textbackslash{}\$fpath | grep <plugin>\\
\color{white}- Ensure the plugin uses as'completion' or adds to fpath\\
\color{white}Keybinding conflicts?\\
\color{white}- Load order matters. Plugins loaded later override earlier ones.\\
\color{white}- ZPWR uses wait'0a' through wait'0e' to control this.\\
\color{white}- zpwrBindOverrideZLE fires in the 0a phase to win conflicts.\\
\vspace{4pt}\textbf{\color{magenta}ADDING CUSTOM PLUGINS}\\
\color{white}Add to ZPWR\_GH\_PLUGINS in \textbackslash{}PWR\_LOCAL/.tokens.sh:\\
\color{white}ZPWR\_GH\_PLUGINS+=( your-user/your-plugin )\\
\color{white}The .tokens.sh file is sourced before plugin loading starts\\
\color{white}(via zpwrTokenPre), so your additions join the array in time.\\
\color{white}For conditional loading:\\
\texttt{\color{green}zpwrCommandExists myapp \&\& ZPWR\_GH\_PLUGINS+=( user/zsh-myapp )}\\
\vspace{4pt}\textbf{\color{magenta}UPDATING PLUGINS}\\
\texttt{\color{green}zpwr updatezinit}\\
\color{white}Runs zinit update --all to pull latest from all repos.\\
\texttt{\color{green}After updating, run zpwr regenzsh to rebuild completions.}\\
\vspace{4pt}\textbf{\color{magenta}NUCLEAR OPTION}\\
\color{white}If zinit is truly broken, nuke and rebuild:\\
\color{white}rm -rf \textbackslash{}PWR\_PLUGIN\_MANAGER\_HOME\\
\color{white}Next shell start auto-clones zinit and re-downloads everything.\\
\color{white}Slow but thorough. The cockroach approach to debugging.\\

\vspace{8pt}


\section{CHAPTER 28: ZSH INTERNALS}

\subsection{Page 137: Autoloading: Lazy Function Loading}

\vspace{4pt}\textbf{\color{cyan}=== AUTOLOADING: LAZY FUNCTION LOADING ===}\vspace{2pt}\\
\textit{\color{gray}// 443 functions. zero parsing until you call one. //}\\
\color{white}Autoloading is zsh's lazy-loading mechanism. Register a name,\\
\color{white}and zsh only reads the function body from disk on first call.\\
\color{white}This is why ZPWR can have 443+ functions with negligible startup cost.\\
\vspace{4pt}\textbf{\color{magenta}HOW autoload -z WORKS}\\
\color{white}1. autoload -z funcname registers 'funcname' as autoloadable.\\
\color{white}2. Zsh records the name but does NOT read the file yet.\\
\color{white}3. On first call to funcname, zsh searches \textbackslash{}\$fpath for a file\\
\color{white}named 'funcname' (exact match, no extension).\\
\color{white}4. The file is sourced. The function is now defined in memory.\\
\color{white}5. The function executes with the original arguments.\\
\color{white}6. Subsequent calls use the in-memory version. No more disk.\\
\vspace{4pt}\textbf{\color{magenta}THE GLOB IN .zshrc}\\
\color{white}builtin autoload -z \textbackslash{}PWR\_AUTOLOAD\_COMMON/*(.:t)\\
\color{white}The glob *(.:t) means:\\
\color{white}*    All files in the directory\\
\color{white}(.)  Qualifier: regular files only (not dirs or symlinks)\\
\color{white}:t   Modifier: tail only (strip directory, keep filename)\\
\color{white}So it registers every regular file's name as a function.\\
\color{white}Same pattern runs for ZPWR\_AUTOLOAD\_DARWIN or ZPWR\_AUTOLOAD\_LINUX.\\
\vspace{4pt}\textbf{\color{magenta}THE fpath}\\
\color{white}fpath is the function search path (like PATH but for functions):\\
\color{white}fpath=( \textbackslash{}PWR\_AUTOLOAD\_COMMON \textbackslash{}PWR\_AUTOLOAD\_COMP\_UTILS \textbackslash{}PWR\_COMPS \textbackslash{}\$fpath )\\
\color{white}Order matters: first match wins. ZPWR dirs come first so they\\
\color{white}can override system completions if needed.\\
\color{white}fpath is NOT exported (typeset +x FPATH) -- child shells don't inherit it.\\
\vspace{4pt}\textbf{\color{magenta}THE FUNCTION FILE FORMAT}\\
\color{white}Each file in the autoload dir contains one function. The file\\
\color{white}is named exactly like the function. Inside:\\
\texttt{\color{green}function zpwrSomething() \{}\\
\color{white}\# function body\\
\color{white}\}\\
\texttt{\color{green}zpwrSomething \textbackslash{}"\textbackslash{}\$@\textbackslash{}"}\\
\color{white}The last line calls the function with args passed through.\\
\color{white}This is the ZPWR convention. The function wraps itself.\\
\vspace{4pt}\textbf{\color{magenta}WHY AUTOLOADING IS FAST}\\
\color{white}443 function registrations = 443 hash table entries.\\
\color{white}Cost: microseconds. No file I/O, no parsing, no execution.\\
\color{white}Compare to sourcing 443 files: hundreds of milliseconds.\\
\color{white}Autoloading turns O(n) startup into O(1) with lazy O(1) per call.\\
\vspace{4pt}\textbf{\color{magenta}autoload -Uz}\\
\color{white}-U suppresses alias expansion during loading (safety).\\
\color{white}-z forces zsh-style function loading (vs ksh-style).\\
\color{white}ZPWR uses -z for its functions, -Uz for system utilities like\\
\color{white}zrecompile, zmv, zargs, and compinit.\\

\vspace{8pt}

\subsection{Page 138: ZLE: The Zsh Line Editor}

\vspace{4pt}\textbf{\color{cyan}=== ZLE: THE ZSH LINE EDITOR ===}\vspace{2pt}\\
\textit{\color{gray}// the engine behind every keystroke. readline on steroids. //}\\
\color{white}ZLE (Zsh Line Editor) is the subsystem that handles your\\
\color{white}command line input. Every key press triggers a ZLE widget.\\
\color{white}ZPWR rewires dozens of these widgets for custom behavior.\\
\vspace{4pt}\textbf{\color{magenta}WIDGETS: NAMED EDITOR ACTIONS}\\
\color{white}A widget is a named action that ZLE can execute:\\
\color{white}.accept-line          Built-in: execute the current line\\
\color{white}.backward-delete-char  Built-in: backspace\\
\color{white}.self-insert           Built-in: insert the typed character\\
\color{white}The dot prefix means 'built-in widget'. User-defined widgets\\
\color{white}can override them. ZPWR defines custom widgets like:\\
\texttt{\color{green}zpwrMagicEnter         Custom Enter key behavior}\\
\texttt{\color{green}zpwrExpandTerminateSpace Custom Space key behavior}\\
\texttt{\color{green}zpwrAcceptLine          Custom line acceptance}\\
\vspace{4pt}\textbf{\color{magenta}bindkey: MAPPING KEYS TO WIDGETS}\\
\color{white}bindkey connects key sequences to widgets:\\
\texttt{\color{green}bindkey '\textasciicircum{}M' zpwrMagicEnter    Enter key}\\
\texttt{\color{green}bindkey ' ' zpwrExpandTerminateSpace  Space key}\\
\color{white}bindkey -v                     Switch to vi mode\\
\color{white}bindkey -e                     Switch to emacs mode\\
\color{white}ZPWR\_BINDKEY\_VI controls which mode: true=vi, false=emacs.\\
\vspace{4pt}\textbf{\color{magenta}HOW zpwrMagicEnter WORKS}\\
\color{white}On empty line (just pressing Enter):\\
\color{white}- If in a git dir: runs git status -u .\\
\color{white}- If not: runs ls -lh .\\
\color{white}- Resets the prompt via \_p9k\_precmd\_impl + zle reset-prompt\\
\color{white}On non-empty line: delegates to zle accept-line (zpwrAcceptLine).\\
\vspace{4pt}\textbf{\color{magenta}HOW zpwrAcceptLine WORKS}\\
\color{white}The accept-line replacement does serious work:\\
\color{white}1. Updates zconvey name with current command + TTY + PID\\
\color{white}2. Detects commands that modify files (rm, mv, cp, mkdir...)\\
\color{white}and sets ZPWR\_WILL\_CLEAR=true for post-exec ls refresh\\
\color{white}3. Handles sudo + grc colorization for aliased commands\\
\color{white}4. Prevents global alias expansion in first position\\
\color{white}5. Optionally expands aliases before execution\\
\color{white}6. Calls .accept-line (the real built-in)\\
\vspace{4pt}\textbf{\color{magenta}THE Ctrl-F CHORD PREFIX SYSTEM}\\
\color{white}ZPWR binds Ctrl-F as a prefix key for multi-key chords:\\
\color{white}Ctrl-F Ctrl-D   FZF file search\\
\color{white}Ctrl-F Ctrl-I   FZF word search (ag/grep)\\
\color{white}Ctrl-F Ctrl-G   FZF git operations\\
\color{white}This gives you an entire second keymap layer without\\
\color{white}conflicting with normal editing keys. Like tmux prefix\\
\color{white}but for your command line.\\

\vspace{8pt}

\subsection{Page 139: Hook Functions: precmd, preexec, chpwd}

\vspace{4pt}\textbf{\color{cyan}=== HOOK FUNCTIONS: precmd, preexec, chpwd ===}\vspace{2pt}\\
\textit{\color{gray}// zsh calls your code at key moments. event-driven shell. //}\\
\color{white}Zsh provides hook points where you can inject code that\\
\color{white}runs automatically before prompts, before commands, and\\
\color{white}when directories change. ZPWR uses all three.\\
\vspace{4pt}\textbf{\color{magenta}precmd: BEFORE EACH PROMPT}\\
\color{white}Runs after a command finishes, right before the prompt draws.\\
\color{white}p10k's \_p9k\_precmd is the main user -- it updates git status,\\
\color{white}execution time, and every prompt segment.\\
\color{white}ZPWR adds zpwrPrecmd to the hook:\\
\texttt{\color{green}function zpwrPrecmd()\{ export ZPWR\_OPTS=\textbackslash{}"\textbackslash{}\$-\textbackslash{}" \}}\\
\color{white}This captures the current shell options before each prompt.\\
\vspace{4pt}\textbf{\color{magenta}preexec: BEFORE EACH COMMAND}\\
\color{white}Runs after Enter but BEFORE the command executes.\\
\color{white}Receives the command string as \textbackslash{}\$1.\\
\color{white}ZPWR's zpwrPreexec is minimal (a placeholder hook).\\
\color{white}p10k uses preexec to record the start timestamp for\\
\color{white}command\_execution\_time segment.\\
\vspace{4pt}\textbf{\color{magenta}chpwd: WHEN DIRECTORY CHANGES}\\
\color{white}Runs whenever \textbackslash{}WD changes (cd, pushd, popd, auto\_cd).\\
\color{white}ZPWR's zpwrChpwd hook handles directory-change actions.\\
\vspace{4pt}\textbf{\color{magenta}THE \_functions ARRAYS}\\
\color{white}Zsh supports multiple hooks via arrays:\\
\color{white}precmd\_functions+=( myHook )      Add a precmd hook\\
\color{white}preexec\_functions+=( myHook )     Add a preexec hook\\
\color{white}chpwd\_functions+=( myHook )       Add a chpwd hook\\
\color{white}Every function in the array runs at the hook point.\\
\color{white}Order matters: first in array = first to execute.\\
\vspace{4pt}\textbf{\color{magenta}HOW ZPWR REGISTERS HOOKS}\\
\color{white}zpwrBindPrecmd (runs during p10k atload):\\
\texttt{\color{green}precmd\_functions=( zpwrPrecmd \textbackslash{}\$\{(@)precmd\_functions:\#zpwrPrecmd\} )}\\
\color{white}This puts zpwrPrecmd FIRST, removing any duplicate. The\\
\color{white}\textbackslash{}\$\{(@)array:\#pattern\} syntax removes matching elements.\\
\color{white}zpwrBindPreexecChpwd (runs during syntax-highlighting atload):\\
\texttt{\color{green}preexec\_functions=( \textbackslash{}\$\{(@)preexec\_functions:\#zpwrPreexec\} zpwrPreexec )}\\
\texttt{\color{green}chpwd\_functions=( \textbackslash{}\$\{(@)chpwd\_functions:\#zpwrChpwd\} zpwrChpwd )}\\
\color{white}These put zpwr hooks LAST (after p10k and other plugins).\\
\vspace{4pt}\textbf{\color{magenta}THE HOOK COUNT IN zpwr top}\\
\texttt{\color{green}The zpwr top dashboard shows the total hook count:}\\
\color{white}nhooks=\textbackslash{}\$((\textbackslash{}\$\{\#precmd\_functions\}+\textbackslash{}\$\{\#preexec\_functions\}+\textbackslash{}\$\{\#chpwd\_functions\}))\\
\color{white}Typical value: 5-8 hooks total across all three arrays.\\
\color{white}If this number grows unexpectedly, a plugin is adding hooks\\
\color{white}you may not want. Check with: print -l \textbackslash{}\$precmd\_functions\\

\vspace{8pt}

\subsection{Page 140: Emulate \& Options}

\vspace{4pt}\textbf{\color{cyan}=== EMULATE \& OPTIONS ===}\vspace{2pt}\\
\textit{\color{gray}// zsh has 200+ options. emulate -L zsh resets them all locally. //}\\
\color{white}Zsh options control fundamental behavior: globbing, word\\
\color{white}splitting, quoting, history. ZPWR functions use emulate -L zsh\\
\color{white}to guarantee a known environment regardless of user settings.\\
\vspace{4pt}\textbf{\color{magenta}emulate -L zsh}\\
\color{white}builtin emulate -L zsh\\
\color{white}Resets ALL options to zsh defaults for the current function scope.\\
\color{white}-L means local: options revert when the function returns.\\
\color{white}Without this, a function might inherit noglob from the caller\\
\color{white}and break globbing, or inherit shwordsplit and break arrays.\\
\color{white}ZPWR functions start with this to be immune to the caller's state.\\
\vspace{4pt}\textbf{\color{magenta}KEY OPTIONS IN ZPWR}\\
\color{white}extendedglob      \#\# = one-or-more, (\#i) = case-insensitive glob\\
\color{white}\textasciicircum{}pattern = negation. Essential for ZPWR globs.\\
\color{white}glob\_dots         Dot files matched by * (no need for .*)\\
\color{white}globstarshort     **.c = **/*.c (recursive glob shorthand)\\
\color{white}braceccl          \{abcd0-9\} expands as character class\\
\color{white}rc\_quotes         '' inside single quotes = literal single quote\\
\color{white}rc\_expand\_param   Array expansion includes prefix: xx\textbackslash{}\$\{arr\} = xxA xxB\\
\color{white}nocaseglob        Globs are case-insensitive globally\\
\color{white}numeric\_glob\_sort  file1 file2 file10 sort numerically\\
\color{white}prompt\_subst       Allow \textbackslash{}\$(...) and \textbackslash{}\$\{...\} in prompt strings\\
\color{white}transient\_rprompt  Right prompt only on current line\\
\vspace{4pt}\textbf{\color{magenta}THE noshwordsplit DEFAULT}\\
\color{white}Bash splits unquoted variables on whitespace by default.\\
\color{white}Zsh does NOT (noshwordsplit). This means:\\
\color{white}x=\textbackslash{}"hello world\textbackslash{}"; print \textbackslash{}\$x  -> one argument in zsh, two in bash\\
\color{white}If a function inherits shwordsplit (e.g., from emulate sh),\\
\color{white}arrays and parameter expansions break silently. Another reason\\
\color{white}for emulate -L zsh.\\
\vspace{4pt}\textbf{\color{magenta}THE local-IN-LOOP BUG}\\
\color{white}In zsh, 'local' inside a loop prints the variable value as\\
\color{white}a side effect when the variable already exists:\\
\color{white}for i in 1 2 3; do local x=\textbackslash{}\$i; done  \# prints x=1 x=2 x=3\\
\color{white}This is because 'local' on an existing local acts like typeset,\\
\color{white}which prints. Declare locals BEFORE loops, not inside them.\\
\vspace{4pt}\textbf{\color{magenta}auto\_name\_dirs: REMOVED}\\
\color{white}This option auto-registers any parameter set to a directory\\
\color{white}path as a named directory. Sounds nice, but p10k/gitstatus\\
\color{white}set internal variables to paths, polluting the named dir table.\\
\color{white}ZPWR removed it. Explicit hash -d in zpwrBindDirs handles it.\\
\color{white}The comment in .zshrc: \# auto\_name\_dirs removed: pollutes named\\
\color{white}\# dirs with p10k/gitstatus internals\\

\vspace{8pt}

\subsection{Page 141: Parameter Expansion Magic}

\vspace{4pt}\textbf{\color{cyan}=== PARAMETER EXPANSION MAGIC ===}\vspace{2pt}\\
\textit{\color{gray}// zsh parameter expansion is a programming language in itself. //}\\
\color{white}Zsh's \textbackslash{}\$\{...\} syntax can do transformations that would\\
\color{white}require sed, awk, basename, dirname, tr, and cut in bash.\\
\color{white}No forks. No subshells. Pure in-process string manipulation.\\
\vspace{4pt}\textbf{\color{magenta}PARAMETER EXPANSION FLAGS (the parenthesized letters)}\\
\color{white}\textbackslash{}\$\{(k)assoc\}          Keys of associative array\\
\color{white}\textbackslash{}\$\{(v)assoc\}          Values of associative array\\
\color{white}\textbackslash{}\$\{(kv)assoc\}         Key-value pairs interleaved\\
\color{white}\textbackslash{}\$\{(s:/:)var\}         Split string on / into array\\
\color{white}\textbackslash{}\$\{(j:,:)array\}       Join array elements with ,\\
\color{white}\textbackslash{}\$\{(L)var\}            Lowercase the entire string\\
\color{white}\textbackslash{}\$\{(U)var\}            Uppercase the entire string\\
\color{white}\textbackslash{}\$\{(f)var\}            Split on newlines (lines to array)\\
\color{white}\textbackslash{}\$\{(F)array\}          Join with newlines (array to lines)\\
\color{white}\textbackslash{}\$\{(t)var\}            Type of variable (scalar, array, etc.)\\
\color{white}\textbackslash{}\$\{(q)var\}            Quote the value for re-use in eval\\
\color{white}\textbackslash{}\$\{(Q)var\}            Remove one level of quoting\\
\color{white}\textbackslash{}\$\{(z)var\}            Split into words like the shell parser\\
\color{white}\textbackslash{}\$\{(o)array\}          Sort ascending\\
\color{white}\textbackslash{}\$\{(O)array\}          Sort descending\\
\vspace{4pt}\textbf{\color{magenta}MODIFIERS (the colon letters)}\\
\color{white}\textbackslash{}\$\{var:t\}             Tail: basename (/a/b/c.txt -> c.txt)\\
\color{white}\textbackslash{}\$\{var:h\}             Head: dirname (/a/b/c.txt -> /a/b)\\
\color{white}\textbackslash{}\$\{var:r\}             Root: remove extension (file.txt -> file)\\
\color{white}\textbackslash{}\$\{var:e\}             Extension: (file.txt -> txt)\\
\color{white}\textbackslash{}\$\{var:A\}             Absolute path (resolve symlinks)\\
\color{white}\textbackslash{}\$\{var:l\}             Lowercase (zsh 5.0+)\\
\color{white}\textbackslash{}\$\{var:u\}             Uppercase (zsh 5.0+)\\
\vspace{4pt}\textbf{\color{magenta}SUBSTITUTION AND PATTERN MATCHING}\\
\textit{\color{gray}\textbackslash{}\$\{var//pat/rep\}       Global substitution}\\
\color{white}\textbackslash{}\$\{var/pat/rep\}        First match substitution\\
\color{white}\textbackslash{}\$\{var\#pat\}            Remove shortest prefix match\\
\color{white}\textbackslash{}\$\{var\#\#pat\}           Remove longest prefix match\\
\color{white}\textbackslash{}\$\{var\%pat\}            Remove shortest suffix match\\
\color{white}\textbackslash{}\$\{var\%\%pat\}           Remove longest suffix match\\
\vspace{4pt}\textbf{\color{magenta}ZPWR EXAMPLES}\\
\color{white}From zpwrBindDirs (writing hash dir cache):\\
\color{white}print -r \textbackslash{}"hash -d \textbackslash{}\$\{k\}=\textbackslash{}\$\{(q)v\}\textbackslash{}"  -- (q) quotes the value safely\\
\color{white}From zpwrAcceptLine (parsing command buffer):\\
\color{white}mywords=( \textbackslash{}\$\{(z)cmdstr\} )           -- (z) shell-word splits\\
\color{white}cmd=\textbackslash{}\$\{mywords[1]\#=\}                -- strip leading = from =cmd\\
\color{white}From zpwrBindPrecmd (removing duplicates from array):\\
\texttt{\color{green}\textbackslash{}\$\{(@)precmd\_functions:\#zpwrPrecmd\}  -- remove matching element}\\
\color{white}These are not academic exercises. They replace forks.\\

\vspace{8pt}

\subsection{Page 142: The .zwc Compilation System}

\vspace{4pt}\textbf{\color{cyan}=== THE .zwc COMPILATION SYSTEM ===}\vspace{2pt}\\
\textit{\color{gray}// bytecode for your shell. because even scripts deserve a compiler. //}\\
\color{white}Zsh can compile scripts to .zwc (Zsh Word Code) bytecode.\\
\color{white}These files skip the parsing phase on load, going straight\\
\color{white}to execution. ZPWR compiles 50+ files for speed.\\
\vspace{4pt}\textbf{\color{magenta}WHAT .zwc FILES ARE}\\
\color{white}A .zwc file is a binary representation of parsed zsh code.\\
\color{white}When zsh sources a .zsh file, it checks for a .zwc sibling.\\
\color{white}If the .zwc exists and is newer than the source, zsh loads\\
\color{white}the bytecode directly. No lexing, no parsing.\\
\color{white}Speedup: 2-5x for large files. Marginal for tiny ones.\\
\vspace{4pt}\textbf{\color{magenta}zrecompile}\\
\color{white}The zrecompile utility compiles only if the source is newer:\\
\color{white}autoload -Uz zrecompile\\
\color{white}zrecompile -p -U \textbackslash{}PWR\_ENV\_FILE\\
\color{white}-p means 'check pairs' (source vs .zwc timestamps).\\
\color{white}-U suppresses alias expansion during compilation.\\
\color{white}If the .zwc is current, zrecompile is a no-op.\\
\vspace{4pt}\textbf{\color{magenta}zpwrRecompileFiles: WHAT GETS COMPILED}\\
\color{white}The function compiles two categories of files:\\
\color{white}Local files (\textasciitilde{}47):\\
\color{white}\textasciitilde{}/.zshrc, \textasciitilde{}/.zshenv, \textasciitilde{}/.zlogout, \textasciitilde{}/.zlogin\\
\color{white}p10k-instant-prompt-*.zsh, p10k-dump-*.zsh\\
\color{white}\textbackslash{}SH\_COMPDUMP, \textbackslash{}PWR\_ALIAS\_FILE, \textbackslash{}PWR\_PROMPT\_FILE\\
\color{white}\textbackslash{}PWR\_LOCAL/.tokens.sh, \textbackslash{}PWR\_ENV\_FILE, \textbackslash{}PWR\_RE\_ENV\_FILE\\
\color{white}All .zsh/.sh in \textbackslash{}PWR\_ENV/, all .zsh in \textbackslash{}PWR\_SCRIPTS/\\
\color{white}All plugin .zsh files under \textbackslash{}PWR\_PLUGIN\_MANAGER\_HOME/bin/\\
\color{white}OMZ lib, plugin, theme, and custom files\\
\color{white}Global files (\textasciitilde{}8):\\
\color{white}/etc/profile, /etc/zshrc, /etc/zshenv, /etc/zprofile, etc.\\
\color{white}These need sudo (compiled with sudo zsh -c zrecompile).\\
\vspace{4pt}\textbf{\color{magenta}ZPWR VERBS FOR COMPILATION}\\
\texttt{\color{green}zpwr recompile       Compile all known files (zpwrRecompileFiles)}\\
\texttt{\color{green}zpwr decompile       Remove .zwc files (zpwrUncompile)}\\
\texttt{\color{green}zpwr refreshzwc      Remove + recompile (clean rebuild)}\\
\vspace{4pt}\textbf{\color{magenta}THE zsh/files MODULE}\\
\color{white}Zsh can load builtin file operations via zmodload:\\
\color{white}zmodload -F zsh/files b:zf\_ln b:zf\_rm b:zf\_mkdir\\
\color{white}These are shell builtins that replace ln, rm, mkdir.\\
\color{white}No fork, no exec -- pure kernel syscalls from the shell process.\\
\color{white}ZPWR uses zf\_ln for the OMZ completion symlink hack.\\
\color{white}In a loop creating 12 symlinks, this saves 12 fork+exec pairs.\\
\vspace{4pt}\textbf{\color{magenta}OTHER USEFUL MODULES}\\
\color{white}zsh/datetime    \textbackslash{}POCHREALTIME, \textbackslash{}POCHSECONDS\\
\color{white}zsh/complist    Menu-select completion mode\\
\color{white}zsh/zprof       Profiler (ZPWR\_PROFILING=true enables it)\\
\color{white}Modules are loaded with: zmodload zsh/datetime\\

\vspace{8pt}


\section{CHAPTER 29: FZF ARCHITECTURE}

\subsection{Page 143: FZF: The Neural Interface}

\vspace{4pt}\textbf{\color{cyan}=== FZF: THE NEURAL INTERFACE ===}\vspace{2pt}\\
\textit{\color{gray}// the universal adapter between data and decisions. //}\\
\color{white}fzf (fuzzy finder) is a general-purpose command-line filter.\\
\color{white}Pipe anything into it, fuzzy-search interactively, get results.\\
\color{white}ZPWR wraps fzf into dozens of workflows: files, words, git, verbs.\\
\vspace{4pt}\textbf{\color{magenta}THE PATTERN: PIPE | FZF | ACTION}\\
\color{white}Every fzf integration follows the same architecture:\\
\color{white}<data source> | fzf [options] | <action on selection>\\
\color{white}The data source generates lines. fzf presents them. The action\\
\color{white}consumes the selected line(s). Simple. Composable. Powerful.\\
\vspace{4pt}\textbf{\color{magenta}FZF\_DEFAULT\_OPTS}\\
\color{white}ZPWR sets global fzf defaults in zpwrBindFZFLate:\\
\color{white}FZF\_DEFAULT\_OPTS=\textbackslash{}"\textbackslash{}PWR\_COMMON\_FZF\_ELEM --reverse --border --height 100\%\textbackslash{}"\\
\color{white}Plus a Solarized Dark color scheme layered on top:\\
\color{white}--color fg:-1,bg:-1,hl:\textbackslash{}\$blue,fg+:\textbackslash{}\$base2,bg+:\textbackslash{}\$base02,hl+:\textbackslash{}\$blue\\
\color{white}ZPWR\_COMMON\_FZF\_ELEM adds the logo prompt and navigation binds:\\
\color{white}--prompt='\textbackslash{}PWR\_FZF\_LOGO ' --bind=ctrl-n:page-down,...\\
\vspace{4pt}\textbf{\color{magenta}THE PREVIEW WINDOW}\\
\color{white}fzf's --preview flag runs a command for each highlighted item:\\
\color{white}fzf --preview 'bat --color=always \{\}'\\
\color{white}ZPWR configures rich previews per context:\\
\color{white}Files:   bat with syntax highlighting, line numbers, grid\\
\color{white}Dirs:    stat output\\
\color{white}Packages: dpkg -I or rpm -qi\\
\color{white}Git:     diff preview, commit log\\
\color{white}If bat is installed, it's the default colorizer. Falls back\\
\color{white}to pygmentize with terminal256 output.\\
\vspace{4pt}\textbf{\color{magenta}fzf-completion (Tab)}\\
\color{white}fzf replaces the default Tab completion with a fuzzy finder.\\
\color{white}Type a command, press Tab, and fzf opens with completions:\\
\color{white}cd **<Tab>     fzf through subdirectories\\
\color{white}kill <Tab>      fzf through running processes\\
\color{white}ssh **<Tab>     fzf through known hosts\\
\color{white}FZF\_COMPLETION\_OPTS controls the preview for these.\\
\vspace{4pt}\textbf{\color{magenta}fzf-tab PLUGIN}\\
\color{white}MenkeTechnologies/fzf-tab replaces zsh's menu completion with fzf.\\
\color{white}Any time the completion system would show a menu, fzf-tab\\
\color{white}intercepts it and shows fzf instead. This means:\\
\color{white}- Fuzzy matching for every completion, not just ** triggers\\
\color{white}- Previews for file completions, git branches, everything\\
\color{white}- FZF\_TAB\_OPTS configures the behavior (ansi, nth, delimiter)\\
\color{white}- FZF\_TAB\_COMMAND sets which fzf binary to use\\
\vspace{4pt}\textbf{\color{magenta}GLOBAL ALIASES FOR FZF}\\
\color{white}\textbackslash{}\$\{ZPWR\_GLOBAL\_ALIAS\_PREFIX\}z   Pipe into fzf: ls | ,,z\\
\color{white}\textbackslash{}\$\{ZPWR\_GLOBAL\_ALIAS\_PREFIX\}zb  Inline fzf selection: vim ,,zb\\

\vspace{8pt}

\subsection{Page 144: FZF in Practice}

\vspace{4pt}\textbf{\color{cyan}=== FZF IN PRACTICE ===}\vspace{2pt}\\
\textit{\color{gray}// the theory is simple. the practice is where it gets good. //}\\
\color{white}ZPWR ships dozens of fzf-powered workflows. Each follows\\
\color{white}the generate | fzf | action pattern. Here are the heaviest hitters.\\
\vspace{4pt}\textbf{\color{magenta}FILE SEARCH (zpwr filesearch / Ctrl-F Ctrl-D)}\\
\color{white}Pattern: find files | fzf --preview 'bat' | open in editor\\
\color{white}1. find/fd generates a file list (respecting .gitignore)\\
\color{white}2. fzf shows files with bat-powered syntax-highlighted preview\\
\color{white}3. Selected file opens in \textbackslash{}DITOR at the matched location\\
\color{white}FZF\_CTRL\_T\_OPTS configures the preview command.\\
\vspace{4pt}\textbf{\color{magenta}WORD SEARCH (zpwr wordsearch / Ctrl-F Ctrl-I)}\\
\color{white}Pattern: ag/rg search | fzf --preview 'context' | open at line\\
\color{white}1. ag (silver searcher) or rg searches file contents\\
\color{white}2. fzf shows matches with file:line:content format\\
\color{white}3. Selection opens the file at the exact line number\\
\color{white}FZF\_AG\_OPTS: delimiter ':', match on 3rd+ field (skip file:line).\\
\vspace{4pt}\textbf{\color{magenta}GIT ADD (forgitadd)}\\
\color{white}Pattern: git diff --stat | fzf --preview 'git diff' | git add\\
\color{white}1. Lists modified files with diff stats\\
\color{white}2. Preview pane shows the actual diff for each file\\
\color{white}3. Toggle files with Tab, confirm to stage them\\
\color{white}Part of forgit: the git + fzf integration plugin.\\
\vspace{4pt}\textbf{\color{magenta}VERB SEARCH (zpwr verbsfzf)}\\
\color{white}Pattern: ZPWR\_VERBS keys | fzf --preview 'help' | eval\\
\color{white}1. All 408+ zpwr verbs piped into fzf\\
\color{white}2. Preview shows the verb description and command\\
\color{white}3. Select a verb to execute it immediately\\
\color{white}FZF\_ENV\_OPTS\_VERBS configures the preview for verbs.\\
\vspace{4pt}\textbf{\color{magenta}CUSTOM FZF KEY BINDINGS}\\
\color{white}The Ctrl-F prefix opens ZPWR's fzf keymap layer:\\
\color{white}Ctrl-F Ctrl-D   File search (find + fzf + editor)\\
\color{white}Ctrl-F Ctrl-I   Word search (ag + fzf + editor)\\
\color{white}Ctrl-F Ctrl-G   Git search (git log + fzf)\\
\color{white}Ctrl-F Ctrl-R   History search with fzf\\
\color{white}These are ZLE widgets bound via zpwrBindOverrideZLE.\\
\vspace{4pt}\textbf{\color{magenta}zpwrIntoFzf AND zpwrIntoFzfAg}\\
\color{white}Two core ZLE widgets that bridge the command line with fzf:\\
\color{white}zpwrIntoFzf     File-oriented fzf (Ctrl-T style)\\
\color{white}zpwrIntoFzfAg   Content-oriented fzf (ag/grep style)\\
\color{white}Both capture fzf output and insert it into the command line\\
\color{white}buffer (\textbackslash{}UFFER). You select files/matches, they appear\\
\color{white}at your cursor. Seamless integration between search and input.\\
\vspace{4pt}\textbf{\color{magenta}THE COLORIZER CHAIN}\\
\color{white}ZPWR\_COLORIZER preference: bat > pygmentize.\\
\color{white}bat uses themes (BAT\_THEME=\textbackslash{}PWR\_BAT\_THEME).\\
\color{white}pygmentize uses PYGMENTIZE\_COLOR style.\\
\color{white}Language-specific colorizers: FZF\_COLORIZER\_C, \_SH, \_YAML, \_JAVA.\\

\vspace{8pt}

\subsection{Page 145: Building Your Own FZF Verbs}

\vspace{4pt}\textbf{\color{cyan}=== BUILDING YOUR OWN FZF VERBS ===}\vspace{2pt}\\
\textit{\color{gray}// once you grok the pattern, everything becomes fzf-able. //}\\
\color{white}The generate | fzf | action pattern is infinitely extensible.\\
\color{white}Here's how to build your own fzf-powered verbs.\\
\vspace{4pt}\textbf{\color{magenta}THE PATTERN}\\
\color{white}<generator> | fzf [--preview '<preview-cmd>'] | <consumer>\\
\color{white}Generator: anything that outputs lines (find, git, docker, etc.)\\
\color{white}Preview: command that shows detail for highlighted line (\textbackslash{}\{\textbackslash{}\})\\
\color{white}Consumer: action on selected lines (xargs, while read, eval)\\
\vspace{4pt}\textbf{\color{magenta}EXAMPLE: FZF THROUGH DOCKER CONTAINERS}\\
\color{white}function fzf-docker() \{\\
\color{white}docker ps -a --format '\{\{.ID\}\} \{\{.Names\}\} \{\{.Status\}\}' |\\
\color{white}fzf --preview 'docker inspect \{1\} | head -50' |\\
\color{white}awk '\{print \textbackslash{}\$1\}' |\\
\color{white}xargs -r docker logs --tail 100\\
\color{white}\}\\
\color{white}Generator: docker ps. Preview: docker inspect. Consumer: docker logs.\\
\vspace{4pt}\textbf{\color{magenta}EXAMPLE: FZF THROUGH TMUX SESSIONS}\\
\color{white}function fzf-tmux-session() \{\\
\color{white}tmux list-sessions -F '\#\{session\_name\}: \#\{session\_windows\} windows' |\\
\color{white}fzf --preview 'tmux list-windows -t \{1\}' |\\
\color{white}cut -d: -f1 |\\
\color{white}xargs tmux switch-client -t\\
\color{white}\}\\
\color{white}Generator: tmux list-sessions. Consumer: tmux switch-client.\\
\vspace{4pt}\textbf{\color{magenta}THE ACCEPT vs EDIT PATTERN}\\
\color{white}ZPWR verbs come in two flavors:\\
\color{white}Accept verbs  Execute the selection immediately\\
\color{white}Example: verbsfzf selects a verb and runs it.\\
\color{white}Edit verbs    Insert the selection into \textbackslash{}UFFER for editing\\
\color{white}Example: filesearch inserts the path into your command line.\\
\color{white}The difference is in the consumer. Accept uses eval/exec.\\
\color{white}Edit uses BUFFER=\textbackslash{}\$selection + zle redisplay.\\
\vspace{4pt}\textbf{\color{magenta}MAKING IT A ZLE WIDGET}\\
\color{white}To bind your fzf verb to a key:\\
\color{white}function my-fzf-widget() \{\\
\color{white}BUFFER=\textbackslash{}\$(my-generator | fzf)\\
\color{white}CURSOR=\textbackslash{}\$\{\#BUFFER\}\\
\color{white}zle redisplay\\
\color{white}\}\\
\color{white}zle -N my-fzf-widget\\
\color{white}bindkey '\textasciicircum{}F\textasciicircum{}X' my-fzf-widget\\
\color{white}Now Ctrl-F Ctrl-X triggers your custom fzf workflow.\\
\vspace{4pt}\textbf{\color{magenta}TIPS FOR GOOD FZF EXPERIENCES}\\
\color{white}- Use --ansi if your generator outputs color codes\\
\color{white}- Use --delimiter and --nth to control what fzf matches on\\
\color{white}- Use -m for multi-select (toggle with Tab)\\
\color{white}- Use --header for context hints at the top\\
\color{white}- Preview window size: --preview-window=right:50\%:wrap\\
\color{white}- Test with: echo 'a\textbackslash{}nb\textbackslash{}nc' | fzf --preview 'echo \{\}'\\

\vspace{8pt}


\section{CHAPTER 30: TMUX DEEP DIVE}

\subsection{Page 146: Tmux Architecture in ZPWR}

\vspace{4pt}\textbf{\color{cyan}=== TMUX ARCHITECTURE IN ZPWR ===}\vspace{2pt}\\
\textit{\color{gray}// the terminal multiplexer. one screen, infinite shells. //}\\
\color{white}Tmux is the backbone of the ZPWR workflow. Multiple shells,\\
\color{white}persistent sessions, split panes, and code execution across\\
\color{white}panes -- all orchestrated from \textbackslash{}PWR\_TMUX.\\
\vspace{4pt}\textbf{\color{magenta}SESSION / WINDOW / PANE MODEL}\\
\color{white}Session    Named group of windows (e.g., 'dev', 'zpwr-study')\\
\color{white}Window     A tab inside a session (like browser tabs)\\
\color{white}Pane       A split inside a window (like vim splits)\\
\color{white}Typical ZPWR layout: 1 session, 2-4 windows, 2-4 panes each.\\
\color{white}Prefix key: Ctrl-B (default). Followed by a command key.\\
\vspace{4pt}\textbf{\color{magenta}SESSION PERSISTENCE}\\
\color{white}tmux-resurrect   Manual save/restore of sessions\\
\color{white}Save: prefix + Ctrl-s\\
\color{white}Restore: prefix + Ctrl-r\\
\color{white}Saves: pane layouts, working directories, running programs.\\
\color{white}@resurrect-processes ':all:' -- saves ALL running programs.\\
\color{white}tmux-continuum    Auto-save every 15 minutes\\
\color{white}Works silently in the background. No manual intervention.\\
\color{white}Combined with resurrect: your tmux state survives reboots.\\
\vspace{4pt}\textbf{\color{magenta}THE POWERLINE STATUS BAR}\\
\color{white}Tmux sources powerline bindings for the status bar:\\
\color{white}source-file \textasciitilde{}/.tmux/powerline/bindings/tmux/powerline.conf\\
\color{white}The symlink \textasciitilde{}/.tmux/powerline points to pip's powerline-status.\\
\color{white}powerline-daemon runs in the background for performance.\\
\color{white}Status bar shows: session name, window list, hostname, time,\\
\color{white}battery level (via tmux-battery plugin).\\
\vspace{4pt}\textbf{\color{magenta}\textbackslash{}PWR\_TMUX: THE CONFIG DIRECTORY}\\
\color{white}All tmux config lives under \textbackslash{}PWR\_TMUX:\\
\color{white}init.conf               Main config (bindings, options)\\
\color{white}tmux-mac.conf           macOS-specific overrides\\
\color{white}tmux-linux.conf         Linux-specific overrides\\
\color{white}tmux\_ge\_21.conf         Tmux >= 2.1 features\\
\color{white}tmux\_ge\_29.conf         Tmux >= 2.9 features\\
\color{white}tmux\_ge\_31.conf         Tmux >= 3.1 features\\
\color{white}four-panes.conf         Layout: 4 panes\\
\color{white}sixteen-panes.conf      Layout: 16 panes\\
\color{white}config-files.conf       Opens brc, vrc, trc, zrc in 4 panes\\
\color{white}learn.conf              Study layout\\
\vspace{4pt}\textbf{\color{magenta}init.conf HIGHLIGHTS}\\
\color{white}- 30,000 line scrollback (ZPWR\_TMUX\_HISTORY\_LIMIT)\\
\color{white}- Vi keybindings in copy mode\\
\color{white}- Prefix S: synchronize-panes (type in all panes at once)\\
\color{white}- Prefix q: display pane numbers for 5 seconds\\
\color{white}- 3-second repeat-time for directional keys\\

\vspace{8pt}

\subsection{Page 147: Vim-Tmux Code Execution}

\vspace{4pt}\textbf{\color{cyan}=== VIM-TMUX CODE EXECUTION ===}\vspace{2pt}\\
\textit{\color{gray}// select code in vim. execute in the pane next door. magic. //}\\
\color{white}The killer feature of the ZPWR vim+tmux integration:\\
\color{white}write code in vim, send it to an adjacent tmux pane for\\
\color{white}execution. No copy-paste. No switching windows.\\
\vspace{4pt}\textbf{\color{magenta}THE TmuxRepeat FUNCTION}\\
\color{white}Defined in .vimrc, TmuxRepeat has three modes:\\
\color{white}'file'    Save the file, execute it in the right pane\\
\color{white}'visual'  Save visual selection to temp file, execute it\\
\color{white}'repl'    Paste visual selection directly into right pane\\
\vspace{4pt}\textbf{\color{magenta}SUPPORTED LANGUAGES}\\
\color{white}TmuxRepeat detects file type by extension:\\
\color{white}sh, zsh, py, rb, pl, clj, tcl, vim, lisp, hs, ml,\\
\color{white}coffee, swift, lua, java, f90, cr\\
\color{white}For known types, it sends the file to zpwrRunner.sh which\\
\color{white}dispatches to the correct interpreter. For unknown types,\\
\color{white}it falls back to TmuxRepeatGeneric (re-run previous command).\\
\vspace{4pt}\textbf{\color{magenta}HOW IT WORKS}\\
\color{white}1. Ctrl-V in normal/insert mode (or leader-vj):\\
\color{white}Saves the file, then sends:\\
\texttt{\color{green}tmux send-keys -t right 'bash \textbackslash{}"\textbackslash{}PWR\_SCRIPTS/zpwrRunner.sh\textbackslash{}" \textbackslash{}"file\textbackslash{}"' C-m}\\
\color{white}2. leader-vk in visual mode:\\
\color{white}Yanks selection to temp file \textbackslash{}PWR\_HIDDEN\_DIR\_TEMP/.vimTempFile.<ext>,\\
\color{white}then sends that temp file to the right pane for execution.\\
\color{white}3. leader-vl in visual mode (REPL mode):\\
\color{white}Yanks selection, loads it into tmux buffer 'buffer0099',\\
\color{white}then pastes directly into the right pane. For REPLs like\\
\color{white}python, irb, node -- code runs line-by-line in the REPL.\\
\vspace{4pt}\textbf{\color{magenta}THE PANE TARGETING}\\
\color{white}Standard tmux: sends to pane named 'right' (the next pane).\\
\color{white}MacVim/GUI: sends to session 'vimmers:0.' (external tmux).\\
\color{white}The pane\_count check ensures a target pane exists.\\
\vspace{4pt}\textbf{\color{magenta}allPanes.zsh: BROADCAST TO ALL PANES}\\
\color{white}\textbackslash{}PWR\_SCRIPTS/allPanes.zsh sends commands to all tmux panes:\\
\color{white}prefix + Space   Run command in all panes (single mode)\\
\color{white}prefix + v       Run command in all panes (multi mode)\\
\color{white}prefix + b       Open URL in all panes\\
\vspace{4pt}\textbf{\color{magenta}SYNCHRONIZE-PANES}\\
\color{white}prefix + S   Toggle synchronize-panes\\
\color{white}When enabled, every keystroke goes to ALL panes in the window.\\
\color{white}Useful for running identical commands on multiple servers,\\
\color{white}or setting up parallel environments simultaneously.\\
\vspace{4pt}\textbf{\color{magenta}GOOGLE FROM TMUX}\\
\color{white}prefix + g   Google search from tmux\\
\color{white}Runs \textbackslash{}PWR\_TMUX/google.sh with the current pane content.\\
\color{white}Select text, press prefix+g, and your browser opens\\
\color{white}with a Google search for the selection.\\

\vspace{8pt}

\subsection{Page 148: Tmux Plugins \& Extensions}

\vspace{4pt}\textbf{\color{cyan}=== TMUX PLUGINS \& EXTENSIONS ===}\vspace{2pt}\\
\textit{\color{gray}// tpm: the plugin manager for your multiplexer's plugins. yo dawg. //}\\
\color{white}ZPWR manages 9 tmux plugins via tpm (Tmux Plugin Manager).\\
\color{white}They live in \textasciitilde{}/.tmux/plugins/ and add features that tmux\\
\color{white}does not provide out of the box.\\
\vspace{4pt}\textbf{\color{magenta}TPM: TMUX PLUGIN MANAGER}\\
\color{white}Installed to \textasciitilde{}/.tmux/plugins/tpm. Managed via prefix keys:\\
\color{white}prefix + I       Install plugins listed in .tmux.conf\\
\color{white}prefix + U       Update all installed plugins\\
\color{white}prefix + Alt-u   Clean plugins not in config\\
\color{white}Plugins are declared in \textbackslash{}PWR\_INSTALL/.tmux.conf:\\
\color{white}set-option -g @plugin 'tmux-plugins/tpm'\\
\color{white}set-option -g @plugin 'tmux-plugins/tmux-resurrect'\\
\vspace{4pt}\textbf{\color{magenta}THE PLUGIN ROSTER (.tmuxplugins)}\\
\color{white}tmux-plugins/tpm                Plugin manager itself\\
\color{white}tmux-plugins/tmux-sensible      Sane defaults for tmux\\
\color{white}tmux-plugins/tmux-resurrect     Session save/restore\\
\color{white}tmux-plugins/tmux-continuum      Auto-save sessions\\
\color{white}tmux-plugins/tmux-battery        Battery indicator\\
\color{white}tmux-plugins/tmux-prefix-highlight Highlight when prefix active\\
\color{white}tmux-plugins/tmux-yank           System clipboard integration\\
\color{white}MenkeTechnologies/tmux-thumbs     Quick copy with hints\\
\color{white}MenkeTechnologies/tmux-fzf-url    FZF through URLs in pane\\
\vspace{4pt}\textbf{\color{magenta}tmux-resurrect: WHAT IT SAVES}\\
\color{white}- Pane layouts (splits, sizes, positions)\\
\color{white}- Working directories for each pane\\
\color{white}- Running programs (@resurrect-processes ':all:')\\
\color{white}- Command history is NOT saved (that's zsh's job)\\
\color{white}Save: prefix+Ctrl-s. Restore: prefix+Ctrl-r.\\
\vspace{4pt}\textbf{\color{magenta}tmux-continuum: AUTO-SAVE}\\
\color{white}Piggybacks on resurrect. Auto-saves every 15 minutes.\\
\color{white}@continuum-restore is commented out in ZPWR by default.\\
\color{white}Enable it to auto-restore on tmux server start.\\
\vspace{4pt}\textbf{\color{magenta}tmux-fzf-url: URL PICKER}\\
\color{white}Scans the visible pane for URLs, opens them in fzf:\\
\color{white}prefix + e   FZF through URLs, open selected in browser\\
\color{white}prefix + u   FZF through URLs, search mode\\
\color{white}Uses ZPWR\_TMUX\_HISTORY\_LIMIT for scrollback depth.\\
\vspace{4pt}\textbf{\color{magenta}CUSTOM SCRIPTS IN \textbackslash{}PWR\_TMUX}\\
\color{white}Layout scripts create multi-pane environments:\\
\color{white}prefix + D   control-window.conf (IDE layout)\\
\color{white}prefix + F   four-panes.conf\\
\color{white}prefix + G   eight-panes.conf\\
\color{white}prefix + O   sixteen-panes.conf\\
\color{white}prefix + T   config-files.conf (brc, vrc, trc, zrc)\\
\color{white}prefix + M   learn.conf (study layout)\\
\vspace{4pt}\textbf{\color{magenta}zpwr study: THE SPLIT-PANE LEARNING ENVIRONMENT}\\
\color{white}Creates a dedicated tmux session 'zpwr-study' with the\\
\color{white}encyclopedia on the left and a live shell on the right.\\
\color{white}Read the docs, try commands. Knowledge through practice.\\

\vspace{8pt}


\section{CHAPTER 31: VIM DEEP DIVE}

\subsection{Page 149: Vim Plugin Architecture}

\vspace{4pt}\textbf{\color{cyan}=== VIM PLUGIN ARCHITECTURE ===}\vspace{2pt}\\
\textit{\color{gray}// 15 plugins in .vimrc. each one a force multiplier. //}\\
\color{white}ZPWR's vim config uses Vundle as the plugin manager and\\
\color{white}ships 15 curated plugins that transform vim into a\\
\color{white}full IDE with completion, snippets, and tmux integration.\\
\vspace{4pt}\textbf{\color{magenta}VUNDLE: THE PLUGIN MANAGER}\\
\color{white}Plugins live in \textasciitilde{}/.vim/bundle/. Managed by Vundle:\\
\color{white}set runtimepath+=\textasciitilde{}/.vim/bundle/Vundle.vim\\
\color{white}call vundle\#begin()\\
\color{white}Plugin 'VundleVim/Vundle.vim'\\
\color{white}... more plugins ...\\
\color{white}call vundle\#end()\\
\color{white}Install: :PluginInstall  Update: :PluginUpdate  Clean: :PluginClean\\
\vspace{4pt}\textbf{\color{magenta}YouCompleteMe: SEMANTIC COMPLETION}\\
\color{white}Plugin 'ycm-core/YouCompleteMe'\\
\color{white}The heavyweight champion. Provides semantic completion for:\\
\color{white}C/C++ (via clangd), Python (via jedi), and more.\\
\color{white}Config: completion after 1 char, enabled in all file types,\\
\color{white}completes in comments. Uses C-n/C-p for selection.\\
\vspace{4pt}\textbf{\color{magenta}UltiSnips + vim-snippets: SNIPPET ENGINE}\\
\color{white}Plugin 'SirVer/ultisnips'\\
\color{white}Plugin 'honza/vim-snippets'\\
\color{white}Tab expands snippets, Tab/Shift-Tab navigate tab stops.\\
\color{white}Hundreds of snippets for every language: type 'for<Tab>'\\
\color{white}in Python and get a complete for loop with cursor placement.\\
\vspace{4pt}\textbf{\color{magenta}SuperTab: TAB COMPLETION GLUE}\\
\color{white}Plugin 'ervandew/supertab'\\
\color{white}Unifies Tab behavior: C-n cycles forward through completions.\\
\color{white}Coordinates with YCM and UltiSnips so Tab does the right\\
\color{white}thing in every context (snippet vs completion vs indent).\\
\vspace{4pt}\textbf{\color{magenta}vim-auto-save: NEVER :w AGAIN}\\
\color{white}Plugin '907th/vim-auto-save'\\
\color{white}Saves on InsertLeave, TextChanged, and TextChangedI.\\
\color{white}Every edit auto-saves silently. No more 'Unsaved changes' panics.\\
\color{white}g:auto\_save\_silent=1 suppresses the notification.\\
\vspace{4pt}\textbf{\color{magenta}rainbow: MATCHING BRACKET COLORS}\\
\color{white}Plugin 'luochen1990/rainbow'\\
\color{white}Nested brackets get different colors: parentheses in blue,\\
\color{white}inner brackets in orange, deeper in cyan, deepest in magenta.\\
\color{white}g:rainbow\_active=1 enables it globally on startup.\\
\vspace{4pt}\textbf{\color{magenta}vim-multiple-cursors: MULTI-CURSOR EDITING}\\
\color{white}Plugin 'TerryMa/vim-multiple-cursors'\\
\color{white}Sublime Text-style multi-cursor. Ctrl-N selects next match,\\
\color{white}then edit all occurrences simultaneously.\\
\vspace{4pt}\textbf{\color{magenta}tmux-complete.vim: CROSS-PANE COMPLETION}\\
\color{white}Plugin 'wellle/tmux-complete.vim'\\
\color{white}Completes words from ALL tmux panes. Writing a command in vim\\
\color{white}that references output in another pane? Tab-complete it.\\
\color{white}Uses async.vim and asyncomplete.vim for non-blocking operation.\\

\vspace{8pt}

\subsection{Page 150: Vim Sessions \& Workflow}

\vspace{4pt}\textbf{\color{cyan}=== VIM SESSIONS \& WORKFLOW ===}\vspace{2pt}\\
\textit{\color{gray}// your vim state, frozen in amber. thaw it when you need it. //}\\
\color{white}Vim sessions save your entire editing state: open buffers,\\
\color{white}window layouts, cursor positions. ZPWR integrates sessions\\
\color{white}with tmux for a seamless development environment.\\
\vspace{4pt}\textbf{\color{magenta}VIM SESSION MANAGEMENT}\\
\color{white}ZPWR configures vim sessions as manual-save only:\\
\color{white}let g:session\_autosave = 'no'\\
\color{white}let g:session\_autoload = 'no'\\
\color{white}No surprises. Sessions save when you tell them to.\\
\color{white}Save: Ctrl-D q (mapped to :SaveSession! in normal/insert mode).\\
\vspace{4pt}\textbf{\color{magenta}CONFIG SESSIONS: vrc, brc, zrc, trc}\\
\color{white}ZPWR defines verbs for editing its own config files:\\
\color{white}vrc   Open .vimrc in vim\\
\color{white}brc   Open shell aliases file (.zpwr\_alias.sh)\\
\color{white}zrc   Open .zshrc in vim\\
\color{white}trc   Open tmux.conf in vim\\
\color{white}The tmux layout script config-files.conf opens all four\\
\color{white}in a 4-pane arrangement. Prefix+T triggers it.\\
\color{white}Each pane runs one of these verbs automatically.\\
\vspace{4pt}\textbf{\color{magenta}THE AUXILIARY VIM PLUGINS}\\
\color{white}vim-autotag           Auto-regenerate ctags on save\\
\color{white}Plugin 'craigemery/vim-autotag'\\
\color{white}g:autotagStartMethod='fork' -- non-blocking tag generation.\\
\color{white}vim-online-thesaurus   Look up synonyms (visual mode \textbackslash{}\textbackslash{}K)\\
\color{white}Plugin 'beloglazov/vim-online-thesaurus'\\
\color{white}vim-zsh-completion     Zsh syntax completion in vim\\
\color{white}Plugin 'Valodim/vim-zsh-completion'\\
\color{white}vim-minimap            Sublime-style code minimap\\
\color{white}Plugin 'severin-lemaignan/vim-minimap'\\
\vspace{4pt}\textbf{\color{magenta}THE zpwr snapshot/restore INTEGRATION}\\
\color{white}ZPWR has verbs for full environment snapshots:\\
\texttt{\color{green}zpwr snapshot    Capture terminal environment state}\\
\texttt{\color{green}zpwr restore     Restore from snapshot}\\
\color{white}zpwrRestore rebuilds shell state: directory stack, aliases,\\
\color{white}parameters, and more. This complements tmux-resurrect:\\
\color{white}resurrect restores the tmux layout, snapshot/restore\\
\color{white}restores the shell state within each pane.\\
\vspace{4pt}\textbf{\color{magenta}THE FULL ZPWR VIM WORKFLOW}\\
\color{white}1. Open a tmux session with your project layout\\
\color{white}2. Vim in left pane, shell in right pane\\
\color{white}3. Edit code in vim, Ctrl-V to execute in right pane\\
\color{white}4. Visual select + leader-vl to paste into a REPL\\
\color{white}5. tmux-complete.vim grabs words from the REPL output\\
\color{white}6. Auto-save means you never lose work\\
\color{white}7. Prefix+Ctrl-s saves the entire tmux state\\
\color{white}This is not an IDE. This is a weapon system.\\

\vspace{8pt}


\section{CHAPTER 32: TEMPRS — TEMPORARY FILE STACK MANAGER}

\subsection{Page 151: temprs: The Stack Machine}

\vspace{4pt}\textbf{\color{cyan}=== TEMPRS: THE STACK MACHINE ===}\vspace{2pt}\\
\textit{\color{gray}// your filesystem is a stack. you just didn't know it yet. //}\\
\color{white}temprs is a Rust-based temporary file stack manager by\\
\color{white}MenkeTechnologies. Think of it as a LIFO/FIFO clipboard that\\
\color{white}persists across shell sessions using actual tempfiles on disk.\\
\color{white}Version 2.9.3. Every operation is instant. No databases, no\\
\color{white}daemons — just a master record file and a directory of tempfiles.\\
\vspace{4pt}\textbf{\color{magenta}THE STACK PARADIGM}\\
\color{white}temprs models your scratch data as a stack with both ends open:\\
\color{white}push (default), pop, shift, unshift — like a Perl array on disk.\\
\vspace{4pt}\textbf{\color{magenta}CORE OPERATIONS}\\
\color{white}temprs              \# no args: create new tempfile, push to top\\
\color{white}temprs -o INDEX     \# read contents of tempfile at INDEX to stdout\\
\color{white}temprs -i INDEX     \# write stdin into tempfile at INDEX\\
\color{white}temprs -p           \# pop from top of stack (print + remove)\\
\color{white}temprs -s           \# shift from bottom of stack (print + remove)\\
\color{white}temprs -u           \# unshift: push stdin to bottom (silent)\\
\color{white}temprs -a INDEX     \# insert new tempfile at INDEX\\
\color{white}temprs -r INDEX     \# remove tempfile at INDEX\\
\vspace{4pt}\textbf{\color{magenta}THE MASTER RECORD}\\
\color{white}temprs keeps a master record file that maps stack positions to\\
\color{white}actual tempfile paths on disk. The temprs directory holds the\\
\color{white}files themselves. Both are in your system's temp directory.\\
\color{white}temprs -d           \# list the temprs directory contents\\
\color{white}temprs -m           \# list the master record file path\\
\vspace{4pt}\textbf{\color{magenta}PIPELINE INTEGRATION}\\
\color{white}temprs reads stdin when available — perfect for pipelines:\\
\color{white}echo 'stash this' | temprs -u  \# silent push to bottom\\
\color{white}curl -s api.io/data | temprs    \# push response to top\\
\color{white}temprs -o 0 | jq .             \# read top, pipe to jq\\
\color{white}The -u unshift flag reads stdin and produces no stdout —\\
\color{white}ideal for mid-pipeline stashing without breaking the flow.\\
\vspace{4pt}\textbf{\color{magenta}POSITIONAL ARGUMENT}\\
\color{white}temprs myfile.txt   \# read file into new tempfile on stack\\
\color{white}If both stdin and a file arg are present, stdin takes priority.\\
\color{white}Your data flows in, gets a stack address, and waits for recall.\\

\vspace{8pt}

\subsection{Page 152: temprs: Advanced Operations}

\vspace{4pt}\textbf{\color{cyan}=== TEMPRS: ADVANCED OPERATIONS ===}\vspace{2pt}\\
\textit{\color{gray}// the stack goes deeper. every tempfile has a story. //}\\
\vspace{4pt}\textbf{\color{magenta}LISTING \& INSPECTION}\\
\color{white}temprs -l           \# list all tempfile paths on the stack\\
\color{white}temprs -n           \# list numbered (index + path)\\
\color{white}temprs -L           \# list all tempfile contents\\
\color{white}temprs -N           \# list numbered with full contents\\
\color{white}temprs -k           \# count tempfiles on the stack\\
\color{white}temprs -I INDEX     \# show metadata for tempfile at INDEX\\
\vspace{4pt}\textbf{\color{magenta}EDITING \& TAGGING}\\
\color{white}temprs -e INDEX     \# open tempfile at INDEX in \textbackslash{}DITOR\\
\color{white}temprs -w NAME      \# tag a tempfile with an alias name\\
\color{white}temprs -R OLD NEW   \# rename tag from OLD to NEW\\
\color{white}Tags let you retrieve tempfiles by name instead of index.\\
\color{white}Your stack positions shift; your tags don't.\\
\vspace{4pt}\textbf{\color{magenta}SEARCH \& CONCATENATE}\\
\color{white}temprs -g PATTERN   \# grep through all tempfile contents\\
\color{white}temprs -C 0 2 5    \# concatenate tempfiles at indices to stdout\\
\color{white}The grep flag searches across every tempfile on the stack —\\
\color{white}like a personal knowledge base with instant full-text search.\\
\vspace{4pt}\textbf{\color{magenta}STACK MANIPULATION}\\
\color{white}temprs -D 0 1      \# diff two tempfiles by index or name\\
\color{white}temprs -M 0 3      \# move tempfile from position 0 to 3\\
\color{white}temprs -x INDEX     \# duplicate tempfile onto top of stack\\
\color{white}temprs -S 0 2      \# swap two tempfiles by index or name\\
\color{white}temprs --rev        \# reverse the entire stack order\\
\color{white}temprs -A INDEX     \# append stdin to existing tempfile\\
\vspace{4pt}\textbf{\color{magenta}HOUSEKEEPING}\\
\color{white}temprs -c           \# clear/purge ALL tempfiles from stack\\
\color{white}temprs --expire 24  \# purge tempfiles older than 24 hours\\
\color{white}temprs --wc INDEX   \# print line count of tempfile\\
\color{white}temprs --size INDEX \# print byte size of tempfile\\
\color{white}temprs --path INDEX \# print the actual filesystem path\\
\color{white}temprs --sort mtime \# sort stack by name, size, or mtime\\
\vspace{4pt}\textbf{\color{magenta}MODE FLAGS}\\
\color{white}temprs -q           \# quiet mode: suppress output on create\\
\color{white}temprs -v           \# verbose mode: extra diagnostic output\\
\vspace{4pt}\textbf{\color{magenta}THE CYBERPUNK -h}\\
\color{white}Run temprs -h for the full ASCII art banner with the\\
\color{white}TEMPRS logo, status bar, and signal meter. Because even\\
\color{white}help text should look like a console from a Gibson novel.\\
\color{white}Also: --head INDEX N and --tail INDEX N for partial reads,\\
\color{white}and --replace INDEX PAT REPL for in-place sed-style edits.\\

\vspace{8pt}


\section{CHAPTER 33: LSOFRS — ENHANCED LSOF}

\subsection{Page 153: lsofrs: Network Intelligence}

\vspace{4pt}\textbf{\color{cyan}=== LSOFRS: NETWORK INTELLIGENCE ===}\vspace{2pt}\\
\textit{\color{gray}// every open file descriptor is a signal. lsofrs reads them all. //}\\
\color{white}lsofrs is a Rust-based enhanced lsof wrapper v4.7.0 by\\
\color{white}MenkeTechnologies. It wraps the venerable lsof with colorized\\
\color{white}output, structured formatting, and cyberpunk-grade aesthetics.\\
\vspace{4pt}\textbf{\color{magenta}WHY NOT RAW LSOF?}\\
\color{white}Raw lsof dumps walls of monochrome text. lsofrs adds color\\
\color{white}coding by type, structured columns, and filtering that doesn't\\
\color{white}require piping through 5 layers of awk and grep.\\
\vspace{4pt}\textbf{\color{magenta}SELECTION FLAGS}\\
\color{white}lsofrs -i [ADDR]   \# select internet connections\\
\color{white}ADDR format: [4|6][proto][@host|addr][:svc|port]\\
\color{white}lsofrs -c COMMAND   \# select by command name (prefix match)\\
\color{white}lsofrs -p PID       \# select by PID (comma-separated, \textasciicircum{}excludes)\\
\color{white}lsofrs -u USER      \# select by login name or UID\\
\color{white}lsofrs -d FD        \# select by file descriptor set\\
\color{white}lsofrs -g [PGID]    \# select by process group ID\\
\color{white}lsofrs -a           \# AND selections (default is OR)\\
\vspace{4pt}\textbf{\color{magenta}NETWORK FLAGS}\\
\color{white}lsofrs -n           \# inhibit hostname resolution (fast mode)\\
\color{white}lsofrs -P           \# inhibit port-to-name conversion\\
\color{white}lsofrs -N           \# select NFS files\\
\color{white}lsofrs -U           \# select UNIX domain socket files\\
\vspace{4pt}\textbf{\color{magenta}ZPWR INTEGRATION}\\
\color{white}The zpwr lsof verb pipes lsof output through fzf for\\
\color{white}interactive process selection and killing:\\
\texttt{\color{green}zpwr lsof           \# fzf over lsof, select PIDs to kill}\\
\texttt{\color{green}zpwr lsofedit       \# same but edit the kill command first}\\
\color{white}The zpwrLsoffzf widget (bound to \textasciicircum{}F\textasciicircum{}H in viins and vicmd)\\
\color{white}runs sudo lsof -i | fzf -m, extracts PIDs with awk,\\
\color{white}and inserts them into your command line. Scan the local\\
\color{white}network like a deck jockey, pick your targets with fzf.\\
\vspace{4pt}\textbf{\color{magenta}FILES \& DIRECTORIES}\\
\color{white}lsofrs FILE...      \# list processes using these files\\
\color{white}lsofrs --dir DIR    \# open files in DIR (one level, like +d)\\
\color{white}lsofrs --dir-recurse DIR \# recursive (like +D)\\

\vspace{8pt}

\subsection{Page 154: lsofrs: Practical Usage}

\vspace{4pt}\textbf{\color{cyan}=== LSOFRS: PRACTICAL USAGE ===}\vspace{2pt}\\
\textit{\color{gray}// who's listening on port 8080? lsofrs knows. lsofrs always knows. //}\\
\vspace{4pt}\textbf{\color{magenta}FINDING WHAT'S ON A PORT}\\
\color{white}lsofrs -i :8080        \# anything on port 8080\\
\color{white}lsofrs -i TCP:443      \# TCP connections on 443\\
\color{white}lsofrs -i 4            \# all IPv4 connections\\
\color{white}lsofrs -i 6TCP         \# IPv6 TCP only\\
\vspace{4pt}\textbf{\color{magenta}DISPLAY \& OUTPUT}\\
\color{white}lsofrs -F [FIELDS]  \# machine-readable output fields\\
\color{white}lsofrs -J / --json  \# JSON output for piping to jq\\
\color{white}lsofrs -R           \# list parent PIDs\\
\color{white}lsofrs -t           \# terse: PID only (for scripting)\\
\color{white}lsofrs -0           \# NUL field terminator (xargs -0 friendly)\\
\color{white}lsofrs --color MODE \# auto, always, or never\\
\vspace{4pt}\textbf{\color{magenta}LIVE MONITORING}\\
\color{white}lsofrs +r 2         \# repeat every 2 seconds\\
\color{white}lsofrs -W / --monitor \# live full-screen like top\\
\color{white}lsofrs --delta       \# highlight new/gone FDs in repeat mode\\
\color{white}lsofrs --follow PID  \# watch one process, highlight opens/closes\\
\color{white}lsofrs --top 20     \# live top-N processes by FD count\\
\color{white}lsofrs --watch FILE  \# watch who opens/closes a file over time\\
\vspace{4pt}\textbf{\color{magenta}DIAGNOSTICS}\\
\color{white}lsofrs --leak-detect \# detect FD leaks (poll every 5s, 3 rounds)\\
\color{white}lsofrs --stale       \# find FDs pointing to deleted files\\
\color{white}lsofrs --ports       \# listening ports summary (like ss -tlnp)\\
\color{white}lsofrs --tree        \# process tree with FD counts (pstree + lsof)\\
\color{white}lsofrs --summary     \# aggregate FD summary by type and user\\
\vspace{4pt}\textbf{\color{magenta}COMPARISON}\\
\color{white}lsof    — the OG. every Unix has it. output is a wall of text.\\
\color{white}ss      — socket statistics. fast but sockets only, no files.\\
\color{white}netstat — deprecated on Linux. still works on macOS/BSD.\\
\color{white}lsofrs  — wraps lsof with color, JSON, live monitoring,\\
\color{white}leak detection, tree views, and the cyberpunk banner\\
\color{white}you get with lsofrs -h. Built for operators who\\
\color{white}stare at FD tables and like what they see.\\
\vspace{4pt}\textbf{\color{magenta}QUICK RECIPES}\\
\color{white}lsofrs -t -i :3000 | xargs kill  \# kill whatever's on 3000\\
\color{white}lsofrs -c node -a -i TCP          \# node's TCP connections\\
\color{white}lsofrs -u \textbackslash{}\$(whoami) -a -i         \# your network connections\\

\vspace{8pt}


\section{CHAPTER 34: EZA — THE MODERN LS}

\subsection{Page 155: eza: ls Evolved}

\vspace{4pt}\textbf{\color{cyan}=== EZA: LS EVOLVED ===}\vspace{2pt}\\
\textit{\color{gray}// ls is dead. long live ls. //}\\
\color{white}eza is the modern Rust replacement for ls, formerly known as\\
\color{white}exa. Maintained by eza-community after exa's author went AWOL.\\
\color{white}In zpwr, it's the default file lister for everything.\\
\vspace{4pt}\textbf{\color{magenta}\textbackslash{}PWR\_EXA\_COMMAND — THE CONFIGURED INVOCATION}\\
\color{white}command eza --git -il --classify=always -H --color-scale -g -a --colour=always\\
\color{white}Every flag earns its place:\\
\color{white}--git              \# git status per file (M modified, N new, I ignored)\\
\color{white}-i                 \# show inode numbers\\
\color{white}-l                 \# long listing format\\
\color{white}--classify=always  \# / after dirs, * after executables, always\\
\color{white}-H                 \# show hard link count\\
\color{white}--color-scale      \# gradient colors for file sizes and ages\\
\color{white}-g                 \# show group ownership\\
\color{white}-a                 \# show hidden files (dotfiles)\\
\color{white}--colour=always    \# force color even when piped\\
\vspace{4pt}\textbf{\color{magenta}WHY --classify=always?}\\
\color{white}Not just -F. Under p10k instant prompt, stdout gets redirected\\
\color{white}during init. eza's --classify auto-detects terminal capabilities\\
\color{white}and when stdout is redirected, it drops the indicators. The\\
\color{white}=always suffix forces classification regardless of fd state.\\
\color{white}Without it, your login banner loses all the / and * markers.\\
\vspace{4pt}\textbf{\color{magenta}WHERE IT'S USED}\\
\color{white}zpwrClearList and zpwrListNoClear both eval \textbackslash{}PWR\_EXA\_COMMAND.\\
\color{white}The fallback chain: eza -> grc + gls -> raw ls.\\
\color{white}zpwrListNoClear checks for eza first, falls back to gls on\\
\color{white}macOS (with grc for colorization) or plain ls on Linux.\\
\texttt{\color{green}alias eza=\textbackslash{}"\textbackslash{}PWR\_EXA\_COMMAND\textbackslash{}"  \# bound in zpwrBindAliasesLate}\\
\color{white}alias lr=\textbackslash{}"\textbackslash{}PWR\_EXA\_COMMAND -R --tree | less -rMN\textbackslash{}"\\
\color{white}The lr alias gives you a recursive tree piped through less\\
\color{white}with line numbers — a full filesystem X-ray.\\
\vspace{4pt}\textbf{\color{magenta}EXTENDED ATTRIBUTES}\\
\texttt{\color{green}Set ZPWR\_EXA\_EXTENDED=true in .zpwr\_env.sh to add --extended,}\\
\color{white}which shows xattr metadata on macOS (com.apple.quarantine etc).\\

\vspace{8pt}

\subsection{Page 156: eza: Colors \& Integration}

\vspace{4pt}\textbf{\color{cyan}=== EZA: COLORS \& INTEGRATION ===}\vspace{2pt}\\
\textit{\color{gray}// a file listing without color is just noise. //}\\
\vspace{4pt}\textbf{\color{magenta}EZA\_COLORS ENVIRONMENT VARIABLE}\\
\color{white}eza uses EZA\_COLORS (or EXA\_COLORS for compat) to customize\\
\color{white}colors beyond the defaults. Format is colon-separated key=value:\\
\color{white}export EZA\_COLORS=\textbackslash{}"di=34:ln=35:so=32:pi=33:ex=31\textbackslash{}"\\
\color{white}Keys match LS\_COLORS extensions. eza also adds its own:\\
\color{white}da=date, sn=file size number, sb=file size unit,\\
\color{white}uu=current user, un=other user, gu=current group.\\
\vspace{4pt}\textbf{\color{magenta}LS\_COLORS COMPATIBILITY}\\
\color{white}eza reads LS\_COLORS as a fallback. If you have a vivid or\\
\color{white}dircolors setup, eza inherits it. EZA\_COLORS overrides.\\
\vspace{4pt}\textbf{\color{magenta}TREE VIEW}\\
\color{white}eza --tree          \# recursive tree (like tree command)\\
\color{white}eza --tree -L 3     \# limit depth to 3 levels\\
\color{white}eza --tree --git-ignore \# respect .gitignore in tree\\
\color{white}The tree integrates with --git, so you see git status markers\\
\color{white}inline with the tree structure. New files glow green.\\
\vspace{4pt}\textbf{\color{magenta}THE FALLBACK CHAIN IN zpwrListNoClear}\\
\color{white}1. If eza exists: eval \textbackslash{}PWR\_EXA\_COMMAND\\
\color{white}2. macOS + grc: grc -c \textasciitilde{}/conf.gls gls -iFlhA --color=always\\
\color{white}3. macOS bare: ls -iFlhAO (O = macOS file flags)\\
\color{white}4. Linux + grc: grc -c \textasciitilde{}/conf.gls ls -iFlhA --color=always\\
\color{white}5. Linux bare: ls -ifhlA\\
\color{white}grc (Generic Colouriser) wraps ls output with conf.gls rules.\\
\color{white}It's the middle ground between raw ls and eza.\\
\vspace{4pt}\textbf{\color{magenta}ZPWR VERBS USING EZA}\\
\color{white}zpwrClearList clears the terminal then runs the listing.\\
\color{white}It's called on every cd via the chpwd hook — so you always\\
\color{white}see what's in the directory you just landed in.\\
\color{white}zpwrListNoClear does the same listing without clearing.\\
\vspace{4pt}\textbf{\color{magenta}SORTING \& FILTERING}\\
\color{white}eza --sort=modified  \# sort by modification time\\
\color{white}eza --sort=size      \# sort by file size\\
\color{white}eza -r               \# reverse sort order\\
\color{white}eza --group-directories-first \# dirs at top\\
\color{white}eza -I '*.pyc|\_\_pycache\_\_'   \# ignore patterns\\
\color{white}Combine with --git-ignore to hide everything in .gitignore.\\
\color{white}Your file listings become as curated as your code.\\

\vspace{8pt}


\section{CHAPTER 35: BAT — CAT WITH WINGS}

\subsection{Page 157: bat: Syntax Highlighting for Everything}

\vspace{4pt}\textbf{\color{cyan}=== BAT: SYNTAX HIGHLIGHTING FOR EVERYTHING ===}\vspace{2pt}\\
\textit{\color{gray}// cat is for reading files. bat is for understanding them. //}\\
\color{white}bat is a cat clone with syntax highlighting, git integration,\\
\color{white}and automatic paging. Written in Rust. In zpwr it's the\\
\texttt{\color{green}default colorizer: ZPWR\_COLORIZER=bat in .zpwr\_env.sh.}\\
\vspace{4pt}\textbf{\color{magenta}CORE FEATURES}\\
\color{white}Line numbers in the gutter. Git diff markers (+/-/\textasciitilde{}) in the\\
\color{white}margin. Syntax highlighting for 200+ languages. Automatic\\
\color{white}paging via less when output exceeds terminal height.\\
\vspace{4pt}\textbf{\color{magenta}DISCOVERY}\\
\color{white}bat --list-themes    \# see all available color themes\\
\color{white}bat --list-languages \# every syntax bat can highlight\\
\vspace{4pt}\textbf{\color{magenta}HOW ZPWR USES BAT}\\
\color{white}bat is the preview engine for fzf throughout zpwr:\\
\color{white}fzf --preview 'bat --color=always \{\}'\\
\color{white}The full zpwr colorizer invocation:\\
\color{white}bat --paging never --wrap character --color always\\
\color{white}--style=\textbackslash{}"numbers,grid,changes,header\textbackslash{}"\\
\color{white}This is stored in \textbackslash{}PWR\_COLORIZER and expanded into:\\
\color{white}FZF\_COLORIZER          — preview for generic fzf\\
\color{white}FZF\_COLORIZER\_FILE     — preview for file-specific fzf\\
\color{white}FZF\_COLORIZER\_FILE\_TEXT — text file preview (uses -l ASP)\\
\color{white}When bat is unavailable, zpwr falls back to pygmentize.\\
\vspace{4pt}\textbf{\color{magenta}ZPWR VERBS}\\
\color{white}zpwr filesearch and zpwr wordsearch both use bat in their\\
\color{white}fzf preview windows. vimfilesearch, vimwordsearch, and the\\
\color{white}emacs variants (emacsfilesearch, emacswordsearch) all inherit\\
\color{white}the same bat-powered preview pipeline.\\
\vspace{4pt}\textbf{\color{magenta}STANDALONE USAGE}\\
\color{white}bat src/main.rs              \# highlighted source view\\
\color{white}bat -l json data.txt         \# force language detection\\
\color{white}bat --diff                    \# show only git-changed lines\\
\color{white}bat -A                        \# show non-printable chars\\
\color{white}bat --plain                   \# no line numbers, no border\\
\color{white}bat -r 10:20 file.zsh        \# show only lines 10-20\\
\color{white}The -r (--line-range) flag is perfect for surgical code review.\\

\vspace{8pt}

\subsection{Page 158: bat: Configuration \& Tricks}

\vspace{4pt}\textbf{\color{cyan}=== BAT: CONFIGURATION \& TRICKS ===}\vspace{2pt}\\
\textit{\color{gray}// configure your cat. teach it new tricks. make it purr. //}\\
\vspace{4pt}\textbf{\color{magenta}BAT\_THEME}\\
\color{white}export BAT\_THEME=\textbackslash{}"Dracula\textbackslash{}"\\
\color{white}Set this in your env to change the default theme globally.\\
\color{white}Popular themes: Dracula, Monokai Extended, Nord, OneHalfDark,\\
\color{white}Solarized (Dark), TwoDark, ansi. Preview them all:\\
\color{white}bat --list-themes | fzf --preview='bat --theme=\{\} src/main.rs'\\
\vspace{4pt}\textbf{\color{magenta}STYLE COMPONENTS}\\
\color{white}bat --style=numbers,changes,header\\
\color{white}Components: full, auto, plain, numbers, changes, header,\\
\color{white}header-filename, header-filesize, grid, rule, snip.\\
\color{white}Combine with commas. 'full' enables everything.\\
\color{white}'plain' disables everything — just colored text, no chrome.\\
\vspace{4pt}\textbf{\color{magenta}BAT AS MANPAGE COLORIZER}\\
\color{white}export MANPAGER=\textbackslash{}"sh -c 'col -bx | bat -l man -p'\textbackslash{}"\\
\color{white}This makes every man page syntax-highlighted. The -l man forces\\
\color{white}manpage language detection, -p is --plain (no line numbers\\
\color{white}cluttering man output). col -bx strips backspace formatting.\\
\vspace{4pt}\textbf{\color{magenta}SYNTAX CACHE}\\
\color{white}bat cache --build    \# rebuild syntax/theme cache\\
\color{white}bat cache --clear    \# clear and rebuild from scratch\\
\color{white}After adding custom .sublime-syntax files to\\
\color{white}\textasciitilde{}/.config/bat/syntaxes/, rebuild the cache to pick them up.\\
\vspace{4pt}\textbf{\color{magenta}DEBIAN/UBUNTU GOTCHA}\\
\color{white}On Debian-based systems, the binary is batcat not bat\\
\color{white}due to a name conflict with an unrelated package. Fix:\\
\color{white}alias bat=batcat     \# or symlink: ln -s batcat \textasciitilde{}/.local/bin/bat\\
\vspace{4pt}\textbf{\color{magenta}BAT + DELTA FOR GIT DIFFS}\\
\color{white}delta is a syntax-highlighting pager for git diffs. It uses\\
\color{white}bat's syntax themes internally. Configure in .gitconfig:\\
\color{white}[core] pager = delta\\
\color{white}[delta] syntax-theme = Dracula\\
\color{white}delta adds line numbers, side-by-side view, and word-level\\
\color{white}diff highlighting. bat does files, delta does diffs.\\
\vspace{4pt}\textbf{\color{magenta}BAT\_PAGER}\\
\color{white}export BAT\_PAGER=\textbackslash{}"less -RFX\textbackslash{}"\\
\color{white}Customize the pager bat uses. -R for color, -F to quit if one\\
\color{white}screen, -X to not clear screen on exit. Or set to '' to disable\\
\color{white}paging entirely — zpwr does this with --paging never.\\

\vspace{8pt}


\section{CHAPTER 36: FD-FIND — FIND REPLACEMENT}

\subsection{Page 159: fd: find for Humans}

\vspace{4pt}\textbf{\color{cyan}=== FD: FIND FOR HUMANS ===}\vspace{2pt}\\
\textit{\color{gray}// find is powerful. fd is powerful AND usable. //}\\
\color{white}fd is a fast, user-friendly alternative to find, written in\\
\color{white}Rust by David Peter (sharkdp — also the author of bat).\\
\color{white}It respects .gitignore by default and uses regex, not globs.\\
\vspace{4pt}\textbf{\color{magenta}BASIC USAGE}\\
\texttt{\color{green}fd PATTERN           \# search current dir recursively}\\
\texttt{\color{green}fd PATTERN path/     \# search specific directory}\\
\texttt{\color{green}fd '\textasciicircum{}test.*\textbackslash{}\textbackslash{}.zsh\textbackslash{}\$'  \# full regex support}\\
\color{white}No -name, no -type f, no quoting gymnastics. Just the pattern.\\
\vspace{4pt}\textbf{\color{magenta}KEY DIFFERENCES FROM FIND}\\
\color{white}1. Regex by default (find uses glob with -name)\\
\color{white}2. Ignores .git, node\_modules, .gitignore entries automatically\\
\color{white}3. 5-10x faster than find due to parallel directory traversal\\
\color{white}4. Smart case: case-insensitive unless pattern has uppercase\\
\color{white}5. Colorized output by default\\
\vspace{4pt}\textbf{\color{magenta}COMMON FLAGS}\\
\color{white}fd -e zsh           \# filter by extension\\
\color{white}fd -t f             \# files only\\
\color{white}fd -t d             \# directories only\\
\color{white}fd -t l             \# symlinks only\\
\color{white}fd -t x             \# executables only\\
\color{white}fd -H               \# include hidden files\\
\color{white}fd -I               \# don't respect .gitignore\\
\color{white}fd -u               \# unrestricted (-H -I combined)\\
\color{white}fd -g GLOB          \# use glob pattern instead of regex\\
\vspace{4pt}\textbf{\color{magenta}EXECUTION}\\
\color{white}fd -x cmd \{\}        \# execute command per result (parallel)\\
\color{white}fd -X cmd\{\}         \# execute once with all results as args\\
\textit{\color{gray}Placeholders: \{\} = path, \{/\} = basename, \{//\} = parent dir,}\\
\color{white}\{.\} = path without extension, \{/.\} = basename without extension.\\
\texttt{\color{green}fd -e jpg -x convert \{\} \{.\}.png  \# convert all jpgs to png}\\
\color{white}The -x flag runs in parallel by default — your cores earn their\\
\color{white}keep. -X collects all results and passes them as one batch.\\
\color{white}This is the find | xargs pattern without the pipe.\\

\vspace{8pt}

\subsection{Page 160: fd: Integration with zpwr \& fzf}

\vspace{4pt}\textbf{\color{cyan}=== FD: INTEGRATION WITH ZPWR \& FZF ===}\vspace{2pt}\\
\textit{\color{gray}// fd feeds fzf. fzf feeds your brain. the pipeline of knowing. //}\\
\vspace{4pt}\textbf{\color{magenta}FZF\_DEFAULT\_COMMAND}\\
\color{white}zpwr sets FZF\_DEFAULT\_COMMAND to use find + ag by default:\\
\color{white}export FZF\_DEFAULT\_COMMAND='find * | ag -v \textbackslash{}".git/\textbackslash{}"'\\
\color{white}You can override this with fd for faster results:\\
\texttt{\color{green}export FZF\_DEFAULT\_COMMAND='fd --type f --hidden --follow'}\\
\color{white}fd is faster because it parallelizes traversal and skips\\
\color{white}.gitignore entries without needing ag to filter them out.\\
\vspace{4pt}\textbf{\color{magenta}FD + FZF PATTERNS}\\
\texttt{\color{green}fd -e zsh | fzf                 \# fuzzy find zsh files}\\
\texttt{\color{green}fd -t f | fzf --preview 'bat \{\}'  \# find + preview}\\
\texttt{\color{green}vim \textbackslash{}\$(fd -t f | fzf)             \# find, select, edit}\\
\color{white}This three-tool combo (fd + fzf + bat) is the holy trinity\\
\color{white}of file navigation. fd finds, fzf filters, bat previews.\\
\vspace{4pt}\textbf{\color{magenta}TEMPORAL \& SIZE FILTERS}\\
\color{white}fd --changed-within 1d  \# modified in last 24 hours\\
\color{white}fd --changed-before 1w  \# modified more than a week ago\\
\color{white}fd --size +1m           \# files larger than 1 megabyte\\
\color{white}fd --size -10k          \# files smaller than 10 kilobytes\\
\color{white}Units: b (bytes), k (kilobytes), m (megabytes), g (gigabytes).\\
\color{white}Time: s (seconds), m (minutes), h (hours), d (days), w (weeks).\\
\vspace{4pt}\textbf{\color{magenta}FD + BAT}\\
\texttt{\color{green}fd -e rs --exec bat\{\}   \# colorized find results}\\
\texttt{\color{green}fd -e md -X bat          \# cat all markdown files at once}\\
\vspace{4pt}\textbf{\color{magenta}FD vs FIND vs LOCATE}\\
\color{white}find   — universal. every Unix has it. slow on large trees.\\
\color{white}supports -exec, -newer, -perm, -prune. the Swiss Army.\\
\color{white}fd     — fast parallel traversal, sane defaults, .gitignore\\
\color{white}aware. best for daily use in code repositories.\\
\color{white}locate — instant results from a pre-built database (updatedb).\\
\color{white}stale until the next cron run. no real-time accuracy.\\
\vspace{4pt}\textbf{\color{magenta}DEPTH \& EXCLUSION}\\
\color{white}fd -d 3              \# max depth of 3 levels\\
\color{white}fd -E node\_modules   \# exclude a directory by name\\
\color{white}fd -E '*.min.js'     \# exclude by glob pattern\\
\color{white}Stack -E flags to build precise exclusion lists on the fly.\\

\vspace{8pt}


\section{CHAPTER 37: RIPGREP — GREP AT LIGHT SPEED}

\subsection{Page 161: ripgrep: grep Reimagined}

\vspace{4pt}\textbf{\color{cyan}=== RIPGREP: GREP REIMAGINED ===}\vspace{2pt}\\
\textit{\color{gray}// grep -r is a campfire. ripgrep is a directed energy weapon. //}\\
\color{white}ripgrep (rg) is a blazing-fast recursive grep written in Rust\\
\color{white}by Andrew Gallant (BurntSushi). It respects .gitignore, skips\\
\color{white}binary files, and uses Rust's regex engine for raw speed.\\
\color{white}5-100x faster than grep -r on large codebases. Not a typo.\\
\vspace{4pt}\textbf{\color{magenta}BASIC USAGE}\\
\texttt{\color{green}rg PATTERN           \# search current dir recursively}\\
\texttt{\color{green}rg PATTERN path/     \# search specific directory}\\
\texttt{\color{green}rg 'fn main'         \# literal or regex — rg auto-detects}\\
\color{white}Output is colorized by default: filenames in magenta, line\\
\color{white}numbers in green, matches highlighted. Grouped by file.\\
\vspace{4pt}\textbf{\color{magenta}KEY FLAGS}\\
\color{white}rg -i PATTERN       \# case insensitive search\\
\color{white}rg -l PATTERN       \# files with matches only (paths)\\
\color{white}rg -c PATTERN       \# count matches per file\\
\color{white}rg -w PATTERN       \# whole word matches only\\
\color{white}rg -F PATTERN       \# fixed string (no regex interpretation)\\
\color{white}rg -v PATTERN       \# invert match (lines NOT matching)\\
\vspace{4pt}\textbf{\color{magenta}FILE TYPE FILTERING}\\
\color{white}rg -t zsh PATTERN   \# search only zsh files\\
\color{white}rg -T js PATTERN    \# exclude javascript files\\
\color{white}rg -g '*.zsh' PAT   \# glob filter (include)\\
\color{white}rg -g '!*.min.js' PAT \# glob filter (exclude)\\
\color{white}rg --type-list shows all built-in type definitions.\\
\vspace{4pt}\textbf{\color{magenta}CONTEXT LINES}\\
\color{white}rg -A 3 PATTERN     \# 3 lines after each match\\
\color{white}rg -B 3 PATTERN     \# 3 lines before each match\\
\color{white}rg -C 3 PATTERN     \# 3 lines before AND after\\
\vspace{4pt}\textbf{\color{magenta}MULTILINE MATCHING}\\
\color{white}rg -U PATTERN       \# --multiline: match across lines\\
\texttt{\color{green}rg -U 'struct \textbackslash{}\textbackslash{}\{[\textbackslash{}\textbackslash{}s\textbackslash{}\textbackslash{}S]*?\textbackslash{}\textbackslash{}\}' \# match entire struct blocks}\\
\color{white}Without -U, patterns match within single lines only.\\
\vspace{4pt}\textbf{\color{magenta}THE COMPETITION}\\
\color{white}grep -r  — universal, slow, no .gitignore, no parallelism\\
\color{white}ag      — the silver searcher. fast, .gitignore aware.\\
\color{white}rg is 2-5x faster on most benchmarks.\\
\color{white}ack     — Perl-based, pioneered the \textbackslash{}"better grep\textbackslash{}" category.\\
\color{white}slower than both ag and rg. still has loyal users.\\

\vspace{8pt}

\subsection{Page 162: ripgrep: Advanced \& zpwr Integration}

\vspace{4pt}\textbf{\color{cyan}=== RIPGREP: ADVANCED \& ZPWR INTEGRATION ===}\vspace{2pt}\\
\textit{\color{gray}// rg --json | jq — structured grep for the data-obsessed. //}\\
\vspace{4pt}\textbf{\color{magenta}STRUCTURED OUTPUT}\\
\color{white}rg --json PATTERN   \# JSON output: one object per line\\
\color{white}Each line is a JSON object with type (begin, match, end,\\
\color{white}summary), path, line number, match offsets. Pipe to jq:\\
\texttt{\color{green}rg --json TODO | jq 'select(.type==\textbackslash{}"match\textbackslash{}") | .data.path'}\\
\vspace{4pt}\textbf{\color{magenta}SEARCH \& REPLACE PREVIEW}\\
\color{white}rg --replace 'NEW' PATTERN  \# preview replacements\\
\color{white}This prints what the output WOULD look like with replacements\\
\color{white}but does NOT modify files. It's a dry-run. Capture groups work:\\
\texttt{\color{green}rg 'foo(\textbackslash{}\textbackslash{}w+)' --replace 'bar\textbackslash{}\$1'  \# foo\_x -> bar\_x}\\
\vspace{4pt}\textbf{\color{magenta}STATISTICS}\\
\color{white}rg --stats PATTERN   \# show match stats after results\\
\color{white}Prints: N matches, N matched lines, N files searched,\\
\color{white}N files matched, bytes searched, duration. Quick audit.\\
\vspace{4pt}\textbf{\color{magenta}CUSTOM FILE TYPES}\\
\texttt{\color{green}rg --type-add 'zpwr:*.zsh' -t zpwr PATTERN}\\
\color{white}Define a custom type inline and use it immediately.\\
\color{white}Persist custom types in \textasciitilde{}/.ripgreprc (set RIPGREP\_CONFIG\_PATH).\\
\vspace{4pt}\textbf{\color{magenta}ZPWR INTEGRATION}\\
\color{white}zpwr grep uses ag (silver searcher) piped into fzf:\\
\texttt{\color{green}zpwr grep            \# ag into fzf for interactive search}\\
\color{white}zpwr wordsearch searches file contents with fzf preview:\\
\texttt{\color{green}zpwr wordsearch      \# search by word in subdirectory files}\\
\color{white}Both use bat-powered previews in the fzf window.\\
\vspace{4pt}\textbf{\color{magenta}FZF + RG LIVE RELOAD}\\
\color{white}The ultimate interactive grep: rg feeds fzf, bat previews:\\
\texttt{\color{green}rg --line-number PATTERN |}\\
\color{white}fzf --delimiter : --preview 'bat --highlight-line \{2\} \{1\}'\\
\color{white}rg outputs file:line:match, fzf splits on :, bat highlights\\
\color{white}the matching line. Three tools, one seamless experience.\\
\vspace{4pt}\textbf{\color{magenta}CONFIGURATION}\\
\color{white}export RIPGREP\_CONFIG\_PATH=\textasciitilde{}/.ripgreprc\\
\color{white}Put default flags in this file, one per line:\\
\color{white}--smart-case, --hidden, --glob=!.git, --max-columns=200\\
\color{white}Also: .rgignore in any directory works like .gitignore —\\
\color{white}project-specific patterns that rg will skip.\\
\color{white}.ignore files (used by ag too) are also respected by rg.\\

\vspace{8pt}


\section{CHAPTER 38: NEOVIM}

\subsection{Page 163: Neovim: The Next Generation}

\vspace{4pt}\textbf{\color{cyan}=== NEOVIM: THE NEXT GENERATION ===}\vspace{2pt}\\
\textit{\color{gray}// vim, but the timeline where Bram had unlimited funding. //}\\
\color{white}Neovim is the modernized fork of vim. Refactored internals,\\
\color{white}async everything, built-in LSP, Lua plugin API. In zpwr,\\
\texttt{\color{green}the default editor: ZPWR\_VIM='nvim' in .zpwr\_env.sh.}\\
\vspace{4pt}\textbf{\color{magenta}THE INIT.VIM BRIDGE}\\
\color{white}\textbackslash{}PWR\_INSTALL/init.vim is just three lines:\\
\color{white}set runtimepath\textasciicircum{}=\textasciitilde{}/.vim runtimepath+=\textasciitilde{}/.vim/after\\
\color{white}let \&packpath = \&runtimepath\\
\color{white}source \textasciitilde{}/.vimrc\\
\color{white}This means neovim sources the same .vimrc as vim. One config\\
\color{white}to rule them both. All 80+ vim plugins work in either editor.\\
\color{white}The symlink is created by plugins\_install.sh:\\
\color{white}ln -sf \textbackslash{}PWR\_INSTALL/init.vim \textasciitilde{}/.config/nvim/init.vim\\
\vspace{4pt}\textbf{\color{magenta}NEOVIM-SPECIFIC FEATURES}\\
\color{white}Async job control  — run external commands without blocking\\
\color{white}the editor. vim 8 added this too, but\\
\color{white}nvim's API is cleaner (jobstart/jobstop).\\
\color{white}Built-in LSP       — nvim 0.5+ has a native LSP client.\\
\color{white}Configure with lua: vim.lsp.start().\\
\color{white}No plugin needed for basic LSP.\\
\color{white}Lua plugin API     — write plugins in Lua, not just VimScript.\\
\color{white}Lua runs 10-100x faster than VimScript.\\
\color{white}:terminal          — built-in terminal emulator. Better than\\
\color{white}vim's :terminal — proper job control,\\
\color{white}scrollback, and buffer integration.\\
\color{white}True color         — 24-bit color out of the box. No termguicolors\\
\color{white}hacks needed (though the option exists).\\
\vspace{4pt}\textbf{\color{magenta}PROVIDER SUPPORT}\\
\color{white}nvim uses providers for python3 and node integration:\\
\color{white}pip3 install pynvim          \# python3 provider\\
\color{white}npm install -g neovim        \# node.js provider\\
\color{white}Run :checkhealth to verify all providers are working.\\
\vspace{4pt}\textbf{\color{magenta}ZPWR VIM INFO}\\
\color{white}zpwr catviminfo and zpwr catrecentfviminfo read the nvim\\
\color{white}shada (shared data) file — the nvim equivalent of .viminfo.\\

\vspace{8pt}

\subsection{Page 164: Neovim: Configuration \& Tips}

\vspace{4pt}\textbf{\color{cyan}=== NEOVIM: CONFIGURATION \& TIPS ===}\vspace{2pt}\\
\textit{\color{gray}// one vimrc. two editors. zero compromises. //}\\
\vspace{4pt}\textbf{\color{magenta}THE SHARED CONFIG PHILOSOPHY}\\
\color{white}zpwr's init.vim just sources \textasciitilde{}/.vimrc. This means:\\
\color{white}1. Every plugin in .vimbundle works in both vim and nvim\\
\color{white}2. All keybindings, settings, and autocmds are shared\\
\color{white}3. One file to edit: zpwr vrc opens .vimrc\\
\color{white}4. No divergent configs to maintain across editors\\
\color{white}The .vimrc is the single source of truth. init.vim is a shim.\\
\vspace{4pt}\textbf{\color{magenta}NEOVIM-SPECIFIC SETTINGS}\\
\color{white}Guard nvim-only settings with has('nvim'):\\
\color{white}if has('nvim')\\
\color{white}set inccommand=nosplit   \textbackslash{}" live substitution preview\\
\color{white}set pumblend=15          \textbackslash{}" popup menu transparency\\
\color{white}set winblend=15          \textbackslash{}" floating window transparency\\
\color{white}endif\\
\color{white}inccommand is a killer feature: :s/old/new shows replacements\\
\color{white}live as you type. No more guessing what your regex will do.\\
\vspace{4pt}\textbf{\color{magenta}:checkhealth}\\
\color{white}Run this in nvim to diagnose your setup. It checks:\\
\color{white}- python3/node providers (pynvim, neovim npm package)\\
\color{white}- clipboard support (pbcopy on macOS, xclip on Linux)\\
\color{white}- terminal capabilities, true color support\\
\color{white}- plugin health (many plugins register health checks)\\
\color{white}First thing to run when something feels wrong.\\
\vspace{4pt}\textbf{\color{magenta}SESSION TRACKING}\\
\color{white}tpope/vim-obsession (in .vimbundle) tracks sessions:\\
\color{white}:Obsession           \# start tracking session to Session.vim\\
\color{white}:Obsession!          \# stop tracking and delete Session.vim\\
\color{white}nvim -S Session.vim  \# restore session\\
\color{white}Works in both vim and nvim since it's a VimScript plugin.\\
\vspace{4pt}\textbf{\color{magenta}WHY ZPWR SUPPORTS BOTH}\\
\color{white}nvim is better locally — async, LSP, Lua, true color. But:\\
\color{white}- Remote servers often only have vim (or vi)\\
\color{white}- Docker containers ship vim, not nvim\\
\color{white}- macOS ships vim in /usr/bin/vim\\
\color{white}The shared .vimrc means your muscle memory works everywhere.\\
\color{white}nvim when you can, vim when you must, same config either way.\\
\vspace{4pt}\textbf{\color{magenta}ZPWR EDITOR VERBS}\\
\texttt{\color{green}zpwr vrc             \# edit .vimrc (shared by both)}\\
\color{white}Both editors, one config, zero excuses. The way it should be.\\

\vspace{8pt}


\section{CHAPTER 39: PERL IN ZPWR}

\subsection{Page 165: Perl Oneliners: The Secret Weapon}

\vspace{4pt}\textbf{\color{cyan}=== PERL ONELINERS: THE SECRET WEAPON ===}\vspace{2pt}\\
\color{white}sed wishes it could do this. awk cries in a corner.\\
\color{white}Perl is zpwr's scalpel for text surgery. Where sed and awk\\
\color{white}stumble on complex patterns, perl oneliners slice clean.\\
\color{white}zpwr uses perl throughout for text transformation, parsing,\\
\color{white}delimiter handling, and file manipulation.\\
\vspace{4pt}\textbf{\color{magenta}THE KEY FLAGS}\\
\color{white}-e 'code'     execute code\\
\color{white}-n            loop over input lines (like sed)\\
\color{white}-p            loop and print each line (like sed -p)\\
\color{white}-l            chomp newlines, add them back on print\\
\color{white}-a            autosplit into @F array (like awk)\\
\color{white}-F'sep'       set field separator for -a\\
\color{white}-i            in-place edit (like sed -i)\\
\color{white}-E            enable say, state, switch (modern perl)\\
\vspace{4pt}\textbf{\color{magenta}WHITESPACE CLEANUP (scripts/lib.sh)}\\
\texttt{\color{green}perl -pi -e 's@\textbackslash{}\textbackslash{}s+\textbackslash{}\$@\textbackslash{}\textbackslash{}n@g; s@\textbackslash{}\textbackslash{}x09\textbackslash{}\$@    @g'}\\
\color{white}Strips trailing whitespace, converts trailing tabs to spaces.\\
\color{white}The -pi combo: -p loops, -i edits in place. zpwr's file cleaner.\\
\vspace{4pt}\textbf{\color{magenta}PRINT SEPARATOR LINES}\\
\texttt{\color{green}perl -le \textbackslash{}"print '\#'x80\textbackslash{}"}\\
\color{white}Prints 80 \# characters. The x operator repeats a string.\\
\color{white}Used in zpwrPrettyPrint for visual separators.\\
\vspace{4pt}\textbf{\color{magenta}DELIMITER TRANSFORMATION}\\
\texttt{\color{green}perl -F\textbackslash{}"\textbackslash{}PWR\_DELIMITER\_CHAR\textbackslash{}" -anE 'say ...'}\\
\color{white}Splits lines on ZPWR\_DELIMITER\_CHAR into @F fields.\\
\color{white}Used heavily in zpwrContribCount, zpwrTotalLines for\\
\color{white}reformatting git contribution data into tabular output.\\
\vspace{4pt}\textbf{\color{magenta}FILTER BLANK LINES}\\
\texttt{\color{green}perl -ne 'print if /\textbackslash{}\textbackslash{}S/'}\\
\color{white}Only prints lines containing non-whitespace.\\
\color{white}Used in zpwrRegenAllKeybindingsCache to clean keybinding dumps.\\
\vspace{4pt}\textbf{\color{magenta}THE PERL REPL ALIAS}\\
\texttt{\color{green}plr (alias for rlwrap perl -wnE'say eval()//\textbackslash{}\$@')}\\
\color{white}Interactive perl REPL with readline. Type expressions, get results.\\

\vspace{8pt}

\subsection{Page 166: Perl Oneliners: Advanced Patterns}

\vspace{4pt}\textbf{\color{cyan}=== PERL ONELINERS: ADVANCED PATTERNS ===}\vspace{2pt}\\
\color{white}when you need a chainsaw, not a butter knife\\
\vspace{4pt}\textbf{\color{magenta}PATH QUOTING (zpwrGetFound)}\\
\texttt{\color{green}perl -pe 's@\textasciicircum{}([\textasciitilde{}]*)([\textasciicircum{}\textasciitilde{}].*)\textbackslash{}\$@\textbackslash{}\$1\textbackslash{}"\textbackslash{}\$2\textbackslash{}"@;s@\textbackslash{}\textbackslash{}s+@ @g'}\\
\color{white}Wraps non-tilde paths in double quotes and collapses spaces.\\
\color{white}Handles filenames with spaces for safe shell consumption.\\
\vspace{4pt}\textbf{\color{magenta}FZF RESULT FORMATTING (zpwrAgIntoFzf)}\\
\texttt{\color{green}perl -pe 's@\textasciicircum{}(.*?):(\textbackslash{}\textbackslash{}d+):(.*)@+\textbackslash{}\$2 \textbackslash{}"PWD/\textbackslash{}\$1\textbackslash{}"@'}\\
\color{white}Converts grep output (file:line:match) into vim args (+line \textbackslash{}"file\textbackslash{}").\\
\color{white}The brain behind zpwr grep -> fzf -> vim pipeline.\\
\vspace{4pt}\textbf{\color{magenta}VERB LIST FORMATTING (zpwrVerbsFZF)}\\
\texttt{\color{green}perl -e '@a=<>;for(@a)\{print \textbackslash{}"zpwr \textbackslash{}\$1\textbackslash{}" if m\{\textasciicircum{}(\textbackslash{}\textbackslash{}S+)\textbackslash{}\textbackslash{}s+\}\}'}\\
\color{white}Reads verb list, extracts first word, prepends 'zpwr '.\\
\color{white}Builds the command string that fzf executes on selection.\\
\vspace{4pt}\textbf{\color{magenta}IN-PLACE RENAME (zpwrChangeNameZpwr)}\\
\texttt{\color{green}perl -i -pe 's@\textbackslash{}\textbackslash{}bOLDNAME\textbackslash{}\textbackslash{}b@NEWNAME@g' **/*(.)'}\\
\color{white}Renames all occurrences across all files. The \textbackslash{}\textbackslash{}b word boundary\\
\color{white}prevents partial matches. The **(.) glob hits all regular files.\\
\vspace{4pt}\textbf{\color{magenta}LOGIN COUNTING (zpwrLoginCount)}\\
\texttt{\color{green}perl -e '"'"'print qx(last -f "\$\_") for </var/log/wtmp*>'"'"'}\\
\color{white}Reads all wtmp files and runs last on each. The diamond\\
\color{white}operator <> globs and iterates. Then pipes to:\\
\texttt{\color{green}perl -lane '"'"'print [0] if /\textbackslash{}S+/ \&\& !/wtmp/'"'"'}\\
\color{white}Extracts username (first field) from non-empty, non-header lines.\\
\vspace{4pt}\textbf{\color{magenta}FILE COUNT DEDUP (zpwrEnvCounts)}\\
\texttt{\color{green}perl -lne '@l=<>;@u=do\{my \%seen;grep\{!\textbackslash{}\$seen\{\textbackslash{}\$\_\}++\}@l\};'}\\
\color{white}Reads all lines, deduplicates with a hash, then checks if\\
\color{white}each unique path is a real file. The densest oneliner in zpwr.\\
\color{white}Counts unique zpwr files across all tracked locations.\\
\vspace{4pt}\textbf{\color{magenta}IP ADDRESS EXTRACTION (ipa alias)}\\
\texttt{\color{green}ifconfig | perl -lane 'do \{print \textbackslash{}[1]=\textasciitilde{}s/addr//r;exit0\} if /inet\textbackslash{}\textbackslash{}s/ and !/127/'}\\
\textit{\color{gray}Finds the first non-loopback inet address. The s///r flag}\\
\color{white}returns the modified string without changing the original.\\
\vspace{4pt}\textbf{\color{magenta}WHY PERL OVER SED/AWK?}\\
\color{white}1. Perl regex is more powerful (lookahead, backrefs, /x)\\
\color{white}2. -i flag for in-place edit is reliable across OS\\
\color{white}3. Complex logic in one line (if/unless/for/do blocks)\\
\color{white}4. @F autosplit is awk-like but with perl's full stdlib\\
\color{white}5. Perl is installed everywhere (even minimal containers)\\

\vspace{8pt}


\section{CHAPTER 40: HIDDEN SUBSYSTEMS}

\subsection{Page 167: Zconvey: Terminal Communication}

\vspace{4pt}\textbf{\color{cyan}=== ZCONVEY: TERMINAL COMMUNICATION ===}\vspace{2pt}\\
\color{white}every terminal screams its identity into the void\\
\vspace{4pt}\textbf{\color{magenta}WHAT IS ZCONVEY?}\\
\color{white}MenkeTechnologies/zconvey - a zinit plugin\\
\color{white}Allows sending commands between terminal sessions.\\
\color{white}Your terminals can talk to each other. Yes, really.\\
\color{white}Loaded via zinit in \textbackslash{}PWR\_INSTALL/.zshplugins\\
\vspace{4pt}\textbf{\color{magenta}HOW ZPWR USES IT}\\
\color{white}Every time you press Enter, zpwrAcceptLine fires.\\
\color{white}It calls zc-rename to label the current terminal with:\\
\color{white}TTY:\textbackslash{} PID:\textbackslash{}\$\{\$\} CMD:\textbackslash{}UFFER PWD:\textbackslash{} DATE:\textbackslash{}\$(date)\\
\color{white}This broadcasts your terminal's state to the zconvey bus.\\
\color{white}Every pane, every tab, every session - all labeled.\\
\vspace{4pt}\textbf{\color{magenta}THE CONVEY VARIABLES}\\
\color{white}ZPWR\_CONVEY\_NAME      the full label string sent to zconvey\\
\color{white}ZPWR\_CONVEY\_LAST\_CMD  the last command (for zconvey commands)\\
\color{white}When you run a zc-* command, it preserves the previous label\\
\color{white}so your terminal doesn't get renamed to \textbackslash{}"zc-send ls\textbackslash{}".\\
\vspace{4pt}\textbf{\color{magenta}THE RENAME LOGIC}\\
\texttt{\color{green}if ! [[ \textbackslash{}\$(zpwrExpandAliases \textbackslash{}UFFER) = zc* ]]; then}\\
\color{white}ZPWR\_CONVEY\_NAME=\textbackslash{}"TTY:\textbackslash{} PID:\textbackslash{}\$\{\$\} CMD:...\textbackslash{}"\\
\color{white}zc-rename \textbackslash{}PWR\_CONVEY\_NAME\\
\color{white}If BUFFER starts with zc*, keep the old label.\\
\color{white}Otherwise, stamp the terminal with the new command.\\
\vspace{4pt}\textbf{\color{magenta}WHY THIS MATTERS}\\
\color{white}tmux status bar can show what each pane is running.\\
\color{white}You can identify panes by their last command and PWD.\\
\color{white}Remote debugging: see exactly what terminal ran what.\\
\color{white}It's like process accounting, but for your shell sessions.\\
\vspace{4pt}\textbf{\color{magenta}SENDING COMMANDS BETWEEN TERMINALS}\\
\color{white}zc-send 2 'cd /tmp \&\& ls'  send command to terminal 2\\
\color{white}zc-all 'echo hello'        broadcast to all terminals\\
\color{white}Your terminals are a mesh network. Act accordingly.\\
\color{white}SIGNAL: zpwrAcceptLine is the heartbeat of zpwr.\\
\color{white}Every Enter keypress: rename terminal, check for clear,\\
\color{white}detect file-modifying commands, send tmux keys.\\
\color{white}zconvey makes that heartbeat visible across sessions.\\

\vspace{8pt}

\subsection{Page 168: The 24-Hour Updater Loop}

\vspace{4pt}\textbf{\color{cyan}=== THE 24-HOUR UPDATER LOOP ===}\vspace{2pt}\\
\color{white}your system updates itself while you sleep\\
\vspace{4pt}\textbf{\color{magenta}THE INFINITE LOOP}\\
\color{white}\textbackslash{}PWR\_SCRIPTS/autoUpdater.sh runs forever.\\
\color{white}Every 24 hours it wakes up and updates everything.\\
\color{white}while true; do\\
\color{white}bash -l updater.sh -e\\
\color{white}\# sleep in 10-minute chunks until 24h elapsed\\
\color{white}done\\
\color{white}Uses Perl Time::Piece to calculate the next update time.\\
\color{white}Logs the time diff on each wake to \textbackslash{}PWR\_LOGFILE.\\
\vspace{4pt}\textbf{\color{magenta}WHAT GETS UPDATED}\\
\color{white}updater.sh orchestrates the full update:\\
\color{white}zpwrUpdater      runs the main zpwr update sequence\\
\color{white}zpwrUpdateDeps   updates all package manager dependencies\\
\color{white}Homebrew, npm, pip, gem, go - everything gets refreshed.\\
\color{white}History files backed up. Logs rotated. Deps pinned.\\
\vspace{4pt}\textbf{\color{magenta}SYSTEMD SERVICE (LINUX)}\\
\color{white}On Linux, autoUpdater.sh runs as a systemd service.\\
\color{white}zpwr restart      restart the zpwr systemd service\\
\color{white}zpwr serviceup    start and enable the systemd service\\
\color{white}zpwr servicedown  stop and disable the systemd service\\
\color{white}On macOS, \textbackslash{}PWR\_SCRIPTS/updaterForLaunchCtl.sh is used.\\
\vspace{4pt}\textbf{\color{magenta}EMAIL NOTIFICATIONS}\\
\color{white}\textbackslash{}PWR\_SCRIPTS/updaterEmail.sh\\
\color{white}Pipes the update log to email via mutt.\\
\color{white}Wake up to a full report of what changed overnight.\\
\color{white}Your system admin is a bash script. It never sleeps.\\
\vspace{4pt}\textbf{\color{magenta}THE SLEEP STRATEGY}\\
\color{white}Doesn't sleep for 24h straight (that would drift).\\
\color{white}Sleeps in 10-minute intervals, checks elapsed time:\\
\color{white}((timediff > \textbackslash{}\$((3600 * 24)))) \&\& break\\
\color{white}This handles laptop suspends and clock adjustments.\\
\color{white}WARNING: The updater keeps zpwr and all deps current.\\
\color{white}If something breaks after an update, check the log:\\
\color{white}zpwr taillog  to see what the updater did last.\\
\color{white}The machine maintains itself. You just write code.\\

\vspace{8pt}

\subsection{Page 169: The 32-Pane REPL Layout}

\vspace{4pt}\textbf{\color{cyan}=== THE 32-PANE REPL LAYOUT ===}\vspace{2pt}\\
\color{white}every language, one screen, zero excuses\\
\vspace{4pt}\textbf{\color{magenta}THE TMUX KEYBINDINGS}\\
\color{white}prefix R  source \textbackslash{}PWR\_TMUX/thirtytwo-panes-repl.conf\\
\color{white}prefix O  source \textbackslash{}PWR\_TMUX/sixteen-panes.conf\\
\color{white}prefix G  source \textbackslash{}PWR\_TMUX/eight-panes.conf\\
\color{white}prefix F  source \textbackslash{}PWR\_TMUX/four-panes.conf\\
\color{white}One keystroke splits your terminal into a REPL army.\\
\vspace{4pt}\textbf{\color{magenta}THE 32 REPLs (prefix R)}\\
\color{white}pane 0:  lua          pane 1:  swift        pane 2:  ghci\\
\color{white}pane 3:  clisp        pane 4:  ocaml        pane 5:  clojure\\
\color{white}pane 6:  pry (ruby)   pane 7:  bpython      pane 8:  php -a\\
\color{white}pane 9:  groovysh     pane 10: csharp       pane 11: kotlinc\\
\color{white}pane 12: node         pane 13: coffee       pane 14: tclsh\\
\color{white}pane 15: ksh          pane 16: bash         pane 17: csh\\
\color{white}pane 18: perl -dE1    pane 19: gore (go)    pane 20: jshell\\
\color{white}pane 21: mycli        pane 22: pgcli        pane 23: python2\\
\color{white}pane 24: nslookup     pane 25: mongo        pane 26: irb\\
\color{white}pane 27: scala        pane 28: utop         pane 29: iex\\
\color{white}pane 30: bc           pane 31: gforth\\
\color{white}32 languages. 32 panes. One tmux window. Absolute madness.\\
\vspace{4pt}\textbf{\color{magenta}THE VIM-TMUX EXECUTION PIPELINE}\\
\color{white}Write code in vim, send it to a REPL pane:\\
\color{white}Supported languages for vim-to-tmux execution:\\
\color{white}zsh, bash, python, ruby, perl, node, go\\
\color{white}java, scala, R, octave\\
\color{white}Select code in vim, fire it into the matching pane.\\
\color{white}Edit-run-edit-run without ever leaving vim.\\
\vspace{4pt}\textbf{\color{magenta}PROGRESSIVE LAYOUTS}\\
\color{white}Don't need 32 languages? Scale down:\\
\color{white}prefix G  8 panes   - quick polyglot session\\
\color{white}prefix O  16 panes  - serious multi-language work\\
\color{white}prefix R  32 panes  - full REPL matrix (needs big screen)\\
\color{white}Each layout is a tmux conf file in \textbackslash{}PWR\_TMUX/\\
\color{white}NOTE: prefix R on a 13\textbackslash{}" laptop is an act of hubris.\\
\color{white}You'll need at least a 27\textbackslash{}" monitor to read anything.\\
\color{white}But hey, it works. And it's glorious.\\

\vspace{8pt}

\subsection{Page 170: Cross-Platform Installer}

\vspace{4pt}\textbf{\color{cyan}=== CROSS-PLATFORM INSTALLER ===}\vspace{2pt}\\
\color{white}one repo to rule them all, one script to install them\\
\vspace{4pt}\textbf{\color{magenta}INSTALLER SCRIPTS IN \textbackslash{}PWR\_INSTALL}\\
\color{white}formulae.sh        Homebrew formulae (macOS/Linux)\\
\color{white}casks.sh           Homebrew casks (macOS GUI apps)\\
\color{white}mincasks.sh        minimal cask set for lean installs\\
\color{white}taps.sh            Homebrew tap repositories\\
\color{white}npm\_install.sh     Node.js packages via npm\\
\color{white}npm.sh             npm package list\\
\color{white}gems.sh            Ruby gems\\
\color{white}pip\_install.sh     Python packages via pip\\
\color{white}pip.sh / pip3.sh   pip package lists\\
\color{white}go\_install.sh      Go packages\\
\color{white}rustupinstall.sh   Rust toolchain via rustup\\
\vspace{4pt}\textbf{\color{magenta}EDITOR \& PLUGIN INSTALLERS}\\
\color{white}neovim\_install.sh       Neovim setup\\
\color{white}vim\_install.sh          Vim from source\\
\color{white}vim\_plugins\_install.sh  Vim plugin manager + plugins\\
\color{white}zsh\_plugins\_install.sh  Zsh plugins via zinit\\
\color{white}tmux\_plugins\_install.sh Tmux plugin manager + plugins\\
\color{white}spacemacs\_install.sh    Spacemacs (Emacs) setup\\
\vspace{4pt}\textbf{\color{magenta}SHARED INFRASTRUCTURE}\\
\color{white}common.sh  shared installer functions used by all scripts\\
\color{white}zpwrInstall.sh  the master installer orchestrator\\
\color{white}Supports 20+ Linux distros, macOS, and FreeBSD.\\
\color{white}Detects your OS and calls the right package manager.\\
\vspace{4pt}\textbf{\color{magenta}ZPWR VERBS FOR INSTALLATION}\\
\color{white}zpwr install            run configure, make, make install\\
\color{white}zpwr uninstall          uninstall all zpwr configs\\
\color{white}zpwr regenconfiglinks   regen symlinks from \textbackslash{}PWR\_INSTALL to \textbackslash{}OME\\
\color{white}regenconfiglinks recreates dotfile symlinks after a reinstall.\\
\vspace{4pt}\textbf{\color{magenta}THE SYMLINK STRATEGY}\\
\color{white}zpwr doesn't copy config files to \textbackslash{}OME. It symlinks them.\\
\color{white}.vimrc, .tmux.conf, .gitconfig - all point back to \textbackslash{}PWR\_INSTALL.\\
\color{white}Edit once in the repo, git commit, and every machine gets it.\\
\texttt{\color{green}zpwr regenconfiglinks rebuilds them if they break.}\\
\color{white}PHILOSOPHY: Package lists are just text files.\\
\color{white}Add a line to formulae.sh, and your next machine gets it.\\
\color{white}Your entire dev environment is version-controlled.\\

\vspace{8pt}


\section{CHAPTER 41: QUALITY ASSURANCE}

\subsection{Page 171: The Test Suite}

\vspace{4pt}\textbf{\color{cyan}=== THE TEST SUITE ===}\vspace{2pt}\\
\color{white}60 test files standing between you and chaos\\
\vspace{4pt}\textbf{\color{magenta}TESTING FRAMEWORK}\\
\color{white}zpwr uses zunit - a testing framework written in zsh.\\
\color{white}Config lives in \textbackslash{}PWR/.zunit.yml\\
\color{white}Tests dir: \textbackslash{}PWR\_TEST  Output: \textbackslash{}PWR\_TEST/\_output\\
\color{white}Support helpers: \textbackslash{}PWR\_TEST/\_support\\
\vspace{4pt}\textbf{\color{magenta}TEST CATEGORIES}\\
\color{white}t-autoload-*       verify all autoloaded functions\\
\color{white}load, defined, existence, fpath, whence, type, body\\
\color{white}Every autoload function must: exist, load, have a body.\\
\color{white}t-autoload-subdir-* same checks for subdirectory autoloads\\
\color{white}t-env-*            verify environment variables\\
\color{white}defined, exported, comprehensive, type checks\\
\color{white}Every ZPWR\_* var must be defined and exported.\\
\color{white}t-scripts-*        verify script files\\
\color{white}executable, readable, non-empty, have shebangs, valid syntax\\
\color{white}t-verbs-*          verify zpwr verb definitions\\
\color{white}key format, non-empty values, callable, type checks\\
\vspace{4pt}\textbf{\color{magenta}MORE TEST FILES}\\
\color{white}t-functions.zsh     test core shell functions\\
\color{white}t-callable-fns.zsh  verify functions are callable\\
\color{white}t-navigation.zsh    test directory navigation\\
\color{white}t-install-files.zsh verify installer file integrity\\
\color{white}t-repo-files.zsh    verify repo structure\\
\color{white}t-plugin-counts.zsh verify plugin load counts\\
\vspace{4pt}\textbf{\color{magenta}RUNNING TESTS}\\
\color{white}zpwr test     run all zpwr tests\\
\color{white}zpwr testall  run ALL env tests (comprehensive sweep)\\
\color{white}fail\_fast: true means first failure stops the run.\\
\color{white}verbose: true means you see every assertion.\\
\color{white}GUARANTEE: Every autoload loads. Every env var defined.\\
\color{white}Every script executable. Every verb callable.\\
\color{white}60 test files verify zpwr won't break your shell.\\
\color{white}If tests pass, your environment is sound. Ship it.\\

\vspace{8pt}


\section{CHAPTER 42: COLORIZATION STACK}

\subsection{Page 172: GRC: Generic Colorizer}

\vspace{4pt}\textbf{\color{cyan}=== GRC: GENERIC COLORIZER ===}\vspace{2pt}\\
\color{white}because monochrome ls is a crime against humanity\\
\vspace{4pt}\textbf{\color{magenta}WHAT IS GRC?}\\
\color{white}grc wraps commands and adds color output via regex rules.\\
\color{white}It reads a config file that maps patterns to colors.\\
\color{white}Any command's output can be colorized. Zero modifications.\\
\vspace{4pt}\textbf{\color{magenta}GRC IN zpwrListNoClear}\\
\color{white}The fallback chain when listing files:\\
\color{white}1. eza              preferred (color-scale, git status)\\
\color{white}2. grc + ls         fallback (colorized via conf.gls)\\
\color{white}3. plain ls         last resort (no color config)\\
\vspace{4pt}\textbf{\color{magenta}PLATFORM-SPECIFIC COMMANDS}\\
\color{white}macOS:\\
\color{white}grc -c \textbackslash{}OME/conf.gls gls -iFlhA --color=always\\
\color{white}Linux:\\
\color{white}grc -c \textbackslash{}OME/conf.gls ls -iFlhA --color=always\\
\color{white}macOS uses gls (GNU ls from coreutils) because BSD ls\\
\color{white}doesn't support all the flags. Linux uses native ls.\\
\vspace{4pt}\textbf{\color{magenta}THE conf.gls CONFIG}\\
\color{white}Custom color configuration for ls output, lives at:\\
\color{white}\textbackslash{}OME/conf.gls\\
\color{white}Maps file types, permissions, sizes to ANSI colors.\\
\color{white}Directories, executables, symlinks - all distinct colors.\\
\vspace{4pt}\textbf{\color{magenta}OTHER GRC TARGETS}\\
\color{white}grc can colorize almost any command:\\
\color{white}ping         latency numbers in color\\
\color{white}traceroute   hop-by-hop color coding\\
\color{white}make         errors red, warnings yellow\\
\color{white}gcc          compiler output colorized\\
\color{white}netstat      connection states highlighted\\
\color{white}diff         additions green, deletions red\\
\vspace{4pt}\textbf{\color{magenta}GRC IN LOG COLORIZATION}\\
\color{white}\textbackslash{}PWR\_SCRIPTS/colorLogger.sh uses grc for system logs.\\
\color{white}tmux control windows show colorized live log output.\\
\color{white}TIP: If eza is missing and ls looks dull, install grc.\\
\color{white}It's the middle tier of zpwr's color fallback chain.\\
\color{white}One package, instant color for dozens of commands.\\

\vspace{8pt}

\subsection{Page 173: Ponysay, Lolcat \& ASCII Art}

\vspace{4pt}\textbf{\color{cyan}=== PONYSAY, LOLCAT \& ASCII ART ===}\vspace{2pt}\\
\color{white}enterprise-grade cartoon horse infrastructure\\
\vspace{4pt}\textbf{\color{magenta}THE ASCII ART TOOLCHAIN}\\
\color{white}ponysay    My Little Pony ASCII art (cowsay but fabulous)\\
\color{white}lolcat     rainbow gradient text output\\
\color{white}figlet     large ASCII text banners from font files\\
\color{white}cowsay     classic ASCII art speech bubbles\\
\color{white}toilet     another ASCII art text renderer\\
\color{white}zpwr combines all of these for its startup banner.\\
\vspace{4pt}\textbf{\color{magenta}ZPWR BANNER VERBS}\\
\color{white}zpwr bannerpony      hostname banner with ponysay\\
\color{white}zpwr bannerlolcat    ponysay + lolcat rainbow gradient\\
\color{white}zpwr bannernopony    banner without ponysay\\
\color{white}zpwr figletfonts     preview all installed figlet fonts\\
\color{white}Each verb calls a different autoloaded function:\\
\texttt{\color{green}zpwrPonyBanner, zpwrBannerLolcat, zpwrNoPonyBanner}\\
\vspace{4pt}\textbf{\color{magenta}THE STARTUP BANNER}\\
\color{white}On shell startup, zpwrDarwinBanner / zpwrLinuxBanner fires.\\
\color{white}It picks a random figlet font each time via:\\
\color{white}\textbackslash{}PWR\_SCRIPTS/macOnly/figletRandomFontOnce.sh\\
\color{white}ZPWR\_BANNER\_TYPE controls pony vs no-pony:\\
\color{white}ponies  -> zpwrBannerLolcat (rainbow horse)\\
\color{white}other   -> zpwrNoPonyBanner (figlet only)\\
\vspace{4pt}\textbf{\color{magenta}BANNER CONTROL}\\
\color{white}ZPWR\_BANNER\_CLEARLIST  controls startup banner + ls display\\
\color{white}When true: clear screen, show banner, then zpwrListNoClear.\\
\color{white}When false: skip the banner entirely (fast startup).\\
\vspace{4pt}\textbf{\color{magenta}FONT LOGGING}\\
\color{white}The random figlet font name is logged to:\\
\color{white}\textbackslash{}PWR\_LOGFILE\\
\color{white}Want to know which font made that gorgeous banner?\\
\color{white}zpwr taillog  and look for the font name.\\
\color{white}Found one you love? Set it permanently in your config.\\
\color{white}REALITY CHECK: A rainbow pony greeting you on every\\
\color{white}terminal open is either peak engineering or peak insanity.\\
\color{white}zpwr does not distinguish between the two.\\
\color{white}Your hostname, rendered in ASCII art, by a cartoon horse.\\
\color{white}This is what unlimited power looks like.\\

\vspace{8pt}

\subsection{Page 174: GRC, ccze \& the Color Pipeline}

\vspace{4pt}\textbf{\color{cyan}=== GRC, CCZE \& THE COLOR PIPELINE ===}\vspace{2pt}\\
\color{white}every byte of output deserves to be beautiful\\
\vspace{4pt}\textbf{\color{magenta}THE FULL COLORIZATION STACK}\\
\color{white}zpwr has a color tool for every type of output:\\
\color{white}FILE LISTING\\
\color{white}eza (color-scale, git status) -> grc+ls -> plain ls\\
\color{white}FILE VIEWING\\
\color{white}bat (syntax highlighting) -> cat\\
\texttt{\color{green}ZPWR\_COLORIZER=bat in \textbackslash{}PWR\_ENV/.zpwr\_env.sh}\\
\color{white}LOG VIEWING\\
\color{white}ccze (log colorizer) -> plain tail\\
\color{white}zpwrTailLog pipes \textbackslash{}PWR\_LOGFILE through ccze.\\
\color{white}GIT DIFFS\\
\color{white}delta (syntax-aware diff) -> plain diff\\
\color{white}delta adds line numbers, word-level highlighting,\\
\color{white}and side-by-side view with Dracula theme.\\
\vspace{4pt}\textbf{\color{magenta}CCZE: LOG FILE COLORIZER}\\
\color{white}ccze parses log formats and applies semantic colors:\\
\color{white}timestamps, severity levels, PIDs, IP addresses.\\
\color{white}zpwr taillog  uses ccze for colorized log tailing.\\
\color{white}System logs become readable at a glance.\\
\vspace{4pt}\textbf{\color{magenta}COLOR TESTING VERBS}\\
\color{white}zpwr colortest       test terminal basic 16 colors\\
\color{white}zpwr colortest256    test full 256-color palette\\
\color{white}zpwr pygmentcolors   show all pygment color schemes\\
\color{white}zpwr colorsdiff      colorized side-by-side diff\\
\color{white}Use these to verify your terminal supports the full palette.\\
\vspace{4pt}\textbf{\color{magenta}THE FALLBACK PHILOSOPHY}\\
\color{white}Every color tool has a plain fallback:\\
\color{white}No eza? grc+ls. No grc? plain ls. No bat? cat.\\
\color{white}zpwr degrades gracefully. Color is never a hard dependency.\\
\color{white}SSH into a minimal server? Everything still works.\\
\color{white}Just less pretty. Function over form, always.\\
\color{white}THE RESULT: In a fully-equipped zpwr terminal:\\
\color{white}File listings glow. Logs are semantic. Diffs are surgical.\\
\color{white}Your eyes never have to parse raw monochrome again.\\
\color{white}The color pipeline is zpwr's cosmetic surgery department.\\

\vspace{8pt}


\section{CHAPTER 43: THE END}

\subsection{Page 175: The Epilogue: You Are the Wizard}

\vspace{4pt}\textbf{\color{cyan}=== THE EPILOGUE: YOU ARE THE WIZARD ===}\vspace{2pt}\\
\color{white}175 pages later, the terminal bows to you\\
\vspace{4pt}\textbf{\color{magenta}WHAT YOU'VE LEARNED}\\
\color{white}175 pages across 43 chapters. From keybindings to kernel.\\
\color{white}Autoloaded functions. Tmux pipelines. Fzf integration.\\
\color{white}Vim execution. Perl oneliners. Systemd services.\\
\color{white}32-pane REPL layouts. Cross-platform installers.\\
\color{white}A 24-hour updater loop. A 60-file test suite.\\
\color{white}A colorization stack five tools deep.\\
\color{white}And a startup banner delivered by a rainbow pony.\\
\vspace{4pt}\textbf{\color{magenta}THE 5 COMMANDS TO MEMORIZE}\\
\color{white}zpwr verbs    list every zpwr subcommand\\
\color{white}zpwr doctor   run environment health check diagnostic\\
\color{white}zpwr bench    benchmark startup with percentiles\\
\color{white}zpwr wizard   interactive guided tour (you're here)\\
\color{white}zpwr study    split-pane learning: book left, shell right\\
\color{white}Five commands. That's your entire onboarding.\\
\vspace{4pt}\textbf{\color{magenta}HOW TO ADD YOUR OWN VERBS}\\
\color{white}1. Create an autoload function in \textbackslash{}PWR/autoload/common/\\
\texttt{\color{green}2. Add to ZPWR\_VERBS in \textbackslash{}PWR\_ENV/zpwr.zsh}\\
\color{white}Format: ZPWR\_VERBS[myverb]='myFunc=description'\\
\color{white}3. Reload: exec zsh. Your verb is live.\\
\vspace{4pt}\textbf{\color{magenta}HOW TO EXTEND THE WIZARD}\\
\color{white}Add a new page\_NNN.zsh to \textbackslash{}PWR\_INSTALL/wizard\_pages/\\
\color{white}Follow the format: PAGE\_TITLE, PAGE\_CHAPTER, PAGE\_CONTENT().\\
\color{white}The wizard auto-discovers pages. No registration needed.\\
\vspace{4pt}\textbf{\color{magenta}HOW TO CONTRIBUTE}\\
\color{white}Fork github.com/MenkeTechnologies/zpwr\\
\color{white}Add your feature. Write tests. Submit a PR.\\
\color{white}The codebase is the documentation. Read it.\\
\vspace{4pt}\textbf{\color{magenta}THE ZPWR PHILOSOPHY}\\
\color{white}Verbs, not scripts.  Discoverability over memorization.\\
\color{white}Speed over bloat.    Fallbacks over hard dependencies.\\
\color{white}Convention over config.  Your shell should just work.\\
\vspace{4pt}\textbf{\color{magenta}ONE LAST THING}\\
\color{white}Run zpwr about to see the zpwr banner.\\
\color{white}Then open a new terminal. Watch the pony greet you.\\
\color{white}You are the wizard now. The terminal is yours.\\
\vspace{4pt}\textbf{\color{cyan}>>> END OF TRANSMISSION <<<}\vspace{2pt}\\
\textit{\color{gray}// ZPWR WIZARD ENCYCLOPEDIA v1.0 // SIGNAL LOST //}\\
\color{white}May your shell be fast, your pipes be clean,\\
\color{white}and your exit codes always zero.\\

\vspace{8pt}


\end{document}
